% !TEX program = pdflatex
\documentclass[12pt]{article}

\usepackage[utf8]{inputenc}
\usepackage{CJKutf8}
\usepackage{amsmath,amssymb}
\usepackage{amsthm}
\usepackage{geometry}
\geometry{
  a4paper,
  top=25mm,
  bottom=25mm,
  left=20mm,
  right=20mm
}
\usepackage{xcolor}
\usepackage{url}
\usepackage[unicode,colorlinks=true,linkcolor=blue,urlcolor=blue]{hyperref}
\usepackage{xurl}
\Urlmuskip=0mu plus 2mu
\sloppy
\usepackage{tocloft}

\renewcommand{\cftsecleader}{\cftdotfill{\cftdotsep}}

\newtheorem{definition}{定义}[section]
\newtheorem{axiom}{公理}[section]
\newtheorem{assumption}{假设}[section]
\newtheorem{theorem}{定理}[section]
\newtheorem{proposition}{命题}[section]
\newtheorem{corollary}{推论}[section]
\newtheorem{lemma}{引理}[section]
\newtheorem{remark}{备注}[section]
\newtheorem{Certificate}{证书}[section]

\newcommand{\NN}{\mathbb{N}}
\newcommand{\code}[1]{\path{#1}}
\newcommand{\GFramework}{\textsf{G-Framework}}
\newcommand{\Ob}{\operatorname{Ob}}
\newcommand{\Sem}{\operatorname{Sem}}
\newcommand{\Norm}{\operatorname{Norm}}
\newcommand{\Tw}{\operatorname{Tw}}
\newcommand{\PL}{\operatorname{PL}}
\newcommand{\Parse}{\operatorname{Parse}}
\newcommand{\Audit}{\operatorname{Audit}}
\newcommand{\Cert}{\operatorname{Cert}}
\newcommand{\Auth}{\operatorname{Auth}}
\newcommand{\Conv}{\operatorname{Conv}}
\newcommand{\Downgrade}{\operatorname{Downgrade}}
\newcommand{\Refine}{\operatorname{Refine}}
\newcommand{\Close}{\operatorname{Close}}
\newcommand{\Issue}{\operatorname{Issue}}
\newcommand{\ExtAssump}{\operatorname{ExtAssump}}
\newcommand{\Render}{\operatorname{Render}}
\newcommand{\Select}{\operatorname{Select}}
\newcommand{\Promote}{\operatorname{Promote}}
\newcommand{\Style}{\operatorname{Style}}
\newcommand{\Engine}{\operatorname{Engine}}
\newcommand{\Log}{\operatorname{Log}}
\newcommand{\Img}{\operatorname{Im}}
\newcommand{\id}{\mathrm{id}}

\setlength{\emergencystretch}{6em}
\setlength{\parindent}{0pt}
\setlength{\parskip}{0.6em}

\begin{document}
\begin{CJK*}{UTF8}{gkai}

\title{G 框架正则语义规范场到自然语言语义逻辑占位规范场的\\高阶同伦扭曲双射、LLM 运行时降级与逻辑引擎\\（工作笔记：形式系统版）}
\author{Gao Zheng（高政）}
\date{2026-03-10}
\maketitle

\section*{前言}
\addcontentsline{toc}{section}{前言}

本文的目标不是孤立地给出一组语义规范场、自然语言占位、LLM 运行时或逻辑引擎
术语，而是在 \GFramework{} 的整体脉络中，把正则语义规范场、高阶同伦扭曲双射、
自然语言语义逻辑占位、确定性审计器、LLM 表面表达、上下文修复塔、领域细化插件
与签发收束协议组织成一条可追踪的形式链条。若这些对象分散出现，它们容易被误解为
一组围绕大模型提示词工程的局部技巧；本文将其压合为一个可以定义、可以规范化、
可以审计、可以降级运行、可以插件化扩张、可以回收求解决策主权的语义运行形式系统。

本文的第一层立场是语义主权的回收。自然语言在本文中不被提升为语义本体，LLM 也
不承担证明义务或最终裁决权。语义源头保留在 \GFramework{} 的正则语义规范场内；
自然语言只承担可读投影，LLM 只承担受控运行时表达。换言之，本文的核心链条不是
\[
  \text{自然语言生成语义},
\]
而是
\[
  \text{规范语义}
  \Longrightarrow
  \text{逻辑占位}
  \Longrightarrow
  \text{可审计自然语言表面表达}.
\]
这一层承接仓内“判定 + 规范化”接口与“LLM/NN 只提案、确定性审计器裁决”的口径
\cite{RepoTex20SemanticGaugeNormalization,RepoTex21AuditedWrapper,RepoTex23SemanticBase}。

本文的第二层立场是占位层的严格化。自然语言语义逻辑占位不是自然语言句子本身，
而是包含量词骨架、谓词依赖、证书槽、失败槽、引用接口与风格参数的逻辑容器。
因此，表达任务被拆成两步：先把 \GFramework{} 侧的规范代表搬运到占位层，再由运行时
渲染为自然语言。这样，表面文本可以变化，量词结构、证书引用、失败口径与审计日志
不能随之漂移。本文由此把“可读性”与“语义正确性”分离开来。

本文的第三层立场是高阶同伦扭曲双射的中介作用。正则语义规范场与自然语言逻辑
占位规范场之间，并不是靠一次性翻译表连接，而是靠逐阶兼容的同伦扭曲塔连接。
每一阶都必须保持语义、证书、失败与审计接口相容，最终形成
\[
  \mathsf{NF}^{(\infty)}_G
  \simeq
  \mathsf{PL}^{(\infty)}_{\mathrm{NL}}
\]
意义下的可审计对应。由此，本文把“翻译自然语言”改写为“在同伦扭曲保持下迁移规范
代表”，这与仓内关于路径积分支撑类、同伦扭曲正规形与终端编码的对象链
\cite{RepoTex36HomotopyTwistBijection} 相互衔接。

本文的第四层立场是运行时降级。LLM 在本文中不是认知主体、证明机或语义裁判，而是
\[
  \Engine_{\mathrm{LLM}}
\]
这一类可替换运行时组件。它的输出必须被解析、审计、证书检查与失败闭合规则约束；
当外部证书不足、上下文不闭合、引用不可追踪或谓词链不完整时，运行时必须降级为
拒绝、补证、细化或显式失败。于是，“智能生成”被降级为“在白盒合同下执行表面渲染”，
这使本文能够在使用自然语言表达能力的同时，保留形式系统的主权边界。

本文的第五层立场是三层分离：规范语义塔、逻辑占位塔与表面文本层必须保持不同角色。
规范语义塔负责语义定义、证明链、正规代表与证书闭合；逻辑占位塔负责量词、谓词、
引用、失败与风格参数的结构化承载；表面文本层只负责可读输出。任何把表面文本直接
当作语义源头的做法，都会破坏这一分离。本文因此把“自然语言逻辑性”重写为占位层
与审计层的性质，而不是重写为语言模型输出的自发可靠性。

本文的第六层立场是上下文修复塔与问答内核的对象化。问答系统的内核不是一段回答，
也不是一次 prompt 结果，而是上下文基底上的同伦修复塔路径类。问题、上下文、证据、
引用、证书、障碍与失败槽共同构成可修复对象；回答只是修复路径到表面文本层的投影。
这一路径继续承接扩展签发协议、可修复障碍与结构表示模板
\cite{RepoTex35ExtendedCertificationTower,RepoTex29RepairableObstruction}，并为后续开放系统博弈
与 Jacobi--Bianchi 修复稿件提供语义接口 \cite{RepoTex38JacobiBianchiRepair}。

本文的第七层立场是领域细化插件与白盒主权合同。领域知识不应绕过主基底直接支配
输出，而应作为扇区插件、认证纤维、局部联络与签发包进入系统。主基底保留规则、
证书、失败闭合与审计权；插件提供领域约束、局部模型、数据接口与专用证书。由此，
本文把领域扩展写成保守扩张，而不是让外部工具、经验语料或局部专家系统接管语义
主权。

本文的第八层立场是隐式解空间与持久化原则。系统不必每次从零生成答案；它可以在
上下文条件化命中之后，把可审计解空间、路径支撑类、占位结构与证书包持久化。后续
调用不是简单复读文本，而是在已持久化的隐式解空间中重新投影、补证、细化或降级。
这样，知识积累不被降为语料记忆，而被提升为可审计语义结构的持久化。

本文的第九层立场是为后续 \GFramework{} 主链提供语义底座。后续关于开放系统博弈、
Jacobi--Bianchi 障碍修复、统计晶格同伦隧穿、高阶正交确定性收敛与 PDE 求解范式
替代的稿件，都需要一个更早的语义事实：对象、路径、证书与命中集合必须先能在
白盒语义规范场中被命名、占位、解析、审计和降级。本文提供的正则语义规范场、
自然语言逻辑占位、LLM 运行时降级、领域插件与签发收束协议，正是这条主链的语义
运行底座。

因此，本文在整体 \GFramework{} 中承担的是“语义主权与运行时降级锚定”作用。它
向前连接语义函子场、语义规范化接口、逻辑占位与证书治理，向中间连接同伦扭曲双射、
上下文修复塔、白盒审计合同与领域细化插件，向后连接开放系统博弈、离散障碍修复、
微观隧穿命中与 PDE 残差求解范式。本文所说的自然语言不是语义本体，而是规范语义
的可读壳；本文所说的智能运行时不是裁决主体，而是受审计合同约束的表达组件。

概括地说，本文把 \GFramework{} 中“如何让形式语义进入自然语言运行时而不丧失
主权”这一问题，从提示词式经验推进为正则语义规范场、逻辑占位塔、同伦扭曲双射、
审计运行时、领域插件与签发收束协议共同构成的形式系统。它提供的不是单一表达
策略，而是一条由语义源头、占位骨架、运行时降级、白盒审计、插件扩张与持久化
解空间组成的方法论主干。

\setcounter{tocdepth}{3}
\setcounter{secnumdepth}{3}
\tableofcontents
\clearpage

%%%%%%%%%%%%%%%%%%%%%%%%%%%%%%%%%%%%%%%%%%%%%%%%%%%%%%%%%%%%
\section{形式逻辑：G 框架正则语义规范场经高阶同伦扭曲双射变换为自然语言语义逻辑占位规范场，并由 LLM 作为可审计确定性运行时表达}
\label{sec:tex37-ch1}
%%%%%%%%%%%%%%%%%%%%%%%%%%%%%%%%%%%%%%%%%%%%%%%%%%%%%%%%%%%%

\begin{remark}[核心陈述（严格化后）]
将 \GFramework{} 的语义及符号体系视为一个\textbf{正则语义规范场}，并把自然语言的可审计语义骨架视为一个
\textbf{自然语言语义逻辑占位规范场}。若二者之间存在逐阶兼容的高阶同伦扭曲双射，
则所有被 \GFramework{} 规范化成功的表达都可先无损搬运到自然语言逻辑占位，再经受控运行时输出为自然语言。
在此口径下，LLM 不再承担语义定义权或证明义务，而只承担“在既定占位与审计合同下生成自然语言表面表达”的运行时职责。
\end{remark}

本章将如下自然语言链条严格化：
\[
  \begin{aligned}
    &\text{G 框架正则语义规范场}
    \Longrightarrow
    \text{高阶同伦扭曲双射}\\
    \Longrightarrow\ &
    \text{自然语言语义逻辑占位规范场}
    \Longrightarrow
    \text{LLM 可审计确定性运行时}.
  \end{aligned}
\]
其目标不是把自然语言提升为语义本体，而是把自然语言降为
\[
  \text{占位的可读投影}+\text{审计日志的可视外壳}。
\]

\begin{remark}[阅读导航：仓内文稿与外部参照]
本章的接口风格直接参考仓内关于语义函子/语义规范场与“判定+规范化”接口的整理稿 \cite{RepoTex20SemanticGaugeNormalization}，
关于“LLM/NN 只提案，确定性审计器裁决”的形式化口径，参考 \cite{RepoTex21AuditedWrapper,RepoTex23SemanticBase}；
关于同伦扭曲正规形与双射集合写法，参考 \cite{RepoTex36HomotopyTwistBijection}；
关于 \GFramework{}、语义规范场、逻辑占位与自然语言规范子空间的总背景，参考
\cite{GaoZhengAppVol2,GaoZhengAppVol3,GaoZheng2025GFramework}。

外部标准参照方面，本章默认借用范畴/函子语言、同伦理论与同调证书写法，可参考
\cite{MacLane1998Categories,Hatcher2002AT,Weibel1994HomologicalAlgebra,KobayashiNomizu1963Foundations}。
\end{remark}

\subsection{对象域、正则语义规范场与自然语言逻辑占位规范场}
\label{subsec:tex37-objects}

\begin{definition}[G 框架表达域、语义解释与语义等价]
令
\[
  \mathsf{Expr}_G
\]
表示 \GFramework{} 允许的表达域，它可同时容纳数学符号、形式证明片段、规则化文本与受控说明文字。
给定语义解释映射
\[
  \Sem_G:\mathsf{Expr}_G\to\mathsf{Obj}_G,
\]
定义语义等价关系
\[
  e_1\sim_G e_2
  \quad:\Longleftrightarrow\quad
  \Sem_G(e_1)=\Sem_G(e_2).
\]
\end{definition}

\begin{definition}[正则语义规范场与总规范化映射]
称一族映射
\[
  \mathcal{F}^{\mathrm{reg}}_G=\{F_i\}_{i\in I},
  \qquad
  F_i:\mathsf{Expr}_G\to \mathsf{NF}_{G,i}\sqcup\{\bot\},
\]
为 \GFramework{} 的一个\textbf{正则语义规范场}，若存在总规范化映射
\[
  \Norm_G:\mathsf{Expr}_G\to \mathsf{NF}_G\sqcup\{\bot\},
\]
满足：
\begin{enumerate}
  \item \textbf{语义保持：}若 $\Norm_G(e)=n\neq\bot$，则 $e\sim_G n$；
  \item \textbf{规范唯一：}若 $\Norm_G(e_1)=\Norm_G(e_2)\neq\bot$，则 $e_1\sim_G e_2$；
  \item \textbf{失败显式：}若无法在既定规则库与证书库内完成规范化，则返回 $\bot$。
\end{enumerate}
\end{definition}

\begin{definition}[自然语言语义逻辑占位空间]
\label{def:tex37-placeholder-space}
定义
\[
  \mathsf{PL}_{\mathrm{NL}}
\]
为自然语言语义逻辑占位的集合。每个
\[
  p\in \mathsf{PL}_{\mathrm{NL}}
\]
至少包含如下槽位：
\begin{enumerate}
  \item 量词骨架与自由/束缚变量约束；
  \item 谓词依赖链、因果链与引用接口；
  \item 证书槽、证明槽与失败槽；
  \item 风格参数与严格层级参数的可投影接口。
\end{enumerate}
它不是自然语言句子本身，而是自然语言可复原的语义逻辑占位。
\end{definition}

\begin{definition}[自然语言语义逻辑占位规范场]
称一族映射
\[
  \mathcal{F}^{\mathrm{pl}}_{\mathrm{NL}}=\{H_k\}_{k\in K},
  \qquad
  H_k:\mathsf{PL}_{\mathrm{NL}}\to \mathsf{NF}_{\mathrm{PL},k}\sqcup\{\bot\},
\]
为\textbf{自然语言语义逻辑占位规范场}，若其共同作用把每个占位
\[
  p\in\mathsf{PL}_{\mathrm{NL}}
\]
约束为一个唯一可审计的逻辑骨架，并且不同的风格输出只允许改变表面表达，
不允许改变量词、谓词依赖、证书引用与失败口径。
\end{definition}

\begin{remark}[两类规范场的分工]
\GFramework{} 侧负责“语义定义、证明链与规范代表”；自然语言逻辑占位侧负责“把该规范代表压成自然语言可表达但不失真的逻辑骨架”。
前者是语义源头，后者是语义投影面；两者都必须保持 fail-closed。
\end{remark}

\subsection{高阶同伦扭曲塔与双射的定义}
\label{subsec:tex37-twist}

\begin{definition}[高阶同伦扭曲扩张塔]
\label{def:tex37-twist-tower}
定义两列递归对象：
\[
  \mathsf{NF}_G^{(0)}:=\mathsf{NF}_G,
  \qquad
  \mathsf{PL}_{\mathrm{NL}}^{(0)}:=\mathsf{PL}_{\mathrm{NL}}.
\]
对每个 $n\in\NN$，给定纤维化扭曲数据族
\[
  \Tw_G^{(n)}(x)\quad (x\in \mathsf{NF}_G^{(n)}),
  \qquad
  \Tw_{\mathrm{NL}}^{(n)}(p)\quad (p\in \mathsf{PL}_{\mathrm{NL}}^{(n)}),
\]
并递归定义
\[
  \mathsf{NF}_G^{(n+1)}
  :=
  \bigsqcup_{x\in \mathsf{NF}_G^{(n)}}\{x\}\times \Tw_G^{(n)}(x),
\]
\[
  \mathsf{PL}_{\mathrm{NL}}^{(n+1)}
  :=
  \bigsqcup_{p\in \mathsf{PL}_{\mathrm{NL}}^{(n)}}\{p\}\times \Tw_{\mathrm{NL}}^{(n)}(p).
\]
称其为 G 侧与自然语言占位侧的\textbf{高阶同伦扭曲扩张塔}。
\end{definition}

\begin{definition}[高阶同伦扭曲双射]
\label{def:tex37-high-twist-bijection}
一族映射
\[
  \Theta=\{\Theta_n\}_{n\in\NN},
  \qquad
  \Theta_n:\mathsf{NF}_G^{(n)}\to \mathsf{PL}_{\mathrm{NL}}^{(n)},
\]
称为一个\textbf{高阶同伦扭曲双射}，若满足：
\begin{enumerate}
  \item 每个 $\Theta_n$ 都是双射；
  \item 对每个 $n$、每个 $x\in\mathsf{NF}_G^{(n)}$，存在纤维双射
  \[
    \Psi_x^{(n)}:\Tw_G^{(n)}(x)\to \Tw_{\mathrm{NL}}^{(n)}(\Theta_n(x));
  \]
  \item 对任意 $(x,\eta)\in \mathsf{NF}_G^{(n+1)}$，都有递归兼容式
  \[
    \Theta_{n+1}(x,\eta)=\bigl(\Theta_n(x),\Psi_x^{(n)}(\eta)\bigr).
  \]
\end{enumerate}
\end{definition}

\begin{remark}[“高阶同伦扭曲”在本章中的含义]
这里的“扭曲”不是任意修辞上的变形，而是逐阶附着在规范代表之上的附加同伦数据、证明依赖、证书依赖与展开层级。
因此 $\Theta$ 的作用不是自由翻译，而是\textbf{受纤维约束的逐层搬运}。
\end{remark}

\subsection{公理（降级为定理）：高阶同伦扭曲双射可按层递归构造}
\label{subsec:tex37-axiom-downgrade}

\begin{theorem}[公理降级：基阶双射与纤维双射存在 $\Rightarrow$ 全塔高阶同伦扭曲双射存在]
\label{thm:tex37-axiom-downgrade}
设满足下列条件：
\begin{enumerate}
  \item 存在基阶双射
  \[
    \Theta_0:\mathsf{NF}_G^{(0)}\to \mathsf{PL}_{\mathrm{NL}}^{(0)};
  \]
  \item 对任意 $n\in\NN$ 与任意 $x\in\mathsf{NF}_G^{(n)}$，存在纤维双射
  \[
    \Psi_x^{(n)}:\Tw_G^{(n)}(x)\to \Tw_{\mathrm{NL}}^{(n)}(\Theta_n(x)).
  \]
\end{enumerate}
则存在唯一一族映射
\[
  \Theta=\{\Theta_n\}_{n\in\NN}
\]
满足定义~\ref{def:tex37-high-twist-bijection} 中高阶同伦扭曲双射的全部条件。
\end{theorem}

\begin{Certificate}[证书：基步 + 归纳步（数学归纳法证书）]
\label{cert:tex37-axiom-downgrade}
定义谓词
\[
  P(n):\ \text{“存在唯一双射 }\Theta_n:\mathsf{NF}_G^{(n)}\to \mathsf{PL}_{\mathrm{NL}}^{(n)}
  \text{，且与所有已定义的纤维双射兼容”}.
\]
证书登记为
\[
  \pi_{\mathrm{ind}}
  =
  \bigl(P(0),\ \forall n\in\NN,\ P(n)\Rightarrow P(n+1)\bigr).
\]
其中
\[
  P(n)\Rightarrow P(n+1)
\]
的具体递推式为
\[
  \Theta_{n+1}(x,\eta)
  :=
  \bigl(\Theta_n(x),\Psi_x^{(n)}(\eta)\bigr).
\]
\end{Certificate}

\begin{proof}[证明（数学归纳法证书；谓词因果链）]
\textbf{基步 $P(0)$：}
由条件(1)，$\Theta_0$ 已被假定为双射，故 $P(0)$ 成立。

\textbf{归纳步：}
设 $P(n)$ 成立，即已存在唯一双射
\[
  \Theta_n:\mathsf{NF}_G^{(n)}\to \mathsf{PL}_{\mathrm{NL}}^{(n)}.
\]
对任意 $(x,\eta)\in\mathsf{NF}_G^{(n+1)}$，定义
\[
  \Theta_{n+1}(x,\eta):=\bigl(\Theta_n(x),\Psi_x^{(n)}(\eta)\bigr).
\]
现证其为双射。

\textbf{单射性：}
若
\[
  \Theta_{n+1}(x_1,\eta_1)=\Theta_{n+1}(x_2,\eta_2),
\]
则有
\[
  \Theta_n(x_1)=\Theta_n(x_2),
  \qquad
  \Psi_{x_1}^{(n)}(\eta_1)=\Psi_{x_2}^{(n)}(\eta_2).
\]
由 $\Theta_n$ 的单射性得 $x_1=x_2$；再由对应纤维双射 $\Psi_{x_1}^{(n)}$ 的单射性得 $\eta_1=\eta_2$。
故 $(x_1,\eta_1)=(x_2,\eta_2)$。

\textbf{满射性：}
任取
\[
  (p,\zeta)\in \mathsf{PL}_{\mathrm{NL}}^{(n+1)}.
\]
由 $\Theta_n$ 的满射性，存在唯一 $x\in \mathsf{NF}_G^{(n)}$ 使得 $\Theta_n(x)=p$。
再由纤维双射
\[
  \Psi_x^{(n)}:\Tw_G^{(n)}(x)\to \Tw_{\mathrm{NL}}^{(n)}(p)
\]
的满射性，存在唯一 $\eta\in \Tw_G^{(n)}(x)$ 使得 $\Psi_x^{(n)}(\eta)=\zeta$。
于是
\[
  \Theta_{n+1}(x,\eta)=(p,\zeta),
\]
故 $\Theta_{n+1}$ 满射。

\textbf{唯一性：}
一旦要求
\[
  \Theta_{n+1}(x,\eta)=\bigl(\Theta_n(x),\Psi_x^{(n)}(\eta)\bigr),
\]
其取值便被递推式唯一确定，故无第二种兼容定义。

综上，$P(n)\Rightarrow P(n+1)$ 成立。由数学归纳法，$\forall n\in\NN,\ P(n)$。
故全塔高阶同伦扭曲双射存在且唯一。证毕。
\end{proof}

\subsection{引理：占位搬运保持 G 框架规范代表的唯一性}
\label{subsec:tex37-lemma}

\begin{lemma}[自然语言逻辑占位不丢失 G 框架规范代表]
\label{lem:tex37-placeholder-uniqueness}
定义基阶占位搬运映射
\[
  \Pi_G:=\Theta_0\circ \Norm_G:\mathsf{Expr}_G\to \mathsf{PL}_{\mathrm{NL}}^{(0)}\sqcup\{\bot\}.
\]
若对两个表达 $e_1,e_2\in\mathsf{Expr}_G$ 有
\[
  \Pi_G(e_1)=\Pi_G(e_2)\neq\bot,
\]
则
\[
  \Norm_G(e_1)=\Norm_G(e_2)\neq\bot
  \qquad\text{且}\qquad
  e_1\sim_G e_2.
\]
\end{lemma}

\begin{Certificate}[证书：占位相同 $\Rightarrow$ 规范形相同 $\Rightarrow$ 语义等价]
\label{cert:tex37-placeholder-uniqueness}
证书登记谓词链：
\[
\begin{aligned}
  L_1 &: \Pi_G(e_1)=\Pi_G(e_2)=p\neq\bot,\\
  L_2 &: p=\Theta_0(\Norm_G(e_1))=\Theta_0(\Norm_G(e_2)),\\
  L_3 &: \Theta_0\ \text{为双射，故特别是单射},\\
  L_4 &: L_2\wedge L_3\Rightarrow \Norm_G(e_1)=\Norm_G(e_2),\\
  L_5 &: \Norm_G\ \text{的语义保持性}\Rightarrow e_1\sim_G \Norm_G(e_1),\ e_2\sim_G \Norm_G(e_2),\\
  L_6 &: L_4\wedge L_5\Rightarrow e_1\sim_G e_2.
\end{aligned}
\]
\end{Certificate}

\begin{proof}[证明（谓词因果链）]
由 $L_1$，两式的占位搬运值同为非失败元 $p$。
展开 $\Pi_G$ 的定义得 $L_2$。
由定理~\ref{thm:tex37-axiom-downgrade} 中的基阶双射条件，$\Theta_0$ 为单射，即 $L_3$。
因此由 $L_2$ 与 $L_3$ 得
\[
  \Norm_G(e_1)=\Norm_G(e_2)\neq\bot,
\]
即 $L_4$。
再由正则语义规范场的语义保持性，得到 $L_5$。
把 $L_4$ 代回 $L_5$，可得
\[
  e_1\sim_G \Norm_G(e_1)=\Norm_G(e_2)\sim_G e_2,
\]
于是 $e_1\sim_G e_2$，即 $L_6$。证毕。
\end{proof}

\subsection{定理：LLM 被降级为自然语言表达的可审计确定性运行时}
\label{subsec:tex37-runtime}

\begin{definition}[风格模式、严格层级与自然语言候选池]
令
\[
  \Sigma_{\mathrm{sty}}
\]
表示风格参数集合（如简洁、教材、法律、论文、说明书等），
令
\[
  \Lambda
\]
表示严格层级集合（如口语说明、技术说明、形式推演、逐步证明等）。
记
\[
  \Xi:=\Sigma_{\mathrm{sty}}\times\Lambda
\]
为输出模式集合。对每个 $\xi\in\Xi$，记自然语言候选集合为 $\mathsf{Txt}_\xi$。
\end{definition}

\begin{definition}[解析器、审计器与候选采集协议]
对每个 $\xi\in\Xi$，定义：
\begin{enumerate}
  \item 确定性解析器
  \[
    \Parse_\xi:\mathsf{Txt}_\xi\to \mathsf{PL}_{\mathrm{NL}}^{(0)}\sqcup\{\bot\};
  \]
  \item 风格/层级审计谓词
  \[
    \Style_\xi:\mathsf{Txt}_\xi\to\{\top,\bot\};
  \]
  \item 固定模型版本、固定提示模板、固定解码策略与固定随机种子的候选采集协议
  \[
    \mathcal{C}^{\mathrm{LLM}}_\xi:
    \mathsf{PL}_{\mathrm{NL}}^{(0)}\to
    2^{\mathsf{Txt}_\xi}_{\mathrm{fin}}.
  \]
\end{enumerate}
称
\[
  \mathcal{C}^{\mathrm{LLM}}_\xi(p)
\]
为模式 $\xi$ 下由 LLM 运行时给出的候选集。
\end{definition}

\begin{definition}[可审计接受集与确定性运行时]
固定每个 $\xi\in\Xi$ 上的全序
\[
  \prec_\xi.
\]
对任意 $p\in \mathsf{PL}_{\mathrm{NL}}^{(0)}$，定义可审计接受集
\[
  \Audit_\xi(p)
  :=
  \bigl\{
    t\in \mathcal{C}^{\mathrm{LLM}}_\xi(p)
    \mid
    \Parse_\xi(t)=p\ \wedge\ \Style_\xi(t)=\top
  \bigr\}.
\]
定义确定性运行时
\[
  \Render^{\mathrm{LLM}}_\xi(p)
  :=
  \begin{cases}
    \min\nolimits_{\prec_\xi}\Audit_\xi(p), & \Audit_\xi(p)\neq\varnothing,\\
    \bot, & \Audit_\xi(p)=\varnothing.
  \end{cases}
\]
\end{definition}

\begin{theorem}[LLM 运行时降级定理]
\label{thm:tex37-runtime-deterministic}
对任意 $\xi\in\Xi$ 与任意占位 $p\in\mathsf{PL}_{\mathrm{NL}}^{(0)}$，若
\[
  t=\Render^{\mathrm{LLM}}_\xi(p)\neq\bot,
\]
则有
\[
  \Parse_\xi(t)=p
  \qquad\text{且}\qquad
  \Style_\xi(t)=\top.
\]
并且在固定候选采集协议 $\mathcal{C}^{\mathrm{LLM}}_\xi$ 与固定全序 $\prec_\xi$ 下，$t$ 唯一确定。
因此，LLM 对外只表现为
\[
  \mathsf{PL}_{\mathrm{NL}}^{(0)}\longrightarrow \mathsf{Txt}_\xi\sqcup\{\bot\}
\]
的\textbf{可审计确定性运行时}，而不再是语义定义者。
\end{theorem}

\begin{Certificate}[证书：候选采集固定 + 解析回投影 + 全序最小接受元]
\label{cert:tex37-runtime-deterministic}
证书登记谓词链：
\[
\begin{aligned}
  R_1 &: \mathcal{C}^{\mathrm{LLM}}_\xi\ \text{由固定模型版本、固定提示模板、固定策略、固定种子定义},\\
  R_2 &: t=\Render^{\mathrm{LLM}}_\xi(p)\neq\bot,\\
  R_3 &: R_2\Rightarrow t\in \Audit_\xi(p),\\
  R_4 &: t\in \Audit_\xi(p)\Rightarrow \Parse_\xi(t)=p\ \wedge\ \Style_\xi(t)=\top,\\
  R_5 &: \prec_\xi\ \text{为全序}\Rightarrow \min\nolimits_{\prec_\xi}\Audit_\xi(p)\ \text{唯一}.
\end{aligned}
\]
\end{Certificate}

\begin{proof}[证明（谓词因果链）]
由 $R_1$，候选集
\[
  \mathcal{C}^{\mathrm{LLM}}_\xi(p)
\]
对每个给定的 $p$ 与 $\xi$ 都是确定的有限集合。
由 $R_2$ 与 $\Render^{\mathrm{LLM}}_\xi$ 的定义知，只有在 $\Audit_\xi(p)\neq\varnothing$ 时才有非失败输出，
且此时
\[
  t=\min\nolimits_{\prec_\xi}\Audit_\xi(p),
\]
故 $R_3$ 成立。
再由 $\Audit_\xi(p)$ 的定义，$R_3$ 立即推出
\[
  \Parse_\xi(t)=p
  \qquad\text{且}\qquad
  \Style_\xi(t)=\top,
\]
即 $R_4$。
最后，由 $R_5$，在全序上最小元一旦存在便唯一，因此 $t$ 唯一确定。

于是，外部系统真正依赖的是：
\[
  \text{固定候选采集协议}
  +\text{解析回投影}
  +\text{风格审计}
  +\text{确定性最小接受元选择},
\]
而不是 LLM 的自由叙述能力本身。故 LLM 被降级为自然语言表达运行时。证毕。
\end{proof}

\subsection{命题：不同风格或严格层级的表达属于同一逻辑占位类}
\label{subsec:tex37-style}

\begin{definition}[自然语言占位等价]
对不同模式 $\xi_1,\xi_2\in\Xi$ 下的两个自然语言输出 $t_1\in\mathsf{Txt}_{\xi_1}$、$t_2\in\mathsf{Txt}_{\xi_2}$，
定义
\[
  t_1\sim_{\mathrm{pl}} t_2
  \quad:\Longleftrightarrow\quad
  \Parse_{\xi_1}(t_1)=\Parse_{\xi_2}(t_2)\neq\bot.
\]
称其为\textbf{自然语言占位等价}。
\end{definition}

\begin{proposition}[同一占位可生成多风格、多严格层级表达而不改语义骨架]
\label{prop:tex37-style-invariance}
任取 $p\in \mathsf{PL}_{\mathrm{NL}}^{(0)}$ 与两个模式
\[
  \xi_1=(\sigma_1,\lambda_1),\qquad \xi_2=(\sigma_2,\lambda_2)\in\Xi.
\]
若
\[
  t_1=\Render^{\mathrm{LLM}}_{\xi_1}(p)\neq\bot,
  \qquad
  t_2=\Render^{\mathrm{LLM}}_{\xi_2}(p)\neq\bot,
\]
则
\[
  \Parse_{\xi_1}(t_1)=\Parse_{\xi_2}(t_2)=p,
\]
从而
\[
  t_1\sim_{\mathrm{pl}} t_2.
\]
因此，不同风格或严格层级只能改变表面措辞与展开粒度，不能改变由占位确定的逻辑骨架。
\end{proposition}

\begin{Certificate}[证书：同源占位 + 解析回投影 $\Rightarrow$ 占位等价]
\label{cert:tex37-style-invariance}
证书登记谓词链：
\[
\begin{aligned}
  S_1 &: t_1=\Render^{\mathrm{LLM}}_{\xi_1}(p)\neq\bot,\\
  S_2 &: t_2=\Render^{\mathrm{LLM}}_{\xi_2}(p)\neq\bot,\\
  S_3 &: S_1\Rightarrow \Parse_{\xi_1}(t_1)=p\ \text{（定理~\ref{thm:tex37-runtime-deterministic}）},\\
  S_4 &: S_2\Rightarrow \Parse_{\xi_2}(t_2)=p\ \text{（定理~\ref{thm:tex37-runtime-deterministic}）},\\
  S_5 &: S_3\wedge S_4\Rightarrow t_1\sim_{\mathrm{pl}} t_2.
\end{aligned}
\]
\end{Certificate}

\begin{proof}[证明（谓词因果链）]
由 $S_1$ 与 $S_2$，两个输出都不是失败符号。
分别对 $\xi_1$ 与 $\xi_2$ 应用定理~\ref{thm:tex37-runtime-deterministic}，得到
\[
  \Parse_{\xi_1}(t_1)=p,
  \qquad
  \Parse_{\xi_2}(t_2)=p.
\]
于是两者的解析结果相同且非 $\bot$，根据定义便有
\[
  t_1\sim_{\mathrm{pl}} t_2.
\]
故命题成立。证毕。
\end{proof}

\subsection{推论：基于 G 框架的逻辑引擎由此建立}
\label{subsec:tex37-engine}

\begin{corollary}[G 框架逻辑引擎的可审计闭环]
\label{cor:tex37-logic-engine}
对任意模式 $\xi\in\Xi$，定义映射
\[
  \Engine_{G,\xi}(e)
  :=
  \begin{cases}
    \Render^{\mathrm{LLM}}_\xi\bigl(\Theta_0(\Norm_G(e))\bigr), & \Norm_G(e)\neq\bot,\\
    \bot, & \Norm_G(e)=\bot.
  \end{cases}
\]
则对任意 $e\in\mathsf{Expr}_G$，若
\[
  \Engine_{G,\xi}(e)=t\neq\bot,
\]
则有
\[
  \Parse_\xi(t)=\Theta_0(\Norm_G(e)),
\]
并进一步有
\[
  \Theta_0^{-1}\bigl(\Parse_\xi(t)\bigr)=\Norm_G(e).
\]
故
\[
  e
  \xrightarrow{\Norm_G}
  \Norm_G(e)
  \xrightarrow{\Theta_0}
  \Theta_0(\Norm_G(e))
  \xrightarrow{\Render^{\mathrm{LLM}}_\xi}
  t
\]
构成一个可审计、可回投影、可 fail-closed 的逻辑引擎闭环。
\end{corollary}

\begin{Certificate}[证书：规范化成功 + 占位唯一 + 运行时回投影]
\label{cert:tex37-logic-engine}
证书登记谓词链：
\[
\begin{aligned}
  E_1 &: \Norm_G(e)=n\neq\bot,\\
  E_2 &: p:=\Theta_0(n)\in \mathsf{PL}_{\mathrm{NL}}^{(0)},\\
  E_3 &: t=\Render^{\mathrm{LLM}}_\xi(p)\neq\bot,\\
  E_4 &: E_3\Rightarrow \Parse_\xi(t)=p\ \text{（定理~\ref{thm:tex37-runtime-deterministic}）},\\
  E_5 &: E_2\wedge E_4\Rightarrow \Theta_0^{-1}(\Parse_\xi(t))=n,\\
  E_6 &: E_1\wedge E_5\Rightarrow \Theta_0^{-1}(\Parse_\xi(t))=\Norm_G(e).
\end{aligned}
\]
\end{Certificate}

\begin{proof}[证明（谓词因果链）]
若 $\Norm_G(e)=\bot$，则按定义 $\Engine_{G,\xi}(e)=\bot$，推论仅讨论非失败情形。
故设 $E_1$ 成立，记
\[
  n:=\Norm_G(e),\qquad p:=\Theta_0(n).
\]
由 $\Theta_0$ 的值域定义，得 $E_2$。
再设
\[
  t=\Render^{\mathrm{LLM}}_\xi(p)\neq\bot,
\]
即 $E_3$。由定理~\ref{thm:tex37-runtime-deterministic}，得
\[
  \Parse_\xi(t)=p,
\]
即 $E_4$。
由 $E_2$ 与 $E_4$，
\[
  \Theta_0^{-1}\bigl(\Parse_\xi(t)\bigr)=\Theta_0^{-1}(p)=n,
\]
即 $E_5$。
最后将 $n=\Norm_G(e)$ 代回，得到
\[
  \Theta_0^{-1}\bigl(\Parse_\xi(t)\bigr)=\Norm_G(e),
\]
即 $E_6$。
因此从 G 框架表达出发，可经规范形、占位与自然语言运行时得到一个全程可回投影的闭环逻辑引擎。证毕。
\end{proof}

\subsection{备注：边界、解释与工程含义}
\label{subsec:tex37-remarks}

\begin{remark}[LLM 的权力边界]
本章并未把 LLM 视为“证明机器”或“语义裁判”。
真正承担语义定义权的是
\[
  \mathcal{F}^{\mathrm{reg}}_G,\ \Norm_G,\ \Theta,\ \Parse_\xi,\ \Style_\xi,
\]
而 LLM 只在固定协议下承担候选采集与自然语言表面展开职责。
这与仓内关于“LLM 只提案，证书/审计器裁决”的结论完全一致 \cite{RepoTex21AuditedWrapper,RepoTex23SemanticBase,GaoZhengAppVol2}。
\end{remark}

\begin{remark}[多风格与多严格层级的真正自由度]
风格变化、说明粒度变化、数学符号占比变化、解释性文字增减，都只能发生在
\[
  \Render^{\mathrm{LLM}}_\xi
\]
之后的表面表达层。
只要解析器能把输出回投影回同一个占位
\[
  p\in\mathsf{PL}_{\mathrm{NL}}^{(0)},
\]
这些变化就是合法风格自由度；若无法回投影，则必须 fail-closed，而不能把自由叙述冒充为语义正确。
\end{remark}

\begin{remark}[逻辑引擎的工程读法]
工程上，这个逻辑引擎可被读作：
\[
  \begin{aligned}
    &\text{G 框架规则/符号体系}
    \to
    \text{规范语义}
    \to
    \text{逻辑占位}\\
    \to\ &
    \text{不同模式的自然语言输出}
    \to
    \text{解析回投影与审计}.
  \end{aligned}
\]
因此系统可以同时支持
\[
  \text{严格数学稿、技术说明稿、教学展开稿、审计报告稿}
\]
等多种外观，但它们都必须共享同一规范语义内核。
\end{remark}

\begin{remark}[失败显式与扩张原则]
若某个 $e\in\mathsf{Expr}_G$ 无法被 $\Norm_G$ 规范化，或某个占位 $p$ 在某种模式 $\xi$ 下无可审计接受元，
系统必须返回 $\bot$。这意味着缺口应被记账为：
\[
  \text{规则库不足、占位槽不足、解析器不足或风格合同不足},
\]
而不是被交给 LLM 以自由生成方式“补脑”。这正是逻辑引擎得以可审计的根本边界。
\end{remark}

\clearpage

%%%%%%%%%%%%%%%%%%%%%%%%%%%%%%%%%%%%%%%%%%%%%%%%%%%%%%%%%%%%
\section{形式逻辑：第一章的增益评价、工程桥接地位与逻辑引擎化意义}
\label{sec:tex37-ch2}
%%%%%%%%%%%%%%%%%%%%%%%%%%%%%%%%%%%%%%%%%%%%%%%%%%%%%%%%%%%%

\begin{remark}[核心结论（严格化后）]
就本稿当前结构而言，第一章的总体增益应判为\textbf{高}，但其强项不是新增一个更大的本体宇宙，
而是把
\[
  \text{规范语义}
  \to
  \text{自然语言占位}
  \to
  \text{可审计运行时}
  \to
  \text{解析回投影}
\]
这条链条写成了严格闭环。故第一章的主增益满足
\[
  \boxed{
    \text{接口闭环增益}
    >
    \text{本体扩张增益}
  }.
\]
\end{remark}

本章不再构造新的运行时接口，而是对第~\ref{sec:tex37-ch1} 章已经证明的结构进行\textbf{增益评价的定理化整理}，
即把“第一章究竟增益在哪里”压成定义、定理、引理、命题、推论与备注。

\subsection{定义：第一章增益函数}
\label{subsec:tex37-gain-definition}

\begin{definition}[第一章综合增益函数]
\label{def:tex37-gain-function}
定义第一章的综合增益为
\[
  \mathrm{Gain}_{\mathrm{Ch1}}
  :=
  \frac{
    G_{\mathrm{closure}}
    +
    G_{\mathrm{runtime}}
    +
    G_{\mathrm{audit}}
    +
    G_{\mathrm{style}}
  }{
    C_{\mathrm{unproved}}
  },
\]
其中约定 $C_{\mathrm{unproved}}>0$，并解释如下：
\begin{enumerate}
  \item $G_{\mathrm{closure}}$：理论闭环增益，即从 $\Norm_G$、$\Theta_0$ 到 $\Parse_\xi$ 的结构闭合度；
  \item $G_{\mathrm{runtime}}$：运行时降级增益，即把 LLM 从语义裁判降为表达运行时；
  \item $G_{\mathrm{audit}}$：可审计回投影增益，即输出可经 $\Parse_\xi$ 与 $\Theta_0^{-1}$ 回到规范代表；
  \item $G_{\mathrm{style}}$：多风格不变性增益，即同一占位可被多模式输出而不改语义骨架；
  \item $C_{\mathrm{unproved}}$：尚未完全构造的接口成本，如 $\Theta_0$ 的完全可计算实现、$\Psi_x^{(n)}$ 的统一构造、
  $\Parse_\xi$ 的开放域完备性与 $\Style_\xi$ 的全边界形式化等。
\end{enumerate}
\end{definition}

\begin{remark}[评价口径]
在上述定义中，若分子主要来自“接口闭环、运行时、审计与风格不变性”，
则说明第一章的增益重心偏向工程桥接；若主要来自新宇宙、新基本算子或新修复塔，
则说明增益重心偏向本体扩张。本章将证明前者占主导。
\end{remark}

\subsection{公理（降级为定理）：接口闭环增益可按已认证接口数递归分解}
\label{subsec:tex37-gain-axiom}

\begin{theorem}[公理降级：第一章增益分子可按接口链长递归累加]
\label{thm:tex37-gain-decomposition}
定义接口增益增量
\[
  \Delta_1:=G_{\mathrm{closure}},\qquad
  \Delta_2:=G_{\mathrm{runtime}},\qquad
  \Delta_3:=G_{\mathrm{audit}},\qquad
  \Delta_4:=G_{\mathrm{style}}.
\]
再定义部分和
\[
  N(1):=\Delta_1,
  \qquad
  N(m+1):=N(m)+\Delta_{m+1}
  \quad (1\le m\le 3).
\]
则对所有 $m\in\{1,2,3,4\}$，有
\[
  N(m)=\sum_{k=1}^{m}\Delta_k.
\]
特别地，
\[
  N(4)
  =
  G_{\mathrm{closure}}
  +
  G_{\mathrm{runtime}}
  +
  G_{\mathrm{audit}}
  +
  G_{\mathrm{style}}
\]
正是定义~\ref{def:tex37-gain-function} 中第一章综合增益的分子。
\end{theorem}

\begin{Certificate}[证书：基步 + 归纳步（数学归纳法证书）]
\label{cert:tex37-gain-decomposition}
定义谓词
\[
  P(m):\ N(m)=\sum_{k=1}^{m}\Delta_k.
\]
归纳证书登记为
\[
  \pi_{\mathrm{ind}}
  =
  \bigl(P(1),\ \forall m\in\{1,2,3\},\ P(m)\Rightarrow P(m+1)\bigr).
\]
其中
\[
  P(m)\Rightarrow P(m+1)
\]
只使用递推关系
\[
  N(m+1)=N(m)+\Delta_{m+1}.
\]
\end{Certificate}

\begin{proof}[证明（数学归纳法证书；谓词因果链）]
\textbf{基步 $P(1)$：}
由定义，
\[
  N(1)=\Delta_1=\sum_{k=1}^{1}\Delta_k.
\]
故 $P(1)$ 成立。

\textbf{归纳步：}
设 $P(m)$ 成立，即
\[
  N(m)=\sum_{k=1}^{m}\Delta_k.
\]
则由递推定义，
\[
  N(m+1)=N(m)+\Delta_{m+1}.
\]
将归纳假设代入得
\[
  N(m+1)
  =
  \sum_{k=1}^{m}\Delta_k+\Delta_{m+1}
  =
  \sum_{k=1}^{m+1}\Delta_k.
\]
故 $P(m)\Rightarrow P(m+1)$。

由数学归纳法，$P(m)$ 对 $m=1,2,3,4$ 全部成立。
于是 $N(4)$ 即为四类增益的总和，也就是第一章综合增益的分子。证毕。
\end{proof}

\subsection{定理：第一章的最大增益是把语义主权从 LLM 收回到 G 框架}
\label{subsec:tex37-gain-semantic-sovereignty}

\begin{theorem}[语义主权回收定理]
\label{thm:tex37-semantic-sovereignty}
对任意 $e\in\mathsf{Expr}_G$ 与任意模式 $\xi\in\Xi$，若
\[
  t=\Engine_{G,\xi}(e)\neq\bot,
\]
则该输出的语义决定权完全由
\[
  \Norm_G,\ \Theta_0,\ \Parse_\xi,\ \Style_\xi
\]
共同确定，而不由 LLM 的自然语言生成自由度决定。换言之，第一章的首要增益在于：
\[
  \boxed{
    \text{LLM 失去语义裁判权，退化为受审计的表达运行时}
  }.
\]
\end{theorem}

\begin{Certificate}[证书：规范化先行 + 占位搬运 + 回投影闭环]
\label{cert:tex37-semantic-sovereignty}
证书登记谓词链：
\[
\begin{aligned}
  T_1 &: e\xrightarrow{\Norm_G} n:=\Norm_G(e)\neq\bot,\\
  T_2 &: n\xrightarrow{\Theta_0} p:=\Theta_0(n),\\
  T_3 &: p\xrightarrow{\Render^{\mathrm{LLM}}_\xi} t\neq\bot,\\
  T_4 &: T_3\Rightarrow \Parse_\xi(t)=p\ \wedge\ \Style_\xi(t)=\top\ \text{（定理~\ref{thm:tex37-runtime-deterministic}）},
    \\
  T_5 &: T_1\wedge T_2\wedge T_4\Rightarrow \Theta_0^{-1}(\Parse_\xi(t))=\Norm_G(e)\ \text{（推论~\ref{cor:tex37-logic-engine}）},\\
  T_6 &: T_5\Rightarrow \text{语义裁判权属于上游规范系统而非 LLM}.
\end{aligned}
\]
\end{Certificate}

\begin{proof}[证明（谓词因果链）]
由 $T_1$，原始表达先被规约到规范代表 $n$；由 $T_2$，该规范代表再被搬运到占位 $p$。
由 $T_3$，LLM 运行时只是在给定占位 $p$ 的条件下输出表面文本 $t$。
由定理~\ref{thm:tex37-runtime-deterministic}，$T_3$ 推出
\[
  \Parse_\xi(t)=p
  \qquad\text{且}\qquad
  \Style_\xi(t)=\top,
\]
即 $T_4$。
再由推论~\ref{cor:tex37-logic-engine}，$T_1$、$T_2$ 与 $T_4$ 蕴含
\[
  \Theta_0^{-1}(\Parse_\xi(t))=\Norm_G(e),
\]
即 $T_5$。

这说明：
\[
  t\ \text{是否被接受、被解释为何占位、最终对应哪个规范代表}
\]
都由 $\Norm_G,\Theta_0,\Parse_\xi,\Style_\xi$ 决定；
LLM 只负责在固定协议下提供候选与表面表达，不再拥有语义裁判权。
故定理成立。证毕。
\end{proof}

\begin{remark}[评价]
这一增益的实质不是“让 LLM 更强”，而是把系统架构从
\[
  \text{模型中心}
  \to
  \text{规范语义中心}
\]
重写为可审计闭环。它是控制权迁移，而非措辞优化。
\end{remark}

\subsection{引理：第一章的对象层新增不以本体爆发为主，而以中介层分离为主}
\label{subsec:tex37-gain-middle-layer}

\begin{lemma}[自然语言逻辑占位空间是关键中介层对象]
\label{lem:tex37-middle-layer}
第一章新增并明确化的关键对象是
\[
  \mathsf{PL}_{\mathrm{NL}},
\]
它把结构
\[
  \text{形式规范代表}
  \to
  \text{自然语言输出}
\]
严格分解为
\[
  \text{形式规范代表}
  \xrightarrow{\Theta_0}
  \text{自然语言逻辑占位}
  \xrightarrow{\Render^{\mathrm{LLM}}_\xi}
  \text{多风格表面表达}.
\]
因此第一章的对象层增益主要体现为\textbf{中介层分离}，而非大规模新增本体宇宙。
\end{lemma}

\begin{Certificate}[证书：占位层插入 + 风格不变性 + 回投影]
\label{cert:tex37-middle-layer}
证书登记谓词链：
\[
\begin{aligned}
  L_1 &: \Theta_0:\mathsf{NF}_G\to \mathsf{PL}_{\mathrm{NL}}^{(0)}\ \text{在第一章被显式定义},\\
  L_2 &: \Render^{\mathrm{LLM}}_\xi:\mathsf{PL}_{\mathrm{NL}}^{(0)}\to \mathsf{Txt}_\xi\sqcup\{\bot\}\ \text{在第一章被显式定义},
    \\
  L_3 &: \Parse_\xi:\mathsf{Txt}_\xi\to \mathsf{PL}_{\mathrm{NL}}^{(0)}\sqcup\{\bot\}\ \text{在第一章被显式定义},\\
  L_4 &: t_1,t_2\ \text{若来自同一 }p,\ \text{则 }t_1\sim_{\mathrm{pl}} t_2\ \text{（命题~\ref{prop:tex37-style-invariance}）},\\
  L_5 &: L_1\wedge L_2\wedge L_3\wedge L_4\Rightarrow \text{语义核与表面文本被严格分层}.
\end{aligned}
\]
\end{Certificate}

\begin{proof}[证明（谓词因果链）]
由 $L_1$、$L_2$ 与 $L_3$，第一章已经把
\[
  \mathsf{NF}_G
  \to
  \mathsf{PL}_{\mathrm{NL}}^{(0)}
  \to
  \mathsf{Txt}_\xi
  \to
  \mathsf{PL}_{\mathrm{NL}}^{(0)}
\]
写成显式可组合的结构。
因此形式规范代表与自然语言输出之间不再是硬跳跃，而是经过占位层中介。

再由命题~\ref{prop:tex37-style-invariance}，对同一占位 $p$ 的不同模式输出，
都满足
\[
  \Parse_{\xi_1}(t_1)=\Parse_{\xi_2}(t_2)=p.
\]
故风格变化、解释粒度变化与术语变换不再直接污染语义核，而被限制在占位层之后的表面投影层。
这正说明 $\mathsf{PL}_{\mathrm{NL}}$ 是关键中介层对象。证毕。
\end{proof}

\begin{remark}[评价]
所以第一章不是“本体爆发章”，而是“接口分层章”。
其重要性在于把
\[
  \text{定义语义}
  \neq
  \text{组织逻辑槽位}
  \neq
  \text{输出自然语言}
\]
这三件事正式拆开。
\end{remark}

\subsection{定理：第一章的工程增益显著高于其纯本体增益}
\label{subsec:tex37-gain-engineering}

\begin{theorem}[工程增益主导定理]
\label{thm:tex37-engineering-dominates}
设
\[
  G_{\mathrm{eng}}
  :=
  G_{\mathrm{closure}}
  +
  G_{\mathrm{runtime}}
  +
  G_{\mathrm{audit}}
  +
  G_{\mathrm{style}},
\]
并以
\[
  G_{\mathrm{ontology}}
\]
表示“新增基础本体宇宙、基本算子族或独立修复塔”的增益。
则在第一章当前已证明的范围内，有
\[
  G_{\mathrm{eng}} > G_{\mathrm{ontology}}.
\]
其原因是：第一章明确构造并证明的是
\[
  \Norm_G,\ \Theta_0,\ \Parse_\xi,\ \Style_\xi,\ \Audit_\xi,\ \Render^{\mathrm{LLM}}_\xi,\ \Engine_{G,\xi},
\]
而非新的大规模本体宇宙或新的基础修复塔。
\end{theorem}

\begin{Certificate}[证书：已证明模块的性质分类]
\label{cert:tex37-engineering-dominates}
证书登记谓词链：
\[
\begin{aligned}
  E_1 &: \Norm_G,\Theta_0,\Parse_\xi,\Style_\xi,\Render^{\mathrm{LLM}}_\xi,\Engine_{G,\xi}\ \text{均在第~\ref{sec:tex37-ch1}
    章被显式定义并证明},\\
  E_2 &: E_1\Rightarrow \text{这些对象全部属于“输入规约/输出规约/审计/回投影/运行时”接口类},\\
  E_3 &: \text{第一章新增的对象层核心是 } \mathsf{PL}_{\mathrm{NL}}\ \text{（引理~\ref{lem:tex37-middle-layer}）},\\
  E_4 &: E_3\Rightarrow \text{第一章没有像部分仓内文稿那样引入新的主丛宇宙、障碍修复塔或新的对象大类},\\
  E_5 &: E_2\wedge E_4\Rightarrow G_{\mathrm{eng}} > G_{\mathrm{ontology}}.
\end{aligned}
\]
\end{Certificate}

\begin{proof}[证明（谓词因果链）]
由 $E_1$，第一章的主要成果集中在规范化、占位搬运、自然语言回投影、风格审计与逻辑引擎闭环。
这些结构直接对应软件系统的模块边界，因此由 $E_2$ 可知它们属于工程接口类增益。

另一方面，由引理~\ref{lem:tex37-middle-layer}，第一章对象层的关键新增主要是
\[
  \mathsf{PL}_{\mathrm{NL}},
\]
即逻辑占位中介层；它的重要性很高，但并不等价于引入全新的本体宇宙、全新基本算子族或独立修复塔。
结合与本仓库其他偏本体/偏几何/偏障碍修复文稿的对照 \cite{RepoTex36HomotopyTwistBijection,GaoZhengAppVol3}，
可知第一章的已证明重量主要落在工程桥接侧而非本体扩张侧，即 $E_4$。

因此由 $E_2$ 与 $E_4$，得到
\[
  G_{\mathrm{eng}} > G_{\mathrm{ontology}}.
\]
证毕。
\end{proof}

\begin{remark}[评价]
若以“如何真正做成逻辑引擎、审计引擎、受控自然语言系统”为尺度，第一章很重；
若以“是否提出新的本体宇宙”为尺度，第一章并非最重。故其精确定位应是：
\[
  \boxed{
    \text{极强的工程桥接章}
  }.
\]
\end{remark}

\subsection{定理：第一章完成了“多风格输出而不改语义骨架”的严格封装}
\label{subsec:tex37-gain-style}

\begin{theorem}[同核多外观封装定理]
\label{thm:tex37-same-core-many-surfaces}
任取
\[
  p\in\mathsf{PL}_{\mathrm{NL}}^{(0)},
  \qquad
  \xi_1,\xi_2\in\Xi.
\]
若
\[
  t_1=\Render^{\mathrm{LLM}}_{\xi_1}(p)\neq\bot,
  \qquad
  t_2=\Render^{\mathrm{LLM}}_{\xi_2}(p)\neq\bot,
\]
则
\[
  \Parse_{\xi_1}(t_1)=\Parse_{\xi_2}(t_2)=p
\]
并且
\[
  \Theta_0^{-1}\bigl(\Parse_{\xi_1}(t_1)\bigr)
  =
  \Theta_0^{-1}\bigl(\Parse_{\xi_2}(t_2)\bigr).
\]
因此第一章已经把“严格版、解释版、教学版共享同一语义内核”提升为可审计机制，而不再只是经验直觉。
\end{theorem}

\begin{Certificate}[证书：同一占位 + 两次回投影 + 同一规范代表]
\label{cert:tex37-same-core-many-surfaces}
证书登记谓词链：
\[
\begin{aligned}
  S_1 &: t_1=\Render^{\mathrm{LLM}}_{\xi_1}(p)\neq\bot,\\
  S_2 &: t_2=\Render^{\mathrm{LLM}}_{\xi_2}(p)\neq\bot,\\
  S_3 &: S_1\wedge S_2\Rightarrow \Parse_{\xi_1}(t_1)=\Parse_{\xi_2}(t_2)=p\ \text{（命题~\ref{prop:tex37-style-invariance}）
    },\\
  S_4 &: S_3\Rightarrow \Theta_0^{-1}(\Parse_{\xi_1}(t_1))=\Theta_0^{-1}(p)=\Theta_0^{-1}(\Parse_{\xi_2}(t_2)),\\
  S_5 &: S_4\Rightarrow \text{多风格外观共享同一规范代表}.
\end{aligned}
\]
\end{Certificate}

\begin{proof}[证明（谓词因果链）]
由 $S_1$ 与 $S_2$，两个输出都为非失败输出。
根据命题~\ref{prop:tex37-style-invariance}，可得
\[
  \Parse_{\xi_1}(t_1)=\Parse_{\xi_2}(t_2)=p,
\]
即 $S_3$。
对上式两端同时施加 $\Theta_0^{-1}$，得到
\[
  \Theta_0^{-1}(\Parse_{\xi_1}(t_1))
  =
  \Theta_0^{-1}(p)
  =
  \Theta_0^{-1}(\Parse_{\xi_2}(t_2)),
\]
即 $S_4$。
因此不同模式只改变表面措辞与展开层级，而不会改变被回投影到的规范代表，故 $S_5$ 成立。证毕。
\end{proof}

\begin{remark}[评价]
这一点的实际意义在于：
\[
  \boxed{
    \text{“同核多外观”已从经验直觉提升为可审计结构}
  }.
\]
自然语言层的差异被约束为同一占位类的不同投影，而不再自动意味着语义漂移风险。
\end{remark}

\subsection{命题：第一章不是最大的本体扩张章，但属于重型工程桥接章}
\label{subsec:tex37-gain-bridge}

\begin{proposition}[第一章的章节定位]
\label{prop:tex37-bridge-chapter}
若以“新增宇宙本体、基本算子族、独立障碍修复塔”为主要评价尺度，则第一章不是本稿中最重的本体扩张章；
但若以“把已有理论压成可执行闭环、可审计运行时、可回投影逻辑接口”为尺度，则第一章属于重型工程桥接章。
\end{proposition}

\begin{Certificate}[证书：两套评价尺度上的比较]
\label{cert:tex37-bridge-chapter}
证书登记谓词链：
\[
\begin{aligned}
  B_1 &: \text{第一章的主成果集中于 } \Norm_G,\Theta_0,\Parse_\xi,\Style_\xi,\Render^{\mathrm{LLM}}_\xi,\Engine_{G,\xi},\\
  B_2 &: B_1\Rightarrow \text{第一章在工程接口尺度上具有高权重（定理~\ref{thm:tex37-engineering-dominates}）},\\
  B_3 &: \text{第一章并未主打一整套新的宇宙本体或新修复塔（引理~\ref{lem:tex37-middle-layer}）},\\
  B_4 &: B_3\Rightarrow \text{第一章在本体扩张尺度上不是峰值},\\
  B_5 &: B_2\wedge B_4\Rightarrow \text{第一章的准确定位是重型工程桥接章}.
\end{aligned}
\]
\end{Certificate}

\begin{proof}[证明（谓词因果链）]
由 $B_1$，第一章的证明重心在接口闭环而非宇宙扩张；
由定理~\ref{thm:tex37-engineering-dominates}，得 $B_2$；
由引理~\ref{lem:tex37-middle-layer}，得 $B_3$ 与 $B_4$。
将工程侧高权重与本体侧非峰值同时合并，便得到 $B_5$。命题成立。证毕。
\end{proof}

\subsection{推论：第一章是体系迈向“逻辑引擎化”的关键一步}
\label{subsec:tex37-gain-corollary}

\begin{corollary}[第一章的逻辑引擎化推论]
\label{cor:tex37-logic-engine-step}
对任意 $e\in\mathsf{Expr}_G$ 与任意 $\xi\in\Xi$，若存在
\[
  e
  \xrightarrow{\Norm_G}
  n
  \xrightarrow{\Theta_0}
  p
  \xrightarrow{\Render^{\mathrm{LLM}}_\xi}
  t
  \xrightarrow{\Parse_\xi}
  p
  \xrightarrow{\Theta_0^{-1}}
  n,
\]
则该系统同时具有
\[
  \text{可验证输入}
  +
  \text{可验证输出}
  +
  \text{可追责失败}
\]
三种性质。因此，第一章已经把体系从“语义理论”推进到“逻辑引擎架构”的门槛之上。
\end{corollary}

\begin{Certificate}[证书：输入规范化 + 输出回投影 + fail-closed]
\label{cert:tex37-logic-engine-step}
证书登记谓词链：
\[
\begin{aligned}
  C_1 &: e\xrightarrow{\Norm_G} n\neq\bot,\\
  C_2 &: n\xrightarrow{\Theta_0} p,\quad p\xrightarrow{\Render^{\mathrm{LLM}}_\xi} t\neq\bot,\\
  C_3 &: C_2\Rightarrow \Parse_\xi(t)=p\ \text{（定理~\ref{thm:tex37-runtime-deterministic}）},\\
  C_4 &: C_1\wedge C_2\wedge C_3\Rightarrow \Theta_0^{-1}(\Parse_\xi(t))=n\ \text{（推论~\ref{cor:tex37-logic-engine}）},\\
  C_5 &: \Norm_G(e)=\bot\ \text{或}\ \Audit_\xi(p)=\varnothing\Rightarrow \bot\ \text{被显式返回},\\
  C_6 &: C_4\wedge C_5\Rightarrow \text{系统已具备逻辑引擎的基本闭环}.
\end{aligned}
\]
\end{Certificate}

\begin{proof}[证明（谓词因果链）]
由 $C_1$，输入先被规范化，因此输入是可验证的；
由 $C_2$ 与 $C_3$，输出不是自由漂移文本，而是可被解析回原占位的自然语言表述，因此输出是可验证的；
由 $C_5$，一切失败都以 $\bot$ 的方式显式暴露，因此失败是可追责的。

再由推论~\ref{cor:tex37-logic-engine}，$C_1$、$C_2$ 与 $C_3$ 共同推出
\[
  \Theta_0^{-1}(\Parse_\xi(t))=n,
\]
即 $C_4$。
于是输入、输出与失败三端全部进入统一闭环，故 $C_6$ 成立。
这正是逻辑引擎架构的基本门槛。证毕。
\end{proof}

\subsection{备注：分层总评、边界与一句话结论}
\label{subsec:tex37-gain-remarks}

\begin{remark}[分层总评]
若把第一章增益分层记账，则可给出如下结论：
\[
  \text{公理级增益}=\text{中},
  \qquad
  \text{定理级增益}=\text{高},
  \qquad
  \text{工程级增益}=\text{很高},
  \qquad
  \text{体系级增益}=\text{高}.
\]
其理由是：第一章没有大规模扩充新的本体宇宙，但它把若干原本偏直觉的接口关系压成了可运行、可审计、可回投影的定理链。
\end{remark}

\begin{remark}[第一章最核心的新增不是“生成内容”，而是“约束内容”]
通常人们会把 LLM 的价值理解为
\[
  \text{它能生成很多文字},
\]
而第一章最重要的改写是把这一点降为
\[
  \text{它只能在既定占位与审计合同内表达}.
\]
因此，这一章的深层增益并非“增强生成”，而是
\[
  \text{把生成能力从主权地位降为执行地位}.
\]
\end{remark}

\begin{remark}[边界：闭环架构已被严格写出，但若干关键接口仍待构造性落地]
需要诚实记账的未完全构造部分包括：
\begin{enumerate}
  \item $\Theta_0$ 的完全可计算实现是否已经在所有目标域上给出；
  \item 各阶纤维双射 $\Psi_x^{(n)}$ 是否已有统一而可执行的构造；
  \item $\Parse_\xi$ 对开放自然语言的完备性与鲁棒边界是否已经完全刻画；
  \item $\Style_\xi$ 的合同边界与冲突裁决是否已全部形式化；
  \item 候选采集协议在真实系统中的鲁棒性是否已完成系统化实测。
\end{enumerate}
因此，第一章当前最准确的定位应写成：
\[
  \boxed{
    \text{闭环架构已被严格写出，但若干关键接口仍需构造性落地}
  }.
\]
\end{remark}

\begin{remark}[一句话结论]
若只用一句话概括第一章的最大增益，则可写成：
\[
  \boxed{
    \text{第一章把 G 框架的规范语义正式接成了“自然语言可审计运行时”的闭环逻辑引擎。}
  }.
\]
进一步压缩为更短的一句，则是：
\[
  \boxed{
    \text{它把“理论能说清楚”推进成了“理论能被受控地说出来，并且能回查是否说对”。}
  }.
\]
\end{remark}

\clearpage

%%%%%%%%%%%%%%%%%%%%%%%%%%%%%%%%%%%%%%%%%%%%%%%%%%%%%%%%%%%%
\section{形式逻辑：规范语义塔、逻辑占位塔与表面文本层的三层严格分离}
\label{sec:tex37-ch3}
%%%%%%%%%%%%%%%%%%%%%%%%%%%%%%%%%%%%%%%%%%%%%%%%%%%%%%%%%%%%

\begin{remark}[核心结论（收紧版）]
第一章的核心并不是
\[
  \text{G 侧规范语义塔}
  \leftrightarrow
  \text{LLM 最终输出文本},
\]
而是
\[
  \boxed{
    \text{G 侧规范语义塔}
    \leftrightarrow
    \text{自然语言语义逻辑占位塔}
  }.
\]
其后才由 LLM 在既定候选协议、解析器与风格合同约束下，把给定占位展开成模式 $\xi$ 下的表面文本。
因此必须严格区分：
\[
  \text{规范语义层}
  \neq
  \text{逻辑占位层}
  \neq
  \text{表面文本层}.
\]
\end{remark}

本章对第~\ref{sec:tex37-ch1} 章作两步收紧：
\begin{enumerate}
  \item 把“高阶同伦扭曲双射的终点对象”严格限定为逻辑占位塔，而不是自然语言表面文本层；
  \item 把“LLM 固定占位”改写为“占位由规范化、搬运、解析与审计共同固定，LLM 只服从该占位并给出受筛选的表面表达”。
\end{enumerate}

\subsection{对象链、三层分离与混合表面表达口径}
\label{subsec:tex37-ch3-objects}

\begin{definition}[三层对象链]
第一章的基础对象链应严格写为
\[
  \mathsf{Expr}_G
  \xrightarrow{\Norm_G}
  \mathsf{NF}_G^{(0)}
  \xrightarrow{\Theta_0}
  \mathsf{PL}_{\mathrm{NL}}^{(0)}
  \xrightarrow{\Render^{\mathrm{LLM}}_\xi}
  \mathsf{Txt}_\xi.
\]
其中：
\begin{enumerate}
  \item $\mathsf{NF}_G^{(0)}$ 是 G 侧规范语义的 $0$ 阶规范代表空间；
  \item $\mathsf{PL}_{\mathrm{NL}}^{(0)}$ 是自然语言语义逻辑占位空间；
  \item $\mathsf{Txt}_\xi$ 是模式 $\xi$ 下的表面文本空间。
\end{enumerate}
\end{definition}

\begin{definition}[高阶扩张塔与双射塔]
记第一章中已构造的高阶扩张塔与双射塔为
\[
  \Theta=\{\Theta_n\}_{n\in\NN},
  \qquad
  \Theta_n:\mathsf{NF}_G^{(n)}\to \mathsf{PL}_{\mathrm{NL}}^{(n)}.
\]
其中每个 $\Theta_n$ 由第~\ref{sec:tex37-ch1} 章中的高阶同伦扭曲递推规则定义，
而其对象端点始终分别属于
\[
  \mathsf{NF}_G^{(n)}
  \quad\text{与}\quad
  \mathsf{PL}_{\mathrm{NL}}^{(n)}.
\]
\end{definition}

\begin{definition}[自然语言/数学符号混合表面表达口径]
在本章中，把
\[
  \mathsf{Txt}_\xi
\]
更精确地解释为\textbf{面向自然语言主骨架、并允许嵌入数学符号串的混合表面表达空间}。
因此这里的“自然语言”不是排斥数学符号，而是指：
\[
  \text{自然语言主骨架}
  +
  \text{数学符号可嵌入}.
\]
\end{definition}

\begin{remark}[与仓内三层架构文稿的对齐]
这种三层分离口径与仓内关于 HACA/语义规范场三层结构的总背景一致 \cite{GaoZhengAppVol3}，
也与“LLM 只提案、审计器裁决”的工程包装口径一致 \cite{RepoTex21AuditedWrapper,RepoTex23SemanticBase}。
本章只是把这种分离进一步压成针对第~\ref{sec:tex37-ch1} 章的严格对象论断。
\end{remark}

\subsection{公理（降级为定理）：双射对象的类型边界在 N 阶同伦塔上递归保持}
\label{subsec:tex37-ch3-axiom}

\begin{theorem}[公理降级：双射对象始终是占位塔而非表面文本层]
\label{thm:tex37-ch3-layer-induction}
定义谓词
\[
  P(n):\
  \Bigl(
    \Theta_n:\mathsf{NF}_G^{(n)}\to \mathsf{PL}_{\mathrm{NL}}^{(n)}
    \text{ 是双射}
  \Bigr)
  \wedge
  \Bigl(
    \Theta_n\text{ 的递推构造不以 }\mathsf{Txt}_\xi\text{ 为端点对象}
  \Bigr).
\]
则对所有 $n\in\NN$，$P(n)$ 成立。
因此，第~\ref{sec:tex37-ch1} 章所证明的规范双射始终是
\[
  \mathsf{NF}_G^{(n)}
  \leftrightarrow
  \mathsf{PL}_{\mathrm{NL}}^{(n)},
\]
而不是
\[
  \mathsf{NF}_G^{(n)}
  \leftrightarrow
  \mathsf{Txt}_\xi.
\]
\end{theorem}

\begin{Certificate}[证书：基步 + 归纳步（数学归纳法证书）]
\label{cert:tex37-ch3-layer-induction}
归纳证书登记为
\[
  \pi_{\mathrm{ind}}
  =
  \bigl(P(0),\ \forall n\in\NN,\ P(n)\Rightarrow P(n+1)\bigr).
\]
其中：
\[
  P(0)
  \ \text{来自基阶双射}\
  \Theta_0:\mathsf{NF}_G^{(0)}\to \mathsf{PL}_{\mathrm{NL}}^{(0)},
\]
而归纳步由递推式
\[
  \Theta_{n+1}(x,\eta)
  =
  \bigl(
    \Theta_n(x),
    \Psi_x^{(n)}(\eta)
  \bigr)
\]
给出。
\end{Certificate}

\begin{proof}[证明（数学归纳法证书；谓词因果链）]
\textbf{基步 $P(0)$：}
由第~\ref{sec:tex37-ch1} 章的定义与定理~\ref{thm:tex37-axiom-downgrade}，
已知
\[
  \Theta_0:\mathsf{NF}_G^{(0)}\to \mathsf{PL}_{\mathrm{NL}}^{(0)}
\]
是基阶双射。
并且其值域对象是占位空间 $\mathsf{PL}_{\mathrm{NL}}^{(0)}$，而不是任何 $\mathsf{Txt}_\xi$。
故 $P(0)$ 成立。

\textbf{归纳步：}
设 $P(n)$ 成立。则由第~\ref{sec:tex37-ch1} 章的递推定义，
\[
  \Theta_{n+1}(x,\eta)
  =
  \bigl(
    \Theta_n(x),
    \Psi_x^{(n)}(\eta)
  \bigr).
\]
该构造只调用：
\begin{enumerate}
  \item 上一层双射 $\Theta_n$；
  \item 纤维双射 $\Psi_x^{(n)}$；
  \item 两侧的纤维化扩张对象 $\mathsf{NF}_G^{(n+1)}$ 与 $\mathsf{PL}_{\mathrm{NL}}^{(n+1)}$。
\end{enumerate}
其中没有任何一步把 $\mathsf{Txt}_\xi$ 作为递推端点对象并入定义。

再由定理~\ref{thm:tex37-axiom-downgrade}，$\Theta_{n+1}$ 仍为双射。
因此
\[
  P(n+1)
\]
成立。

由数学归纳法，对所有 $n\in\NN$ 都有 $P(n)$。
于是高阶同伦扭曲双射的对象边界在每一阶都保持为
\[
  \mathsf{NF}_G^{(n)}
  \leftrightarrow
  \mathsf{PL}_{\mathrm{NL}}^{(n)},
\]
而不是表面文本层。证毕。
\end{proof}

\subsection{引理：第一章中被证明的双射对象必须收紧为逻辑占位塔}
\label{subsec:tex37-ch3-lemma}

\begin{lemma}[双射终点对象不是表面句子空间]
\label{lem:tex37-ch3-bijection-target}
在第~\ref{sec:tex37-ch1} 章的已证明结构中，被规范证明为双射的一端只能是
\[
  \mathsf{PL}_{\mathrm{NL}}^{(n)},
\]
而不能直接写成模式 $\xi$ 下的表面文本空间 $\mathsf{Txt}_\xi$。
\end{lemma}

\begin{Certificate}[证书：总双射 vs. 模式依赖部分运行时]
\label{cert:tex37-ch3-bijection-target}
证书登记谓词链：
\[
\begin{aligned}
  L_1 &: \Theta_n:\mathsf{NF}_G^{(n)}\to \mathsf{PL}_{\mathrm{NL}}^{(n)}\ \text{是双射（定理~\ref{thm:tex37-ch3-layer-induction}）},\\
  L_2 &: \Render^{\mathrm{LLM}}_\xi:\mathsf{PL}_{\mathrm{NL}}^{(0)}\to \mathsf{Txt}_\xi\sqcup\{\bot\}\ \text{是模式依赖的部分运行时}
    ,\\
  L_3 &: L_2\Rightarrow \Render^{\mathrm{LLM}}_\xi\ \text{可能失败返回 }\bot,\\
  L_4 &: L_1\wedge L_3\Rightarrow \mathsf{Txt}_\xi\ \text{不是第一章中被证明的双射端点对象}.
\end{aligned}
\]
\end{Certificate}

\begin{proof}[证明（谓词因果链）]
由 $L_1$，第一章中被证明为双射的对象对是
\[
  \mathsf{NF}_G^{(n)}
  \quad\text{与}\quad
  \mathsf{PL}_{\mathrm{NL}}^{(n)}.
\]
由 $L_2$，从占位到文本的映射是
\[
  \Render^{\mathrm{LLM}}_\xi:\mathsf{PL}_{\mathrm{NL}}^{(0)}\to \mathsf{Txt}_\xi\sqcup\{\bot\},
\]
它依赖模式 $\xi$，并且由 $L_3$ 可能返回 $\bot$。

一个带失败值、且依赖外部模式参数的部分运行时，不能作为第~\ref{sec:tex37-ch1} 章中
“规范语义塔的规范双射终点对象”来表述。
因此，若要严格复述第一章的主结论，双射终点对象必须收紧为逻辑占位塔，而不是表面句子空间。证毕。
\end{proof}

\subsection{命题：应将“最小自然语言”收紧为“自然语言可复原的最小逻辑骨架空间”}
\label{subsec:tex37-ch3-minimal-language}

\begin{proposition}[“最小自然语言”的严格改写]
\label{prop:tex37-ch3-minimal-language}
若要严格描述第~\ref{sec:tex37-ch1} 章中的对象
\[
  \mathsf{PL}_{\mathrm{NL}}^{(0)},
\]
则更准确的说法不是“最小自然语言句子空间”，而应写成：
\[
  \boxed{
    \text{自然语言可复原的最小逻辑骨架空间}
  }
\]
或
\[
  \boxed{
    \text{面向自然语言/数学符号混合表达的规范语义逻辑占位空间}
  }.
\]
\end{proposition}

\begin{Certificate}[证书：占位定义 vs. 句子定义]
\label{cert:tex37-ch3-minimal-language}
证书登记谓词链：
\[
\begin{aligned}
  M_1 &: \mathsf{PL}_{\mathrm{NL}}^{(0)}\ \text{在第~\ref{sec:tex37-ch1} 章中被定义为逻辑占位空间},\\
  M_2 &: M_1\Rightarrow \text{其元素是“可被自然语言复原的语义逻辑槽位”而非最终句子},\\
  M_3 &: \mathsf{Txt}_\xi\ \text{才承担最终表面表达角色},\\
  M_4 &: M_2\wedge M_3\Rightarrow \text{“最小自然语言”若不加限定会误导为句子层对象},\\
  M_5 &: M_4\Rightarrow \text{应改写为“自然语言可复原的最小逻辑骨架空间”}.
\end{aligned}
\]
\end{Certificate}

\begin{proof}[证明（谓词因果链）]
由 $M_1$，$\mathsf{PL}_{\mathrm{NL}}^{(0)}$ 在第~\ref{sec:tex37-ch1} 章中被定义为占位空间，
其内部包含量词骨架、谓词依赖链、证书槽与风格参数接口，而不是完整自然语言句子。
故由 $M_2$，它的元素只能解释为“可复原逻辑骨架”。

另一方面，由 $M_3$，最终句子或混合表面文本位于 $\mathsf{Txt}_\xi$ 层。
因此若把 $\mathsf{PL}_{\mathrm{NL}}^{(0)}$ 直接称为“最小自然语言”，就会把占位层与句子层混在一起，
即 $M_4$。
故最严格的表述只能是 $M_5$ 所示的“自然语言可复原的最小逻辑骨架空间”。
证毕。
\end{proof}

\begin{remark}[补充说明]
若需要把数学符号表达显式并入，则最稳妥的表述是：
\[
  \text{自然语言/数学符号混合表面表达}
\]
服从同一占位骨架。这样既不把输出层缩成纯口语，也不把占位层误写成句子层。
\end{remark}

\subsection{命题：不是 LLM 固定占位，而是 LLM 服从既定占位}
\label{subsec:tex37-ch3-placeholder-fixed}

\begin{proposition}[占位的固定权不属于 LLM]
\label{prop:tex37-ch3-placeholder-fixed}
对任意 $e\in\mathsf{Expr}_G$ 与任意模式 $\xi\in\Xi$，若
\[
  t=\Engine_{G,\xi}(e)\neq\bot,
\]
并记
\[
  p:=\Theta_0(\Norm_G(e)),
\]
则
\[
  \Parse_\xi(t)=p.
\]
因此占位 $p$ 的固定来自
\[
  \Norm_G+\Theta_0+\Parse_\xi
\]
的共同约束，审计接受性来自
\[
  \Style_\xi+\Audit_\xi,
\]
而不是由 LLM 单方面“固定占位”。
\end{proposition}

\begin{Certificate}[证书：规范化先行 + 占位搬运 + 解析回投影]
\label{cert:tex37-ch3-placeholder-fixed}
证书登记谓词链：
\[
\begin{aligned}
  P_1 &: p=\Theta_0(\Norm_G(e)),\\
  P_2 &: t=\Render^{\mathrm{LLM}}_\xi(p)\neq\bot,\\
  P_3 &: P_2\Rightarrow \Parse_\xi(t)=p\ \text{（定理~\ref{thm:tex37-runtime-deterministic}）},\\
  P_4 &: P_1\wedge P_3\Rightarrow p\ \text{由 }\Norm_G,\Theta_0,\Parse_\xi\ \text{共同固定},\\
  P_5 &: t\in\Audit_\xi(p)\Rightarrow \Style_\xi(t)=\top,\\
  P_6 &: P_4\wedge P_5\Rightarrow \text{LLM 只服从既定占位并提供候选表面表达}.
\end{aligned}
\]
\end{Certificate}

\begin{proof}[证明（谓词因果链）]
由 $P_1$，占位首先由规范化与占位搬运确定。
由 $P_2$，LLM 运行时只是接收该占位并生成一个非失败输出。
再由定理~\ref{thm:tex37-runtime-deterministic}，得到
\[
  \Parse_\xi(t)=p,
\]
即 $P_3$。
因此输出文本只能被接受为“这个占位 $p$ 的一个表面展开”，而不能反过来决定 $p$ 本身，
从而 $P_4$ 成立。

另一方面，若 $t$ 被接受，则按 $\Audit_\xi(p)$ 的定义还必须满足 $\Style_\xi(t)=\top$，
即 $P_5$。
合并 $P_4$ 与 $P_5$，可知 LLM 的职责只是：
\[
  \text{在既定占位与审计约束下生成并被筛选出的表面表达},
\]
故 $P_6$ 成立。证毕。
\end{proof}

\subsection{定理：非失败输出与审计可接受集非空构成充要等价}
\label{subsec:tex37-ch3-iff}

\begin{theorem}[运行时非失败输出的充要条件]
\label{thm:tex37-ch3-render-iff-audit}
对任意模式 $\xi\in\Xi$ 与任意占位 $p\in\mathsf{PL}_{\mathrm{NL}}^{(0)}$，都有
\[
  \boxed{
    \Render^{\mathrm{LLM}}_\xi(p)\neq\bot
    \iff
    \Audit_\xi(p)\neq\varnothing
  }.
\]
因此，可审计性不仅是必要条件，而且在固定候选协议与全序选择下也是充分条件。
\end{theorem}

\begin{Certificate}[证书：两向构造]
\label{cert:tex37-ch3-render-iff-audit}
证书登记谓词链：
\[
\begin{aligned}
  A_1 &: \Audit_\xi(p)\neq\varnothing,\\
  A_2 &: A_1\Rightarrow \min\nolimits_{\prec_\xi}\Audit_\xi(p)\ \text{存在},\\
  A_3 &: A_2\Rightarrow \Render^{\mathrm{LLM}}_\xi(p)\neq\bot,\\
  A_4 &: \Render^{\mathrm{LLM}}_\xi(p)\neq\bot,\\
  A_5 &: A_4\Rightarrow \Render^{\mathrm{LLM}}_\xi(p)\in \Audit_\xi(p),\\
  A_6 &: A_5\Rightarrow \Audit_\xi(p)\neq\varnothing.
\end{aligned}
\]
\end{Certificate}

\begin{proof}[证明（谓词因果链）]
\textbf{正向：}
若 $\Audit_\xi(p)\neq\varnothing$，则存在至少一个候选
\[
  t\in\mathcal{C}^{\mathrm{LLM}}_\xi(p)
\]
满足
\[
  \Parse_\xi(t)=p
  \qquad\text{且}\qquad
  \Style_\xi(t)=\top.
\]
由全序 $\prec_\xi$，最小元
\[
  \min\nolimits_{\prec_\xi}\Audit_\xi(p)
\]
存在，因此按定义
\[
  \Render^{\mathrm{LLM}}_\xi(p)\neq\bot.
\]
故 $A_1\Rightarrow A_3$。

\textbf{反向：}
若 $\Render^{\mathrm{LLM}}_\xi(p)\neq\bot$，则按定义
\[
  \Render^{\mathrm{LLM}}_\xi(p)=\min\nolimits_{\prec_\xi}\Audit_\xi(p),
\]
于是该最小元属于 $\Audit_\xi(p)$。因此 $\Audit_\xi(p)$ 必非空，即 $A_4\Rightarrow A_6$。

两向合并，得到
\[
  \Render^{\mathrm{LLM}}_\xi(p)\neq\bot
  \iff
  \Audit_\xi(p)\neq\varnothing.
\]
证毕。
\end{proof}

\begin{remark}[强化解释]
因此更严格的工程化表述不是“可审计性是必要条件”，而是：
\[
  \boxed{
    \text{存在可审计候选}
    \iff
    \text{存在合法表面表达}
  }.
\]
这把 LLM 输出从“自由生成”收紧为“审计可接受集上的确定性选择”。
\end{remark}

\subsection{命题：风格模板改变的是表面投影，不改变占位核}
\label{subsec:tex37-ch3-style}

\begin{proposition}[风格变化属于同一占位类上的投影变化]
\label{prop:tex37-ch3-style}
若
\[
  t_1=\Render^{\mathrm{LLM}}_{\xi_1}(p)\neq\bot,
  \qquad
  t_2=\Render^{\mathrm{LLM}}_{\xi_2}(p)\neq\bot,
\]
则
\[
  \Parse_{\xi_1}(t_1)=\Parse_{\xi_2}(t_2)=p.
\]
因此可选风格模板改变的是表面投影，而不改变占位核。
\end{proposition}

\begin{Certificate}[证书：运行时降级定理的两次实例化]
\label{cert:tex37-ch3-style}
证书登记谓词链：
\[
\begin{aligned}
  S_1 &: t_1=\Render^{\mathrm{LLM}}_{\xi_1}(p)\neq\bot,\\
  S_2 &: t_2=\Render^{\mathrm{LLM}}_{\xi_2}(p)\neq\bot,\\
  S_3 &: S_1\Rightarrow \Parse_{\xi_1}(t_1)=p\ \text{（定理~\ref{thm:tex37-runtime-deterministic}）},\\
  S_4 &: S_2\Rightarrow \Parse_{\xi_2}(t_2)=p\ \text{（定理~\ref{thm:tex37-runtime-deterministic}）},\\
  S_5 &: S_3\wedge S_4\Rightarrow \text{风格变化只改变表面投影}.
\end{aligned}
\]
\end{Certificate}

\begin{proof}[证明（谓词因果链）]
由 $S_1$ 与 $S_2$，两个输出都是非失败输出。
分别对 $\xi_1$ 与 $\xi_2$ 应用定理~\ref{thm:tex37-runtime-deterministic}，得
\[
  \Parse_{\xi_1}(t_1)=p,
  \qquad
  \Parse_{\xi_2}(t_2)=p.
\]
于是两种风格模板只对应同一占位 $p$ 的不同表面截面，而不改变占位核。
证毕。
\end{proof}

\subsection{推论：更严格的总表述}
\label{subsec:tex37-ch3-corollary}

\begin{corollary}[规范语义塔、逻辑占位塔与表面文本层的严格总表述]
\label{cor:tex37-ch3-strict-summary}
第一章的严格总表述应写成：
\[
  \boxed{
    \text{基于 G 框架规范语义纤维丛的基底，逐阶构造 $N$ 阶同伦塔上的纤维化扩张空间，}
  }
\]
\[
  \boxed{
    \begin{aligned}
      &\text{并通过逐阶兼容的纤维双射，建立 G 侧规范语义塔}\\
      &\text{与自然语言语义逻辑占位塔之间的双射；}
    \end{aligned}
  }
\]
\[
  \boxed{
    \begin{aligned}
      &\text{该双射的对象不是自然语言表面句子本身，}\\
      &\text{而是可被自然语言/数学符号混合表达复原的最小逻辑占位骨架。}
    \end{aligned}
  }
\]
\[
  \boxed{
    \begin{aligned}
      &\text{随后，LLM 仅在固定候选协议、解析器与风格合同约束下，}\\
      &\text{把给定占位映射为某一模式 }\xi\text{ 下的表面表达；}
    \end{aligned}
  }
\]
\[
  \boxed{
    \Render^{\mathrm{LLM}}_\xi(p)\neq\bot
    \iff
    \Audit_\xi(p)\neq\varnothing.
  }
\]
\[
  \boxed{
    \text{因此风格模板的可选性只改变表面展开，不改变由占位所固定的规范语义核。}
  }
\]
\end{corollary}

\begin{Certificate}[证书：双射对象收紧 + 审计充要条件 + 风格投影不变]
\label{cert:tex37-ch3-strict-summary}
证书登记谓词链：
\[
\begin{aligned}
  C_1 &: \mathsf{NF}_G^{(n)}\leftrightarrow \mathsf{PL}_{\mathrm{NL}}^{(n)}\ \text{是双射对象对（定理~\ref{thm:tex37-ch3-layer-induction}）},\\
  C_2 &: \mathsf{PL}_{\mathrm{NL}}^{(0)}\ \text{应解释为最小逻辑骨架空间（命题~\ref{prop:tex37-ch3-minimal-language}）},\\
  C_3 &: \Render^{\mathrm{LLM}}_\xi(p)\neq\bot\iff \Audit_\xi(p)\neq\varnothing\ \text{（定理~\ref{thm:tex37-ch3-render-iff-audit}）},\\
  C_4 &: \text{风格变化只改变表面投影，不改变占位核（命题~\ref{prop:tex37-ch3-style}）},\\
  C_5 &: C_1\wedge C_2\wedge C_3\wedge C_4\Rightarrow \text{上面的严格总表述成立}.
\end{aligned}
\]
\end{Certificate}

\begin{proof}[证明（谓词因果链）]
由 $C_1$，高阶双射的对象端点只能写成规范语义塔与逻辑占位塔；
由 $C_2$，占位塔应解释为“可复原逻辑骨架”，而不是句子集合；
由 $C_3$，非失败输出与审计可接受集非空构成充要等价；
由 $C_4$，风格模板只改变表面截面，而不改占位核。
把四点合并，即得推论中的严格总表述。证毕。
\end{proof}

\subsection{备注：一句话压缩与最关键的一刀}
\label{subsec:tex37-ch3-remarks}

\begin{remark}[一句话压缩版]
最短的严格压缩可写成：
\[
  \boxed{
    \begin{aligned}
      &\text{不是“LLM 与规范语义双射”，}\\
      &\text{而是“规范语义与逻辑占位双射；LLM 只是在审计约束下}\\
      &\text{把占位展开成可选风格的表面文本”。}
    \end{aligned}
  }.
\]
\end{remark}

\begin{remark}[最关键的一刀]
本章最关键的收紧，不是新增对象，而是把三层关系彻底钉死：
\[
  \text{规范语义层}
  \neq
  \text{逻辑占位层}
  \neq
  \text{表面文本层}.
\]
第一章最强的地方，正是把 LLM 严格降到第三层；本章只是把这一点从正确直觉再压成不可混淆的形式命题。
\end{remark}

\clearpage

%%%%%%%%%%%%%%%%%%%%%%%%%%%%%%%%%%%%%%%%%%%%%%%%%%%%%%%%%%%%
\section{形式逻辑：问答内核对象是上下文基底上的同伦修复塔路径类，表面文本只是可审计投影壳}
\label{sec:tex37-ch4}
%%%%%%%%%%%%%%%%%%%%%%%%%%%%%%%%%%%%%%%%%%%%%%%%%%%%%%%%%%%%

\begin{remark}[核心陈述（再次收紧后）]
本章把前文的“三层分离”继续向内核动力学侧收紧为：
\[
  \boxed{
    \begin{aligned}
      &\text{外显对话对象是逻辑占位与表面文本；}\\
      &\text{内核求解对象是上下文基底上的同伦修复塔允许路径类。}
    \end{aligned}
  }
\]
更具体地说：
\[
  \boxed{
    \begin{aligned}
      &\text{问答的“解”不是一句话本身，}\\
      &\text{而是上下文延拓后在修复塔上选出的可行且近优的路径积分代表。}
    \end{aligned}
  }
\]
同时还必须加上一条边界限定：
\[
  \boxed{
    \text{“比安基恒等式成立”一般不是上下文延拓的先验事实，而是障碍修复后的目标闭合态。}
  }
\]
\end{remark}

\begin{remark}[阅读导航：仓内文稿与外部参照]
本章直接承接前文关于规范语义塔、逻辑占位塔与表面文本层三层分离的结论，
并把仓内关于主丛联络--路径积分--同伦扭曲集合等价的整理稿 \cite{RepoTex36HomotopyTwistBijection}
继续向“上下文修复动力学”推进；关于扩展签发协议、同伦修复塔与结构表示定理，可参照 \cite{RepoTex35ExtendedCertificationTower}；
关于“可修复障碍”及其严格边界，可参照 \cite{RepoTex29RepairableObstruction}；
更系统的 law-space、Jacobi gap、开放系统与证书治理背景，可参照 \cite{GaoZhengAppVol3}。

外部标准参照方面，本章借用联络、曲率、比安基恒等式、同伦修复与同调证书的标准语言，
可参照 \cite{KobayashiNomizu1963Foundations,Hatcher2002AT,Weibel1994HomologicalAlgebra}。
\end{remark}

\subsection{上下文基底、修复塔、障碍与投影链}
\label{subsec:tex37-ch4-definitions}

\begin{definition}[问答上下文基底与同伦修复塔纤维丛]
\label{def:tex37-ch4-base}
设
\[
  \mathcal{B}_{\mathrm{ctx}}
\]
为问答上下文基底空间。一个基点
\[
  b\in \mathcal{B}_{\mathrm{ctx}}
\]
编码当前会话中已被接受的语义约束、引用接口、变量绑定、证书槽、失败槽与风格合同约束。

在其上给定一个 $N$ 阶同伦修复塔纤维丛
\[
  \pi_N:\mathcal{E}^{(N)}_{\mathrm{ctx}}\to \mathcal{B}_{\mathrm{ctx}},
\]
其纤维
\[
  \mathcal{E}^{(N)}_b:=\pi_N^{-1}(b)
\]
表示在上下文基点 $b$ 上所有允许的规范形、扭曲数据、证书依赖、修复分支与展开层级。

若把第一章的高阶同伦扭曲双射
\[
  \Theta_n:\mathsf{NF}_G^{(n)}\to \mathsf{PL}_{\mathrm{NL}}^{(n)}
\]
嵌入该结构，则可把 $\mathsf{NF}_G^{(n)}$ 视为 $\mathcal{E}^{(N)}_{\mathrm{ctx}}$ 的规范语义截面族，
把 $\mathsf{PL}_{\mathrm{NL}}^{(n)}$ 视为其自然语言逻辑占位投影族。
\end{definition}

\begin{definition}[上下文延拓、联络、曲率、比安基缺陷与总障碍]
\label{def:tex37-ch4-obstruction}
给定当前问题 $q$，定义上下文延拓算子
\[
  \operatorname{Ext}_q:\mathcal{B}_{\mathrm{ctx}}\to \mathcal{B}_{\mathrm{ctx}},
  \qquad
  b\mapsto b^+:=\operatorname{Ext}_q(b).
\]
对每个基点 $b$，赋予局部联络
\[
  A_b,
\]
曲率
\[
  F_b,
\]
以及由局部括号失配、依赖失配、证书失配所诱导的雅可比障碍
\[
  J_b.
\]
定义比安基缺陷
\[
  \mathcal{B}\!i_b:=D_{A_b}F_b,
\]
并定义总障碍向量
\[
  \operatorname{Ob}_b
  :=
  \bigl(
    |\mathcal{B}\!i_b|,
    |J_b|,
    |\Delta_{\mathrm{cert},b}|,
    |\Delta_{\mathrm{ref},b}|
  \bigr)
  \in \mathbb{R}_{\ge 0}^{4}.
\]
其中
\[
  |\mathcal{B}\!i_b|
\]
衡量“联络--曲率--上下文延拓”是否闭合，
\[
  |J_b|
\]
衡量“语义括号--推理拼接--结构依赖”是否出现雅可比残差，
其余两项分别衡量证书与引用接口是否断裂。
\end{definition}

\begin{definition}[修复算子、允许路径类与作用泛函]
\label{def:tex37-ch4-action}
对每个基点 $b$，设
\[
  \Gamma_b
\]
为在修复塔 $\mathcal{E}^{(N)}_{\mathrm{ctx}}$ 中、以 $b$ 为基底的允许路径类集合，其元素记为
\[
  [\gamma].
\]
定义修复算子
\[
  \mathcal{R}:\mathcal{E}^{(N)}_{\mathrm{ctx}}\to \mathcal{E}^{(N)}_{\mathrm{ctx}},
\]
满足：若 $x$ 仍含不可接受障碍，则
\[
  \operatorname{Ob}(\mathcal{R}(x))
  \preceq_{\mathrm{lex}}
  \operatorname{Ob}(x),
\]
且在尚未进入可行域时优先严格压低可行性分量。

给定可接受边界
\[
  \varepsilon_b\in \mathbb{R}_{\ge 0}^{4},
\]
定义可接受路径类集合
\[
  \Gamma_b^{\mathrm{adm}}
  :=
  \bigl\{
    [\gamma]\in \Gamma_b
    \mid
    \operatorname{Ob}_{\gamma(1)}
    \preceq_{\mathrm{lex}}
    \varepsilon_b
  \bigr\}.
\]

定义作用泛函
\[
  \mathcal{S}_b([\gamma])
  :=
  \mathcal{S}_{\mathrm{feas},b}([\gamma])
  +
  \lambda\,\mathcal{S}_{\mathrm{opt},b}([\gamma]),
  \qquad
  \lambda\ge 0,
\]
其中
\[
  \mathcal{S}_{\mathrm{feas},b}([\gamma])
  :=
  \int_{\gamma}
  \bigl(
    \alpha_1|\mathcal{B}\!i|
    +
    \alpha_2|J|
    +
    \alpha_3|\Delta_{\mathrm{cert}}|
    +
    \alpha_4|\Delta_{\mathrm{ref}}|
  \bigr)\,\mathrm{d}s,
\]
\[
  \mathcal{S}_{\mathrm{opt},b}([\gamma])
\]
衡量在已可行前提下的复杂度、长度、代价、稳定性损耗与说明冗余。
这里默认路径类已按“保留终点与障碍积分”的等价关系取商，故上述作用量良定义。
\end{definition}

\begin{definition}[规范语义代表、逻辑占位截断与表面投影链]
\label{def:tex37-ch4-projection}
对每个 $n\in\mathbb{N}$，设
\[
  \Sigma_n:\Gamma_b^{\mathrm{adm},(n)}\to \mathsf{NF}_G^{(n)}
\]
为把可接受路径类送到其规范语义代表的选择映射，
并设
\[
  \operatorname{pr}_{n\to 0}:\mathsf{PL}_{\mathrm{NL}}^{(n)}\to \mathsf{PL}_{\mathrm{NL}}^{(0)}
\]
为占位塔到基阶逻辑占位的截断映射。
对任意输出模式 $\xi\in\Xi$，定义表面投影链
\[
  \mathcal{P}^{(n)}_{b,\xi}
  :=
  \Render^{\mathrm{LLM}}_\xi
  \circ
  \operatorname{pr}_{n\to 0}
  \circ
  \Theta_n
  \circ
  \Sigma_n.
\]
因此，一个对外显示的文本
\[
  t\in \mathsf{Txt}_\xi\sqcup\{\bot\}
\]
若非失败，则只能作为某个可接受路径类经过
\[
  [\gamma]
  \longmapsto
  \Sigma_n([\gamma])
  \longmapsto
  \Theta_n\bigl(\Sigma_n([\gamma])\bigr)
  \longmapsto
  \operatorname{pr}_{n\to 0}
  \longmapsto
  \Render^{\mathrm{LLM}}_\xi
\]
后的表面代表。
\end{definition}

\begin{remark}[本章的分层口径]
由定义~\ref{def:tex37-ch4-projection} 可见，本章把“求解对象”与“显示对象”再分成四层：
\[
  \text{修复塔路径类/截面}
  \to
  \text{规范语义代表}
  \to
  \text{逻辑占位}
  \to
  \text{表面文本}.
\]
这比第三章的三层分离再向内多压了一层：把“规范语义代表”继续追溯到其在上下文基底上的动态生成机制。
\end{remark}

\subsection{公理（降级为定理）：修复塔与允许路径类可按阶递归构造}
\label{subsec:tex37-ch4-axiom-downgrade}

\begin{theorem}[公理降级：修复塔纤维与允许路径类存在唯一的递归构造（数学归纳法证书）]
\label{thm:tex37-ch4-tower-induction}
固定一个上下文基点 $b\in\mathcal{B}_{\mathrm{ctx}}$。设满足下列条件：
\begin{enumerate}
  \item 存在基阶纤维 $\mathcal{E}_b^{(0)}$、基阶允许路径类集合 $\Gamma_b^{(0)}$，以及规范代表映射
  \[
    \Sigma_0:\Gamma_b^{\mathrm{adm},(0)}\to \mathsf{NF}_G^{(0)};
  \]
  \item 对任意 $n\in\mathbb{N}$ 与任意 $x\in \mathcal{E}_b^{(n)}$，存在修复扩张纤维
  \[
    \operatorname{Tw}_{\mathrm{rep},b}^{(n)}(x);
  \]
  \item 对任意 $n\in\mathbb{N}$ 与任意 $[\gamma]\in \Gamma_b^{(n)}$，存在允许提升纤维
  \[
    \operatorname{Lift}_{\mathrm{rep},b}^{(n)}([\gamma]);
  \]
  \item 对任意 $n\in\mathbb{N}$，存在固定的代表更新规则
  \[
    U_b^{(n)}:
    \mathsf{NF}_G^{(n)}
    \times
    \operatorname{Lift}_{\mathrm{rep},b}^{(n)}([\gamma])
    \to
    \mathsf{NF}_G^{(n+1)}.
  \]
\end{enumerate}
则存在唯一一族递归对象
\[
  \{\mathcal{E}_b^{(n)}\}_{n\in\mathbb{N}},
  \qquad
  \{\Gamma_b^{(n)}\}_{n\in\mathbb{N}},
  \qquad
  \{\Sigma_n\}_{n\in\mathbb{N}},
\]
满足
\[
  \mathcal{E}_b^{(n+1)}
  :=
  \bigsqcup_{x\in\mathcal{E}_b^{(n)}}
  \{x\}\times \operatorname{Tw}_{\mathrm{rep},b}^{(n)}(x),
\]
\[
  \Gamma_b^{(n+1)}
  :=
  \bigsqcup_{[\gamma]\in\Gamma_b^{(n)}}
  \{[\gamma]\}\times \operatorname{Lift}_{\mathrm{rep},b}^{(n)}([\gamma]),
\]
且对任意
\[
  ([\gamma],\ell)\in \Gamma_b^{(n+1)}
\]
有
\[
  \Sigma_{n+1}([\gamma],\ell)
  :=
  U_b^{(n)}\bigl(\Sigma_n([\gamma]),\ell\bigr).
\]
\end{theorem}

\begin{Certificate}[证书：基步 + 归纳步的修复塔构造证书]
\label{cert:tex37-ch4-tower-induction}
定义谓词
\[
  P(n):\ \text{“$\mathcal{E}_b^{(n)}$、$\Gamma_b^{(n)}$ 与 $\Sigma_n$ 已唯一构造完成，且彼此递归兼容”}.
\]
则数学归纳法证书可登记为
\[
  \pi_{\mathrm{rep}}
  =
  \bigl(
    P(0),\
    \forall n\in\mathbb{N},\ P(n)\Rightarrow P(n+1)
  \bigr).
\]
其中归纳步的显式递推式正是
\[
  \mathcal{E}_b^{(n+1)}
  =
  \bigsqcup_{x\in\mathcal{E}_b^{(n)}}
  \{x\}\times \operatorname{Tw}_{\mathrm{rep},b}^{(n)}(x),
\]
\[
  \Gamma_b^{(n+1)}
  =
  \bigsqcup_{[\gamma]\in\Gamma_b^{(n)}}
  \{[\gamma]\}\times \operatorname{Lift}_{\mathrm{rep},b}^{(n)}([\gamma]),
\]
\[
  \Sigma_{n+1}([\gamma],\ell)
  =
  U_b^{(n)}\bigl(\Sigma_n([\gamma]),\ell\bigr).
\]
\end{Certificate}

\begin{proof}[证明（数学归纳法证书；谓词因果链）]
\textbf{基步 $P(0)$：}
由条件(1)，$\mathcal{E}_b^{(0)}$、$\Gamma_b^{(0)}$ 与 $\Sigma_0$ 已给定，故 $P(0)$ 成立。

\textbf{归纳步：}
设 $P(n)$ 成立，即第 $n$ 阶纤维、路径类与规范代表映射已经唯一确定。
由条件(2)，对每个
\[
  x\in \mathcal{E}_b^{(n)}
\]
都有唯一给定的修复扩张纤维
\[
  \operatorname{Tw}_{\mathrm{rep},b}^{(n)}(x),
\]
因此不交并
\[
  \bigsqcup_{x\in\mathcal{E}_b^{(n)}}
  \{x\}\times \operatorname{Tw}_{\mathrm{rep},b}^{(n)}(x)
\]
唯一确定，从而 $\mathcal{E}_b^{(n+1)}$ 唯一确定。

同理，由条件(3)，对每个
\[
  [\gamma]\in \Gamma_b^{(n)}
\]
都有唯一给定的允许提升纤维
\[
  \operatorname{Lift}_{\mathrm{rep},b}^{(n)}([\gamma]),
\]
因此
\[
  \Gamma_b^{(n+1)}
  =
  \bigsqcup_{[\gamma]\in\Gamma_b^{(n)}}
  \{[\gamma]\}\times \operatorname{Lift}_{\mathrm{rep},b}^{(n)}([\gamma])
\]
也唯一确定。

再由条件(4)，一旦给定
\[
  ([\gamma],\ell)\in \Gamma_b^{(n+1)},
\]
则
\[
  \Sigma_{n+1}([\gamma],\ell)
  :=
  U_b^{(n)}\bigl(\Sigma_n([\gamma]),\ell\bigr)
\]
被唯一确定，故规范代表映射亦唯一确定。

因此 $P(n)\Rightarrow P(n+1)$ 成立。由数学归纳法，$\forall n\in\mathbb{N}, P(n)$。故修复塔纤维、允许路径类与规范代表映射存在且唯一。证毕。
\end{proof}

\subsection{定理：真正被求解的对象是修复塔上的允许截面/路径类，而不是表面文本}
\label{subsec:tex37-ch4-object}

\begin{theorem}[问答内核对象定理]
\label{thm:tex37-ch4-object}
设 $t\in \mathsf{Txt}_\xi$ 是某次问答对外显示的非失败文本。若
\[
  t=\mathcal{P}^{(n)}_{b,\xi}([\gamma])\neq\bot
\]
对某个上下文基点 $b$、某个阶数 $n$ 与某个可接受路径类
\[
  [\gamma]\in \Gamma_b^{\mathrm{adm},(n)}
\]
成立，则被求解的对象不是文本 $t$ 本身，而是该路径类
\[
  [\gamma]
\]
或与之等价确定的修复塔截面。
文本只是在审计约束下对该对象的表面投影壳。
\end{theorem}

\begin{Certificate}[证书：路径类 $\Rightarrow$ 规范代表 $\Rightarrow$ 占位 $\Rightarrow$ 文本]
\label{cert:tex37-ch4-object}
证书登记谓词链：
\[
\begin{aligned}
  O_1 &: [\gamma]\in \Gamma_b^{\mathrm{adm},(n)},\\
  O_2 &: n_\gamma=\Sigma_n([\gamma])\in \mathsf{NF}_G^{(n)},\\
  O_3 &: p_\gamma=\operatorname{pr}_{n\to 0}\!\bigl(\Theta_n(n_\gamma)\bigr)\in \mathsf{PL}_{\mathrm{NL}}^{(0)},\\
  O_4 &: t=\Render^{\mathrm{LLM}}_\xi(p_\gamma)\neq\bot,\\
  O_5 &: O_4\Rightarrow \Parse_\xi(t)=p_\gamma\ \text{（定理~\ref{thm:tex37-runtime-deterministic}）},\\
  O_6 &: O_2\wedge O_3\wedge O_5\Rightarrow t\ \text{不引入新语义对象，只表面表达既有路径类代表}.
\end{aligned}
\]
\end{Certificate}

\begin{proof}[证明（谓词因果链）]
由 $O_1$ 与定义~\ref{def:tex37-ch4-projection}，路径类 $[\gamma]$ 先经
\[
  \Sigma_n
\]
送到规范语义代表 $n_\gamma$，再经第一章的高阶同伦扭曲双射
\[
  \Theta_n
\]
送到逻辑占位塔，并截断为基阶占位 $p_\gamma$，即 $O_2$ 与 $O_3$。

由 $O_4$ 与第一章的 LLM 运行时降级定理，非失败文本 $t$ 必须满足
\[
  \Parse_\xi(t)=p_\gamma,
\]
即 $O_5$。
因此文本层并未制造一个新的求解对象；它只把既有的占位代表投影成外显文本。

再由第三章的对象收紧结论，规范语义塔与逻辑占位塔的对象边界必须与表面文本层分离。
于是求解链只能写成
\[
  [\gamma]
  \to
  n_\gamma
  \to
  p_\gamma
  \to
  t,
\]
而不能反过来把 $t$ 当作本体。
故真正被求解的对象是修复塔上的允许路径类/截面，文本只是其表面投影壳。证毕。
\end{proof}

\subsection{引理：比安基闭合一般是障碍修复的目标态，而不是上下文延拓的先验事实}
\label{subsec:tex37-ch4-bianchi}

\begin{lemma}[上下文延拓并不先验保持比安基闭合]
\label{lem:tex37-ch4-bianchi-not-prior}
设当前上下文基点为 $b$，问题 $q$ 诱导延拓
\[
  b^+=\operatorname{Ext}_q(b).
\]
若联络与曲率随延拓发生局部更新
\[
  A_{b^+}=A_b+\delta A_q,
  \qquad
  F_{b^+}=F_b+\delta F_q,
\]
并且雅可比障碍发生更新
\[
  J_{b^+}=J_b+\delta J_q,
\]
则一般不能仅由“发生了上下文延拓”推出
\[
  D_{A_{b^+}}F_{b^+}=0
  \qquad\text{或}\qquad
  J_{b^+}=0.
\]
更严格地说，除非额外假设 $\operatorname{Ext}_q$ 是硬约束闭合保持算子，否则比安基闭合与雅可比闭合都只能视为修复目标，而不是先验事实。
\end{lemma}

\begin{Certificate}[证书：延拓增量项阻断先验闭合]
\label{cert:tex37-ch4-bianchi-not-prior}
证书登记谓词链：
\[
\begin{aligned}
  B_1 &: A_{b^+}=A_b+\delta A_q,\quad F_{b^+}=F_b+\delta F_q,\quad J_{b^+}=J_b+\delta J_q,\\
  B_2 &: \mathcal{B}\!i_{b^+}=D_{A_{b^+}}F_{b^+},\\
  B_3 &: D_{A_{b^+}}F_{b^+}=D_{A_b}F_b + D_{A_b}\delta F_q + [\delta A_q,F_b] + [\delta A_q,\delta F_q],\\
  B_4 &: \delta A_q,\ \delta F_q,\ \delta J_q\ \text{若不被额外约束抵消，则可使}\ \mathcal{B}\!i_{b^+}\neq 0\ \text{或}\ J_{b^+}\neq 0,
    \\
  B_5 &: B_3\wedge B_4\Rightarrow \text{闭合性应作为修复目标，而非上下文延拓的先验结论}.
\end{aligned}
\]
\end{Certificate}

\begin{proof}[证明（谓词因果链）]
由 $B_1$，上下文延拓后的联络、曲率与雅可比障碍都带有新增量项。
代入 $B_2$ 得
\[
  \mathcal{B}\!i_{b^+}
  =
  D_{A_b+\delta A_q}(F_b+\delta F_q).
\]
按协变导数的线性展开，得到 $B_3$：
\[
  \mathcal{B}\!i_{b^+}
  =
  D_{A_b}F_b
  +
  D_{A_b}\delta F_q
  +
  [\delta A_q,F_b]
  +
  [\delta A_q,\delta F_q].
\]
即便旧上下文已满足
\[
  D_{A_b}F_b=0,
\]
新增的三项也未必自动相互抵消。
同理，
\[
  J_{b^+}=J_b+\delta J_q
\]
也表明只要新问题引入新的依赖拼接、证书责任或引用接口，雅可比残差就可能重新出现。

因此，单从
\[
  b^+=\operatorname{Ext}_q(b)
\]
并不能推出
\[
  \mathcal{B}\!i_{b^+}=0
  \qquad\text{或}\qquad
  J_{b^+}=0.
\]
除非进一步强加硬约束
\[
  D_{A_{\operatorname{Ext}_q(b)}}F_{\operatorname{Ext}_q(b)}=0,
  \qquad
  J_{\operatorname{Ext}_q(b)}=0
\]
对所有 $b$ 先验成立。故比安基闭合一般只是修复后的目标态。证毕。
\end{proof}

\subsection{定理：问答的解答应解释为上下文延拓后的障碍修复}
\label{subsec:tex37-ch4-answer}

\begin{theorem}[解答是延拓后的修复结果，而不是旧上下文上的直接读值]
\label{thm:tex37-ch4-answer-as-repair}
设当前上下文基点为 $b$，问题 $q$ 诱导扩张
\[
  b^+=\operatorname{Ext}_q(b).
\]
若某个模式 $\xi$ 下存在非失败回答
\[
  t=\mathcal{P}^{(n)}_{b^+,\xi}([\gamma^\star])\neq\bot
\]
且
\[
  [\gamma^\star]\in \Gamma_{b^+}^{\mathrm{adm},(n)},
\]
则必存在一个经修复达到可接受边界的终点状态
\[
  \tilde b:=\gamma^\star(1)\in \mathcal{B}_{\mathrm{ctx}}
\]
使得
\[
  \operatorname{Ob}_{\tilde b}
  \preceq_{\mathrm{lex}}
  \varepsilon_{b^+}.
\]
因此，问答系统中的“解答”应解释为：由上下文延拓触发、并在修复塔上完成的障碍修复结果。
\end{theorem}

\begin{Certificate}[证书：延拓 $\Rightarrow$ 可接受路径类 $\Rightarrow$ 修复后终点 $\Rightarrow$ 表面回答]
\label{cert:tex37-ch4-answer-as-repair}
证书登记谓词链：
\[
\begin{aligned}
  A_1 &: b^+=\operatorname{Ext}_q(b),\\
  A_2 &: [\gamma^\star]\in \Gamma_{b^+}^{\mathrm{adm},(n)},\\
  A_3 &: A_2\Rightarrow \operatorname{Ob}_{\gamma^\star(1)}\preceq_{\mathrm{lex}}\varepsilon_{b^+},\\
  A_4 &: t=\mathcal{P}^{(n)}_{b^+,\xi}([\gamma^\star])\neq\bot,\\
  A_5 &: A_2\wedge A_4\Rightarrow t\ \text{是修复后路径类的表面投影（定理~\ref{thm:tex37-ch4-object}）},\\
  A_6 &: A_1\wedge A_3\wedge A_5\Rightarrow \text{解答属于延拓后的障碍修复结果}.
\end{aligned}
\]
\end{Certificate}

\begin{proof}[证明（谓词因果链）]
由 $A_1$，新问题并不是在旧基点 $b$ 上取值，而是先把上下文扩张到 $b^+$。
由定义~\ref{def:tex37-ch4-action}，只有当路径类属于
\[
  \Gamma_{b^+}^{\mathrm{adm},(n)}
\]
时，才可被纳入可接受解答域。
因此由 $A_2$ 立得 $A_3$：
\[
  \operatorname{Ob}_{\gamma^\star(1)}
  \preceq_{\mathrm{lex}}
  \varepsilon_{b^+}.
\]
把
\[
  \tilde b:=\gamma^\star(1)
\]
记为修复后的终点状态，则 $\tilde b$ 已位于可接受边界内。

再由 $A_4$ 与定理~\ref{thm:tex37-ch4-object}，文本 $t$ 只是该路径类的表面投影，而不是独立解对象。
于是，新问题的回答链只能写成
\[
  b
  \xrightarrow{\operatorname{Ext}_q}
  b^+
  \xrightarrow{\text{修复塔路径类 }[\gamma^\star]}
  \tilde b
  \xrightarrow{\mathcal{P}^{(n)}_{b^+,\xi}}
  t.
\]
故问答中的“解答”本质上就是上下文延拓后的障碍修复结果，而非对旧上下文状态的直接读取。证毕。
\end{proof}

\subsection{命题：路径积分是选择与收敛准则，而不是对话对象本身}
\label{subsec:tex37-ch4-proposition}

\begin{proposition}[路径积分的角色是对路径类加权并选择代表]
\label{prop:tex37-ch4-path-integral-selector}
在定义~\ref{def:tex37-ch4-action} 的设定下，作用泛函
\[
  \mathcal{S}_b:\Gamma_b^{\mathrm{adm}}\to \mathbb{R}_{\ge 0}\cup\{+\infty\}
\]
只在可接受路径类集合上定义一个“可行性优先、最优性次之”的排序准则，
而不改变路径类本身的对象类型。
因此，更严格的表述不是“对话的是路径积分本身”，而是：
\[
  \boxed{
    \text{对话内核是修复塔上的允许路径类；路径积分是对这些路径类加权并选择代表的机制。}
  }
\]
若某个路径类
\[
  [\gamma_b^\star]\in \Gamma_b^{\mathrm{adm}}
\]
满足
\[
  \mathcal{S}_b([\gamma_b^\star])
  \le
  \inf_{[\gamma]\in \Gamma_b^{\mathrm{adm}}}\mathcal{S}_b([\gamma])
  +
  \eta_b
  \qquad
  (\eta_b\ge 0),
\]
则称其为基点 $b$ 上的一个可行且 $\eta_b$-近优的路径积分代表。
\end{proposition}

\begin{Certificate}[证书：对象域与排序域分离]
\label{cert:tex37-ch4-path-integral-selector}
证书登记谓词链：
\[
\begin{aligned}
  P_1 &: \Gamma_b^{\mathrm{adm}}\ \text{是路径类对象域},\\
  P_2 &: \mathcal{S}_b:\Gamma_b^{\mathrm{adm}}\to \mathbb{R}_{\ge 0}\cup\{+\infty\}\ \text{是数值型作用泛函},\\
  P_3 &: P_2\Rightarrow \mathcal{S}_b\ \text{只定义排序/加权，不定义新的对象},\\
  P_4 &: [\gamma_b^\star]\ \text{近优} \Longleftrightarrow \mathcal{S}_b([\gamma_b^\star])\le \inf \mathcal{S}_b + \eta_b,
    \\
  P_5 &: P_1\wedge P_3\wedge P_4\Rightarrow \text{路径积分是选择机制，而非对话对象本身}.
\end{aligned}
\]
\end{Certificate}

\begin{proof}[证明（谓词因果链）]
由 $P_1$，本章的对象域是修复塔上的可接受路径类；
由 $P_2$，作用泛函的值域是非负实数或 $+\infty$，故其类型是
\[
  \Gamma_b^{\mathrm{adm}}
  \longrightarrow
  \mathbb{R}_{\ge 0}\cup\{+\infty\},
\]
而不是
\[
  \Gamma_b^{\mathrm{adm}}
  \longrightarrow
  \Gamma_b^{\mathrm{adm}}.
\]
因此 $\mathcal{S}_b$ 只能充当排序、筛选与收敛准则，不能把“对象本身”替换掉，即 $P_3$。

再由 $P_4$，所谓“解答近优”只是说某个路径类的作用量接近可接受域上的下确界。
这仍然是在对象域
\[
  \Gamma_b^{\mathrm{adm}}
\]
上选代表，而不是把作用量函数当作被对话对象。
故命题成立。证毕。
\end{proof}

\subsection{定理：路径积分代表在条件下呈现由可行性到最优性的双性收敛}
\label{subsec:tex37-ch4-convergence}

\begin{theorem}[可行性到最优性的双性收敛定理]
\label{thm:tex37-ch4-dual-convergence}
固定上下文基点 $b$。设修复序列
\[
  [\gamma_0],[\gamma_1],[\gamma_2],\dots \subseteq \Gamma_b
\]
满足下列条件：
\begin{enumerate}
  \item[\textnormal{(H1)}] 可行性作用量单调不增，即
  \[
    \mathcal{S}_{\mathrm{feas},b}([\gamma_{k+1}])
    \le
    \mathcal{S}_{\mathrm{feas},b}([\gamma_k]),
  \]
  且当
  \[
    [\gamma_k]\notin \Gamma_b^{\mathrm{adm}}
  \]
  时严格下降；
  \item[\textnormal{(H2)}] 可行域非空，且所选修复序列具有可行完备性：存在 $K\in\mathbb{N}$ 使得
  \[
    \forall k\ge K,\quad [\gamma_k]\in \Gamma_b^{\mathrm{adm}};
  \]
  \item[\textnormal{(H3)}] 总作用量
  \[
    \mathcal{S}_b
  \]
  在 $\Gamma_b^{\mathrm{adm}}$ 上存在极小代表，或至少存在极小化序列；
  \item[\textnormal{(H4)}] 第一章的高阶同伦扭曲双射与本章的投影链保持可行性判定与作用等价类，即若
  \[
    \Theta_n(\Sigma_n([\gamma_1]))=\Theta_n(\Sigma_n([\gamma_2])),
  \]
  则
  \[
    [\gamma_1]\in\Gamma_b^{\mathrm{adm}}
    \iff
    [\gamma_2]\in\Gamma_b^{\mathrm{adm}},
    \qquad
    \mathcal{S}_b([\gamma_1])=\mathcal{S}_b([\gamma_2]).
  \]
\end{enumerate}
则该修复序列呈现出“由可行性到最优性”的双性收敛：
\[
  \boxed{
    \mathcal{S}_{\mathrm{feas},b}([\gamma_k])\to s_b^\star\ge 0,
  }
\]
并且在理想闭合情形下
\[
  s_b^\star=0;
\]
同时在进入可行域后的尾序列上，有
\[
  \boxed{
    \mathcal{S}_b([\gamma_k])
    \to
    \inf_{[\gamma]\in \Gamma_b^{\mathrm{adm}}}\mathcal{S}_b([\gamma])
  }
\]
或至少存在达到该下确界的极小化子序列。
因此，
\[
  \boxed{
    \text{路径积分收敛 = 可行性收敛 + 最优性收敛，且该双性收敛在同伦扭曲双射下保持。}
  }
\]
\end{theorem}

\begin{Certificate}[证书：单调可行性势 + 可行完备性 + 极小化序列]
\label{cert:tex37-ch4-dual-convergence}
证书登记谓词链：
\[
\begin{aligned}
  D_1 &: \mathcal{S}_{\mathrm{feas},b}([\gamma_{k+1}])\le \mathcal{S}_{\mathrm{feas},b}([\gamma_k])\ \text{且下界为 }0,\\
  D_2 &: D_1\Rightarrow \mathcal{S}_{\mathrm{feas},b}([\gamma_k])\ \text{收敛到某个 }s_b^\star\ge 0,\\
  D_3 &: \exists K\ \forall k\ge K,\ [\gamma_k]\in \Gamma_b^{\mathrm{adm}},\\
  D_4 &: D_3\wedge \textnormal{(H3)}\Rightarrow \text{尾序列可在可行域内做极小化},\\
  D_5 &: \textnormal{(H4)}\Rightarrow \text{同伦扭曲双射不改变可行性与作用等价类},\\
  D_6 &: D_2\wedge D_4\wedge D_5\Rightarrow \text{双性收敛成立}.
\end{aligned}
\]
\end{Certificate}

\begin{proof}[证明（谓词因果链）]
先证可行性收敛。由 \textnormal{(H1)}，
\[
  \mathcal{S}_{\mathrm{feas},b}([\gamma_k])
\]
是一个单调不增的非负实数序列。
由实数完备性，其极限
\[
  s_b^\star:=\lim_{k\to\infty}\mathcal{S}_{\mathrm{feas},b}([\gamma_k])
\]
存在，且 $s_b^\star\ge 0$，即得 $D_1\Rightarrow D_2$。
若修复最终实现理想闭合，则
\[
  |\mathcal{B}\!i|=|J|=|\Delta_{\mathrm{cert}}|=|\Delta_{\mathrm{ref}}|=0
\]
沿尾序列成立，从而
\[
  s_b^\star=0.
\]

再证最优性收敛。由 \textnormal{(H2)}，存在 $K$ 使得所有尾项
\[
  [\gamma_k],\quad k\ge K
\]
都已进入可行域 $\Gamma_b^{\mathrm{adm}}$。
因此从第 $K$ 步起，优化问题被限制在可行域内部。
由 \textnormal{(H3)}，要么存在极小代表
\[
  [\gamma^\star]\in \Gamma_b^{\mathrm{adm}}
\]
满足
\[
  \mathcal{S}_b([\gamma^\star])
  =
  \inf_{[\gamma]\in\Gamma_b^{\mathrm{adm}}}\mathcal{S}_b([\gamma]),
\]
要么存在极小化序列逼近该下确界。
因此可在尾序列上得到
\[
  \mathcal{S}_b([\gamma_k])
  \to
  \inf_{[\gamma]\in\Gamma_b^{\mathrm{adm}}}\mathcal{S}_b([\gamma])
\]
或至少提取一个满足该性质的子序列，即得 $D_3\wedge D_4$。

最后，由 \textnormal{(H4)}，同伦扭曲双射与投影链不会改变“哪条路径类可行、哪条路径类具有给定作用量”等价类。
故前述可行性收敛与最优性收敛，在规范语义塔与逻辑占位塔之间保持不变，即 $D_5$。
三者合并即得 $D_6$。故定理成立。证毕。
\end{proof}

\subsection{推论：LLM 只把收敛后的代表投影为文本，而不参与收敛本身}
\label{subsec:tex37-ch4-corollary}

\begin{corollary}[收敛发生在修复塔/规范语义/占位层；LLM 只表达收敛结果]
\label{cor:tex37-ch4-llm-projection}
在定理~\ref{thm:tex37-ch4-answer-as-repair} 与定理~\ref{thm:tex37-ch4-dual-convergence} 的条件下，若
\[
  [\gamma_b^\star]\in \Gamma_b^{\mathrm{adm}}
\]
是一个可行且近优的路径积分代表，且
\[
  t=\mathcal{P}^{(n)}_{b,\xi}([\gamma_b^\star])\neq\bot,
\]
则
\[
  \Parse_\xi(t)
  =
  \operatorname{pr}_{n\to 0}
  \Bigl(
    \Theta_n\bigl(\Sigma_n([\gamma_b^\star])\bigr)
  \Bigr).
\]
因此，LLM 不参与“收敛到哪个代表”这件事；它只在审计合同下把已经被选出的代表投影成文本。
\end{corollary}

\begin{Certificate}[证书：收敛先于渲染]
\label{cert:tex37-ch4-llm-projection}
证书登记谓词链：
\[
\begin{aligned}
  C_1 &: [\gamma_b^\star]\in \Gamma_b^{\mathrm{adm}}\ \text{且是可行近优代表（命题~\ref{prop:tex37-ch4-path-integral-selector}）},\\
  C_2 &: [\gamma_b^\star]\ \text{由修复与双性收敛选出（定理~\ref{thm:tex37-ch4-answer-as-repair},\ref{thm:tex37-ch4-dual-convergence}）
    },\\
  C_3 &: t=\mathcal{P}^{(n)}_{b,\xi}([\gamma_b^\star])\neq\bot,\\
  C_4 &: C_3\Rightarrow \Parse_\xi(t)=\operatorname{pr}_{n\to 0}\!\bigl(\Theta_n(\Sigma_n([\gamma_b^\star]))\bigr)\
    \text{（定理~\ref{thm:tex37-runtime-deterministic}）},\\
  C_5 &: C_1\wedge C_2\wedge C_4\Rightarrow \text{LLM 只表达收敛结果，不制造收敛}.
\end{aligned}
\]
\end{Certificate}

\begin{proof}[证明（谓词因果链）]
由 $C_1$ 与命题~\ref{prop:tex37-ch4-path-integral-selector}，路径积分只是在可行域内选择一个近优代表；
由 $C_2$，这个代表之所以成立，是因为前面已经完成了上下文延拓后的障碍修复，并在可行域内完成了作用量收敛。
因此，在进入渲染阶段之前，内核已经确定了
\[
  [\gamma_b^\star]
  \quad\text{及其规范语义代表}\quad
  \Sigma_n([\gamma_b^\star]).
\]

再由 $C_3$ 与定义~\ref{def:tex37-ch4-projection}，
\[
  t
  =
  \Render^{\mathrm{LLM}}_\xi
  \circ
  \operatorname{pr}_{n\to 0}
  \circ
  \Theta_n
  \circ
  \Sigma_n
  ([\gamma_b^\star]).
\]
由第一章的 LLM 运行时降级定理，非失败输出满足
\[
  \Parse_\xi(t)
  =
  \operatorname{pr}_{n\to 0}
  \Bigl(
    \Theta_n\bigl(\Sigma_n([\gamma_b^\star])\bigr)
  \Bigr),
\]
即 $C_4$。
这表明文本只是既有代表的表面表达，而不是收敛机制本身。
故推论成立。证毕。
\end{proof}

\subsection{备注：严格收紧版表述与一句话结论}
\label{subsec:tex37-ch4-remarks}

\begin{remark}[建议采用的严格总表述]
本章更适合采用下面这组收紧后的表述：
\[
  \boxed{
    \begin{aligned}
      &\text{在该构造下，问答系统的内核求解对象不是表面文本，}\\
      &\text{而是以上下文为基底的 G 框架同伦修复塔上的允许截面/路径类。}
    \end{aligned}
  }
\]
\[
  \boxed{
    \begin{aligned}
      &\text{一个新问题对应于上下文基底的延拓；}\\
      &\text{该延拓一般会诱导联络、曲率、证书依赖与语义括号结构的局部失衡，}
    \end{aligned}
  }
\]
\[
  \boxed{
    \text{从而表现为比安基缺陷、雅可比障碍以及证书/引用接口残差。}
  }
\]
\[
  \boxed{
    \begin{aligned}
      &\text{解答过程就是在同伦修复塔上，通过障碍驱动的修复算子，}\\
      &\text{寻找使障碍收敛到可接受边界、并在可行域内使作用泛函近极小的路径积分代表。}
    \end{aligned}
  }
\]
\[
  \boxed{
    \text{最终 LLM 只是在审计约束下，把该代表的逻辑占位投影成某种风格模板下的表面表达。}
  }
\]
\end{remark}

\begin{remark}[一句话结论]
最短的严格压缩可写成：
\[
  \boxed{
    \text{问答的“解”是修复塔路径类的可行--近优代表，文本只是这个代表的可审计投影壳。}
  }
\]
若需要把“比安基恒等式成立”写进摘要，则更稳妥的版本应是：
\[
  \boxed{
    \text{上下文延拓诱导局部联络更新；修复过程以恢复比安基闭合为目标。}
  }
\]
而不宜直接写成“上下文延拓就使比安基恒等式成立”。
\end{remark}

\clearpage

%%%%%%%%%%%%%%%%%%%%%%%%%%%%%%%%%%%%%%%%%%%%%%%%%%%%%%%%%%%%
\section{形式逻辑：纠偏后显著增益仍来自求解任务与表达任务的严格分离}
\label{sec:tex37-ch5}
%%%%%%%%%%%%%%%%%%%%%%%%%%%%%%%%%%%%%%%%%%%%%%%%%%%%%%%%%%%%

\begin{remark}[核心陈述（收紧版）]
即便按前一章把对象层次继续收紧，本章仍主张下面这条增益链成立：
\[
  \boxed{
    \begin{aligned}
      &\text{G 框架塔纤维“对话”的结果不是文本，}\\
      &\text{而是障碍修复后的算子语义序列或逻辑占位序列。}
    \end{aligned}
  }
\]
\[
  \boxed{
    \begin{aligned}
      &\text{一旦该序列已经在 G 侧被固定，}\\
      &\text{LLM 侧的任务就从“求解”降级为“受合同约束的投影表达与可读化展开”。}
    \end{aligned}
  }
\]
\[
  \boxed{
    \begin{aligned}
      &\text{因此，文本的充分性、可读性、严谨性与风格化，}\\
      &\text{不再依赖模型是否会证明，而依赖占位完整性、合同严格性与反复审计。}
    \end{aligned}
  }
\]
\end{remark}

\begin{remark}[阅读导航：仓内文稿与本稿前文接口]
本章直接承接第一章关于 LLM 运行时降级、解析回投影与同核多外观的结论
\cite{RepoTex21AuditedWrapper,RepoTex23SemanticBase,RepoTex20SemanticGaugeNormalization}，
并把第二章“工程增益主导”定理~\ref{thm:tex37-engineering-dominates} 与第四章“内核对象是修复塔路径类”
定理~\ref{thm:tex37-ch4-object}、推论~\ref{cor:tex37-ch4-llm-projection} 接起来，专门讨论：
\[
  \text{求解任务} \neq \text{表达任务}
\]
以后，为什么增益不仅没有变弱，反而更容易落成工程闭环。
\end{remark}

\subsection{定义：求解任务、表达任务、槽位覆盖与求解侧负载}
\label{subsec:tex37-ch5-definitions}

\begin{definition}[内核求解任务与表达投影任务]
\label{def:tex37-ch5-task-separation}
记一切待回答问题构成集合
\[
  \mathcal{Q}.
\]
给定上下文基点 $b\in\mathcal{B}_{\mathrm{ctx}}$ 与问题 $q\in\mathcal{Q}$，设
\[
  [\gamma^\star_{b,q}]
  \in
  \Gamma_{\operatorname{Ext}_q(b)}^{\mathrm{adm},(n)}
\]
是第四章意义下的一个可行且近优的修复塔路径积分代表。
定义\textbf{内核求解任务}
\[
  \mathcal{T}_{\mathrm{solve}}:
  \mathcal{B}_{\mathrm{ctx}}\times \mathcal{Q}
  \longrightarrow
  \mathsf{PL}_{\mathrm{NL}}^{(0)},
\]
\[
  \mathcal{T}_{\mathrm{solve}}(b,q)
  :=
  \operatorname{pr}_{n\to 0}
  \Bigl(
    \Theta_n\bigl(\Sigma_n([\gamma^\star_{b,q}])\bigr)
  \Bigr)
  =:p^\star.
\]
它把“上下文延拓后的障碍态”送到“已固定的逻辑占位代表”。

再定义\textbf{表达投影任务}
\[
  \mathcal{T}_{\mathrm{render}}:
  \mathsf{PL}_{\mathrm{NL}}^{(0)}\times \Xi
  \longrightarrow
  \mathsf{Txt}_\xi\sqcup\{\bot\},
\]
\[
  \mathcal{T}_{\mathrm{render}}(p^\star,\xi)
  :=
  \Render^{\mathrm{LLM}}_\xi(p^\star).
\]
因此
\[
  \operatorname{Dom}(\mathcal{T}_{\mathrm{solve}})
  \neq
  \operatorname{Dom}(\mathcal{T}_{\mathrm{render}})
  \qquad\text{且}\qquad
  \operatorname{Cod}(\mathcal{T}_{\mathrm{solve}})
  \neq
  \operatorname{Cod}(\mathcal{T}_{\mathrm{render}}),
\]
从而
\[
  \boxed{
    \mathcal{T}_{\mathrm{solve}}\neq \mathcal{T}_{\mathrm{render}}.
  }
\]
\end{definition}

\begin{definition}[占位槽位、覆盖映射与全审计接受集]
\label{def:tex37-ch5-coverage}
设一个已固定占位可写成有限槽位族
\[
  p^\star=(s_1,\dots,s_m),
  \qquad
  \operatorname{Slot}(p^\star):=\{s_1,\dots,s_m\}.
\]
对每个模式 $\xi\in\Xi$，定义：
\begin{enumerate}
  \item 覆盖映射
  \[
    \operatorname{Cov}_\xi:
    \mathsf{Txt}_\xi\to 2^{\operatorname{Slot}(p^\star)},
  \]
  其值表示文本实际显式覆盖到的占位槽位；
  \item 可读性谓词
  \[
    \operatorname{Read}_\xi:\mathsf{Txt}_\xi\to\{\top,\bot\};
  \]
  \item 全审计接受集
  \[
    \Audit^{\mathrm{full}}_\xi(p^\star)
    :=
    \left\{
      t\in \mathcal{C}^{\mathrm{LLM}}_\xi(p^\star)
      \ \middle|\
      \begin{aligned}
        &\Parse_\xi(t)=p^\star,\\
        &\Style_\xi(t)=\top,\\
        &\operatorname{Read}_\xi(t)=\top,\\
        &\operatorname{Cov}_\xi(t)=\operatorname{Slot}(p^\star)
      \end{aligned}
    \right\}.
  \]
\end{enumerate}
若
\[
  t\in\Audit^{\mathrm{full}}_\xi(p^\star),
\]
则称 $t$ 是对占位 $p^\star$ 的一个\textbf{充分、可读、严谨且风格合格}的文本代表。
\end{definition}

\begin{definition}[求解侧义务集与模型依赖重心迁移量]
\label{def:tex37-ch5-load}
定义求解侧义务集
\[
  \mathfrak{S}_{\mathrm{solve}}
  :=
  \{
    u_{\mathrm{ctx}},
    u_{\mathrm{repair}},
    u_{\mathrm{path}},
    u_{\mathrm{cert}},
    u_{\mathrm{verdict}}
  \},
\]
分别表示开放上下文判定、障碍修复、路径选择、证书拼接与最终裁决。
对每个义务赋予正权重
\[
  \omega:\mathfrak{S}_{\mathrm{solve}}\to \mathbb{R}_{>0}.
\]

若一个系统架构 $\mathfrak{A}$ 把模型实际承担的义务子集记为
\[
  \mathfrak{M}(\mathfrak{A})\subseteq \mathfrak{S}_{\mathrm{solve}},
\]
则定义模型的求解侧负载为
\[
  \operatorname{Load}_{\mathrm{solve}}(\mathfrak{A})
  :=
  \sum_{u\in \mathfrak{M}(\mathfrak{A})}\omega(u).
\]
对“分层前”与“分层后”两种架构 $\mathfrak{A}_{\mathrm{before}}$、$\mathfrak{A}_{\mathrm{after}}$，
定义模型依赖重心迁移量
\[
  \Delta_{\mathrm{dep}}
  :=
  \operatorname{Load}_{\mathrm{solve}}(\mathfrak{A}_{\mathrm{before}})
  -
  \operatorname{Load}_{\mathrm{solve}}(\mathfrak{A}_{\mathrm{after}}).
\]
\end{definition}

\begin{remark}[本章的最小口径]
本章不把“写得更像人”视为核心增益，而把下列四件事视为核心：
\[
  \text{求解归内核}
  +
  \text{表达归运行时}
  +
  \text{覆盖可审计}
  +
  \text{最终裁决外置}.
\]
\end{remark}

\subsection{公理（降级为定理）：已固定占位的覆盖式表达可按槽位数递归构造}
\label{subsec:tex37-ch5-axiom-downgrade}

\begin{theorem}[公理降级：固定占位的受约束表达可按槽位数递归分解（数学归纳法证书）]
\label{thm:tex37-ch5-slot-induction}
设
\[
  p^\star=(s_1,\dots,s_m),
  \qquad
  p^{\star,[k]}:=(s_1,\dots,s_k)\quad (1\le k\le m).
\]
固定某个模式 $\xi\in\Xi$。若满足：
\begin{enumerate}
  \item[\textnormal{(B0)}] 存在基步文本
  \[
    t_\xi^{[1]}\in \mathsf{Txt}_\xi
  \]
  使得
  \[
    \Parse_\xi(t_\xi^{[1]})=p^{\star,[1]},
    \quad
    \Style_\xi(t_\xi^{[1]})=\top,
    \quad
    \operatorname{Read}_\xi(t_\xi^{[1]})=\top,
    \quad
    \operatorname{Cov}_\xi(t_\xi^{[1]})=\{s_1\};
  \]
  \item[\textnormal{(Bk)}] 对任意 $1\le k<m$，存在槽位扩张算子
  \[
    \operatorname{ExtRender}^{(k+1)}_\xi:
    \mathsf{Txt}_\xi^{[k]}
    \to
    \mathsf{Txt}_\xi^{[k+1]}
  \]
  满足：若 $t_\xi^{[k]}$ 已覆盖前 $k$ 个槽位，则
  \[
    t_\xi^{[k+1]}
    :=
    \operatorname{ExtRender}^{(k+1)}_\xi(t_\xi^{[k]})
  \]
  同时满足
  \[
    \Parse_\xi(t_\xi^{[k+1]})=p^{\star,[k+1]},
    \qquad
    \Style_\xi(t_\xi^{[k+1]})=\top,
  \]
  \[
    \operatorname{Read}_\xi(t_\xi^{[k+1]})=\top,
    \qquad
    \operatorname{Cov}_\xi(t_\xi^{[k+1]})=\{s_1,\dots,s_{k+1}\}.
  \]
\end{enumerate}
则存在唯一一条受约束表达递归链
\[
  t_\xi^{[1]}
  \xrightarrow{\operatorname{ExtRender}^{(2)}_\xi}
  t_\xi^{[2]}
  \xrightarrow{\operatorname{ExtRender}^{(3)}_\xi}
  \cdots
  \xrightarrow{\operatorname{ExtRender}^{(m)}_\xi}
  t_\xi^{[m]},
\]
其中
\[
  t_\xi^{[m]}\in \Audit^{\mathrm{full}}_\xi(p^\star).
\]
因此，一旦占位 $p^\star$ 已固定，表达任务便可按槽位扩张与审计检查递归完成，而不再需要重新执行求解任务。
\end{theorem}

\begin{Certificate}[证书：基步 + 槽位扩张归纳步]
\label{cert:tex37-ch5-slot-induction}
令谓词
\[
  P(k):\ \text{“存在唯一文本 }t_\xi^{[k]}\text{，使其对 }p^{\star,[k]}\text{ 通过解析、风格、可读与覆盖检查”}.
\]
则数学归纳法证书可写成
\[
  \pi_{\mathrm{slot}}
  =
  \bigl(
    P(1),\
    \forall k<m,\ P(k)\Rightarrow P(k+1)
  \bigr).
\]
其中归纳步由扩张算子
\[
  \operatorname{ExtRender}^{(k+1)}_\xi
\]
显式给出。
\end{Certificate}

\begin{proof}[证明（数学归纳法证书；谓词因果链）]
\textbf{基步 $P(1)$：}
由条件 \textnormal{(B0)}，文本 $t_\xi^{[1]}$ 已满足
\[
  \Parse_\xi(t_\xi^{[1]})=p^{\star,[1]},
  \qquad
  \operatorname{Cov}_\xi(t_\xi^{[1]})=\{s_1\},
\]
并通过风格与可读检查，因此 $P(1)$ 成立。

\textbf{归纳步：}
设 $P(k)$ 成立，即已存在唯一文本 $t_\xi^{[k]}$ 覆盖前 $k$ 个槽位且可解析回
\[
  p^{\star,[k]}.
\]
由条件 \textnormal{(Bk)}，施加槽位扩张算子
\[
  \operatorname{ExtRender}^{(k+1)}_\xi
\]
后得到
\[
  t_\xi^{[k+1]}
  :=
  \operatorname{ExtRender}^{(k+1)}_\xi(t_\xi^{[k]}),
\]
并且它满足
\[
  \Parse_\xi(t_\xi^{[k+1]})=p^{\star,[k+1]},
  \qquad
  \operatorname{Cov}_\xi(t_\xi^{[k+1]})=\{s_1,\dots,s_{k+1}\},
\]
同时继续通过风格与可读检查。
由于扩张规则对每个 $k$ 已固定，故 $t_\xi^{[k+1]}$ 的构造是唯一的。
因此 $P(k)\Rightarrow P(k+1)$。

由数学归纳法，对所有 $1\le k\le m$，$P(k)$ 成立。
特别地，
\[
  t_\xi^{[m]}
\]
满足
\[
  \Parse_\xi(t_\xi^{[m]})=p^\star,
  \qquad
  \operatorname{Cov}_\xi(t_\xi^{[m]})=\operatorname{Slot}(p^\star),
\]
且风格、可读性均通过，因此
\[
  t_\xi^{[m]}\in \Audit^{\mathrm{full}}_\xi(p^\star).
\]
故定理成立。证毕。
\end{proof}

\subsection{定理：纠偏后，第一章的核心增益仍然显著}
\label{subsec:tex37-ch5-gain}

\begin{theorem}[显著增益定理]
\label{thm:tex37-ch5-significant-gain}
定义五个工程增益谓词：
\[
\begin{aligned}
  \mathsf{G}_1 &: \mathcal{T}_{\mathrm{solve}}\neq \mathcal{T}_{\mathrm{render}},\\
  \mathsf{G}_2 &: t=\Render^{\mathrm{LLM}}_\xi(p)\neq\bot\Rightarrow \Parse_\xi(t)=p,\\
  \mathsf{G}_3 &: \Audit_\xi(p)=\varnothing\Rightarrow \Render^{\mathrm{LLM}}_\xi(p)=\bot,\\
  \mathsf{G}_4 &: \text{同一 }p\text{ 可对应多风格文本且不改语义核},\\
  \mathsf{G}_5 &: \text{LLM 只投影收敛结果，不承担修复塔求解本身}.
\end{aligned}
\]
再定义分离增益指数
\[
  \Gamma_{\mathrm{sep}}
  :=
  \sum_{i=1}^{5}\mathbf{1}[\mathsf{G}_i].
\]
则在本稿当前已经建立的结构内，
\[
  \Gamma_{\mathrm{sep}}=5.
\]
因此，即便经过对象层次纠偏，第一章及其后续展开所带来的核心增益仍然是显著的，
而且其重心是工程职责边界的重写，而不是“让模型更会推理”。
\end{theorem}

\begin{Certificate}[证书：五个增益谓词全部成立]
\label{cert:tex37-ch5-significant-gain}
证书登记谓词链：
\[
\begin{aligned}
  G_1 &: \mathsf{G}_1\ \text{由定义~\ref{def:tex37-ch5-task-separation} 成立},\\
  G_2 &: \mathsf{G}_2\ \text{由定理~\ref{thm:tex37-runtime-deterministic} 成立},\\
  G_3 &: \mathsf{G}_3\ \text{由 }\Render^{\mathrm{LLM}}_\xi\text{ 的定义直接成立},\\
  G_4 &: \mathsf{G}_4\ \text{由命题~\ref{prop:tex37-style-invariance} 与定理~\ref{thm:tex37-same-core-many-surfaces} 成立},\\
  G_5 &: \mathsf{G}_5\ \text{由定理~\ref{thm:tex37-ch4-object} 与推论~\ref{cor:tex37-ch4-llm-projection} 成立},\\
  G_6 &: G_1\wedge G_2\wedge G_3\wedge G_4\wedge G_5\Rightarrow \Gamma_{\mathrm{sep}}=5.
\end{aligned}
\]
\end{Certificate}

\begin{proof}[证明（谓词因果链）]
由定义~\ref{def:tex37-ch5-task-separation}，
\[
  \operatorname{Dom}(\mathcal{T}_{\mathrm{solve}})
  \neq
  \operatorname{Dom}(\mathcal{T}_{\mathrm{render}})
\]
且值域不同，因此 $\mathsf{G}_1$ 成立。
由定理~\ref{thm:tex37-runtime-deterministic}，每个非失败输出都满足解析回投影与风格通过，因此 $\mathsf{G}_2$ 成立。
由 $\Render^{\mathrm{LLM}}_\xi$ 的定义，若无可审计接受元则输出 $\bot$，故 $\mathsf{G}_3$ 成立。

再由命题~\ref{prop:tex37-style-invariance} 与定理~\ref{thm:tex37-same-core-many-surfaces}，
同一占位可以拥有多个模式下的非失败文本代表，而这些代表都解析回同一语义核，故 $\mathsf{G}_4$ 成立。
由第四章定理~\ref{thm:tex37-ch4-object} 与推论~\ref{cor:tex37-ch4-llm-projection}，
问答内核对象是修复塔路径类，LLM 只负责把已收敛的代表投影成文本，故 $\mathsf{G}_5$ 成立。

于是五个增益谓词全部为真，故
\[
  \Gamma_{\mathrm{sep}}
  =
  \sum_{i=1}^{5}\mathbf{1}[\mathsf{G}_i]
  =5.
\]
这说明对象层次的纠偏并没有削弱前文增益，反而把增益精确定位为“职责分离后的工程可控性”。证毕。
\end{proof}

\subsection{引理：占位一旦固定，文本不再引入新的求解义务}
\label{subsec:tex37-ch5-lemma}

\begin{lemma}[固定占位后的表达阶段不新增求解义务]
\label{lem:tex37-ch5-no-new-solve}
对任意已固定占位 $p^\star$，定义其求解义务集
\[
  \operatorname{Obl}(p^\star)
\]
为其中编码的变量绑定、证书槽、引用接口与逻辑依赖的全集。
对任意非失败文本 $t\in\mathsf{Txt}_\xi$，定义
\[
  \operatorname{Obl}(t)
  :=
  \operatorname{Obl}(\Parse_\xi(t))
\]
只要 $\Parse_\xi(t)\neq\bot$。
若
\[
  t\in \Audit^{\mathrm{full}}_\xi(p^\star),
\]
则
\[
  \operatorname{Obl}(t)=\operatorname{Obl}(p^\star).
\]
因此，表达阶段只是在既有义务集上做外显组织，而不会再新增求解型义务。
\end{lemma}

\begin{Certificate}[证书：解析回投影 + 槽位全覆盖 $\Rightarrow$ 义务集不变]
\label{cert:tex37-ch5-no-new-solve}
证书登记谓词链：
\[
\begin{aligned}
  L_1 &: t\in \Audit^{\mathrm{full}}_\xi(p^\star),\\
  L_2 &: L_1\Rightarrow \Parse_\xi(t)=p^\star,\\
  L_3 &: L_1\Rightarrow \operatorname{Cov}_\xi(t)=\operatorname{Slot}(p^\star),\\
  L_4 &: L_2\Rightarrow \operatorname{Obl}(t)=\operatorname{Obl}(\Parse_\xi(t))=\operatorname{Obl}(p^\star),\\
  L_5 &: L_3\wedge L_4\Rightarrow \text{文本未引入占位外的新义务}.
\end{aligned}
\]
\end{Certificate}

\begin{proof}[证明（谓词因果链）]
由 $L_1$ 与定义~\ref{def:tex37-ch5-coverage}，
\[
  \Parse_\xi(t)=p^\star,
  \qquad
  \operatorname{Cov}_\xi(t)=\operatorname{Slot}(p^\star).
\]
前者给出 $L_2$，后者给出 $L_3$。

根据 $\operatorname{Obl}(t)$ 的定义，
\[
  \operatorname{Obl}(t)
  =
  \operatorname{Obl}(\Parse_\xi(t))
  =
  \operatorname{Obl}(p^\star),
\]
即得 $L_4$。
再由槽位全覆盖可知，文本已经把占位中应出现的义务全数外显出来，但并没有越出占位边界。
因此表达阶段只能重排、润色或展开既有义务，而不能凭空新增求解义务，即 $L_5$。证毕。
\end{proof}

\subsection{命题：所谓“充分输出”，其严格口径是覆盖完整且可解析回原占位}
\label{subsec:tex37-ch5-prop-coverage}

\begin{proposition}[覆盖式充分输出命题]
\label{prop:tex37-ch5-coverage}
设 $p^\star=(s_1,\dots,s_m)$ 是已固定占位，$t\in \mathsf{Txt}_\xi$ 是模式 $\xi$ 下的候选文本。
则“$t$ 对 $p^\star$ 充分”应严格定义为
\[
  \operatorname{Cov}_\xi(t)=\operatorname{Slot}(p^\star)
  \qquad\text{且}\qquad
  \Parse_\xi(t)=p^\star.
\]
换言之，充分性不等价于文本更长、措辞更多或论述更繁，
而等价于：
\[
  \boxed{
    \text{应出现的占位槽全部出现，并且文本仍能被解析回原占位。}
  }
\]
\end{proposition}

\begin{Certificate}[证书：充分性 $\Leftrightarrow$ 覆盖完整且回投影不失真]
\label{cert:tex37-ch5-coverage}
证书登记谓词链：
\[
\begin{aligned}
  C_1 &: \text{“充分”意味着占位信息未缺槽},\\
  C_2 &: C_1\Rightarrow \operatorname{Cov}_\xi(t)=\operatorname{Slot}(p^\star),\\
  C_3 &: \text{“对同一占位充分”意味着语义核未漂移},\\
  C_4 &: C_3\Rightarrow \Parse_\xi(t)=p^\star,\\
  C_5 &: C_2\wedge C_4\Rightarrow \text{充分性的严格判定式成立}.
\end{aligned}
\]
\end{Certificate}

\begin{proof}[证明（谓词因果链）]
若文本对已固定占位 $p^\star$ 是充分的，那么它不能漏掉任何应出现的槽位，否则求解义务无法被完整外显，因此 $C_1\Rightarrow C_2$。
另一方面，若文本虽覆盖很多内容却无法解析回原占位，则它已经脱离既定语义核，不再是对同一占位的充分表达，因此 $C_3\Rightarrow C_4$。
把二者合并，得到
\[
  \operatorname{Cov}_\xi(t)=\operatorname{Slot}(p^\star)
  \quad\wedge\quad
  \Parse_\xi(t)=p^\star.
\]
反过来，若这两个条件同时成立，则文本既无缺槽，又未偏离原占位，故正是严格意义下的充分输出。命题成立。证毕。
\end{proof}

\subsection{命题：文本可以同时做到“充分 + 可读 + 严谨 + 可风格化调节”}
\label{subsec:tex37-ch5-prop-multistyle}

\begin{proposition}[同一占位的全审计多外观命题]
\label{prop:tex37-ch5-multistyle}
固定同一占位 $p^\star$。若存在模式
\[
  \xi_1,\dots,\xi_r\in\Xi
\]
满足
\[
  \Audit^{\mathrm{full}}_{\xi_i}(p^\star)\neq\varnothing
  \qquad
  (1\le i\le r),
\]
并令
\[
  t_i
  :=
  \min\nolimits_{\prec_{\xi_i}}
  \Audit^{\mathrm{full}}_{\xi_i}(p^\star),
\]
则每个 $t_i$ 都同时满足：
\[
  \Parse_{\xi_i}(t_i)=p^\star,
  \qquad
  \operatorname{Cov}_{\xi_i}(t_i)=\operatorname{Slot}(p^\star),
\]
\[
  \operatorname{Read}_{\xi_i}(t_i)=\top,
  \qquad
  \Style_{\xi_i}(t_i)=\top.
\]
因此，
\[
  \boxed{
    \text{充分性与可读性可以并存，严谨性与风格化也可以并存。}
  }
\]
\end{proposition}

\begin{Certificate}[证书：全审计接受集非空 $\Rightarrow$ 同核多外观]
\label{cert:tex37-ch5-multistyle}
证书登记谓词链：
\[
\begin{aligned}
  M_1 &: \Audit^{\mathrm{full}}_{\xi_i}(p^\star)\neq\varnothing,\\
  M_2 &: t_i=\min\nolimits_{\prec_{\xi_i}}\Audit^{\mathrm{full}}_{\xi_i}(p^\star),\\
  M_3 &: M_2\Rightarrow t_i\in \Audit^{\mathrm{full}}_{\xi_i}(p^\star),\\
  M_4 &: M_3\Rightarrow \Parse_{\xi_i}(t_i)=p^\star,\ \operatorname{Cov}_{\xi_i}(t_i)=\operatorname{Slot}(p^\star),\\
  M_5 &: M_3\Rightarrow \operatorname{Read}_{\xi_i}(t_i)=\top,\ \Style_{\xi_i}(t_i)=\top,\\
  M_6 &: M_4\wedge M_5\Rightarrow \text{同一语义核可有多种充分、可读、严谨且风格合格的文本壳}.
\end{aligned}
\]
\end{Certificate}

\begin{proof}[证明（谓词因果链）]
由 $M_1$，每个模式 $\xi_i$ 下至少存在一个通过全审计的候选文本。
由 $M_2$ 与全序最小元的定义，得到 $M_3$：
\[
  t_i\in \Audit^{\mathrm{full}}_{\xi_i}(p^\star).
\]
根据全审计接受集的定义，$M_3$ 立即推出
\[
  \Parse_{\xi_i}(t_i)=p^\star,
  \qquad
  \operatorname{Cov}_{\xi_i}(t_i)=\operatorname{Slot}(p^\star),
\]
以及
\[
  \operatorname{Read}_{\xi_i}(t_i)=\top,
  \qquad
  \Style_{\xi_i}(t_i)=\top.
\]
于是每个 $t_i$ 都在不改变语义核的前提下，同时满足充分、可读、严谨与风格合格。
故命题成立。证毕。
\end{proof}

\subsection{定理：在该架构下，模型可以不承担求解型推理，只承担受约束表达}
\label{subsec:tex37-ch5-theorem-render}

\begin{theorem}[职责降级定理]
\label{thm:tex37-ch5-render-not-solve}
设
\[
  p^\star=\mathcal{T}_{\mathrm{solve}}(b,q)
\]
已经由 G 侧修复塔、规范语义层与占位层固定。
若某个模式 $\xi$ 下存在非失败输出
\[
  t=\mathcal{T}_{\mathrm{render}}(p^\star,\xi)\neq\bot,
\]
则模型在该步中所承担的职责可被限制为
\[
  \{
    u_{\mathrm{candidate}},
    u_{\mathrm{phrase}},
    u_{\mathrm{layout}},
    u_{\mathrm{style}}
  \},
\]
把这个集合记为
\[
  \mathfrak{M}_{\mathrm{render}}.
\]
即候选生成、措辞组织、段落编排与风格调整；
而不再需要承担
\[
  \mathfrak{S}_{\mathrm{solve}}
  =
  \{
    u_{\mathrm{ctx}},
    u_{\mathrm{repair}},
    u_{\mathrm{path}},
    u_{\mathrm{cert}},
    u_{\mathrm{verdict}}
  \}
\]
中的任何义务。
因此，在该前提下，模型可以被严格降级为\textbf{受约束表达器}，而不是求解器。
\end{theorem}

\begin{Certificate}[证书：上游已解 + 下游只做投影表达]
\label{cert:tex37-ch5-render-not-solve}
证书登记谓词链：
\[
\begin{aligned}
  R_1 &: p^\star=\mathcal{T}_{\mathrm{solve}}(b,q)\ \text{已由修复塔路径类固定},\\
  R_2 &: t=\mathcal{T}_{\mathrm{render}}(p^\star,\xi)\neq\bot,\\
  R_3 &: R_2\Rightarrow \Parse_\xi(t)=p^\star\ \text{（定理~\ref{thm:tex37-runtime-deterministic}）},\\
  R_4 &: R_1\wedge R_3\Rightarrow t\ \text{没有新增求解义务（引理~\ref{lem:tex37-ch5-no-new-solve}）},\\
  R_5 &: R_4\Rightarrow \mathfrak{S}_{\mathrm{solve}}\cap \mathfrak{M}_{\mathrm{render}}=\varnothing,\\
  R_6 &: R_5\Rightarrow \text{模型职责降级为候选表达与界面组织}.
\end{aligned}
\]
\end{Certificate}

\begin{proof}[证明（谓词因果链）]
由 $R_1$，真正的障碍修复、路径选择、证书拼接与占位固定都已在 G 侧完成。
因此模型接收到的不是开放问题本身，而是一个已固定占位 $p^\star$。

由 $R_2$ 与定理~\ref{thm:tex37-runtime-deterministic}，非失败文本满足
\[
  \Parse_\xi(t)=p^\star,
\]
即 $R_3$。
再由引理~\ref{lem:tex37-ch5-no-new-solve}，表达阶段不会新增求解义务，只会把 $p^\star$ 已有的槽位和义务外显出来。
因此
\[
  \mathfrak{S}_{\mathrm{solve}}\cap \mathfrak{M}_{\mathrm{render}}=\varnothing,
\]
即 $R_5$。

这说明模型在该阶段只需承担候选生成、措辞组织、段落排版与风格调节等表达型工作，
而不再承担上下文判定、障碍修复、路径选择、证书拼接与最终裁决等求解型工作。
故模型在该架构下可以被严格降级为受约束表达器。证毕。
\end{proof}

\subsection{命题：真正应尽量不交给模型的，不只是推理，还有最终裁决}
\label{subsec:tex37-ch5-prop-verdict}

\begin{proposition}[最终接受判定应由确定性合同接口承担]
\label{prop:tex37-ch5-verdict}
对任意已固定占位 $p^\star$ 与任意模式 $\xi$，定义接受谓词
\[
  \operatorname{Accept}_\xi(t;p^\star)
  :\Longleftrightarrow
  \bigl(
    \Parse_\xi(t)=p^\star
  \bigr)
  \wedge
  \bigl(
    \Style_\xi(t)=\top
  \bigr)
  \wedge
  \bigl(
    \operatorname{Read}_\xi(t)=\top
  \bigr)
  \wedge
  \bigl(
    \operatorname{Cov}_\xi(t)=\operatorname{Slot}(p^\star)
  \bigr).
\]
则最终输出应定义为
\[
  \operatorname{VerdictRender}_\xi(p^\star)
  :=
  \begin{cases}
    \min\nolimits_{\prec_\xi}
    \{t\mid \operatorname{Accept}_\xi(t;p^\star)\}, & \exists t,\ \operatorname{Accept}_\xi(t;p^\star),\\
    \bot, & \text{否则}.
  \end{cases}
\]
因此，严格的工程口径不是“模型自己审计自己”，而是：
\[
  \boxed{
    \text{模型可以参与候选生成；最终裁决应由解析器、审计器与合同接口承担。}
  }
\]
\end{proposition}

\begin{Certificate}[证书：接受条件全由外置接口定义]
\label{cert:tex37-ch5-verdict}
证书登记谓词链：
\[
\begin{aligned}
  V_1 &: \operatorname{Accept}_\xi(t;p^\star)\ \text{只依赖 }\Parse_\xi,\Style_\xi,\operatorname{Read}_\xi,
    \operatorname{Cov}_\xi,\\
  V_2 &: \text{这些对象全部是合同化、可复验的确定性接口},\\
  V_3 &: V_1\wedge V_2\Rightarrow \text{接受与否不依赖模型自我声明},\\
  V_4 &: \operatorname{VerdictRender}_\xi(p^\star)\ \text{由全序最小接受元或 }\bot\text{ 给出},\\
  V_5 &: V_3\wedge V_4\Rightarrow \text{最终裁决权在外置接口而不在模型}.
\end{aligned}
\]
\end{Certificate}

\begin{proof}[证明（谓词因果链）]
由 $\operatorname{Accept}_\xi(t;p^\star)$ 的定义，接受条件完全由
\[
  \Parse_\xi,\ \Style_\xi,\ \operatorname{Read}_\xi,\ \operatorname{Cov}_\xi
\]
决定，即 $V_1$。
这些对象在本稿口径下都是外置的解析器、审计器与覆盖检查器，因此是合同化、可复验的接口，即 $V_2$。

因此，候选文本是否被接受，不取决于模型口头声称“自己是对的”，而取决于它是否通过外置接口的检查，即 $V_3$。
再由 $\operatorname{VerdictRender}_\xi(p^\star)$ 的定义，最终输出由最小接受元或失败符号 $\bot$ 唯一给出，即 $V_4$。
故最终裁决权属于确定性合同接口，而不属于模型本身。证毕。
\end{proof}

\subsection{推论：这一章的显著增益，本质上是把模型依赖重心从求解侧移到表达侧}
\label{subsec:tex37-ch5-corollary}

\begin{corollary}[模型依赖重心迁移推论]
\label{cor:tex37-ch5-dependency-shift}
设分层前架构 $\mathfrak{A}_{\mathrm{before}}$ 让模型承担至少一个求解侧义务，即
\[
  \mathfrak{M}(\mathfrak{A}_{\mathrm{before}})\neq\varnothing,
\]
而分层后架构 $\mathfrak{A}_{\mathrm{after}}$ 满足定理~\ref{thm:tex37-ch5-render-not-solve} 与命题~\ref{prop:tex37-ch5-verdict} 的前提，
从而
\[
  \mathfrak{M}(\mathfrak{A}_{\mathrm{after}})\cap \mathfrak{S}_{\mathrm{solve}}=\varnothing.
\]
则
\[
  \operatorname{Load}_{\mathrm{solve}}(\mathfrak{A}_{\mathrm{after}})=0,
\]
并且
\[
  \Delta_{\mathrm{dep}}
  =
  \operatorname{Load}_{\mathrm{solve}}(\mathfrak{A}_{\mathrm{before}})
  >
  0.
\]
因此，本章所谓“显著增益”的更精确含义是：
\[
  \boxed{
    \text{不是让模型更会想，而是让系统更少依赖模型去想。}
  }
\]
\end{corollary}

\begin{Certificate}[证书：求解义务集合被完全外移]
\label{cert:tex37-ch5-dependency-shift}
证书登记谓词链：
\[
\begin{aligned}
  D_1 &: \mathfrak{M}(\mathfrak{A}_{\mathrm{before}})\subseteq \mathfrak{S}_{\mathrm{solve}}\ \text{且非空},\\
  D_2 &: \mathfrak{M}(\mathfrak{A}_{\mathrm{after}})\cap \mathfrak{S}_{\mathrm{solve}}=\varnothing\ \text{（定理~\ref{thm:tex37-ch5-render-not-solve}）},\\
  D_3 &: \omega(u)>0\ \text{对一切 }u\in\mathfrak{S}_{\mathrm{solve}},\\
  D_4 &: D_2\wedge D_3\Rightarrow \operatorname{Load}_{\mathrm{solve}}(\mathfrak{A}_{\mathrm{after}})=0,\\
  D_5 &: D_1\wedge D_3\Rightarrow \operatorname{Load}_{\mathrm{solve}}(\mathfrak{A}_{\mathrm{before}})>0,\\
  D_6 &: D_4\wedge D_5\Rightarrow \Delta_{\mathrm{dep}}>0.
\end{aligned}
\]
\end{Certificate}

\begin{proof}[证明（谓词因果链）]
由 $D_2$，分层后模型不再承担任何求解侧义务。
由于每个义务权重都严格为正，故
\[
  \operatorname{Load}_{\mathrm{solve}}(\mathfrak{A}_{\mathrm{after}})
  =
  \sum_{u\in \mathfrak{M}(\mathfrak{A}_{\mathrm{after}})}\omega(u)
  =0,
\]
即 $D_4$。

另一方面，由 $D_1$ 与 $D_3$，分层前模型至少承担一个正权重求解义务，因此
\[
  \operatorname{Load}_{\mathrm{solve}}(\mathfrak{A}_{\mathrm{before}})>0,
\]
即 $D_5$。
于是
\[
  \Delta_{\mathrm{dep}}
  =
  \operatorname{Load}_{\mathrm{solve}}(\mathfrak{A}_{\mathrm{before}})
  -
  \operatorname{Load}_{\mathrm{solve}}(\mathfrak{A}_{\mathrm{after}})
  >
  0.
\]
这正是“依赖重心从求解侧移到表达侧”的严格数学写法。证毕。
\end{proof}

\subsection{备注：条件边界、收紧版总结与一句话结论}
\label{subsec:tex37-ch5-remarks}

\begin{remark}[“小模型足够”只能写成条件命题]
本章不主张无条件地说“小模型总是足够”。
更严格的写法应是：
\[
  \boxed{
    \begin{aligned}
      &\text{若 }p^\star\text{ 已固定，且模板族、解析器、覆盖检查器与审计器足够完备，}\\
      &\text{则候选表达这一步可以由小模型承担。}
    \end{aligned}
  }
\]
换言之，小模型的可行性来自\textbf{职责降级}，而不是来自模型本身突然变成求解器。
\end{remark}

\begin{remark}[收紧后的结论表述]
本章更适合采用下面这个严格版本：
\[
  \boxed{
    \text{即便在对象层次纠偏后，该章仍然提供了显著增益。}
  }
\]
\[
  \boxed{
    \begin{aligned}
      &\text{因为 G 框架塔纤维“对话”的真实结果是障碍修复后的算子语义序列或逻辑占位序列，}\\
      &\text{而不是直接的自然语言文本。}
    \end{aligned}
  }
\]
\[
  \boxed{
    \begin{aligned}
      &\text{一旦该序列在 G 侧被确定，则通过 LLM 的投影表达就被降级为}\\
      &\text{受合同约束的表面展开问题：文本是否被接受，不取决于模型是否会证明，}\\
      &\text{而取决于其是否覆盖既定占位、能被解析回原占位，并通过风格、可读与层级审计。}
    \end{aligned}
  }
\]
\[
  \boxed{
    \begin{aligned}
      &\text{因此，文本可以在不改变语义核的前提下，通过反复审计获得}\\
      &\text{充分性、可读性、严谨性与风格化调节；}\\
      &\text{真正的裁决权仍应保留在确定性合同接口，而非模型本身。}
    \end{aligned}
  }
\]
\end{remark}

\begin{remark}[一句话结论]
最短的严格压缩可写成：
\[
  \boxed{
    \begin{aligned}
      &\text{这个增益非常实：它把“模型负责想”改写成}\\
      &\text{“G 框架负责想，模型负责把已经想清楚的东西审计后说清楚”。}
    \end{aligned}
  }
\]
\end{remark}

\clearpage

%%%%%%%%%%%%%%%%%%%%%%%%%%%%%%%%%%%%%%%%%%%%%%%%%%%%%%%%%%%%
\section{形式逻辑：联立 tex36 与本稿前文后，求解与决策主权可回收到 G 侧白盒程序}
\label{sec:tex37-ch6}
%%%%%%%%%%%%%%%%%%%%%%%%%%%%%%%%%%%%%%%%%%%%%%%%%%%%%%%%%%%%

\begin{remark}[核心陈述（条件化后）]
把仓内 `tex36` 给出的全编码联络提升、稳定同伦扭曲正规形提取与有效组合作用映射，
与本稿第一章到第五章给出的占位投影、审计回投影与运行时降级链联立之后，可得到如下结构结论：
\[
  \boxed{
    \begin{aligned}
      &\text{G 框架负责求解、修复、裁决与动作选择；}\\
      &\text{LLM 只负责把既定结果投影成可审计文本。}
    \end{aligned}
  }
\]
因此，“G 框架驱动的 AI 由程序确定性表达、白盒，并把 LLM 降级为插件，回收思考和决策主权”
这一方向是成立的；但“彻底回收”必须写成带任务域条件的严格命题，而不能写成无条件口号。
\end{remark}

\begin{remark}[阅读导航：仓内引用与本稿承接位置]
本章对 `tex36` 的依赖主要集中在三条链上 \cite{RepoTex36HomotopyTwistBijection}：
\[
  \begin{aligned}
    &x_0\longmapsto \Sigma_N(x_0),\\
    &\text{路径积分求解}\equiv \text{稳定同伦扭曲正规形提取},\\
    &\mathfrak S_{n,x,L}^{\mathrm{eff}}=\operatorname{Rep}_n(x)=\operatorname{NF}_{A_n}(x).
  \end{aligned}
\]
对本稿前文的依赖主要集中在：
第一章的 LLM 运行时降级定理~\ref{thm:tex37-runtime-deterministic}、
第二章的工程增益主导定理~\ref{thm:tex37-engineering-dominates}、
第四章的内核对象定理~\ref{thm:tex37-ch4-object} 与第五章的职责降级定理~\ref{thm:tex37-ch5-render-not-solve}。
本章的目标，是把这些链条合成为一个更强的“白盒 AI 主权分配”命题。
\end{remark}

\subsection{定义：三种主权、G 侧全编码求解链与 LLM 插件投影链}
\label{subsec:tex37-ch6-definitions}

\begin{definition}[三种主权]
\label{def:tex37-ch6-sovereignty}
定义系统主权分解为
\[
  \mathsf{Sov}
  :=
  \mathsf{Sov}_{\mathrm{solve}}
  \oplus
  \mathsf{Sov}_{\mathrm{decide}}
  \oplus
  \mathsf{Sov}_{\mathrm{render}}.
\]
其中：
\begin{enumerate}
  \item
  \[
    \mathsf{Sov}_{\mathrm{solve}}
  \]
  表示“谁决定规范语义解、障碍修复链与证明责任”；
  \item
  \[
    \mathsf{Sov}_{\mathrm{decide}}
  \]
  表示“谁决定动作、停止、接受、拒绝与执行分支”；
  \item
  \[
    \mathsf{Sov}_{\mathrm{render}}
  \]
  表示“谁决定把既定语义结果如何表述为外显文本”。
\end{enumerate}
本章要回收的核心，不是第三项，而是前两项。
\end{definition}

\begin{definition}[G 侧全编码求解链与动作接口]
\label{def:tex37-ch6-g-solve}
记问题空间为
\[
  \mathcal{Q}.
\]
定义 G 侧求解入口
\[
  \mathcal{G}_{\mathrm{solve}}:
  \mathcal{B}_{\mathrm{ctx}}\times\mathcal{Q}
  \to
  \mathcal{PI}_{A_0}(b_0,B_0^{\mathrm{tar}};p_0),
\]
并记
\[
  x_0:=\mathcal{G}_{\mathrm{solve}}(b,q).
\]

沿用 `tex36` 的全编码口径 \cite{RepoTex36HomotopyTwistBijection}，
对任意 $N\in\mathbb{N}_{>0}$，记
\[
  \Sigma_N(x_0)
  :=
  \bigl(
    (x_n)_{n=0}^{N},
    (u_n)_{n=0}^{N},
    (\tau_n)_{n=0}^{N}
  \bigr)
\]
为从底层类 $x_0$ 唯一提升得到的 $N$ 级全编码求解链。

对每一层 $n$，定义稳定同伦扭曲正规形
\[
  \tau_n^\star:=\operatorname{NF}_{A_n}(x_n),
\]
并定义修复映射
\[
  \operatorname{Rep}_n(x_n):=\operatorname{NF}_{A_n}(x_n).
\]
若再给定确定性动作接口
\[
  \operatorname{Act}_n:
  \mathcal{T}_{A_n}(b_n,B_n^{\mathrm{tar}};p_n)
  \to
  \mathcal{A}_n,
\]
则第 $n$ 层动作定义为
\[
  a_n^\star:=\operatorname{Act}_n(\tau_n^\star).
\]
\end{definition}

\begin{definition}[从 G 侧正规形到占位与文本插件的投影链]
\label{def:tex37-ch6-plugin}
为把 `tex36` 的正规形求解链接到本稿第一章的占位链，
定义一个接口编码映射
\[
  \operatorname{Enc}_G:
  \mathcal{T}_{A_n}(b_n,B_n^{\mathrm{tar}};p_n)
  \to
  \mathsf{NF}_G^{(0)}.
\]
对某个稳定正规形
\[
  \tau^\star\in \mathcal{T}_{A_n}(b_n,B_n^{\mathrm{tar}};p_n),
\]
定义其逻辑占位为
\[
  p^\star:=\Theta_0\bigl(\operatorname{Enc}_G(\tau^\star)\bigr),
\]
并对任意模式 $\xi\in\Xi$ 定义外显文本
\[
  t^\star:=\Render^{\mathrm{LLM}}_\xi(p^\star)\in \mathsf{Txt}_\xi\sqcup\{\bot\}.
\]
若 $t^\star\neq\bot$，则由第一章可知
\[
  \Parse_\xi(t^\star)=p^\star,
  \qquad
  \Style_\xi(t^\star)=\top.
\]
因此 LLM 在该链条中只充当
\[
  p^\star\mapsto t^\star
\]
的可替换文本投影插件。
\end{definition}

\begin{remark}[联立后的总链]
定义~\ref{def:tex37-ch6-g-solve} 与定义~\ref{def:tex37-ch6-plugin} 合起来给出总链
\[
  (\text{上下文},q)
  \xrightarrow{\mathcal{G}_{\mathrm{solve}}}
  x_0
  \xrightarrow{\Sigma_N}
  \tau^\star
  \xrightarrow{\operatorname{Enc}_G}
  \operatorname{Enc}_G(\tau^\star)
  \xrightarrow{\Theta_0}
  p^\star
  \xrightarrow{\Render^{\mathrm{LLM}}_\xi}
  t^\star.
\]
本章要证明的，不是“LLM 消失”，而是：
\[
  \text{它退化为末端表达插件，而不再占据系统大脑位置。}
\]
\end{remark}

\subsection{公理（降级为定理）：有限层级的主权回收程序链可按联络塔层数递归构造}
\label{subsec:tex37-ch6-axiom-downgrade}

\begin{theorem}[公理降级：主权回收程序链可按层数递归构造（数学归纳法证书）]
\label{thm:tex37-ch6-program-chain}
设给定上下文基点 $b\in\mathcal{B}_{\mathrm{ctx}}$、问题 $q\in\mathcal{Q}$ 与层数 $N\in\mathbb{N}_{>0}$。
若满足：
\begin{enumerate}
  \item 存在基步求解入口
  \[
    x_0=\mathcal{G}_{\mathrm{solve}}(b,q);
  \]
  \item 对每个 $0\le n\le N$，`tex36` 的全编码求解链分量
  \[
    x_n,\ u_n,\ \tau_n
  \]
  唯一确定；
  \item 对每个 $0\le n\le N$，动作接口
  \[
    \operatorname{Act}_n
  \]
  是确定性映射；
  \item 接口编码
  \[
    \operatorname{Enc}_G
  \]
  与占位投影
  \[
    \Theta_0
  \]
  已固定。
\end{enumerate}
则存在唯一的有限层级主权回收程序链
\[
  \Omega_N(b,q)
  :=
  \bigl(
    (x_n)_{n=0}^{N},
    (u_n)_{n=0}^{N},
    (\tau_n)_{n=0}^{N},
    (a_n^\star)_{n=0}^{N},
    p^\star
  \bigr),
\]
其中对每个 $0\le n\le N$，
\[
  \tau_n=\operatorname{NF}_{A_n}(x_n),
  \qquad
  a_n^\star=\operatorname{Act}_n(\tau_n),
\]
并且末端占位满足
\[
  p^\star=\Theta_0\bigl(\operatorname{Enc}_G(\tau_N)\bigr).
\]
因此，在有限层级上，求解、动作选择与占位固定都可被递归压成程序化白盒链条。
\end{theorem}

\begin{Certificate}[证书：基步 + 归纳步的主权回收程序链证书]
\label{cert:tex37-ch6-program-chain}
令谓词
\[
  P(N)
  :
  \Longleftrightarrow
  \forall (b,q)\in \mathcal{B}_{\mathrm{ctx}}\times\mathcal{Q},\
  \exists!\,\Omega_N(b,q).
\]
则数学归纳法证书记为
\[
  \pi_{\mathrm{prog}}
  =
  \bigl(
    P(1),\
    \forall N\in\mathbb{N}_{>0},\ P(N)\Rightarrow P(N+1)
  \bigr).
\]
其归纳步的核心递推式为
\[
  \tau_{N+1}=\operatorname{NF}_{A_{N+1}}(x_{N+1}),
  \qquad
  a_{N+1}^\star=\operatorname{Act}_{N+1}(\tau_{N+1}),
  \qquad
  p^\star=\Theta_0\bigl(\operatorname{Enc}_G(\tau_{N+1})\bigr).
\]
\end{Certificate}

\begin{proof}[证明（数学归纳法证书；基于数学符号推演逻辑的谓词因果链）]
令
\[
  P(N)
  :
  \Longleftrightarrow
  \forall (b,q),\ \exists!\,\Omega_N(b,q).
\]

\textbf{基步：}
取 $N=1$。由条件(1)，
\[
  x_0=\mathcal{G}_{\mathrm{solve}}(b,q)
\]
唯一确定。
由条件(2) 以及 `tex36` 的全编码求解链口径，
\[
  (x_0,u_0,\tau_0,x_1,u_1,\tau_1)
\]
唯一确定。
再由条件(3)，
\[
  a_0^\star=\operatorname{Act}_0(\tau_0),
  \qquad
  a_1^\star=\operatorname{Act}_1(\tau_1)
\]
唯一确定。
最后由条件(4)，
\[
  p^\star=\Theta_0(\operatorname{Enc}_G(\tau_1))
\]
唯一确定。
故
\[
  \Omega_1(b,q)
\]
存在且唯一，即 $P(1)$ 成立。

\textbf{归纳步：}
设 $P(N)$ 成立。则
\[
  \Omega_N(b,q)
  =
  \bigl(
    (x_n)_{n=0}^{N},
    (u_n)_{n=0}^{N},
    (\tau_n)_{n=0}^{N},
    (a_n^\star)_{n=0}^{N},
    p^\star_N
  \bigr)
\]
已唯一确定。
由条件(2) 的全编码递推，下一层分量
\[
  x_{N+1},\ u_{N+1},\ \tau_{N+1}
\]
唯一确定，其中
\[
  \tau_{N+1}=\operatorname{NF}_{A_{N+1}}(x_{N+1}).
\]
再由条件(3)，动作
\[
  a_{N+1}^\star=\operatorname{Act}_{N+1}(\tau_{N+1})
\]
唯一确定。
最后由条件(4)，末端占位
\[
  p^\star_{N+1}
  =
  \Theta_0(\operatorname{Enc}_G(\tau_{N+1}))
\]
唯一确定。
故
\[
  \Omega_{N+1}(b,q)
\]
存在且唯一，即 $P(N)\Rightarrow P(N+1)$。

由数学归纳法，对任意 $N\in\mathbb{N}_{>0}$，$P(N)$ 成立。
因此有限层级的主权回收程序链可按联络塔层数递归构造。证毕。
\end{proof}

\subsection{定理：联立 tex36 与 tex37 后，G 框架驱动的 AI 可构造成程序确定性白盒架构，LLM 被降级为插件}
\label{subsec:tex37-ch6-theorem}

\begin{theorem}[白盒主权回收定理]
\label{thm:tex37-ch6-whitebox-plugin}
在定理~\ref{thm:tex37-ch6-program-chain} 的前提下，若末端文本满足
\[
  t^\star=\Render^{\mathrm{LLM}}_\xi(p^\star)\neq\bot,
\]
则联立 `tex36` 与本稿前文可得到：
\begin{enumerate}
  \item
  \[
    \mathsf{Sov}_{\mathrm{solve}}
  \]
  被回收到 G 侧程序链
  \[
    x_0\mapsto \Sigma_N(x_0)\mapsto \tau^\star;
  \]
  \item
  \[
    \mathsf{Sov}_{\mathrm{decide}}
  \]
  被回收到确定性动作接口
  \[
    a^\star=\operatorname{Act}_N(\tau^\star);
  \]
  \item
  \[
    \mathsf{Sov}_{\mathrm{render}}
  \]
  保留在受约束的文本投影层
  \[
    p^\star\mapsto t^\star.
  \]
\end{enumerate}
因此，G 框架驱动的 AI 可被构造成程序确定性白盒架构，而 LLM 在该架构中只是一类可替换的文本投影插件。
\end{theorem}

\begin{Certificate}[证书：全编码求解链 + 作用映射即修复 + 运行时降级]
\label{cert:tex37-ch6-whitebox-plugin}
证书登记谓词链：
\[
\begin{aligned}
  W_1 &: x_0\mapsto \Sigma_N(x_0)\ \text{由 `tex36` 的全编码求解链唯一确定},\\
  W_2 &: \text{路径积分求解}\equiv \text{稳定同伦扭曲正规形提取}\ \text{（`tex36`）},\\
  W_3 &: \mathfrak S_{n,x,L}^{\mathrm{eff}}=\operatorname{Rep}_n(x)=\operatorname{NF}_{A_n}(x)\ \text{（`tex36`）},\\
  W_4 &: p^\star=\Theta_0(\operatorname{Enc}_G(\tau^\star)),\\
  W_5 &: t^\star=\Render^{\mathrm{LLM}}_\xi(p^\star)\neq\bot\Rightarrow \Parse_\xi(t^\star)=p^\star\ \text{（定理~\ref{thm:tex37-runtime-deterministic}）},\\
  W_6 &: W_1\wedge W_2\wedge W_3\Rightarrow \mathsf{Sov}_{\mathrm{solve}}\oplus\mathsf{Sov}_{\mathrm{decide}}\ \text{位于
    G 侧},\\
  W_7 &: W_4\wedge W_5\Rightarrow \mathsf{Sov}_{\mathrm{render}}\ \text{仅位于受审计插件层},\\
  W_8 &: W_6\wedge W_7\Rightarrow \text{白盒主权回收结论成立}.
\end{aligned}
\]
\end{Certificate}

\begin{proof}[证明（基于数学符号推演逻辑的谓词因果链）]
由 $W_1$，问题进入 G 侧以后，底层类
\[
  x_0
\]
会被唯一提升为
\[
  \Sigma_N(x_0)
  =
  \bigl(
    (x_n)_{n=0}^{N},
    (u_n)_{n=0}^{N},
    (\tau_n)_{n=0}^{N}
  \bigr).
\]
由 $W_2$，路径积分求解不再被理解为把思考外包给黑箱语言模型，
而是被改写为对稳定同伦扭曲正规形
\[
  \tau^\star=\operatorname{NF}_{A_N}(x_N)
\]
的提取。
这说明规范语义解与修复链的决定权位于 G 侧程序链内部，故
\[
  \mathsf{Sov}_{\mathrm{solve}}
\]
已回收到 G 侧。

再由 $W_3$，
\[
  \mathfrak S_{N,x_N,L}^{\mathrm{eff}}
  =
  \operatorname{Rep}_N(x_N)
  =
  \operatorname{NF}_{A_N}(x_N)
  =
  \tau^\star.
\]
也就是说，“采取哪一个有效组合/修复动作”并不是由 LLM 临时拍板，
而是已经包含在 G 侧求解结果之内。
配合确定性动作接口
\[
  a^\star=\operatorname{Act}_N(\tau^\star),
\]
可知动作、停止、执行分支的选择权也位于 G 侧，故
\[
  \mathsf{Sov}_{\mathrm{decide}}
\]
被回收到 G 侧程序链。

最后，由 $W_4$ 与 $W_5$，
\[
  \tau^\star
  \xrightarrow{\operatorname{Enc}_G}
  \operatorname{Enc}_G(\tau^\star)
  \xrightarrow{\Theta_0}
  p^\star
  \xrightarrow{\Render^{\mathrm{LLM}}_\xi}
  t^\star,
\]
且非失败输出满足
\[
  \Parse_\xi(t^\star)=p^\star,
  \qquad
  \Style_\xi(t^\star)=\top.
\]
因此 LLM 既不拥有语义定义权，也不拥有证明裁判权、修复权、动作提交权与停止判定权；
它只拥有受解析器、审计器、合同和全序规则约束的候选表达权。
故
\[
  \mathsf{Sov}_{\mathrm{render}}
\]
仅位于插件化表达层。

综上，
\[
  \mathsf{Sov}_{\mathrm{solve}}
  \oplus
  \mathsf{Sov}_{\mathrm{decide}}
\]
被回收到 G 侧白盒程序，而
\[
  \mathsf{Sov}_{\mathrm{render}}
\]
留在受审计插件层。定理成立。证毕。
\end{proof}

\subsection{引理：这里所谓“白盒”具有严格的状态、变换与裁决可见性含义}
\label{subsec:tex37-ch6-lemma}

\begin{lemma}[白盒可见性引理]
\label{lem:tex37-ch6-whitebox}
在定理~\ref{thm:tex37-ch6-whitebox-plugin} 的设定下，以下对象全部可枚举、可追踪、可审计：
\[
  x_n,\ u_n,\ \tau_n,\ a_n^\star,\ \Sigma_N(x_0),\ \operatorname{Rep}_n(x_n),\ p^\star,\ \Audit_\xi(p^\star),\
    \Parse_\xi(t^\star).
\]
因此，这里“白盒”的严格含义不是修辞，而是：
\[
  \boxed{
    \text{状态可见、变换可见、修复链可见、裁决条件可见、失败边界可见。}
  }
\]
\end{lemma}

\begin{Certificate}[证书：全编码链 + 解析审计链的对象可见性]
\label{cert:tex37-ch6-whitebox}
证书登记谓词链：
\[
\begin{aligned}
  L_1 &: \Sigma_N(x_0)\ \text{唯一给出 }x_n,u_n,\tau_n\ \text{（定理~\ref{thm:tex37-ch6-program-chain}）},\\
  L_2 &: a_n^\star=\operatorname{Act}_n(\tau_n)\ \text{由确定性接口给出},\\
  L_3 &: \operatorname{Rep}_n(x_n)=\operatorname{NF}_{A_n}(x_n)\ \text{是显式可写对象},\\
  L_4 &: p^\star=\Theta_0(\operatorname{Enc}_G(\tau^\star)),\\
  L_5 &: t^\star\neq\bot\Rightarrow \Parse_\xi(t^\star)=p^\star\ \wedge\ t^\star\in\Audit_\xi(p^\star),\\
  L_6 &: L_1\wedge L_2\wedge L_3\wedge L_4\wedge L_5\Rightarrow \text{白盒可见性成立}.
\end{aligned}
\]
\end{Certificate}

\begin{proof}[证明（基于数学符号推演逻辑的谓词因果链）]
由 $L_1$，G 侧全编码求解链已经显式给出
\[
  x_n,\ u_n,\ \tau_n.
\]
由 $L_2$，每一层动作
\[
  a_n^\star
\]
都是确定性接口对正规形的像，因此同样显式可追踪。
由 $L_3$，修复映射
\[
  \operatorname{Rep}_n(x_n)=\operatorname{NF}_{A_n}(x_n)
\]
不是口头推断，而是明确对象。
由 $L_4$ 与 $L_5$，末端占位、审计接受集与解析回投影对象也都显式可见。

于是从求解入口到文本出口的每个关键节点都有对象、映射与判定边界可供复验，
系统不再是“模型说了算”的黑箱，而是一个状态、变换与裁决条件均可显式列举的白盒链条。证毕。
\end{proof}

\subsection{命题：所谓“彻底回收思考和决策主权”，只能在已覆盖任务域内写成条件命题}
\label{subsec:tex37-ch6-proposition}

\begin{proposition}[条件化主权回收命题]
\label{prop:tex37-ch6-conditional-sovereignty}
设以下四个条件成立：
\[
  C_1:\ \text{任务域已被 G 侧状态、动作、障碍、证书与目标函数充分覆盖};
\]
\[
  C_2:\ \Sigma_N,\ \operatorname{NF}_{A_n},\ \operatorname{Rep}_n\ \text{在该域内可实现};
\]
\[
  C_3:\ \operatorname{Act}_n,\ \Parse_\xi,\ \Audit_\xi\ \text{在该域内可判定};
\]
\[
  C_4:\ \text{LLM 无权直接改写占位、规则、动作提交与停止判定}.
\]
则在该任务域内，
\[
  \boxed{
    \mathsf{Sov}_{\mathrm{solve}}
    \oplus
    \mathsf{Sov}_{\mathrm{decide}}
  }
\]
\[
  \boxed{
    \text{可以结构性地、近乎彻底地回收到 G 侧程序。}
  }
\]
更严格地说：
\[
  \boxed{
    \text{在已覆盖任务域内，可彻底回收；在超域开放环境中，仍需继续扩充 G 侧形式系统。}
  }
\]
\end{proposition}

\begin{Certificate}[证书：覆盖条件 + 可实现条件 + 裁决边界条件]
\label{cert:tex37-ch6-conditional-sovereignty}
证书登记谓词链：
\[
\begin{aligned}
  P_1 &: C_1\Rightarrow \text{求解与决策对象全部有 G 侧表示},\\
  P_2 &: C_2\Rightarrow \text{这些对象可沿 }\Sigma_N,\operatorname{NF}_{A_n},\operatorname{Rep}_n\ \text{实际求得},\\
  P_3 &: C_3\Rightarrow \text{动作、接受与失败都可由确定性接口裁决},\\
  P_4 &: C_4\Rightarrow \text{LLM 不能越权篡改求解链与决策链},\\
  P_5 &: P_1\wedge P_2\wedge P_3\wedge P_4\Rightarrow \mathsf{Sov}_{\mathrm{solve}}\oplus\mathsf{Sov}_{\mathrm{decide}}\
    \text{在该域内被结构性回收}.
\end{aligned}
\]
\end{Certificate}

\begin{proof}[证明（基于数学符号推演逻辑的谓词因果链）]
由 $C_1$，任务域中需要被求解和被决策的状态、动作、障碍、证书与目标函数都已经被压入 G 侧对象语言，
因此不存在必须临时外包给黑箱模型解释的“无表示核心量”。
这给出 $P_1$。

由 $C_2$，这些对象并非只有名字，而是能通过
\[
  \Sigma_N,\ \operatorname{NF}_{A_n},\ \operatorname{Rep}_n
\]
实际算出，于是给出 $P_2$。
由 $C_3$，动作接口、解析器与审计器在域内可判定，因此停止、接受、拒绝与执行边界都可由确定性接口给出，这就是 $P_3$。
由 $C_4$，LLM 被限制在表达层，不能越权改写占位、规则和动作提交，这给出 $P_4$。

于是，在该任务域中，求解权与决策权已没有结构性理由继续保留在 LLM 侧，
故
\[
  \mathsf{Sov}_{\mathrm{solve}}
  \oplus
  \mathsf{Sov}_{\mathrm{decide}}
\]
可以被结构性回收到 G 侧程序。
但一旦离开已覆盖任务域，新的状态、规则或动作接口可能尚未被 G 侧形式化，
因此“彻底回收”必须被理解为任务域条件下的命题，而不是超域无条件口号。证毕。
\end{proof}

\subsection{推论：该架构同时具有程序确定性求解与程序确定性表达两层确定性}
\label{subsec:tex37-ch6-corollary}

\begin{corollary}[双层程序确定性推论]
\label{cor:tex37-ch6-double-determinism}
在定理~\ref{thm:tex37-ch6-whitebox-plugin} 与命题~\ref{prop:tex37-ch6-conditional-sovereignty} 的条件下，系统具有两层确定性：
\begin{enumerate}
  \item \textbf{求解确定性：}
  \[
    (b,q)\mapsto x_0\mapsto \Sigma_N(x_0)\mapsto \tau^\star\mapsto a^\star;
  \]
  \item \textbf{表达确定性：}
  \[
    p^\star\mapsto \Render^{\mathrm{LLM}}_\xi(p^\star).
  \]
\end{enumerate}
因此，更准确的总结应写成：
\[
  \boxed{
    \text{系统既可以程序确定性求解，也可以程序确定性表达。}
  }
\]
\end{corollary}

\begin{Certificate}[证书：G 侧确定性 + 运行时确定性]
\label{cert:tex37-ch6-double-determinism}
证书登记谓词链：
\[
\begin{aligned}
  C_1 &: (b,q)\mapsto x_0\mapsto \Sigma_N(x_0)\mapsto \tau^\star\mapsto a^\star\ \text{由定理~\ref{thm:tex37-ch6-program-chain} 与动作接口确定},\\
  C_2 &: p^\star\mapsto \Render^{\mathrm{LLM}}_\xi(p^\star)\ \text{由定理~\ref{thm:tex37-runtime-deterministic} 确定},\\
  C_3 &: C_1\wedge C_2\Rightarrow \text{双层程序确定性成立}.
\end{aligned}
\]
\end{Certificate}

\begin{proof}[证明（基于数学符号推演逻辑的谓词因果链）]
由 $C_1$，在已覆盖任务域内，从上下文问题到稳定正规形再到动作输出的整条 G 侧链是确定性的。
由 $C_2$，在固定模型版本、提示模板、策略、种子和全序的前提下，
\[
  \Render^{\mathrm{LLM}}_\xi
\]
也是确定性的可审计运行时。
故求解层与表达层分别都具有程序确定性，合并即得推论。证毕。
\end{proof}

\subsection{备注：收紧版结论与一句话压缩}
\label{subsec:tex37-ch6-remarks}

\begin{remark}[收紧后的结论表述]
本章更适合采用下面这个严格版本：
\[
  \boxed{
    \text{参考 `tex36` 与本稿前文，可将 G 框架驱动的 AI 严格重写为：}
  }
\]
\[
  \boxed{
    \begin{aligned}
      &\text{G 侧程序负责全编码求解链、障碍修复、正规形提取、动作选择与执行裁决；}\\
      &\text{LLM 只负责在既定逻辑占位与审计合同下输出可读文本，因此被降级为可替换插件。}
    \end{aligned}
  }
\]
\[
  \boxed{
    \begin{aligned}
      &\text{于是系统的思考主权与决策主权，从黑箱语言模型回收到}\\
      &\text{程序化、可审计、可回放、可 fail-closed 的白盒 G 框架。}
    \end{aligned}
  }
\]
\end{remark}

\begin{remark}[一句话结论]
最短的严格压缩可写成：
\[
  \boxed{
    \begin{aligned}
      &\text{联立 `tex36` 与本稿前文之后，AI 的“大脑”不必再是 LLM；}\\
      &\text{LLM 可以被降级为“嘴”和“笔”，而真正的思考与决策回到 G 框架程序本体。}
    \end{aligned}
  }
\]
\end{remark}

\clearpage

%%%%%%%%%%%%%%%%%%%%%%%%%%%%%%%%%%%%%%%%%%%%%%%%%%%%%%%%%%%%
\section{形式逻辑：领域细化塔纤维插件原则与白盒主权合同的保守扩张}
\label{sec:tex37-ch7}
%%%%%%%%%%%%%%%%%%%%%%%%%%%%%%%%%%%%%%%%%%%%%%%%%%%%%%%%%%%%

\begin{remark}[核心陈述（承接第六章）]
一旦第~\ref{sec:tex37-ch6} 章已经证明
\[
  \boxed{
    \text{G 核心塔负责求解、修复、裁决与动作选择，LLM 只负责投影表达}
  },
\]
就可进一步得到如下结构结论：
\[
  \boxed{
    \text{在不交出思考主权与决策主权的前提下，可附加领域细化的塔纤维插件。}
  }
\]
这一步的意义不在于增加若干外置规则，而在于把“单一白盒内核”推进为
\[
  \text{白盒求解内核}
  +
  \text{可插拔领域细化丛}
  +
  \text{受审计表达插件}
\]
的严格分层体系。
\end{remark}

\begin{remark}[阅读导航：仓内引用与外部参照]
本章直接承接第~\ref{sec:tex37-ch6} 章关于主权回收、白盒可见性与双层程序确定性的结论；
对象语言方面继续依赖 `tex36` 关于修复塔、纤维丛空间、稳定正规形与路径类的记号 \cite{RepoTex36HomotopyTwistBijection}，
并参考 `tex35` 关于扩展证书塔与结构表示的口径 \cite{RepoTex35ExtendedCertificationTower}；
关于语义规范场、占位层与审计闭环，继续沿用本稿第一章及仓内整理稿 \cite{RepoTex20SemanticGaugeNormalization,RepoTex21AuditedWrapper,RepoTex23SemanticBase,
  GaoZhengAppVol2,GaoZhengAppVol3,GaoZheng2025GFramework}。

外部标准参照方面，本章对“同一基底上的纤维细化、局部联络与局部几何”语言，采用范畴/纤维丛/联络的通用口径 \cite{MacLane1998Categories,KobayashiNomizu1963Foundations,
  Hatcher2002AT}。
\end{remark}

\subsection{定义：核心塔、领域细化纤维丛与核心合同}
\label{subsec:tex37-ch7-definitions}

\begin{definition}[核心塔与领域细化塔纤维插件]
\label{def:tex37-ch7-core-plugin}
承接第~\ref{sec:tex37-ch6} 章的 G 侧全编码求解链，把其对象层统摄为一个 $N$ 阶核心修复塔
\[
  \pi_N:\mathcal E_{\mathrm{core}}^{(N)}\to \mathcal B_{\mathrm{ctx}},
\]
其中 $\mathcal B_{\mathrm{ctx}}$ 是上下文基底，$\mathcal E_{\mathrm{core}}^{(N)}$ 的纤维
\[
  \mathcal E_{\mathrm{core},b}^{(N)}:=\pi_N^{-1}(b)
\]
编码该上下文上的状态、障碍、修复链、正规形、动作接口与证书接口。

给定一个领域
\[
  D\in\mathfrak D,
\]
定义其 $N$ 阶领域细化塔纤维丛为
\[
  \pi_{N,D}:\mathcal E_D^{(N)}\to \mathcal B_{\mathrm{ctx}},
\]
并要求存在保基底态射
\[
  \iota_D:\mathcal E_{\mathrm{core}}^{(N)}\hookrightarrow \mathcal E_D^{(N)},
  \qquad
  \rho_D:\mathcal E_D^{(N)}\to \mathcal E_{\mathrm{core}}^{(N)},
\]
满足
\[
  \rho_D\circ \iota_D=\id_{\mathcal E_{\mathrm{core}}^{(N)}},
  \qquad
  \pi_N\circ \rho_D=\pi_{N,D}.
\]
称三元组
\[
  \mathcal P_D^{(N)}
  :=
  \bigl(
    \pi_{N,D},
    \iota_D,
    \rho_D
  \bigr)
\]
为一个\textbf{领域细化塔纤维插件}。
\end{definition}

\begin{definition}[领域插件的五类增量]
\label{def:tex37-ch7-five-increments}
称领域 $D$ 的五类标准增量为
\[
  \Delta_D^{\mathrm{obs}},
  \quad
  \Delta_D^{\mathrm{op}},
  \quad
  \Delta_D^{\mathrm{obst}},
  \quad
  \Delta_D^{\mathrm{act}},
  \quad
  \Delta_D^{\mathrm{cert}},
\]
分别表示领域观测变量、领域算子族、领域障碍类型、领域动作空间与领域证书接口。
若把它们并入第 $N$ 层纤维，则记
\[
  \mathcal E_D^{(N)}
  =
  \mathcal E_{\mathrm{core}}^{(N)}
  \oplus
  \Delta_D^{\mathrm{obs},(N)}
  \oplus
  \Delta_D^{\mathrm{op},(N)}
  \oplus
  \Delta_D^{\mathrm{obst},(N)}
  \oplus
  \Delta_D^{\mathrm{act},(N)}
  \oplus
  \Delta_D^{\mathrm{cert},(N)}.
\]
因此，领域插件附加的不是新的自由修辞权，而是
\[
  \boxed{
    \text{新变量、新障碍、新修复算子、新动作空间与新证书接口。}
  }
\]
\end{definition}

\begin{definition}[核心主权合同与插件服从性]
\label{def:tex37-ch7-core-contract}
定义核心主权合同为
\[
  \mathfrak C_{\mathrm{core}}
  :=
  \bigl(
    \mathfrak I_{\mathrm{state}},
    \mathfrak I_{\mathrm{obst}},
    \mathfrak I_{\mathrm{repair}},
    \mathfrak I_{\mathrm{act}},
    \mathfrak I_{\mathrm{cert}},
    \mathfrak I_{\mathrm{pl}},
    \mathfrak I_{\mathrm{audit}}
  \bigr),
\]
并约定：若领域插件 $\mathcal P_D^{(N)}$ 满足以下四项，则记为
\[
  \mathcal P_D^{(N)}\models \mathfrak C_{\mathrm{core}}.
\]
四项条件分别为：
\begin{enumerate}
  \item \textbf{保基底：}
  \[
    \pi_N\circ \rho_D=\pi_{N,D};
  \]
  \item \textbf{保核心正规形：}
  若记核心正规形与领域正规形分别为
  \[
    \operatorname{NF}_{\mathrm{core}}:\mathcal E_{\mathrm{core}}^{(N)}\to \mathcal T_{\mathrm{core}}^{(N)},
    \qquad
    \operatorname{NF}_D:\mathcal E_D^{(N)}\to \mathcal T_D^{(N)},
  \]
  则对任意 $x\in\mathcal E_{\mathrm{core}}^{(N)}$，
  \[
    \rho_D\bigl(\operatorname{NF}_D(\iota_D(x))\bigr)
    =
    \operatorname{NF}_{\mathrm{core}}(x);
  \]
  \item \textbf{保裁决接口：}
  领域动作映射
  \[
    \operatorname{Act}_D:\mathcal T_D^{(N)}\to \mathcal A_D
  \]
  必须经由 G 侧核心裁决协议约束，而不能绕开第~\ref{sec:tex37-ch6} 章的白盒裁决链；
  \item \textbf{保占位审计闭环：}
  存在领域占位映射
  \[
    \Theta_{0,D}:\mathcal T_D^{(N)}\to \mathsf{PL}_{\mathrm{NL},D}^{(0)},
  \]
  使任意非失败领域输出都进入
  \[
    \tau_D^\star
    \xrightarrow{\Theta_{0,D}}
    p_D
    \xrightarrow{\Render^{\mathrm{LLM}}_{\xi,D}}
    t_D,
  \]
  并满足
  \[
    \Parse_{\xi,D}(t_D)=p_D,
    \qquad
    \Style_{\xi,D}(t_D)=\top.
  \]
\end{enumerate}
\end{definition}

\begin{remark}[几何解释]
定义~\ref{def:tex37-ch7-core-plugin}--\ref{def:tex37-ch7-core-contract} 说明：
领域插件不是另起一套“外部系统”，而是在同一上下文基底 $\mathcal B_{\mathrm{ctx}}$ 上对纤维对象作局部细化。
因此它的数学地位首先是\textbf{同基底纤维细化}，其次才是领域知识打包。
\end{remark}

\subsection{公理（降级为定理）：领域细化塔纤维插件可按层数递归构造}
\label{subsec:tex37-ch7-axiom}

\begin{theorem}[公理降级：统一核心合同下的领域插件可按层数递归构造（数学归纳法证书）]
\label{thm:tex37-ch7-induction}
设固定一个领域 $D\in\mathfrak D$。若满足：
\begin{enumerate}
  \item 存在 $0$ 阶领域插件
  \[
    \mathcal P_D^{(0)}
  \]
  且
  \[
    \mathcal P_D^{(0)}\models \mathfrak C_{\mathrm{core}};
  \]
  \item 对任意 $n\in\NN$，只要
  \[
    \mathcal P_D^{(n)}\models \mathfrak C_{\mathrm{core}},
  \]
  就存在唯一的层扩张算子
  \[
    \operatorname{Ext}_{D,n}:\mathcal P_D^{(n)}\mapsto \mathcal P_D^{(n+1)},
  \]
  使得
  \[
    \mathcal E_D^{(n+1)}
    =
    \mathcal E_D^{(n)}
    \oplus
    \Delta_D^{\mathrm{obs},(n+1)}
    \oplus
    \Delta_D^{\mathrm{op},(n+1)}
    \oplus
    \Delta_D^{\mathrm{obst},(n+1)}
    \oplus
    \Delta_D^{\mathrm{act},(n+1)}
    \oplus
    \Delta_D^{\mathrm{cert},(n+1)},
  \]
  且
  \[
    \mathcal P_D^{(n+1)}\models \mathfrak C_{\mathrm{core}}.
  \]
\end{enumerate}
则对任意 $N\in\NN$，存在唯一的 $N$ 阶领域细化塔纤维插件
\[
  \mathcal P_D^{(N)},
\]
并且
\[
  \mathcal P_D^{(N)}\models \mathfrak C_{\mathrm{core}}.
\]
\end{theorem}

\begin{Certificate}[证书：基步 + 归纳步（数学归纳法证书）]
\label{cert:tex37-ch7-induction}
定义谓词
\[
  P(N)
  :
  \Longleftrightarrow
  \exists!\,\mathcal P_D^{(N)}
  \ \text{使得}\
  \mathcal P_D^{(N)}\models \mathfrak C_{\mathrm{core}}.
\]
则数学归纳法证书记为
\[
  \pi_{\mathrm{dom}}
  =
  \bigl(
    P(0),\
    \forall N\in\NN,\ P(N)\Rightarrow P(N+1)
  \bigr).
\]
其中归纳步由层扩张算子
\[
  \operatorname{Ext}_{D,N}
\]
与五类增量
\[
  \Delta_D^{\mathrm{obs},(N+1)},
  \Delta_D^{\mathrm{op},(N+1)},
  \Delta_D^{\mathrm{obst},(N+1)},
  \Delta_D^{\mathrm{act},(N+1)},
  \Delta_D^{\mathrm{cert},(N+1)}
\]
共同见证。
\end{Certificate}

\begin{proof}[证明（数学归纳法证书；基于数学符号推演逻辑的谓词因果链）]
定义谓词
\[
  P(N)
  :
  \Longleftrightarrow
  \exists!\,\mathcal P_D^{(N)}
  \ \text{使得}\
  \mathcal P_D^{(N)}\models \mathfrak C_{\mathrm{core}}.
\]

\textbf{基步：}
由条件(1)，$0$ 阶领域插件
\[
  \mathcal P_D^{(0)}
\]
存在且满足
\[
  \mathcal P_D^{(0)}\models \mathfrak C_{\mathrm{core}}.
\]
故 $P(0)$ 成立。

\textbf{归纳步：}
设 $P(N)$ 成立，则存在唯一的
\[
  \mathcal P_D^{(N)}
\]
满足
\[
  \mathcal P_D^{(N)}\models \mathfrak C_{\mathrm{core}}.
\]
由条件(2)，存在唯一的层扩张算子
\[
  \operatorname{Ext}_{D,N}
\]
把
\[
  \mathcal P_D^{(N)}
\]
送到
\[
  \mathcal P_D^{(N+1)}.
\]
并且该扩张只附加五类领域增量，而不改写上下文基底、核心正规形、核心裁决协议与占位审计闭环。
因此
\[
  \mathcal P_D^{(N+1)}\models \mathfrak C_{\mathrm{core}}.
\]
又由于扩张算子是唯一的，故
\[
  \mathcal P_D^{(N+1)}
\]
亦唯一，于是 $P(N+1)$ 成立。

由数学归纳法，对任意 $N\in\NN$，$P(N)$ 都成立。
因此统一核心合同下的领域细化塔纤维插件可按层数递归构造。证毕。
\end{proof}

\subsection{定理：若领域插件是保守扩张，则不会交出思考与决策主权}
\label{subsec:tex37-ch7-theorem}

\begin{theorem}[保守扩张主权不外包定理]
\label{thm:tex37-ch7-conservative}
在定理~\ref{thm:tex37-ch6-whitebox-plugin} 与命题~\ref{prop:tex37-ch6-conditional-sovereignty} 的前提下，
设某个领域插件 $\mathcal P_D^{(N)}$ 满足
\[
  \mathcal P_D^{(N)}\models \mathfrak C_{\mathrm{core}}.
\]
则附加该插件后：
\[
  \mathsf{Sov}_{\mathrm{solve}}
  \oplus
  \mathsf{Sov}_{\mathrm{decide}}
\]
仍由 G 核心程序持有，而
\[
  \mathsf{Sov}_{\mathrm{render}}
\]
仍被限制在受审计的表达层。
换言之，领域插件只能带来
\[
  \text{求解空间细化}
  +
  \text{障碍类型细化}
  +
  \text{动作空间细化}
  +
  \text{证书接口细化},
\]
而不能把思考主权与决策主权外包出去。
\end{theorem}

\begin{Certificate}[证书：保基底 + 保核心正规形 + 保裁决接口 + 保占位审计闭环]
\label{cert:tex37-ch7-conservative}
证书登记谓词链：
\[
\begin{aligned}
  C_1 &: \pi_N\circ \rho_D=\pi_{N,D},\\
  C_2 &: \rho_D\bigl(\operatorname{NF}_D(\iota_D(x))\bigr)=\operatorname{NF}_{\mathrm{core}}(x)\quad (\forall
    x\in\mathcal E_{\mathrm{core}}^{(N)}),\\
  C_3 &: \operatorname{Act}_D\ \text{受核心裁决协议约束},\\
  C_4 &: \tau_D^\star\xrightarrow{\Theta_{0,D}}p_D\xrightarrow{\Render^{\mathrm{LLM}}_{\xi,D}}t_D\
  \text{且}\ \Parse_{\xi,D}(t_D)=p_D,\ \Style_{\xi,D}(t_D)=\top,\\
  C_5 &: C_1\wedge C_2\wedge C_3\wedge C_4
  \Rightarrow
  \mathsf{Sov}_{\mathrm{solve}}\oplus \mathsf{Sov}_{\mathrm{decide}}\ \text{不外包}.
\end{aligned}
\]
\end{Certificate}

\begin{proof}[证明（基于数学符号推演逻辑的谓词因果链）]
由 $C_1$，领域插件仍附着在同一上下文基底 $\mathcal B_{\mathrm{ctx}}$ 上，
因此它没有把“问题属于谁、求解对象属于谁”改写为外部系统的输入口径。
这说明插件只是在原基底之上增加纤维分辨率，而不是迁移求解本体。

由 $C_2$，领域正规形在回投影后仍与核心正规形一致。
因此领域细化不会推翻第~\ref{sec:tex37-ch6} 章已经建立的 G 侧求解链，
它只会把核心正规形分解得更细、更贴近特定领域对象。
故
\[
  \mathsf{Sov}_{\mathrm{solve}}
\]
仍保留在 G 核心侧。

由 $C_3$，即便领域动作空间 $\mathcal A_D$ 变大了，
最终动作、停止、拒绝与执行分支仍必须经过核心裁决协议。
这说明新增的是动作候选与控制维度，而不是裁决主权。
故
\[
  \mathsf{Sov}_{\mathrm{decide}}
\]
仍保留在 G 核心侧。

由 $C_4$，领域输出的外显阶段仍服从“占位固定 $\rightarrow$ 文本渲染 $\rightarrow$ 解析回投影 $\rightarrow$ 风格审计”的闭环，
因此 LLM 在领域场景下也仍只是表达插件，而不是语义定义者或决策裁判。
这说明
\[
  \mathsf{Sov}_{\mathrm{render}}
\]
依旧位于受审计表达层。

于是，领域插件所做的只是对求解纤维、障碍类别、动作空间与证书接口的局部细化，
并未把思考主权与决策主权从 G 核心程序转移到外部黑箱。
定理成立。证毕。
\end{proof}

\subsection{引理：领域插件的正确数学地位是纤维局部几何细化，而不是外挂规则包}
\label{subsec:tex37-ch7-lemma}

\begin{lemma}[领域插件不是无结构规则列表]
\label{lem:tex37-ch7-geometry}
在定义~\ref{def:tex37-ch7-core-plugin}--\ref{def:tex37-ch7-core-contract} 的设定下，
若领域插件 $\mathcal P_D^{(N)}$ 通过五类增量改变的是每个纤维
\[
  \mathcal E_{\mathrm{core},b}^{(N)}
\]
上允许的观测方向、障碍类型、局部修复算子、动作方向与证书判定，
则其正确的数学地位不是“外挂规则包”，而是
\[
  \boxed{
    \text{同一基底上的纤维局部几何、局部联络与局部动作空间细化。}
  }
\]
\end{lemma}

\begin{Certificate}[证书：同基底纤维变化 $\Rightarrow$ 局部几何变化]
\label{cert:tex37-ch7-geometry}
证书登记谓词链：
\[
\begin{aligned}
  L_1 &: \pi_N\circ \rho_D=\pi_{N,D}\Rightarrow \mathcal E_D^{(N)}\ \text{与}\ \mathcal E_{\mathrm{core}}^{(N)}\
    \text{共享同一基底},\\
  L_2 &: \Delta_D^{\mathrm{obs}},\Delta_D^{\mathrm{op}},\Delta_D^{\mathrm{obst}},\Delta_D^{\mathrm{act}},
    \Delta_D^{\mathrm{cert}}
  \Rightarrow \text{纤维内允许方向与判定结构发生变化},\\
  L_3 &: L_2\Rightarrow \text{变化对象是局部可达路径、局部联络代理、局部控制锥与局部证书判定},\\
  L_4 &: L_1\wedge L_3\Rightarrow \mathcal P_D^{(N)}\ \text{应解释为纤维局部几何细化}.
\end{aligned}
\]
\end{Certificate}

\begin{proof}[证明（基于数学符号推演逻辑的谓词因果链）]
由 $L_1$，领域插件与核心塔共享同一上下文基底，
因此改变发生在纤维内部，而不是换了一个新的底空间或新的求解宇宙。
由 $L_2$，领域插件并非只补充若干文本规则，
而是实际改动了观测变量、可用算子、障碍类型、动作集合与证书接口。

这些变化会直接影响每个纤维上的允许路径、局部修复方向、局部控制锥、局部代价比较以及局部证书判定，
即 $L_3$。
而在纤维丛与联络的标准语言中，
对同一基底上纤维允许结构与局部联络代理的细化，
正应解释为局部几何/局部联络结构的细化 \cite{RepoTex36HomotopyTwistBijection,KobayashiNomizu1963Foundations}。
故由 $L_4$，领域插件的正确数学地位是纤维局部几何细化，而不是外接的无结构规则列表。证毕。
\end{proof}

\subsection{命题：领域插件越多，不代表系统越乱；前提是核心合同统一}
\label{subsec:tex37-ch7-proposition}

\begin{proposition}[统一核心合同下的可扩展性命题]
\label{prop:tex37-ch7-scalable}
设领域索引集为
\[
  \mathfrak D_* \subseteq \mathfrak D.
\]
若对每个
\[
  D\in\mathfrak D_*,
\]
都有
\[
  \mathcal P_D^{(N)}\models \mathfrak C_{\mathrm{core}},
\]
则随着领域插件数量增加，系统增加的是覆盖域
\[
  \operatorname{Cov}(\mathfrak D_*),
\]
而不是主权失控。
更准确地说：
\[
  \boxed{
    \text{统一核心合同}+\text{多领域细化插件}
    \Rightarrow
    \text{覆盖域扩张而非裁决权分裂。}
  }
\]
\end{proposition}

\begin{Certificate}[证书：统一合同 + 跨插件一致性 + 核心裁决保留]
\label{cert:tex37-ch7-scalable}
证书登记谓词链：
\[
\begin{aligned}
  P_1 &: \forall D\in\mathfrak D_*,\ \mathcal P_D^{(N)}\models \mathfrak C_{\mathrm{core}},\\
  P_2 &: P_1\Rightarrow \text{所有插件共享同一状态接口、障碍接口、修复接口、动作接口、}\\
  &\qquad \text{证书接口、占位接口与审计接口},\\
  P_3 &: P_2\Rightarrow \text{跨插件一致性由核心合同维持，而不由插件彼此争夺裁决权决定},\\
  P_4 &: P_3\Rightarrow \text{插件数增加只扩大 }\operatorname{Cov}(\mathfrak D_*),\\
  P_5 &: P_4\Rightarrow \mathsf{Sov}_{\mathrm{solve}}\oplus\mathsf{Sov}_{\mathrm{decide}}\\
  &\qquad \text{不因插件增多而失控}.
\end{aligned}
\]
\end{Certificate}

\begin{proof}[证明（基于数学符号推演逻辑的谓词因果链）]
由 $P_1$，每个领域插件都必须服从同一核心合同，
因此它们对外暴露的接口类型是一致的。
这立即给出 $P_2$。

由 $P_2$，插件之间的组合并不是“谁先说服系统就谁拍板”，
而是先把领域差异压进各自插件内部，再通过统一合同接入核心裁决链。
故跨插件一致性来自核心合同，而不是来自多个黑箱之间的博弈，这就是 $P_3$。

一旦 $P_3$ 成立，新增插件只会给系统带来新的可覆盖领域、新的纤维细化与新的证书接口，
而不会产生新的主权中心。
故得到 $P_4$。
再结合第~\ref{sec:tex37-ch6} 章的主权回收结构，立得 $P_5$。
命题成立。证毕。
\end{proof}

\subsection{推论：系统总体结构提升为“核心核 + 领域丛 + 表达插件”的分层白盒体系}
\label{subsec:tex37-ch7-corollary}

\begin{corollary}[插件化白盒体系推论]
\label{cor:tex37-ch7-architecture}
在定理~\ref{thm:tex37-ch7-conservative} 与命题~\ref{prop:tex37-ch7-scalable} 的条件下，
系统总体结构可写成
\[
  \mathcal G_{\mathrm{sys}}
  =
  \mathcal G_{\mathrm{core}}
  \oplus
  \Bigl(
    \bigoplus_{D\in\mathfrak D_*}\mathcal P_D
  \Bigr)
  \oplus
  \mathcal P_{\mathrm{render}},
\]
其中
\[
  \mathcal G_{\mathrm{core}}
\]
是核心 G 求解/修复/裁决核，
\[
  \mathcal P_D
\]
是领域细化塔纤维插件，
\[
  \mathcal P_{\mathrm{render}}
\]
是全局表达插件。

并且每个领域插件还可继续分层为
\[
  \mathcal P_D
  =
  \mathcal P_D^{\mathrm{obs}}
  \oplus
  \mathcal P_D^{\mathrm{op}}
  \oplus
  \mathcal P_D^{\mathrm{obst}}
  \oplus
  \mathcal P_D^{\mathrm{repair}}
  \oplus
  \mathcal P_D^{\mathrm{act}}
  \oplus
  \mathcal P_D^{\mathrm{cert}}
  \oplus
  \mathcal P_D^{\mathrm{render}},
\]
其中最后一项 $\mathcal P_D^{\mathrm{render}}$ 只表示\textbf{领域术语模板/报告风格适配子插件}，
仍须服从全局的占位--解析--审计闭环，而不是独立语义内核。
\end{corollary}

\begin{Certificate}[证书：主权保留 + 统一合同 + 可分层细化]
\label{cert:tex37-ch7-architecture}
证书登记谓词链：
\[
\begin{aligned}
  C_1 &: \text{定理~\ref{thm:tex37-ch7-conservative}}\Rightarrow
  \mathcal G_{\mathrm{core}}\ \text{保留求解与裁决主权},\\
  C_2 &: \text{命题~\ref{prop:tex37-ch7-scalable}}\Rightarrow
  \bigoplus_{D\in\mathfrak D_*}\mathcal P_D\ \text{只扩张覆盖域而不分裂主权},\\
  C_3 &: \text{定义~\ref{def:tex37-ch7-five-increments}}\Rightarrow
  \mathcal P_D\ \text{可按观测、算子、障碍、修复、动作、证书与术语适配继续分层},\\
  C_4 &: C_1\wedge C_2\wedge C_3\Rightarrow
  \mathcal G_{\mathrm{sys}}
  =
  \mathcal G_{\mathrm{core}}
  \oplus
  \bigl(\bigoplus_{D\in\mathfrak D_*}\mathcal P_D\bigr)
  \oplus
  \mathcal P_{\mathrm{render}}.
\end{aligned}
\]
\end{Certificate}

\begin{proof}[证明（基于数学符号推演逻辑的谓词因果链）]
由 $C_1$，求解与裁决主权已经被固定在核心 G 程序内，
因此系统必须以
\[
  \mathcal G_{\mathrm{core}}
\]
作为唯一白盒主核。
由 $C_2$，多领域插件的并入并不会制造新的主权中心，
它们只是在核心核周围增加领域细化丛。

再由 $C_3$，单个领域插件本身并非不可分的黑箱，
而是可以沿观测、算子、障碍、修复、动作、证书与术语适配继续分层。
因此系统结构自然分解为
\[
  \mathcal G_{\mathrm{core}}
  +
  \text{领域丛插件}
  +
  \text{表达插件},
\]
即 $C_4$ 所给出的分层白盒体系。
推论成立。证毕。
\end{proof}

\subsection{备注：结构意义、边界与一句话结论}
\label{subsec:tex37-ch7-remarks}

\begin{remark}[这一步把 G 框架推进到“操作系统内核”地位]
在本章结构下，更贴切的系统类比不再是
\[
  \text{模型本体}+\text{零散工具},
\]
而是
\[
  \boxed{
    \text{白盒求解内核}+\text{领域驱动插件}+\text{表达层插件}.
  }
\]
其中核心 G 框架类似内核，
各领域塔纤维插件类似驱动/子系统，
而 LLM 及其术语模板适配只类似可替换的显示与交互层。
\end{remark}

\begin{remark}[最重要的边界]
“可附加领域细化的塔纤维插件”这一命题成立，但必须守住下述硬边界：
\[
  \boxed{
    \text{插件可以细化纤维，不能篡改核心主权合同。}
  }
\]
一旦某个插件试图绕开核心正规形、核心裁决协议或占位审计闭环，
它就不再是插件，而是在偷换求解核。
\end{remark}

\begin{remark}[一句话结论]
本章最短的严格压缩可写成：
\[
  \boxed{
    \begin{aligned}
      &\text{在 G 框架的白盒主权结构下，完全可以附加领域细化的塔纤维插件；}\\
      &\text{它们负责把同一核心塔在不同领域上做局部几何、障碍、动作与证书的细化，}\\
      &\text{但不交出思考主权与决策主权。}
    \end{aligned}
  }
\]
\end{remark}

\clearpage

%%%%%%%%%%%%%%%%%%%%%%%%%%%%%%%%%%%%%%%%%%%%%%%%%%%%%%%%%%%%
\section{形式逻辑：领域细化塔纤维插件与证书认证富结构对象的自然对应}
\label{sec:tex37-ch8}
%%%%%%%%%%%%%%%%%%%%%%%%%%%%%%%%%%%%%%%%%%%%%%%%%%%%%%%%%%%%

\begin{remark}[核心陈述（在第七章基础上再收紧）]
第~\ref{sec:tex37-ch7} 章已经证明：领域细化塔纤维插件不是外加规则包，而是同一上下文基底上的纤维局部几何细化。
本章进一步把这一结论收紧为更强的结构命题：
\[
  \boxed{
    \begin{aligned}
      &\text{领域细化的塔纤维插件，不只是“可附加模块”，}\\
      &\text{而是可与“基于证书认证的富结构对象”自然对应。}
    \end{aligned}
  }
\]
因此，插件的更准确数学身份不应只写成
\[
  \mathcal P_D,
\]
而应写成
\[
  \widehat{\mathfrak R}_D,
\]
其中“帽子”表示：该对象不是裸的富结构，而是带证书、带认证、带审计闭环的富结构。
\end{remark}

\begin{remark}[阅读导航：承接位置与引用口径]
本章对富结构对象的裸结构部分，直接承接 `tex36` 关于“主纤维丛版广义非交换李代数的富结构对象”的记号与结论 \cite{RepoTex36HomotopyTwistBijection}；
对认证/签发/补证口径，直接承接 `tex35` 的扩展签发认证协议 \cite{RepoTex35ExtendedCertificationTower}；
对占位、解析、审计与运行时降级链，继续承接本稿第一章与第六章、第七章的结果 \cite{RepoTex20SemanticGaugeNormalization,RepoTex21AuditedWrapper,
  RepoTex23SemanticBase,GaoZhengAppVol2,GaoZhengAppVol3,GaoZheng2025GFramework}。

本章的目标不是再造一套平行对象宇宙，而是把
\[
  \text{富结构}
  +
  \text{证书认证}
  +
  \text{占位审计闭环}
\]
压成同一个严格对象。
\end{remark}

\subsection{定义：裸富结构、认证富结构与认证的严格含义}
\label{subsec:tex37-ch8-definitions}

\begin{definition}[裸富结构与认证富结构]
\label{def:tex37-ch8-bare-rich}
承接 `tex36` 关于富结构对象的整理口径 \cite{RepoTex36HomotopyTwistBijection}，
定义一个裸富结构
\[
  \mathfrak R
  :=
  \Bigl(
    \{B_n,P_n,\pi_n,A_n,\Omega_n\}_{n\in\NN},
    \{\mathcal F_n,\mathcal{PI}_n,\mathcal T_n\}_{n\in\NN},
    \{\varepsilon_n,\Pi_n^m\}_{n<m},
    \{[\cdot,\cdot]_n,\Ob_n\}_{n\in\NN},
    \mathfrak c_{\tau}
  \Bigr).
\]
称其为“裸”的理由是：它虽然已经同时携带
\[
  \text{拓扑数据}
  +
  \text{代数数据}
  +
  \text{语义编码数据},
\]
但尚未把证书合法性、动作合法性与表达合法性并入同一对象。

定义一个认证数据包为
\[
  \Cert
  :=
  \bigl(
    \Cert^{\mathrm{norm}},
    \Cert^{\mathrm{repair}},
    \Cert^{\mathrm{equiv}},
    \Cert^{\mathrm{act}},
    \Cert^{\mathrm{render}}
  \bigr),
\]
其中：
\begin{enumerate}
  \item $\Cert^{\mathrm{norm}}$ 认证对象的类型、来源、量纲、索引、层级、边界条件与正规化入口；
  \item $\Cert^{\mathrm{repair}}$ 认证修复链的合法性与障碍压低性；
  \item $\Cert^{\mathrm{equiv}}$ 认证与核心正规形、核心合同和回投影关系的相容性；
  \item $\Cert^{\mathrm{act}}$ 认证动作来自合法正规形与合法裁决接口；
  \item $\Cert^{\mathrm{render}}$ 认证占位到文本的表达闭环满足解析、风格与审计要求。
\end{enumerate}

再定义认证谓词族
\[
  \Auth
  :=
  \bigl(
    \Auth^{\mathrm{obj}},
    \Auth^{\mathrm{mor}},
    \Auth^{\mathrm{repair}},
    \Auth^{\mathrm{act}},
    \Auth^{\mathrm{render}}
  \bigr),
\]
其中每个分量都是可判定的合法性谓词，取值于 $\{\top,\bot\}$。
于是称
\[
  \widehat{\mathfrak R}
  :=
  (\mathfrak R,\Cert,\Auth)
\]
为一个\textbf{认证富结构}。
\end{definition}

\begin{definition}[领域细化塔纤维插件的认证化]
\label{def:tex37-ch8-plugin-auth}
对任意领域 $D\in\mathfrak D$，承接第~\ref{sec:tex37-ch7} 章，记其领域细化塔纤维插件为
\[
  \mathcal P_D
  :=
  \Bigl(
    \Delta_D^{\mathrm{obs}},
    \Delta_D^{\mathrm{op}},
    \Delta_D^{\mathrm{obst}},
    \Delta_D^{\mathrm{act}},
    \Delta_D^{\mathrm{cert}}
  \Bigr).
\]
记由第~\ref{sec:tex37-ch7} 章核心塔诱导出的核心富结构为
\[
  \mathfrak R_{\mathrm{core}}.
\]
把插件与核心富结构合并，定义领域裸富结构
\[
  \mathfrak R_D
  :=
  \mathfrak R_{\mathrm{core}}
  \oplus
  \mathcal P_D.
\]

进一步，要求插件中的新增部分带有证书附着接口
\[
  \Delta_D^{\mathrm{obs}}\rightsquigarrow \Cert_D^{\mathrm{norm}},
  \qquad
  \Delta_D^{\mathrm{op}}\rightsquigarrow \Cert_D^{\mathrm{equiv}},
  \qquad
  \Delta_D^{\mathrm{obst}}\rightsquigarrow \Cert_D^{\mathrm{repair}},
\]
\[
  \Delta_D^{\mathrm{act}}\rightsquigarrow \Cert_D^{\mathrm{act}},
  \qquad
  \Delta_D^{\mathrm{cert}}\rightsquigarrow \Cert_D^{\mathrm{render}},
\]
并令
\[
  \Cert_D
  :=
  \bigl(
    \Cert_D^{\mathrm{norm}},
    \Cert_D^{\mathrm{repair}},
    \Cert_D^{\mathrm{equiv}},
    \Cert_D^{\mathrm{act}},
    \Cert_D^{\mathrm{render}}
  \bigr),
\]
\[
  \Auth_D
  :=
  \bigl(
    \Auth_D^{\mathrm{obj}},
    \Auth_D^{\mathrm{mor}},
    \Auth_D^{\mathrm{repair}},
    \Auth_D^{\mathrm{act}},
    \Auth_D^{\mathrm{render}}
  \bigr).
\]
于是定义该领域的认证富结构为
\[
  \widehat{\mathfrak R}_D
  :=
  (\mathfrak R_D,\Cert_D,\Auth_D).
\]
这说明：领域插件的正确写法，不是“往纤维里塞几个新对象”，而是
\[
  \boxed{
    \text{把纤维细化为可认证、可签发、可审计的富结构层。}
  }
\]
\end{definition}

\begin{definition}[认证的严格含义]
\label{def:tex37-ch8-auth-meaning}
称领域 $D$ 上的一条执行链
\[
  x
  \xrightarrow{f}
  x'
  \xrightarrow{\operatorname{Rep}_n=\operatorname{NF}_{A_n}}
  \tau^\star
  \xrightarrow{\operatorname{Act}_D}
  a^\star
  \xrightarrow{\Theta_{0,D}}
  p_D
  \xrightarrow{\Render^{\mathrm{LLM}}_{\xi,D}}
  t
\]
为\textbf{严格认证的}，若存在证书
\[
  \begin{aligned}
    &c_x\in \Cert_D^{\mathrm{norm}},
    \qquad
    c_f\in \Cert_D^{\mathrm{equiv}},
    \qquad
    c_{\mathrm{rep}}\in \Cert_D^{\mathrm{repair}},\\
    &c_{\mathrm{act}}\in \Cert_D^{\mathrm{act}},
    \qquad
    c_{\mathrm{render}}\in \Cert_D^{\mathrm{render}}
  \end{aligned}
\]
使得以下可判定链成立：
\[
  \Auth_D^{\mathrm{obj}}(x)=\top,
  \qquad
  \Auth_D^{\mathrm{mor}}(f)=\top,
  \qquad
  \Auth_D^{\mathrm{repair}}(x,\tau^\star)=\top,
\]
\[
  \Auth_D^{\mathrm{act}}(a^\star)=\top,
  \qquad
  \Auth_D^{\mathrm{render}}(t)=\top.
\]
并且末端必须满足
\[
  \Parse_{\xi,D}(t)=p_D,
  \qquad
  \Style_{\xi,D}(t)=\top,
  \qquad
  t\in \Audit_{\xi,D}(p_D).
\]
因此这里的“认证”不是口头背书，而是
\[
  \text{对象证书}
  \Rightarrow
  \text{态射/相容证书}
  \Rightarrow
  \text{修复证书}
  \Rightarrow
  \text{动作证书}
  \Rightarrow
  \text{表达证书}
\]
的可判定闭环。
\end{definition}

\begin{remark}[与第七章的关系]
第~\ref{sec:tex37-ch7} 章说明插件是纤维局部几何细化；
本章则说明：一旦把第六章的主权合同与第一章的占位审计闭环压进去，
这种细化不再只是“几何更细”，而是“几何更细且每一步都有证书与认证”。
故插件的最终对象身份应收紧为认证富结构。
\end{remark}

\subsection{公理（降级为定理）：认证富结构可按塔层数递归构造}
\label{subsec:tex37-ch8-axiom}

\begin{theorem}[公理降级：领域认证富结构可按塔层数递归构造（数学归纳法证书）]
\label{thm:tex37-ch8-induction}
设固定领域 $D\in\mathfrak D$。若满足：
\begin{enumerate}
  \item 存在 $0$ 阶领域裸富结构
  \[
    \mathfrak R_D^{(0)}
    =
    \mathfrak R_{\mathrm{core}}^{(0)}
    \oplus
    \mathcal P_D^{(0)},
  \]
  以及认证数据
  \[
    \Cert_D^{(0)},\ \Auth_D^{(0)},
  \]
  使得
  \[
    \widehat{\mathfrak R}_D^{(0)}
    :=
    \bigl(
      \mathfrak R_D^{(0)},
      \Cert_D^{(0)},
      \Auth_D^{(0)}
    \bigr)
  \]
  已被合法签发；
  \item 对任意 $n\in\NN$，只要
  \[
    \widehat{\mathfrak R}_D^{(n)}
  \]
  已定义并满足核心合同兼容，就存在唯一的层扩张
  \[
    \operatorname{Lift}_{D,n}:
    \widehat{\mathfrak R}_D^{(n)}
    \mapsto
    \widehat{\mathfrak R}_D^{(n+1)},
  \]
  该扩张同时增加
  \[
    \Delta_D^{\mathrm{obs},(n+1)},
    \Delta_D^{\mathrm{op},(n+1)},
    \Delta_D^{\mathrm{obst},(n+1)},
    \Delta_D^{\mathrm{act},(n+1)},
    \Delta_D^{\mathrm{cert},(n+1)},
  \]
  及其对应证书和认证谓词，并保持：
  \[
    \mathcal P_D^{(n+1)}\models \mathfrak C_{\mathrm{core}},
    \qquad
    \Parse_{\xi,D}(t)=p_D,\ \Style_{\xi,D}(t)=\top,\ t\in\Audit_{\xi,D}(p_D).
  \]
\end{enumerate}
则对任意 $N\in\NN$，存在唯一的 $N$ 阶领域认证富结构
\[
  \widehat{\mathfrak R}_D^{(N)}.
\]
\end{theorem}

\begin{Certificate}[证书：基步 + 归纳步（数学归纳法证书）]
\label{cert:tex37-ch8-induction}
定义谓词
\[
  P(N)
  :
  \Longleftrightarrow
  \exists!\,\widehat{\mathfrak R}_D^{(N)}
  \ \text{使其兼容核心合同且具有完整认证闭环}.
\]
则数学归纳法证书记为
\[
  \pi_{\mathrm{auth}}
  =
  \bigl(
    P(0),\
    \forall N\in\NN,\ P(N)\Rightarrow P(N+1)
  \bigr).
\]
其中归纳步的见证对象为
\[
  \operatorname{Lift}_{D,N},
  \quad
  \Delta_D^{\mathrm{obs},(N+1)},
  \Delta_D^{\mathrm{op},(N+1)},
  \Delta_D^{\mathrm{obst},(N+1)},
  \Delta_D^{\mathrm{act},(N+1)},
  \Delta_D^{\mathrm{cert},(N+1)},
\]
及其对应的
\[
  \Cert_D^{(N+1)},
  \quad
  \Auth_D^{(N+1)}.
\]
\end{Certificate}

\begin{proof}[证明（数学归纳法证书；基于数学符号推演逻辑的谓词因果链）]
定义谓词
\[
  P(N)
  :
  \Longleftrightarrow
  \exists!\,\widehat{\mathfrak R}_D^{(N)}
  \ \text{使其兼容核心合同且具有完整认证闭环}.
\]

\textbf{基步：}
由条件(1)，$0$ 阶对象
\[
  \widehat{\mathfrak R}_D^{(0)}
  =
  \bigl(
    \mathfrak R_D^{(0)},
    \Cert_D^{(0)},
    \Auth_D^{(0)}
  \bigr)
\]
已经存在并被合法签发。
因此 $P(0)$ 成立。

\textbf{归纳步：}
设 $P(N)$ 成立，则存在唯一的
\[
  \widehat{\mathfrak R}_D^{(N)}
\]
满足核心合同兼容与认证闭环。
由条件(2)，存在唯一的层扩张
\[
  \operatorname{Lift}_{D,N}
\]
把
\[
  \widehat{\mathfrak R}_D^{(N)}
\]
提升为
\[
  \widehat{\mathfrak R}_D^{(N+1)}.
\]
该提升不仅附加新的领域对象、算子、障碍、动作和证书接口，
还同时附加对应的证书数据与认证谓词，并保持核心合同、核心正规形兼容以及占位--解析--审计闭环。
因此
\[
  \widehat{\mathfrak R}_D^{(N+1)}
\]
存在。
又由于扩张算子是唯一的，故其亦唯一，于是 $P(N+1)$ 成立。

由数学归纳法，对任意 $N\in\NN$，$P(N)$ 都成立。
因此领域认证富结构可按塔层数递归构造。证毕。
\end{proof}

\subsection{定理：领域细化塔纤维插件与证书认证富结构对象存在自然对应}
\label{subsec:tex37-ch8-theorem}

\begin{theorem}[证书认证富结构对应定理]
\label{thm:tex37-ch8-correspondence}
对任意领域 $D\in\mathfrak D$，定义
\[
  \mathsf{Plug}_D^{\mathrm{auth}}
  :=
  \left\{
    \mathcal P_D\
    \middle|\
    \begin{aligned}
      &\mathcal P_D\models \mathfrak C_{\mathrm{core}},\\
      &\mathcal P_D\ \text{给每个新增对象、态射、障碍、动作与输出配置证书接口},\\
      &\mathcal P_D\ \text{服从占位--解析--审计闭环},\\
      &\mathcal P_D\ \text{的动作提交权与裁决权不外包给 LLM 或其他黑箱}
    \end{aligned}
  \right\},
\]
\[
  \widehat{\mathsf{Rich}}_D^{\mathrm{auth}}
  :=
  \left\{
    \widehat{\mathfrak R}_D=(\mathfrak R_D,\Cert_D,\Auth_D)\
    \middle|\
    \begin{aligned}
      &\mathfrak R_D=\mathfrak R_{\mathrm{core}}\oplus \mathcal P_D,\\
      &(\mathfrak R_D,\Cert_D,\Auth_D)\\
      &\qquad \text{兼容核心合同并具有完整认证闭环}
    \end{aligned}
  \right\}.
\]
则存在两条典范映射
\[
  \Phi_D:
  \mathsf{Plug}_D^{\mathrm{auth}}
  \to
  \widehat{\mathsf{Rich}}_D^{\mathrm{auth}},
  \qquad
  \Phi_D(\mathcal P_D)
  :=
  \bigl(
    \mathfrak R_{\mathrm{core}}\oplus \mathcal P_D,\Cert_D,\Auth_D
  \bigr),
\]
\[
  \Psi_D:
  \widehat{\mathsf{Rich}}_D^{\mathrm{auth}}
  \to
  \mathsf{Plug}_D^{\mathrm{auth}},
  \qquad
  \Psi_D(\widehat{\mathfrak R}_D)
  :=
  \mathfrak R_D\ominus \mathfrak R_{\mathrm{core}},
\]
其中
\[
  \mathfrak R_D\ominus \mathfrak R_{\mathrm{core}}
\]
表示按直和分解抽取的非核心增量部分。

并且
\[
  \Psi_D\circ \Phi_D
  =
  \id_{\mathsf{Plug}_D^{\mathrm{auth}}},
  \qquad
  \Phi_D\circ \Psi_D
  =
  \id_{\widehat{\mathsf{Rich}}_D^{\mathrm{auth}}}.
\]
因此，
\[
  \mathsf{Plug}_D^{\mathrm{auth}}
  \cong
  \widehat{\mathsf{Rich}}_D^{\mathrm{auth}},
\]
即领域细化塔纤维插件与证书认证富结构对象存在自然对应。
\end{theorem}

\begin{Certificate}[证书：插件到富结构 + 富结构回抽插件 + 双向恒等]
\label{cert:tex37-ch8-correspondence}
证书登记谓词链：
\[
\begin{aligned}
  C_1 &: \mathcal P_D\models \mathfrak C_{\mathrm{core}}
  \Rightarrow
  \mathfrak R_{\mathrm{core}}\oplus \mathcal P_D\ \text{是合法领域裸富结构};\\
  C_2 &: C_1\wedge \text{完整证书接口}
  \Rightarrow
  \Phi_D(\mathcal P_D)=\widehat{\mathfrak R}_D\in \widehat{\mathsf{Rich}}_D^{\mathrm{auth}};\\
  C_3 &: \widehat{\mathfrak R}_D\in \widehat{\mathsf{Rich}}_D^{\mathrm{auth}}
  \Rightarrow
  \mathfrak R_D\ominus \mathfrak R_{\mathrm{core}}\ \text{是良定义的插件增量};\\
  C_4 &: C_3\Rightarrow \Psi_D(\widehat{\mathfrak R}_D)\in \mathsf{Plug}_D^{\mathrm{auth}};\\
  C_5 &: C_2\wedge C_4\Rightarrow \Psi_D(\Phi_D(\mathcal P_D))=\mathcal P_D;\\
  C_6 &: C_2\wedge C_4\Rightarrow \Phi_D(\Psi_D(\widehat{\mathfrak R}_D))=\widehat{\mathfrak R}_D.
\end{aligned}
\]
\end{Certificate}

\begin{proof}[证明（基于数学符号推演逻辑的谓词因果链）]
\textbf{从插件到认证富结构：}
给定
\[
  \mathcal P_D\in \mathsf{Plug}_D^{\mathrm{auth}}.
\]
由 $C_1$，它与核心富结构
\[
  \mathfrak R_{\mathrm{core}}
\]
合并后得到
\[
  \mathfrak R_D
  =
  \mathfrak R_{\mathrm{core}}
  \oplus
  \mathcal P_D.
\]
由 $\mathcal P_D$ 的定义，它附加了新的观测、算子、障碍、动作与证书接口；
由其属于 $\mathsf{Plug}_D^{\mathrm{auth}}$，这些新增部分又都带有证书接口、认证谓词与占位审计闭环。
故可定义
\[
  \Phi_D(\mathcal P_D)
  :=
  (\mathfrak R_D,\Cert_D,\Auth_D)
  \in
  \widehat{\mathsf{Rich}}_D^{\mathrm{auth}},
\]
这就是 $C_2$。

\textbf{从认证富结构回抽插件：}
给定
\[
  \widehat{\mathfrak R}_D=(\mathfrak R_D,\Cert_D,\Auth_D)
  \in
  \widehat{\mathsf{Rich}}_D^{\mathrm{auth}}.
\]
由其定义，存在核心兼容的直和分解
\[
  \mathfrak R_D
  =
  \mathfrak R_{\mathrm{core}}
  \oplus
  \Delta_D,
\]
其中 $\Delta_D$ 精确收集了所有相对于核心新增的观测变量、算子族、障碍类型、动作空间与证书接口。
于是定义
\[
  \Psi_D(\widehat{\mathfrak R}_D)
  :=
  \mathfrak R_D\ominus \mathfrak R_{\mathrm{core}}
  =
  \Delta_D.
\]
由 $C_3$，这一差分是良定义的；又由核心合同兼容、动作不越权以及占位审计闭环，得
\[
  \Delta_D
  \in
  \mathsf{Plug}_D^{\mathrm{auth}},
\]
即 $C_4$。

\textbf{双向恒等：}
对任意
\[
  \mathcal P_D\in \mathsf{Plug}_D^{\mathrm{auth}},
\]
先应用 $\Phi_D$ 再应用 $\Psi_D$，得到
\[
  \Psi_D(\Phi_D(\mathcal P_D))
  =
  \bigl(\mathfrak R_{\mathrm{core}}\oplus \mathcal P_D\bigr)
  \ominus
  \mathfrak R_{\mathrm{core}}
  =
  \mathcal P_D.
\]
故 $C_5$ 成立。

对任意
\[
  \widehat{\mathfrak R}_D=(\mathfrak R_D,\Cert_D,\Auth_D)
  \in
  \widehat{\mathsf{Rich}}_D^{\mathrm{auth}},
\]
先应用 $\Psi_D$ 再应用 $\Phi_D$，会把
\[
  \mathfrak R_D\ominus \mathfrak R_{\mathrm{core}}
\]
重新与核心富结构拼回，同时保留原有的证书数据与认证谓词。
因此
\[
  \Phi_D(\Psi_D(\widehat{\mathfrak R}_D))
  =
  (\mathfrak R_D,\Cert_D,\Auth_D)
  =
  \widehat{\mathfrak R}_D.
\]
故 $C_6$ 成立。

于是 $\Phi_D$ 与 $\Psi_D$ 互为逆映射，
领域细化塔纤维插件与证书认证富结构对象之间不是松散类比，而是严格的自然对应。证毕。
\end{proof}

\subsection{引理：认证富结构回抽插件的差分算子是良定义的}
\label{subsec:tex37-ch8-lemma}

\begin{lemma}[差分抽取良定义引理]
\label{lem:tex37-ch8-difference}
设
\[
  \widehat{\mathfrak R}_D=(\mathfrak R_D,\Cert_D,\Auth_D)
  \in
  \widehat{\mathsf{Rich}}_D^{\mathrm{auth}}.
\]
若
\[
  \mathfrak R_D
  =
  \mathfrak R_{\mathrm{core}}
  \oplus
  \Delta_D
\]
是与核心合同相容的直和分解，则差分对象
\[
  \mathfrak R_D\ominus \mathfrak R_{\mathrm{core}}
\]
唯一确定，且恰好等于
\[
  \Delta_D
  =
  \Bigl(
    \Delta_D^{\mathrm{obs}},
    \Delta_D^{\mathrm{op}},
    \Delta_D^{\mathrm{obst}},
    \Delta_D^{\mathrm{act}},
    \Delta_D^{\mathrm{cert}}
  \Bigr).
\]
故从认证富结构回抽领域插件不存在“抽错层”或“抽出黑箱主权”的歧义。
\end{lemma}

\begin{Certificate}[证书：直和分解唯一性 + 核心合同边界]
\label{cert:tex37-ch8-difference}
证书登记谓词链：
\[
\begin{aligned}
  L_1 &: \mathfrak R_D=\mathfrak R_{\mathrm{core}}\oplus \Delta_D,\\
  L_2 &: \Delta_D\ \text{仅含相对核心新增的观测、算子、障碍、动作与证书接口},\\
  L_3 &: \Delta_D\ \text{不包含核心基底、核心正规形、核心裁决协议与占位审计闭环的所有权},\\
  L_4 &: L_1\wedge L_2\wedge L_3\Rightarrow \mathfrak R_D\ominus \mathfrak R_{\mathrm{core}}\ \text{唯一良定义}.
\end{aligned}
\]
\end{Certificate}

\begin{proof}[证明（基于数学符号推演逻辑的谓词因果链）]
由 $L_1$，领域认证富结构的裸结构部分已经给出一个与核心富结构的直和分解。
由 $L_2$，非核心部分只收集新增的观测、算子、障碍、动作与证书接口，
因此它是一个纯增量对象，而不是另一套平行核心。

再由 $L_3$，核心基底、核心正规形、核心裁决协议与占位审计闭环的所有权没有被并入差分对象，
所以差分抽取不会把主权性对象错误地纳入插件。
因此从
\[
  \mathfrak R_D
\]
中扣除
\[
  \mathfrak R_{\mathrm{core}}
\]
后，剩余部分唯一确定为
\[
  \Delta_D.
\]
即
\[
  \mathfrak R_D\ominus \mathfrak R_{\mathrm{core}}
  =
  \Delta_D.
\]
故差分算子良定义。证毕。
\end{proof}

\subsection{命题：LLM 不是认证富结构的核心部分，只是末端外显接口插件}
\label{subsec:tex37-ch8-proposition}

\begin{proposition}[LLM 不属于认证富结构本体核心]
\label{prop:tex37-ch8-llm-not-core}
在任意
\[
  \widehat{\mathfrak R}_D
  =
  (\mathfrak R_D,\Cert_D,\Auth_D)
  \in
  \widehat{\mathsf{Rich}}_D^{\mathrm{auth}}
\]
中，LLM 最多只出现在末端映射
\[
  p_D
  \xrightarrow{\Render^{\mathrm{LLM}}_{\xi,D}}
  t
\]
里，并且该映射必须服从
\[
  \Parse_{\xi,D}(t)=p_D,
  \qquad
  \Style_{\xi,D}(t)=\top,
  \qquad
  t\in \Audit_{\xi,D}(p_D).
\]
因此 LLM 不是
\[
  \widehat{\mathfrak R}_D
\]
的核心结构分量，而只是认证富结构的一个外显接口插件。
\end{proposition}

\begin{Certificate}[证书：求解/动作证书在前，渲染证书在后]
\label{cert:tex37-ch8-llm-not-core}
证书登记谓词链：
\[
\begin{aligned}
  P_1 &: \mathfrak R_D\ \text{的核心分量是对象、联络、曲率、括号、障碍、修复、正规形与动作接口},\\
  P_2 &: \Cert_D^{\mathrm{norm}},\Cert_D^{\mathrm{repair}},\Cert_D^{\mathrm{equiv}},\Cert_D^{\mathrm{act}}\\
  &\qquad \text{在文本渲染前已决定对象合法性、修复合法性与动作合法性},\\
  P_3 &: \Render^{\mathrm{LLM}}_{\xi,D}\ \text{只出现在}\ \Cert_D^{\mathrm{render}}\ \text{对应的末端外显阶段},\\
  P_4 &: P_2\wedge P_3\Rightarrow \text{LLM 不参与富结构核心本体的决定},\\
  P_5 &: P_4\Rightarrow \text{LLM 只是外显接口插件}.
\end{aligned}
\]
\end{Certificate}

\begin{proof}[证明（基于数学符号推演逻辑的谓词因果链）]
由 $P_1$，认证富结构的核心对象已经由几何、代数、障碍、修复、正规形与动作接口构成。
由 $P_2$，对象是否合法、修复是否合法、动作是否合法，都在文本渲染之前就已由证书与认证谓词决定。

由 $P_3$，LLM 只在
\[
  p_D\mapsto t
\]
这一末端投影阶段出现，
并且它的输出仍受
\[
  \Parse_{\xi,D},\ \Style_{\xi,D},\ \Audit_{\xi,D}
\]
三重约束。
因此 LLM 并不决定对象本体、修复本体或动作本体，只决定已给定占位的表面文本外观。
于是 $P_4$ 与 $P_5$ 同时成立。
命题成立。证毕。
\end{proof}

\subsection{推论：所谓“富”，不再只是多，而是“多且被认证”}
\label{subsec:tex37-ch8-corollary-rich}

\begin{corollary}[可认证地丰富]
\label{cor:tex37-ch8-certified-rich}
在定理~\ref{thm:tex37-ch8-correspondence} 的条件下，
领域插件对应的富结构其“富”不应理解为
\[
  \text{变量多、规则多、接口多},
\]
而应严格理解为
\[
  \boxed{
    \text{多层对象}
    +
    \text{多层态射}
    +
    \text{多层障碍}
    +
    \text{多层动作}
    +
    \text{多层编码}
  }
\]
并且
\[
  \boxed{
    \text{这些层全部带证书认证与审计闭环。}
  }
\]
故其更准确名称不是“丰富”，而是“可认证地丰富”。
\end{corollary}

\begin{Certificate}[证书：富结构层次 + 全链认证]
\label{cert:tex37-ch8-certified-rich}
证书登记谓词链：
\[
\begin{aligned}
  C_1 &: \mathfrak R_D\ \text{同时包含拓扑、代数、语义、障碍、动作与编码层次},\\
  C_2 &: \widehat{\mathfrak R}_D=(\mathfrak R_D,\Cert_D,\Auth_D),\\
  C_3 &: C_2\Rightarrow \text{上述层次全部进入证书与认证闭环},\\
  C_4 &: C_1\wedge C_3\Rightarrow \text{其“富”是可认证地丰富}.
\end{aligned}
\]
\end{Certificate}

\begin{proof}[证明（基于数学符号推演逻辑的谓词因果链）]
由 $C_1$，领域富结构本身已经是多层对象。
由 $C_2$，其上又附加了证书包与认证谓词族。
于是由 $C_3$，这些层不再只是“存在”，而是全部进入合法性判定、修复合法性判定、动作合法性判定与表达合法性判定。
因此该富结构的“富”不是散漫堆叠，而是带签发边界与审计边界的受控丰富性，即 $C_4$。证毕。
\end{proof}

\subsection{推论：领域插件升级为可签发的领域子宇宙}
\label{subsec:tex37-ch8-corollary-universe}

\begin{corollary}[可签发领域子宇宙推论]
\label{cor:tex37-ch8-subuniverse}
在定理~\ref{thm:tex37-ch8-correspondence} 与命题~\ref{prop:tex37-ch8-llm-not-core} 的条件下，
每个领域插件都可被提升为一个可签发的领域认证富结构对象
\[
  \widehat{\mathfrak R}_D.
\]
从而整个系统可写成
\[
  \widehat{\mathfrak R}_{\mathrm{core}}
  \oplus
  \bigoplus_{D\in\mathfrak D_*}\widehat{\mathfrak R}_D
  \oplus
  \mathcal P_{\mathrm{render}},
\]
其中
\[
  \widehat{\mathfrak R}_{\mathrm{core}}
\]
是核心认证富结构，
\[
  \widehat{\mathfrak R}_D
\]
是各领域认证富结构，
\[
  \mathcal P_{\mathrm{render}}
\]
是统一表达插件层。
因此，系统已不再是“一个模型 + 若干工具”，而是
\[
  \boxed{
    \text{一个核心认证富结构}
    +
    \text{多个领域认证富结构}
    +
    \text{一个表达插件层。}
  }
\]
\end{corollary}

\begin{Certificate}[证书：对应定理 + 签发口径 + LLM 外置]
\label{cert:tex37-ch8-subuniverse}
证书登记谓词链：
\[
\begin{aligned}
  U_1 &: \text{定理~\ref{thm:tex37-ch8-correspondence}}\Rightarrow \mathcal P_D\leftrightarrow \widehat{\mathfrak R}_D,\\
  U_2 &: \text{`tex35` 的扩展签发口径}\Rightarrow \widehat{\mathfrak R}_D\ \text{可被理解为可签发对象},\\
  U_3 &: \text{命题~\ref{prop:tex37-ch8-llm-not-core}}\Rightarrow \mathcal P_{\mathrm{render}}\ \text{位于认证富结构之外的末端表达层},\\
  U_4 &: U_1\wedge U_2\wedge U_3\Rightarrow
  \widehat{\mathfrak R}_{\mathrm{core}}
  \oplus
  \bigoplus_{D\in\mathfrak D_*}\widehat{\mathfrak R}_D
  \oplus
  \mathcal P_{\mathrm{render}}
  \ \text{成立}.
\end{aligned}
\]
\end{Certificate}

\begin{proof}[证明（基于数学符号推演逻辑的谓词因果链）]
由 $U_1$，每个合法领域插件都对应一个认证富结构对象。
由 $U_2$，在 `tex35` 的扩展签发协议口径下，
凡带有完整证书包与认证谓词族、且边界条件明确的对象，都可被理解为可签发结构对象。
因此
\[
  \widehat{\mathfrak R}_D
\]
可以被解释为一个可签发的领域子宇宙。

再由 $U_3$，LLM 不在这些认证富结构的核心本体中，
而只在统一表达插件层中出现。
因此整个系统的最准确结构分解就是 $U_4$ 所示的直和式。
推论成立。证毕。
\end{proof}

\subsection{备注：收紧版结论与一句话压缩}
\label{subsec:tex37-ch8-remarks}

\begin{remark}[建议采用的严格版表述]
本章更适合采用如下压缩后的严格表述：
\[
  \boxed{
    \text{领域细化的塔纤维插件，可严格对应于在核心 G 富结构之上的证书认证富结构对象。}
  }
\]
\[
  \boxed{
    \text{其中“富”指多层拓扑--代数--语义--动作--编码数据的同时存在；}
  }
\]
\[
  \boxed{
    \text{“认证”指这些数据及其态射、修复链、动作提交与表达投影，}
  }
\]
\[
  \boxed{
    \text{都必须经由证书接口与审计闭环获得合法性。}
  }
\]
\end{remark}

\begin{remark}[边界]
本章的“自然对应”不是说插件脱离核心合同后仍自动合法；
恰恰相反，它只在以下边界内成立：
\[
  \text{保核心基底}
  +
  \text{保核心正规形}
  +
  \text{保核心裁决接口}
  +
  \text{保占位审计闭环}.
\]
一旦这些边界被破坏，得到的就不再是
\[
  \widehat{\mathfrak R}_D,
\]
而只是未经签发的结构碎片。
\end{remark}

\begin{remark}[一句话结论]
本章最短的严格压缩可写成：
\[
  \boxed{
    \begin{aligned}
      &\text{在该体系里，“领域细化塔纤维插件”最准确的数学身份，}\\
      &\text{就是“基于证书认证的富结构对象”。}
    \end{aligned}
  }
\]
\end{remark}

\clearpage

%%%%%%%%%%%%%%%%%%%%%%%%%%%%%%%%%%%%%%%%%%%%%%%%%%%%%%%%%%%%
\section{形式逻辑：领域细化塔纤维插件与接口收束的证书认证富结构对象}
\label{sec:tex37-ch9}
%%%%%%%%%%%%%%%%%%%%%%%%%%%%%%%%%%%%%%%%%%%%%%%%%%%%%%%%%%%%

\begin{remark}[核心陈述（在第八章基础上再收紧）]
第~\ref{sec:tex37-ch8} 章已经把领域细化塔纤维插件收紧为
\[
  \text{基于证书认证的富结构对象}.
\]
本章进一步加入“接口收束性”，并把这一对应升级为
\[
  \boxed{
    \text{领域细化塔纤维插件}
    \longleftrightarrow
    \text{接口收束的证书认证富结构对象}.
  }
\]
这里“接口收束”指：任一对外接口要么收束为可签发定理、正规形或已认证动作，
要么进入受控的进一步纤维细化展开；
它不能停留在未签发的悬空公理、未闭合的接口声明或任意扩写的黑箱分支。
\end{remark}

\begin{remark}[阅读导航：承接位置与引用口径]
本章直接承接第~\ref{sec:tex37-ch8} 章关于“领域插件 $\leftrightarrow$ 认证富结构对象”的对应定理，
并继续依赖 `tex35` 的扩展签发认证协议 \cite{RepoTex35ExtendedCertificationTower}、
`tex36` 的富结构与修复塔口径 \cite{RepoTex36HomotopyTwistBijection}，
以及本稿第一章、第六章、第七章给出的占位--解析--审计闭环与白盒主权合同 \cite{RepoTex20SemanticGaugeNormalization,RepoTex21AuditedWrapper,
  RepoTex23SemanticBase,GaoZhengAppVol2,GaoZhengAppVol3,GaoZheng2025GFramework}。

本章新增的对象不是新的本体宇宙，而是给第八章的认证富结构再附加一个
\[
  \text{接口命题签发}
  +
  \text{纤维细化展开}
  +
  \text{当前层闭合}
\]
的受控演算层。
\end{remark}

\subsection{定义：接口收束的认证富结构、公理降级签发与纤维细化层级}
\label{subsec:tex37-ch9-definitions}

\begin{definition}[接口收束的证书认证富结构]
\label{def:tex37-ch9-ic-rich}
设领域 $D$ 的认证富结构为
\[
  \widehat{\mathfrak R}_D
  =
  (\mathfrak R_D,\Cert_D,\Auth_D).
\]
再引入接口收束数据
\[
  \Conv_D
  :=
  \bigl(
    \Downgrade_D,
    \Refine_D,
    \Close_D
  \bigr),
\]
其中：
\begin{enumerate}
  \item $\Downgrade_D$ 负责“公理降级为定理”的认证签发；
  \item $\Refine_D$ 负责“当前接口进入更细纤维层级展开”；
  \item $\Close_D$ 负责“当前层接口收束为正规形、定理或已认证动作”。
\end{enumerate}
于是定义
\[
  \widehat{\mathfrak R}_D^{\mathrm{IC}}
  :=
  (\mathfrak R_D,\Cert_D,\Auth_D,\Conv_D),
\]
称其为\textbf{接口收束的证书认证富结构}。
\end{definition}

\begin{definition}[公理降级为定理认证签发]
\label{def:tex37-ch9-downgrade}
记领域 $D$ 的接口命题空间为
\[
  \mathcal I_D^{\mathrm{stmt}}.
\]
对任意接口命题
\[
  \alpha\in \mathcal I_D^{\mathrm{stmt}},
\]
称其完成了\textbf{公理降级为定理认证签发}，若满足以下二择一条件：
\begin{enumerate}
  \item \textbf{内部签发型：}
  存在定理对象 $T_\alpha$ 与证书 $c_\alpha$，使
  \[
    T_\alpha\equiv \alpha,
    \qquad
    \Issue(c_\alpha,T_\alpha)=\top;
  \]
  \item \textbf{外部导入型：}
  \[
    \ExtAssump(\alpha)=\top.
  \]
\end{enumerate}
因此接口命题集合被分解为
\[
  \mathcal I_D^{\mathrm{stmt}}
  =
  \mathcal I_{D,\mathrm{thm}}^{\mathrm{stmt}}
  \sqcup
  \mathcal I_{D,\mathrm{ext}}^{\mathrm{stmt}},
\]
其中第一类必须有定理签发证书，第二类必须显式标注为外部导入前提。
\end{definition}

\begin{definition}[纤维丛细化展开层级与接口秩函数]
\label{def:tex37-ch9-refine}
设某个领域对象、障碍或动作接口在第 $n$ 层的状态为
\[
  u^{(n)}\in \mathcal E_D^{(n)}.
\]
定义细化展开算子
\[
  \Refine_D:
  u^{(n)}
  \longmapsto
  \bigl\{u_\beta^{(n+1)}\bigr\}_{\beta\in B_n},
\]
并要求存在保基底投影
\[
  \rho_{n+1\to n}\bigl(u_\beta^{(n+1)}\bigr)=u^{(n)}.
\]
这表示更细层不是另起炉灶，而是上一层接口的纤维细化。

再定义接口秩函数
\[
  r:
  \bigsqcup_{n\in\NN}\mathcal E_D^{(n)}
  \to
  \NN,
\]
并要求对任意真正细化分支 $u_\beta^{(n+1)}$ 都有
\[
  r\bigl(u_\beta^{(n+1)}\bigr)<r\bigl(u^{(n)}\bigr).
\]
因此细化不是任意膨胀，而是朝着更小的未解决接口复杂度推进。
\end{definition}

\begin{definition}[层级收束与进一步细化展开]
\label{def:tex37-ch9-close-or-refine}
对任意接口状态 $u^{(n)}\in \mathcal E_D^{(n)}$，只允许两种合法去向：
\begin{enumerate}
  \item \textbf{层级收束：}存在闭合对象使
  \[
    \Close_D\bigl(u^{(n)}\bigr)=\operatorname{NF}_D\bigl(u^{(n)}\bigr),
  \]
  或存在定理签发对象 $T_u$ 与证书 $c_u$，使
  \[
    \Close_D\bigl(u^{(n)}\bigr)=T_u,
    \qquad
    \Issue(c_u,T_u)=\top;
  \]
  \item \textbf{进一步细化展开：}若当前层不能闭合，则必须进入
  \[
    u^{(n)}
    \longmapsto
    \bigl\{u_\beta^{(n+1)}\bigr\}_{\beta\in B_n},
    \qquad
    \rho_{n+1\to n}\bigl(u_\beta^{(n+1)}\bigr)=u^{(n)}.
  \]
\end{enumerate}
于是接口收束性的核心条件可压缩为
\[
  \boxed{
    \text{每一层都必须满足：要么闭合，要么受控细化。}
  }
\]
\end{definition}

\begin{remark}[严格含义]
定义~\ref{def:tex37-ch9-downgrade}--\ref{def:tex37-ch9-close-or-refine} 共同说明：
“接口收束”不只要求“最后能收束”，更要求\textbf{每一层都不存在悬空状态}。
也就是说，系统不允许：
\[
  \text{未签发悬空公理}
  \quad\text{或}\quad
  \text{既不闭合也不继续细化的接口残片}.
\]
\end{remark}

\subsection{公理（降级为定理）：接口状态可按秩递归实现“闭合或继续细化”}
\label{subsec:tex37-ch9-axiom}

\begin{theorem}[公理降级：接口收束程序可按接口秩递归构造（数学归纳法证书）]
\label{thm:tex37-ch9-rank-induction}
设领域 $D$ 已给定
\[
  \widehat{\mathfrak R}_D^{\mathrm{IC}}
  =
  (\mathfrak R_D,\Cert_D,\Auth_D,\Conv_D),
\]
并设接口秩函数 $r$ 满足定义~\ref{def:tex37-ch9-refine} 的下降条件。
若对任意接口状态 $u$ 都满足：
\begin{enumerate}
  \item 若 $r(u)=0$，则 $u$ 可被 $\Close_D$ 闭合；
  \item 若 $r(u)=k+1$，则 $u$ 要么可被 $\Close_D$ 直接闭合，要么经 $\Refine_D$ 送入一族子状态
  \[
    \{u_\beta\}_{\beta\in B}
  \]
  且对每个 $\beta$ 都有
  \[
    r(u_\beta)\le k.
  \]
\end{enumerate}
则对任意接口状态 $u$，存在一个有限接口收束链，
其每一步都满足“闭合或继续细化”，并最终到达闭合对象或显式外部导入前提。
\end{theorem}

\begin{Certificate}[证书：基步 + 归纳步（数学归纳法证书）]
\label{cert:tex37-ch9-rank-induction}
定义谓词
\[
  P(k)
  :
  \Longleftrightarrow
  \forall u,\ r(u)\le k
  \Rightarrow
  \text{$u$ 存在有限接口收束链}.
\]
则数学归纳法证书记为
\[
  \pi_{\mathrm{iface}}
  =
  \bigl(
    P(0),\
    \forall k\in\NN,\ P(k)\Rightarrow P(k+1)
  \bigr).
\]
其中归纳步的见证数据为
\[
  \Close_D,
  \qquad
  \Refine_D,
  \qquad
  \rho_{n+1\to n},
  \qquad
  r(u_\beta)<r(u).
\]
\end{Certificate}

\begin{proof}[证明（数学归纳法证书；基于数学符号推演逻辑的谓词因果链）]
定义谓词
\[
  P(k)
  :
  \Longleftrightarrow
  \forall u,\ r(u)\le k
  \Rightarrow
  \text{$u$ 存在有限接口收束链}.
\]

\textbf{基步：}
若 $r(u)=0$，由条件(1)，$u$ 必可被 $\Close_D$ 闭合。
因此对一切秩不超过 $0$ 的接口状态，有限收束链都存在，故 $P(0)$ 成立。

\textbf{归纳步：}
设 $P(k)$ 成立。取任意满足 $r(u)\le k+1$ 的状态 $u$。
若 $u$ 可直接被 $\Close_D$ 闭合，则收束链立即存在。
否则由条件(2)，$u$ 必经 $\Refine_D$ 送入一族子状态 $\{u_\beta\}_{\beta\in B}$，
并且每个 $u_\beta$ 的秩都不超过 $k$。
由归纳假设 $P(k)$，每个子状态 $u_\beta$ 都存在有限接口收束链。
又由于
\[
  \rho_{n+1\to n}(u_\beta)=u,
\]
这些子链都是对原状态 $u$ 的保投影细化延伸。
故把
\[
  u\longmapsto \{u_\beta\}_{\beta\in B}
\]
接到各个子链前端，即可得到 $u$ 的有限接口收束链。
因此 $P(k+1)$ 成立。

由数学归纳法，对所有 $k\in\NN$，$P(k)$ 成立。
故任意接口状态都满足“闭合或继续细化”的有限收束性质。证毕。
\end{proof}

\subsection{定理：领域细化塔纤维插件与接口收束的证书认证富结构对象自然对应}
\label{subsec:tex37-ch9-theorem}

\begin{theorem}[接口收束的证书认证富结构对应定理]
\label{thm:tex37-ch9-correspondence}
对任意领域 $D\in\mathfrak D$，以下两件事等价：
\begin{enumerate}
  \item \textbf{插件型条件：}
  $\mathcal P_D$ 是一个合法的领域细化塔纤维插件，并且其所有接口都满足：
  \begin{enumerate}
    \item 公理降级为定理认证签发；
    \item 纤维丛展开层级在每一步都“闭合或继续细化”；
    \item 所有细化都保基底、保投影、保认证闭环；
    \item 表达层仍服从占位--解析--审计闭环。
  \end{enumerate}
  \item \textbf{富结构型条件：}
  存在一个接口收束的证书认证富结构对象
  \[
    \widehat{\mathfrak R}_D^{\mathrm{IC}}
    =
    (\mathfrak R_D,\Cert_D,\Auth_D,\Conv_D)
  \]
  与之对应。
\end{enumerate}
并且该对应由两条典范映射给出：
\[
  \Phi_D^{\mathrm{IC}}:
  \mathcal P_D
  \longmapsto
  \widehat{\mathfrak R}_D^{\mathrm{IC}},
  \qquad
  \Psi_D^{\mathrm{IC}}:
  \widehat{\mathfrak R}_D^{\mathrm{IC}}
  \longmapsto
  \mathfrak R_D\ominus \mathfrak R_{\mathrm{core}}.
\]
\end{theorem}

\begin{Certificate}[证书：插件条件 $\Rightarrow$ 收束富结构，收束富结构 $\Rightarrow$ 插件条件]
\label{cert:tex37-ch9-correspondence}
证书登记谓词链：
\[
\begin{aligned}
  C_1 &: \mathcal P_D\models \mathfrak C_{\mathrm{core}},\\
  C_2 &: \text{$\mathcal P_D$ 的接口命题都完成定理签发或显式外部导入标注},\\
  C_3 &: \text{$\mathcal P_D$ 的接口状态都满足“闭合或继续细化”且保基底、保投影},\\
  C_4 &: \text{表达层保持占位--解析--审计闭环},\\
  C_5 &: C_1\wedge C_2\wedge C_3\wedge C_4
  \Rightarrow
  \Phi_D^{\mathrm{IC}}(\mathcal P_D)=\widehat{\mathfrak R}_D^{\mathrm{IC}},\\
  C_6 &: \widehat{\mathfrak R}_D^{\mathrm{IC}}
  \Rightarrow
  \Psi_D^{\mathrm{IC}}(\widehat{\mathfrak R}_D^{\mathrm{IC}})
  =
  \mathfrak R_D\ominus \mathfrak R_{\mathrm{core}}
  \ \text{满足插件型条件}.
\end{aligned}
\]
\end{Certificate}

\begin{proof}[证明（基于数学符号推演逻辑的谓词因果链）]
\textbf{证明 $(i)\Rightarrow(ii)$：}
设 $\mathcal P_D$ 满足插件型条件。
由 $C_1$，它已经是一个服从核心主权合同的合法领域细化塔纤维插件。
由 $C_2$，其接口命题不再允许以“内部悬空公理”的方式滞留，
而必须被压成可签发定理，或被显式标记为外部导入前提。
因此接口命题层具有
\[
  \Downgrade_D
\]
所编码的签发机制。

由 $C_3$，任何当前层尚未闭合的几何、障碍或动作接口，
都必须进入保基底、保投影的进一步细化展开。
因此接口状态层具有
\[
  \Refine_D
\]
与
\[
  \Close_D
\]
所编码的“闭合或继续细化”机制。

再由 $C_4$，表达层仍处于占位--解析--审计闭环内，
故接口收束并不在文本层断裂。
于是把
\[
  \mathfrak R_D,\ \Cert_D,\ \Auth_D
\]
与
\[
  \Conv_D=(\Downgrade_D,\Refine_D,\Close_D)
\]
合并，即得
\[
  \widehat{\mathfrak R}_D^{\mathrm{IC}}
  =
  (\mathfrak R_D,\Cert_D,\Auth_D,\Conv_D).
\]
故 $(ii)$ 成立。

\textbf{证明 $(ii)\Rightarrow(i)$：}
设给定
\[
  \widehat{\mathfrak R}_D^{\mathrm{IC}}
  =
  (\mathfrak R_D,\Cert_D,\Auth_D,\Conv_D).
\]
由第~\ref{sec:tex37-ch8} 章的对应定理，
\[
  \mathfrak R_D\ominus \mathfrak R_{\mathrm{core}}
\]
给出领域插件增量
\[
  \mathcal P_D.
\]
由于 $\Conv_D$ 已被并入该对象，
故从 $\widehat{\mathfrak R}_D^{\mathrm{IC}}$ 回抽出的插件天然携带：
\begin{enumerate}
  \item 由 $\Downgrade_D$ 给出的接口命题签发机制；
  \item 由 $\Refine_D$ 给出的纤维细化展开机制；
  \item 由 $\Close_D$ 给出的当前层闭合机制。
\end{enumerate}
又因为该对象仍带有 $\Cert_D$、$\Auth_D$ 与占位--解析--审计闭环，
所以回抽出的 $\mathcal P_D$ 自动满足插件型条件中的四项要求。
故 $(i)$ 成立。

因此 $(i)$ 与 $(ii)$ 等价。
换言之，合法领域插件与接口收束的证书认证富结构对象之间存在自然对应。证毕。
\end{proof}

\subsection{引理：接口命题不允许以“未签发内部公理”形式悬空}
\label{subsec:tex37-ch9-lemma}

\begin{lemma}[接口命题二分引理]
\label{lem:tex37-ch9-statement-split}
对任意
\[
  \alpha\in \mathcal I_D^{\mathrm{stmt}},
\]
若 $\alpha$ 属于接口收束的证书认证富结构
\[
  \widehat{\mathfrak R}_D^{\mathrm{IC}},
\]
则以下二者必居其一且互斥：
\[
  \alpha\in \mathcal I_{D,\mathrm{thm}}^{\mathrm{stmt}},
  \qquad
  \alpha\in \mathcal I_{D,\mathrm{ext}}^{\mathrm{stmt}}.
\]
因此，接口层不允许存在第三类“未签发的内部悬空公理”。
\end{lemma}

\begin{Certificate}[证书：内部签发 or 外部导入]
\label{cert:tex37-ch9-statement-split}
证书登记谓词链：
\[
\begin{aligned}
  L_1 &: \alpha\in \mathcal I_D^{\mathrm{stmt}},\\
  L_2 &: \alpha\notin \mathcal I_{D,\mathrm{ext}}^{\mathrm{stmt}}
  \Rightarrow
  \ExtAssump(\alpha)\neq \top,\\
  L_3 &: L_2\Rightarrow \alpha\ \text{必须经}\ \Downgrade_D\ \text{进入}\ \mathcal I_{D,\mathrm{thm}}^{\mathrm{stmt}},\\
  L_4 &: \mathcal I_{D,\mathrm{thm}}^{\mathrm{stmt}}\cap \mathcal I_{D,\mathrm{ext}}^{\mathrm{stmt}}=\varnothing,\\
  L_5 &: L_1\wedge L_3\wedge L_4\Rightarrow \text{不存在第三类悬空内部公理}.
\end{aligned}
\]
\end{Certificate}

\begin{proof}[证明（基于数学符号推演逻辑的谓词因果链）]
由 $L_1$，$\alpha$ 是一个领域接口命题。
若 $\alpha$ 属于外部导入型，则它已被
\[
  \ExtAssump(\alpha)=\top
\]
显式标注，其地位不是内部真理，而是外部前提。

若 $\alpha$ 不属于外部导入型，则由接口收束定义，它就不能悬空停留；
因此必须经由
\[
  \Downgrade_D
\]
进入定理签发流程，并取得
\[
  \Issue(c_\alpha,T_\alpha)=\top.
\]
这给出 $L_3$。

再由 $L_4$，内部签发型与外部导入型互斥。
因此接口层只能有这两类命题，而不存在既非外部前提、又未被签发为定理的第三类内部悬空公理。
引理成立。证毕。
\end{proof}

\subsection{命题：接口收束性回收的是接口主权，而该主权不属于 LLM}
\label{subsec:tex37-ch9-proposition}

\begin{proposition}[接口主权回收命题]
\label{prop:tex37-ch9-interface-sovereignty}
定义接口主权为
\[
  \mathsf{Sov}_{\mathrm{interface}}
  :=
  \mathsf{Sov}_{\mathrm{axiom\to theorem}}
  \oplus
  \mathsf{Sov}_{\mathrm{refine/close}},
\]
其中：
\begin{enumerate}
  \item
  \[
    \mathsf{Sov}_{\mathrm{axiom\to theorem}}
  \]
  表示“谁决定某个接口命题被压成可签发定理，还是被标注为外部导入前提”；
  \item
  \[
    \mathsf{Sov}_{\mathrm{refine/close}}
  \]
  表示“谁决定某个接口状态在当前层闭合，还是进入下一层纤维细化”。
\end{enumerate}
在
\[
  \widehat{\mathfrak R}_D^{\mathrm{IC}}
\]
的设定下，这一主权不应交给 LLM，而应与
\[
  \mathsf{Sov}_{\mathrm{solve}}
  \oplus
  \mathsf{Sov}_{\mathrm{decide}}
\]
一起回收到 G 侧程序与认证接口。
\end{proposition}

\begin{Certificate}[证书：签发决定权 + 细化闭合决定权 + 表达层后置]
\label{cert:tex37-ch9-interface-sovereignty}
证书登记谓词链：
\[
\begin{aligned}
  P_1 &: \Downgrade_D,\Refine_D,\Close_D\ \text{都属于}\ \Conv_D\ \text{的白盒接口演算层},\\
  P_2 &: \Conv_D\ \text{在文本渲染前决定接口命题签发与接口状态去向},\\
  P_3 &: \Render^{\mathrm{LLM}}_{\xi,D}\ \text{只位于占位固定之后的末端表达层},\\
  P_4 &: P_1\wedge P_2\wedge P_3\Rightarrow \mathsf{Sov}_{\mathrm{interface}}\ \text{不属于 LLM},\\
  P_5 &: P_4\Rightarrow \mathsf{Sov}_{\mathrm{solve}}\oplus\mathsf{Sov}_{\mathrm{decide}}\oplus\mathsf{Sov}
    _{\mathrm{interface}}
  \ \text{共同回收到 G 侧}.
\end{aligned}
\]
\end{Certificate}

\begin{proof}[证明（基于数学符号推演逻辑的谓词因果链）]
由 $P_1$，接口命题的签发机制与接口状态的细化/闭合机制都内置于
\[
  \Conv_D=(\Downgrade_D,\Refine_D,\Close_D)
\]
中。
由 $P_2$，这些机制发生在文本渲染之前，
因此它们决定的是“系统如何理解、签发、分解和闭合接口”，而不是“如何把既定结果写成文本”。

由 $P_3$，LLM 只在占位固定之后执行末端外显，
所以它既不能决定某个命题是否已被签发为定理，
也不能决定某个状态是应闭合还是应继续细化。
故 $P_4$ 成立。

再结合第~\ref{sec:tex37-ch6} 章的求解主权回收与决策主权回收结论，
可知接口收束性额外回收的就是
\[
  \mathsf{Sov}_{\mathrm{interface}}.
\]
因此三类主权应共同留在 G 侧白盒程序与认证接口内。命题成立。证毕。
\end{proof}

\subsection{推论：富结构从静态对象升级为“可演化但不失控”的富结构}
\label{subsec:tex37-ch9-corollary-evolve}

\begin{corollary}[可演化但不失控的富结构]
\label{cor:tex37-ch9-evolve}
在定理~\ref{thm:tex37-ch9-correspondence} 与定理~\ref{thm:tex37-ch9-rank-induction} 的条件下，
领域富结构不再只是
\[
  \text{对象多、层级多、证书多},
\]
而是升级为
\[
  \boxed{
    \text{对象可继续细化，但每一步细化都可投影、可认证、可收束。}
  }
\]
因此它避免了两类坏情形：
\[
  \text{接口悬空}
  \quad\text{与}\quad
  \text{细化失控}.
\]
\end{corollary}

\begin{Certificate}[证书：秩下降 + 保投影 + 闭合机制]
\label{cert:tex37-ch9-evolve}
证书登记谓词链：
\[
\begin{aligned}
  C_1 &: \Refine_D\ \text{伴随}\ \rho_{n+1\to n}\ \text{与秩下降条件}\ r(u_\beta)<r(u),\\
  C_2 &: \Close_D\ \text{提供当前层闭合机制},\\
  C_3 &: C_1\wedge C_2\Rightarrow \text{细化既可继续演化，又不会脱离回投影与收束准则},\\
  C_4 &: C_3\Rightarrow \text{接口悬空与细化失控同时被封死}.
\end{aligned}
\]
\end{Certificate}

\begin{proof}[证明（基于数学符号推演逻辑的谓词因果链）]
由 $C_1$，每一次细化都带有保投影约束与秩下降约束，
所以系统不会沿着无法回投影、也无法降低复杂度的分支任意膨胀。
由 $C_2$，当前层一旦满足闭合条件，就可立即收束为正规形、定理或已认证动作。

故由 $C_3$，富结构既是可演化的，又是受控的：
不能闭合时才继续细化，能够闭合时立即收束。
于是既没有“提出接口但无法签发或分解”的悬空状态，也没有“层级越分越多但没有准则”的失控状态，
即 $C_4$。推论成立。证毕。
\end{proof}

\subsection{推论：表达层同样满足接口收束性}
\label{subsec:tex37-ch9-corollary-render}

\begin{corollary}[表达层接口收束推论]
\label{cor:tex37-ch9-render}
设文本输出为
\[
  t^{(n)}.
\]
则在接口收束口径下，表达层同样只允许两种合法状态：
\begin{enumerate}
  \item \textbf{已闭合输出：}
  \[
    \Parse_\xi\bigl(t^{(n)}\bigr)=p,
    \qquad
    \Style_\xi\bigl(t^{(n)}\bigr)=\top,
    \qquad
    t^{(n)}\in \Audit_\xi(p);
  \]
  \item \textbf{继续细化展开：}
  若当前文本尚不足以覆盖占位槽，则必须进入更细模板展开
  \[
    t^{(n)}
    \longmapsto
    \bigl\{t_\beta^{(n+1)}\bigr\}_{\beta\in B_n^{\mathrm{render}}},
  \]
  直到各槽位被完整覆盖，或显式返回 $\bot$。
\end{enumerate}
因此表达层也不允许“看起来像写完了，但其实没有闭合”的半成品接口。
\end{corollary}

\begin{Certificate}[证书：占位固定 + 审计接受 + 覆盖不足则继续细化]
\label{cert:tex37-ch9-render}
证书登记谓词链：
\[
\begin{aligned}
  R_1 &: p\ \text{已由上游求解与占位链固定},\\
  R_2 &: \Parse_\xi(t)=p,\ \Style_\xi(t)=\top,\ t\in\Audit_\xi(p)\Rightarrow t\ \text{已闭合},\\
  R_3 &: \text{若当前文本对占位槽覆盖不足，则必须进入下一层模板细化或返回 }\bot,\\
  R_4 &: R_2\vee R_3\Rightarrow \text{表达层满足“闭合或继续细化”}.
\end{aligned}
\]
\end{Certificate}

\begin{proof}[证明（基于数学符号推演逻辑的谓词因果链）]
由 $R_1$，表达层并不决定占位本身，它只负责对既定占位做外显。
若某个文本 $t$ 已满足
\[
  \Parse_\xi(t)=p,\ \Style_\xi(t)=\top,\ t\in\Audit_\xi(p),
\]
则由 $R_2$，它已经构成闭合输出。

若当前文本还不能完整覆盖占位槽，
则根据接口收束原则，它不能伪装成已完成输出，
而必须继续进入下一层模板细化，或显式返回 $\bot$，这就是 $R_3$。
于是表达层同样满足“闭合或继续细化”的接口收束性，即 $R_4$。证毕。
\end{proof}

\subsection{备注：建议采用的严格版表述}
\label{subsec:tex37-ch9-remarks}

\begin{remark}[建议采用的压缩表述]
本章更适合采用如下严格版表述：
\[
  \boxed{
    \text{领域细化塔纤维插件不仅对应于基于证书认证的富结构对象，}
  }
\]
\[
  \boxed{
    \text{还必须对应于接口收束的证书认证富结构对象。}
  }
\]
\[
  \boxed{
    \text{所谓接口收束性，是指任何接口命题都必须由“公理降级为定理”的认证签发机制闭合，}
  }
\]
\[
  \boxed{
    \begin{aligned}
      &\text{而任何纤维细化层级，都必须满足“当前层闭合”，}\\
      &\text{或“进入保基底、保投影、保认证的进一步细化展开”。}
    \end{aligned}
  }
\]
\[
  \boxed{
    \text{因此，系统扩张不是任意膨胀，而是受控的、可回投影的、可签发的层级演化。}
  }
\]
\end{remark}

\begin{remark}[一句话结论]
本章最短的严格压缩可写成：
\[
  \boxed{
    \begin{aligned}
      &\text{领域插件的准确身份不只是“认证富结构”，}\\
      &\text{而是“可将公理压为定理，并使细化层级闭合或继续受控展开的}\\
      &\text{接口收束认证富结构”。}
    \end{aligned}
  }
\]
\end{remark}

\clearpage

%%%%%%%%%%%%%%%%%%%%%%%%%%%%%%%%%%%%%%%%%%%%%%%%%%%%%%%%%%%%
\section{形式逻辑：本次讨论的综合增益评价与架构级重大增益判定}
\label{sec:tex37-ch10}
%%%%%%%%%%%%%%%%%%%%%%%%%%%%%%%%%%%%%%%%%%%%%%%%%%%%%%%%%%%%

\begin{remark}[核心结论（总评版）]
对本次从第~\ref{sec:tex37-ch1} 章到第~\ref{sec:tex37-ch9} 章的整轮讨论，可给出如下明确判断：
\[
  \boxed{
    \begin{aligned}
      &\text{本次讨论的总体增益极高，}\\
      &\text{其重心在“体系闭环化、主权回收化、插件化、可签发化”四个方向。}
    \end{aligned}
  }
\]
若再压缩一句，则可写成：
\[
  \boxed{
    \begin{aligned}
      &\text{这次讨论不是单纯新增一个局部结论，}\\
      &\text{而是把若干分散的强直觉压成了一条可部署的白盒 AI 总架构链。}
    \end{aligned}
  }
\]
\end{remark}

\begin{remark}[阅读导航：本章的评价对象]
本章不再新增新的求解对象或新的认证对象，而是把前九章得到的关键结论统一记账：
第~\ref{sec:tex37-ch4} 章给出问答内核对象与修复塔路径类口径；
第~\ref{sec:tex37-ch6} 章给出主权回收到 G 侧白盒程序；
第~\ref{sec:tex37-ch7} 章给出领域细化塔纤维插件；
第~\ref{sec:tex37-ch8} 章给出证书认证富结构对象；
第~\ref{sec:tex37-ch9} 章给出接口收束性。
本章要评价的，正是这些结论联立后所形成的总体架构增益。
\end{remark}

\subsection{定义：本次讨论的综合增益函数与定性评级集}
\label{subsec:tex37-ch10-definitions}

\begin{definition}[本次讨论的综合增益函数]
\label{def:tex37-ch10-gain}
定义本次讨论的综合增益为
\[
  \mathrm{Gain}_{\mathrm{disc}}
  :=
  \frac{
    G_{\mathrm{arch}}
    +
    G_{\mathrm{sov}}
    +
    G_{\mathrm{plugin}}
    +
    G_{\mathrm{cert}}
    +
    G_{\mathrm{conv}}
    +
    G_{\mathrm{render}}
  }{
    C_{\mathrm{pending}}
  },
\]
其中：
\begin{enumerate}
  \item $G_{\mathrm{arch}}$ 表示总架构闭环增益；
  \item $G_{\mathrm{sov}}$ 表示思考/决策/接口主权回收增益；
  \item $G_{\mathrm{plugin}}$ 表示领域细化塔纤维插件化增益；
  \item $G_{\mathrm{cert}}$ 表示证书认证富结构化增益；
  \item $G_{\mathrm{conv}}$ 表示接口收束性增益；
  \item $G_{\mathrm{render}}$ 表示 LLM 降级为表达插件后的工程增益；
  \item $C_{\mathrm{pending}}$ 表示仍需进一步形式化与实现的剩余成本。
\end{enumerate}
\end{definition}

\begin{definition}[定性评级集]
\label{def:tex37-ch10-rank}
定义本章使用的定性评级集为
\[
  \Gamma_{\mathrm{gain}}
  :=
  \{
    \text{低},
    \text{中},
    \text{中高},
    \text{高},
    \text{很高},
    \text{极高}
  \},
\]
并按书写顺序赋予全序
\[
  \text{低}
  \prec
  \text{中}
  \prec
  \text{中高}
  \prec
  \text{高}
  \prec
  \text{很高}
  \prec
  \text{极高}.
\]
记
\[
  \operatorname{Rank}(G)\in \Gamma_{\mathrm{gain}}
\]
为某个增益项的定性评级。
\end{definition}

\begin{remark}[评价口径]
本章中的“极高/很高/高”不是修辞，而是针对以下问题的定性结论：
\[
  \begin{aligned}
    &\text{是否打通全链路闭环},
    \quad
    \text{是否完成主权边界重写},\\
    &\text{是否得到可部署的插件化白盒架构},
    \quad
    \text{是否把扩张对象纳入签发与收束机制}.
  \end{aligned}
\]
\end{remark}

\subsection{公理（降级为定理）：综合增益分子可按六类结构增量递归累加}
\label{subsec:tex37-ch10-axiom}

\begin{theorem}[公理降级：综合增益分子可按六类结构增量递归累加（数学归纳法证书）]
\label{thm:tex37-ch10-decompose}
定义
\[
  \Delta_1:=G_{\mathrm{arch}},
  \quad
  \Delta_2:=G_{\mathrm{sov}},
  \quad
  \Delta_3:=G_{\mathrm{plugin}},
  \quad
  \Delta_4:=G_{\mathrm{cert}},
  \quad
  \Delta_5:=G_{\mathrm{conv}},
  \quad
  \Delta_6:=G_{\mathrm{render}}.
\]
再定义部分和
\[
  N(1):=\Delta_1,
  \qquad
  N(m+1):=N(m)+\Delta_{m+1}
  \quad (1\le m\le 5).
\]
则对所有 $m\in\{1,2,3,4,5,6\}$，都有
\[
  N(m)=\sum_{k=1}^{m}\Delta_k.
\]
特别地，
\[
  N(6)
  =
  G_{\mathrm{arch}}
  +
  G_{\mathrm{sov}}
  +
  G_{\mathrm{plugin}}
  +
  G_{\mathrm{cert}}
  +
  G_{\mathrm{conv}}
  +
  G_{\mathrm{render}},
\]
正是定义~\ref{def:tex37-ch10-gain} 中综合增益函数的分子。
\end{theorem}

\begin{Certificate}[证书：基步 + 归纳步（数学归纳法证书）]
\label{cert:tex37-ch10-decompose}
定义谓词
\[
  P(m):\ N(m)=\sum_{k=1}^{m}\Delta_k.
\]
则数学归纳法证书记为
\[
  \pi_{\mathrm{gain}}
  =
  \bigl(
    P(1),\
    \forall m\in\{1,2,3,4,5\},\ P(m)\Rightarrow P(m+1)
  \bigr).
\]
其中归纳步只依赖递推关系
\[
  N(m+1)=N(m)+\Delta_{m+1}.
\]
\end{Certificate}

\begin{proof}[证明（数学归纳法证书；基于数学符号推演逻辑的谓词因果链）]
\textbf{基步：}
由定义，
\[
  N(1)=\Delta_1=\sum_{k=1}^{1}\Delta_k.
\]
故 $P(1)$ 成立。

\textbf{归纳步：}
设 $P(m)$ 成立，即
\[
  N(m)=\sum_{k=1}^{m}\Delta_k.
\]
则由递推定义，
\[
  N(m+1)=N(m)+\Delta_{m+1}.
\]
代入归纳假设得
\[
  N(m+1)
  =
  \sum_{k=1}^{m}\Delta_k+\Delta_{m+1}
  =
  \sum_{k=1}^{m+1}\Delta_k.
\]
故 $P(m+1)$ 成立。

由数学归纳法，$P(m)$ 对所有 $m=1,\dots,6$ 都成立。
于是 $N(6)$ 恰好是六类结构增益的总和，也就是综合增益函数的分子。证毕。
\end{proof}

\subsection{定理：本次讨论最大的增益，不是新增单个对象，而是打通了完整求解链}
\label{subsec:tex37-ch10-theorem-arch}

\begin{theorem}[总架构闭环增益定理]
\label{thm:tex37-ch10-arch}
联立第~\ref{sec:tex37-ch4} 章、第~\ref{sec:tex37-ch6} 章到第~\ref{sec:tex37-ch9} 章的结果，
本次讨论已把下列链条明确压成连续闭环：
\[
  \begin{aligned}
    &\text{上下文基底}
    \to
    \text{G 框架同伦修复塔}
    \to
    \text{路径积分/正规形/有效作用选择}\\
    &\to
    \text{逻辑占位}
    \to
    \text{可审计文本表达}.
  \end{aligned}
\]
并且严格分离了三层对象：
\[
  \text{修复塔路径/截面层}
  \neq
  \text{逻辑占位层}
  \neq
  \text{表面文本层}.
\]
因此，本次讨论关于总体架构的主要增益落在“链的打通”而非“点的新增”，即
\[
  \operatorname{Rank}(G_{\mathrm{arch}})=\text{极高}.
\]
\end{theorem}

\begin{Certificate}[证书：修复塔对象 + 占位层 + 表达闭环]
\label{cert:tex37-ch10-arch}
证书登记谓词链：
\[
\begin{aligned}
  A_1 &: \text{第~\ref{sec:tex37-ch4} 章给出问答内核对象是修复塔路径/截面层对象},\\
  A_2 &: \text{第~\ref{sec:tex37-ch6} 章给出求解/决策主权回收到 G 侧白盒程序},\\
  A_3 &: \text{第~\ref{sec:tex37-ch7}--\ref{sec:tex37-ch9} 章给出插件化、认证化与接口收束化扩张},\\
  A_4 &: \text{第~\ref{sec:tex37-ch1} 章给出逻辑占位与可审计文本表达闭环},\\
  A_5 &: A_1\wedge A_2\wedge A_3\wedge A_4
  \Rightarrow
  \text{从求解对象到对外表达的完整闭环已建立}.
\end{aligned}
\]
\end{Certificate}

\begin{proof}[证明（基于数学符号推演逻辑的谓词因果链）]
若一个体系只有若干局部强命题，而没有从“求解对象”到“对外表达”的连续闭环，
则其整体状态仍接近理论碎片。
由 $A_1$，本稿已经把真正被求解的对象压到修复塔路径/截面层。
由 $A_2$，这些对象的求解与动作裁决被压回 G 侧白盒程序。
由 $A_3$，领域扩张又被统一压成插件化、认证化与接口收束化结构。
由 $A_4$，这些结果最终进入逻辑占位层与可审计表达层。

故由 $A_5$，本次讨论的核心增益不在“再发明一个单点对象”，而在“把已有强对象接成连续闭环”。
这正是总架构闭环增益的来源。证毕。
\end{proof}

\begin{remark}[评价]
很多理论稿的强项是“局部深”，而本次讨论的重量在于把局部深度接成了系统闭环。
故其增益首先是架构级而不是局部命题级。
\end{remark}

\subsection{定理：本次讨论把“AI 的大脑属于谁”给出了严格答案}
\label{subsec:tex37-ch10-theorem-sov}

\begin{theorem}[主权回收总定理]
\label{thm:tex37-ch10-sovereignty}
在本次讨论的最终结构中，以下四类主权被严格分离：
\[
  \mathsf{Sov}_{\mathrm{solve}},
  \qquad
  \mathsf{Sov}_{\mathrm{decide}},
  \qquad
  \mathsf{Sov}_{\mathrm{interface}},
  \qquad
  \mathsf{Sov}_{\mathrm{render}}.
\]
并满足：
\begin{enumerate}
  \item
  \[
    \mathsf{Sov}_{\mathrm{solve}}
  \]
  属于 G 侧修复塔、正规形与有效作用链；
  \item
  \[
    \mathsf{Sov}_{\mathrm{decide}}
  \]
  属于 G 侧动作裁决协议；
  \item
  \[
    \mathsf{Sov}_{\mathrm{interface}}
  \]
  属于定理签发与接口收束协议；
  \item
  \[
    \mathsf{Sov}_{\mathrm{render}}
  \]
  只保留在末端表达投影层。
\end{enumerate}
因此 AI 的“主脑”从黑箱模型迁回白盒程序，故
\[
  \operatorname{Rank}(G_{\mathrm{sov}})=\text{极高}.
\]
\end{theorem}

\begin{Certificate}[证书：求解主权 + 决策主权 + 接口主权 + 渲染后置]
\label{cert:tex37-ch10-sovereignty}
证书登记谓词链：
\[
\begin{aligned}
  S_1 &: \text{第~\ref{sec:tex37-ch6} 章回收}\ \mathsf{Sov}_{\mathrm{solve}}\oplus \mathsf{Sov}_{\mathrm{decide}},\\
  S_2 &: \text{第~\ref{sec:tex37-ch9} 章回收}\ \mathsf{Sov}_{\mathrm{interface}},\\
  S_3 &: \Render^{\mathrm{LLM}}_\xi\ \text{只位于占位固定后的末端表达层},\\
  S_4 &: S_1\wedge S_2\wedge S_3
  \Rightarrow
  \text{除表达投影外的核心主权全部回收到 G 侧}.
\end{aligned}
\]
\end{Certificate}

\begin{proof}[证明（基于数学符号推演逻辑的谓词因果链）]
由 $S_1$，求解权与决策权已经不属于 LLM。
由 $S_2$，接口命题的签发权与接口状态的闭合/细化权也已回收到 G 侧程序和认证接口。
由 $S_3$，LLM 只负责
\[
  p^\star\mapsto t
\]
这一步末端表达投影。
因此，求解、决策、接口三类核心主权不再属于黑箱模型，而属于白盒程序对象链，即 $S_4$。
这正给出了“AI 的大脑属于谁”的严格答案。证毕。
\end{proof}

\begin{remark}[评价]
这一点是本次讨论最重的核心之一，因为它把
\[
  \text{G 框架负责想与定；LLM 负责说}
\]
从口头判断压成了可引用的体系结论。
\end{remark}

\subsection{引理：本次讨论把领域扩张从“外挂知识”提升为“塔纤维插件”}
\label{subsec:tex37-ch10-lemma-plugin}

\begin{lemma}[插件化升级引理]
\label{lem:tex37-ch10-plugin}
本次讨论中，领域扩张不再被理解为术语表、规则包或提示词堆叠，
而被重写为同一上下文基底上的塔纤维插件
\[
  \mathcal P_D.
\]
其新增内容进入的是求解纤维本身：
\[
  \Delta_D^{\mathrm{obs}},
  \quad
  \Delta_D^{\mathrm{op}},
  \quad
  \Delta_D^{\mathrm{obst}},
  \quad
  \Delta_D^{\mathrm{act}},
  \quad
  \Delta_D^{\mathrm{cert}}.
\]
因此，本次讨论把“领域知识”从语言层外挂提升为几何/语义/动作层插件，故
\[
  \operatorname{Rank}(G_{\mathrm{plugin}})=\text{很高}.
\]
\end{lemma}

\begin{Certificate}[证书：同基底纤维细化 + 五类增量]
\label{cert:tex37-ch10-plugin}
证书登记谓词链：
\[
\begin{aligned}
  L_1 &: \text{第~\ref{sec:tex37-ch7} 章定义}\ \mathcal P_D\ \text{为同基底塔纤维插件},\\
  L_2 &: \mathcal P_D\ \text{附加的是观测变量、障碍类型、修复算子、动作空间与证书接口},\\
  L_3 &: L_1\wedge L_2\Rightarrow \text{领域扩张进入求解纤维本身而非只进入表面语言层}.
\end{aligned}
\]
\end{Certificate}

\begin{proof}[证明（基于数学符号推演逻辑的谓词因果链）]
若领域扩张只是增加术语表或提示词，则它只改变表面语言层。
由 $L_1$ 与 $L_2$，本次讨论中的领域扩张却被写成了对同一基底上纤维对象的细化，
并把新增内容直接压到求解、障碍、动作与证书接口层。
因此由 $L_3$，领域知识进入的是内核求解纤维，而不是外围补丁。
这就是插件化增益的实质。证毕。
\end{proof}

\begin{remark}[评价]
这一步是体系从“通用核”走向“多领域工业核”的关键桥梁；
没有它，体系仍偏抽象；有了它，体系开始具备工业系统架构感。
\end{remark}

\subsection{定理：本次讨论把插件升级为证书认证且接口收束的富结构对象}
\label{subsec:tex37-ch10-theorem-cert-conv}

\begin{theorem}[认证富结构与接口收束增益定理]
\label{thm:tex37-ch10-cert-conv}
联立第~\ref{sec:tex37-ch8} 章与第~\ref{sec:tex37-ch9} 章，可得：
\begin{enumerate}
  \item 领域插件对应于证书认证富结构对象
  \[
    \widehat{\mathfrak R}_D
    =
    (\mathfrak R_D,\Cert_D,\Auth_D);
  \]
  \item 该对象进一步收紧为接口收束的证书认证富结构对象
  \[
    \widehat{\mathfrak R}_D^{\mathrm{IC}}
    =
    (\mathfrak R_D,\Cert_D,\Auth_D,\Conv_D).
  \]
\end{enumerate}
因此，本次讨论不只把插件变为“可认证地丰富”，还把它变为“可签发、可闭合、可继续受控细化的层级结构”，并满足
\[
  \operatorname{Rank}(G_{\mathrm{cert}})=\text{很高},
  \qquad
  \operatorname{Rank}(G_{\mathrm{conv}})=\text{极高}.
\]
\end{theorem}

\begin{Certificate}[证书：认证化对应 + 接口收束对应]
\label{cert:tex37-ch10-cert-conv}
证书登记谓词链：
\[
\begin{aligned}
  C_1 &: \text{第~\ref{sec:tex37-ch8} 章给出}\ \mathcal P_D\leftrightarrow \widehat{\mathfrak R}_D,\\
  C_2 &: \text{第~\ref{sec:tex37-ch9} 章给出}\ \mathcal P_D\leftrightarrow \widehat{\mathfrak R}_D^{\mathrm{IC}},\\
  C_3 &: C_1\Rightarrow \text{插件扩张从功能模块升级为可签发富结构层},\\
  C_4 &: C_2\Rightarrow \text{该富结构层进一步获得接口命题签发与层级收束机制},\\
  C_5 &: C_3\wedge C_4\Rightarrow G_{\mathrm{cert}}\ \text{很高且}\ G_{\mathrm{conv}}\ \text{极高}.
\end{aligned}
\]
\end{Certificate}

\begin{proof}[证明（基于数学符号推演逻辑的谓词因果链）]
由 $C_1$，领域插件不再只是“功能模块”，而是带证书、带认证、带审计闭环的富结构对象。
这说明领域扩张被纳入了可签发、可回放、可审计的结构层，故 $G_{\mathrm{cert}}$ 很高。

由 $C_2$，该对象又进一步附加了
\[
  \Conv_D=(\Downgrade_D,\Refine_D,\Close_D),
\]
从而任何接口都必须“当前层闭合”或“进入受控进一步细化”。
这一步封死了“接口悬空”和“细化失控”两类坏情形，
因此 $G_{\mathrm{conv}}$ 的权重高于单纯认证化，达到极高。
故 $C_5$ 成立。证毕。
\end{proof}

\begin{remark}[评价]
这一步非常漂亮，因为它把“富结构”从静态多层对象推进成了“可演化但不失控的富结构”。
\end{remark}

\subsection{命题：本次讨论对 LLM 的重新定位带来极强工程增益，但仍未等于全部完成}
\label{subsec:tex37-ch10-proposition}

\begin{proposition}[表达插件化与剩余成本命题]
\label{prop:tex37-ch10-render-pending}
本次讨论最终把
\[
  \text{LLM}\rightsquigarrow \text{可替换表达插件}
\]
这一重定位压成了严格结构，
故
\[
  \operatorname{Rank}(G_{\mathrm{render}})=\text{极高}.
\]
但与此同时，剩余成本
\[
  C_{\mathrm{pending}}
\]
仍严格大于 $0$，其主要来源包括：
\begin{enumerate}
  \item 各领域插件的正式接口清单；
  \item 公理降级为定理的统一签发协议；
  \item 细化层级收束的可计算准则；
  \item 正规形到动作执行器的工程实现；
  \item 占位覆盖度与文本充分性的形式化度量。
\end{enumerate}
因此，本次讨论的准确定位不是“全部完成”，而是：
\[
  \boxed{
    \text{完成了总架构定型与主权边界定型。}
  }
\]
\end{proposition}

\begin{Certificate}[证书：LLM 降级 + 成本残留]
\label{cert:tex37-ch10-render-pending}
证书登记谓词链：
\[
\begin{aligned}
  P_1 &: \text{第~\ref{sec:tex37-ch6} 章与第~\ref{sec:tex37-ch9} 章共同给出：LLM 只承担末端表达},\\
  P_2 &: P_1\Rightarrow \text{系统可控性上升、模型替换成本下降、fail-closed 边界清晰},\\
  P_3 &: \text{仍存在若干未完全构造的正式协议、可计算准则与执行接口},\\
  P_4 &: P_2\wedge P_3\Rightarrow G_{\mathrm{render}}\ \text{极高但}\ C_{\mathrm{pending}}>0.
\end{aligned}
\]
\end{Certificate}

\begin{proof}[证明（基于数学符号推演逻辑的谓词因果链）]
由 $P_1$，LLM 的职责已经从开放式求解降级为受审计约束的表面展开。
于是由 $P_2$，系统的可控性、可替换性与失败显式边界都显著提升，
这正构成表达插件化带来的工程增益，因此 $G_{\mathrm{render}}$ 为极高。

但由 $P_3$，体系仍有若干正式接口与实现环节尚待补齐。
故 $C_{\mathrm{pending}}>0$。
将两者合并即得 $P_4$：本次讨论已经完成总架构与主权边界定型，但并不等于所有实现细节均已完成。证毕。
\end{proof}

\subsection{推论：本次讨论属于“架构级重大增益”，其重量高于普通局部命题增益}
\label{subsec:tex37-ch10-corollary}

\begin{corollary}[本次讨论的总体评级]
\label{cor:tex37-ch10-overall}
在定理~\ref{thm:tex37-ch10-arch}、定理~\ref{thm:tex37-ch10-sovereignty}、引理~\ref{lem:tex37-ch10-plugin}、
定理~\ref{thm:tex37-ch10-cert-conv} 与命题~\ref{prop:tex37-ch10-render-pending} 的条件下，
本次讨论满足：
\[
  \operatorname{Rank}(G_{\mathrm{arch}})=\text{极高},
  \qquad
  \operatorname{Rank}(G_{\mathrm{sov}})=\text{极高},
  \qquad
  \operatorname{Rank}(G_{\mathrm{conv}})=\text{极高},
\]
\[
  \operatorname{Rank}(G_{\mathrm{render}})=\text{极高},
  \qquad
  \operatorname{Rank}(G_{\mathrm{plugin}})=\text{很高},
  \qquad
  \operatorname{Rank}(G_{\mathrm{cert}})=\text{很高},
\]
且
\[
  C_{\mathrm{pending}}>0.
\]
因此，本次讨论的准确评价应写成：
\[
  \boxed{
    \text{本次讨论属于“架构级重大增益”，其重量明显高于普通局部命题增益。}
  }
\]
\end{corollary}

\begin{Certificate}[证书：六类增益评级 + 剩余成本约束]
\label{cert:tex37-ch10-overall}
证书登记谓词链：
\[
\begin{aligned}
  O_1 &: \text{定理~\ref{thm:tex37-ch10-arch}}\Rightarrow \operatorname{Rank}(G_{\mathrm{arch}})=\text{极高},\\
  O_2 &: \text{定理~\ref{thm:tex37-ch10-sovereignty}}\Rightarrow \operatorname{Rank}(G_{\mathrm{sov}})=\text{极高},\\
  O_3 &: \text{引理~\ref{lem:tex37-ch10-plugin}}\Rightarrow \operatorname{Rank}(G_{\mathrm{plugin}})=\text{很高},\\
  O_4 &: \text{定理~\ref{thm:tex37-ch10-cert-conv}}\Rightarrow \operatorname{Rank}(G_{\mathrm{cert}})=\text{很高}\ \wedge\
    \operatorname{Rank}(G_{\mathrm{conv}})=\text{极高},\\
  O_5 &: \text{命题~\ref{prop:tex37-ch10-render-pending}}\Rightarrow \operatorname{Rank}(G_{\mathrm{render}})=\text{极高}\
    \wedge\ C_{\mathrm{pending}}>0,\\
  O_6 &: O_1\wedge O_2\wedge O_3\wedge O_4\wedge O_5
  \Rightarrow
  \text{总体评价为架构级重大增益}.
\end{aligned}
\]
\end{Certificate}

\begin{proof}[证明（基于数学符号推演逻辑的谓词因果链）]
由 $O_1$ 与 $O_2$，总架构闭环化与主权回收化都已达到极高权重。
由 $O_3$，领域扩张已被稳定压成塔纤维插件，而不是语言层外挂。
由 $O_4$，这些插件进一步升级为证书认证且接口收束的富结构对象。
由 $O_5$，表达层插件化带来极强工程增益，但体系仍保留诚实记账的剩余成本。

因此，本次讨论并非“全部完成”，却已完成了决定体系命运的主梁与承重墙：
\[
  \begin{aligned}
    &\text{修复塔求解链闭环},
    \quad
    \text{主权边界回收},
    \quad
    \text{领域插件化},\\
    &\text{可签发化},
    \quad
    \text{接口收束化},
    \quad
    \text{表达插件化}.
  \end{aligned}
\]
这正对应“架构级重大增益”的评价。证毕。
\end{proof}

\subsection{备注：统一口径、关键桥梁与边界}
\label{subsec:tex37-ch10-remarks}

\begin{remark}[本次讨论最重要的不是再证明一个局部强命题，而是统一口径]
本次讨论真正的重量在于，它把许多原本分散但很强的判断，统一压成了一句总口径：
\[
  \boxed{
    \begin{aligned}
      &\text{G 框架是白盒求解内核；领域知识是塔纤维插件；证书负责签发；}\\
      &\text{接口必须收束；LLM 只是表达插件。}
    \end{aligned}
  }
\]
这句话本身，就是一次体系压缩。
\end{remark}

\begin{remark}[若以“距离自动化科研/自动化设计主架构还差多少”来衡量]
本次讨论已经不再只是
\[
  \text{“有理论”},
\]
而是进一步到达
\[
  \text{“知道系统该怎么搭”}.
\]
这通常比单个局部证明更难，也更有结构价值，因此它应被视为关键桥梁而非普通增量。
\end{remark}

\begin{remark}[边界与诚实记账]
本章的总评并不意味着一切都已完成。
仍需继续补齐的环节包括：
\begin{enumerate}
  \item 各领域插件的正式接口清单；
  \item 公理降级为定理的统一签发协议；
  \item 细化层级收束的可计算准则；
  \item 正规形到动作执行器的工程实现；
  \item 占位覆盖度与文本充分性的形式化度量。
\end{enumerate}
因此最准确的收尾口径是：
\[
  \boxed{
    \text{本次讨论完成了总架构定型与主权边界定型，但仍保留若干关键接口的实现成本。}
  }
\]
\end{remark}

\begin{remark}[一句话结论]
若只用一句话给出总判断，则可写成：
\[
  \boxed{
    \begin{aligned}
      &\text{这次讨论不是“补了一块砖”，}\\
      &\text{而是把白盒 AI/自动化科研内核的主梁和承重墙一起立起来了。}
    \end{aligned}
  }
\]
\end{remark}

\clearpage

%%%%%%%%%%%%%%%%%%%%%%%%%%%%%%%%%%%%%%%%%%%%%%%%%%%%%%%%%%%%
\section{形式逻辑：从 \texttt{tex36} 的求解法到主基底--扇区插件--签发收束协议的可签发方程系统}
\label{sec:tex37-ch11}
%%%%%%%%%%%%%%%%%%%%%%%%%%%%%%%%%%%%%%%%%%%%%%%%%%%%%%%%%%%%

\begin{remark}[核心陈述（系统编写版）]
若把 `tex36` 严格理解为给出
\[
  x_0\mapsto \Sigma_N(x_0)\mapsto \operatorname{NF}_{A_N}(x_N)
\]
这一类白盒求解链，而把本稿第一章、第六章到第九章理解为给出
\[
  \mathcal E_{\mathrm{core}}^{(N)},
  \qquad
  \mathcal P_D,
  \qquad
  \widehat{\mathfrak R}_D^{\mathrm{IC}},
  \qquad
  p^\star\mapsto t
\]
这一类系统对象边界与表达闭环，
则当前剩余主任务已不再是“再想一个总框架”，而是把这些对象压成
\[
  \boxed{
    \text{方程系统 / 闭环求解系统 / 可签发的求解--裁决--表达协议}
  }
\]
并逐一签发。
\end{remark}

\begin{remark}[阅读导航：本章的引用口径与新增目标]
本章关于“如何求解”的对象链，直接承接 `tex36` 的全编码求解链、稳定同伦扭曲正规形与有效作用代表口径 \cite{RepoTex36HomotopyTwistBijection}；
关于“占位--解析--审计”的表达闭环，直接承接本稿第一章与仓内整理稿 \cite{RepoTex20SemanticGaugeNormalization,RepoTex21AuditedWrapper,
  RepoTex23SemanticBase}；
关于“主权回收、领域插件、认证富结构与接口收束”，直接承接本稿第~\ref{sec:tex37-ch6} 章到第~\ref{sec:tex37-ch9} 章，以及 `tex35` 的扩展签发协议
  \cite{RepoTex35ExtendedCertificationTower}。

外部标准参照方面，本章继续使用范畴、纤维丛、同调与联络的标准语言 \cite{MacLane1998Categories,Hatcher2002AT,Weibel1994HomologicalAlgebra,
  KobayashiNomizu1963Foundations}。
本章新增的不是新的本体宇宙，而是把“求解法已经给出、系统边界已经给出”这一判断，进一步压成可执行的方程化任务书。
\end{remark}

\subsection{定义：闭环求解系统、主基底方程系统与扇区插件方程系统}
\label{subsec:tex37-ch11-definitions}

\begin{definition}[闭环求解系统与可签发方程系统]
\label{def:tex37-ch11-signable-system}
称一个数据组
\[
  \mathfrak E
  =
  \bigl(
    \mathcal B,
    \mathcal X,
    \mathcal O,
    \mathcal R,
    \operatorname{NF},
    \operatorname{Act},
    \mathcal I^{\mathrm{stmt}},
    \Pi_{\mathrm{expr}}
  \bigr)
\]
为一个\textbf{闭环求解系统}，若满足：
\begin{enumerate}
  \item $\mathcal B$ 是问题所依附的基底，$\mathcal X$ 是状态空间；
  \item $\mathcal O$ 是障碍空间；
  \item
  \[
    \mathcal R\subseteq \mathcal X\times \mathcal O\times \mathcal X
  \]
  是修复关系，编码“给定状态与障碍后，如何到达下一可接受状态”；
  \item
  \[
    \operatorname{NF}:\mathcal X\to \mathcal T\sqcup\{\bot\}
  \]
  给出正规形或失败符号；
  \item
  \[
    \operatorname{Act}:\mathcal T\to \mathcal A
  \]
  给出动作选择律；
  \item
  \[
    \Pi_{\mathrm{expr}}
    :=
    \bigl(
      \Theta_0,
      \Render^{\mathrm{LLM}}_\xi,
      \Parse_\xi,
      \Audit_\xi
    \bigr)
  \]
  给出从正规形/占位到文本外显的投影闭环。
\end{enumerate}

若进一步满足：对任意接口命题
\[
  \alpha\in \mathcal I^{\mathrm{stmt}},
\]
都满足
\[
  \exists\,(c_\alpha,T_\alpha),\ \Issue(c_\alpha,T_\alpha)=\top,\ T_\alpha\equiv \alpha
  \qquad\text{或}\qquad
  \ExtAssump(\alpha)=\top,
\]
则称 $\mathfrak E$ 为一个\textbf{可签发方程系统}。
\end{definition}

\begin{definition}[主基底方程系统、扇区插件方程系统与跨扇区耦合]
\label{def:tex37-ch11-core-sector}
定义主基底为
\[
  \mathcal B_{\mathrm{core}}:=\mathcal B_{\mathrm{ctx}},
\]
并把第~\ref{sec:tex37-ch6} 章与第~\ref{sec:tex37-ch7} 章中已经固定的核心对象压成主基底方程系统
\[
  \mathfrak E_{\mathrm{core}}
  :=
  \bigl(
    \mathcal B_{\mathrm{core}},
    \mathcal X_{\mathrm{core}},
    \mathcal O_{\mathrm{core}},
    \mathcal R_{\mathrm{core}},
    \operatorname{NF}_{\mathrm{core}},
    \operatorname{Act}_{\mathrm{core}},
    \mathcal I_{\mathrm{core}}^{\mathrm{stmt}},
    \Pi_{\mathrm{expr}}
  \bigr),
\]
其中
\[
  \mathcal X_{\mathrm{core}}\subseteq \mathcal E_{\mathrm{core}}^{(N)}.
\]

对每个领域（扇区）$D\in\mathfrak D_*$，定义其扇区插件方程系统为
\[
  \mathfrak E_D
  :=
  \bigl(
    \mathcal B_{\mathrm{core}},
    \mathcal X_D,
    \mathcal O_D,
    \mathcal R_D,
    \operatorname{NF}_D,
    \operatorname{Act}_D,
    \mathcal I_D^{\mathrm{stmt}},
    \Pi_{\mathrm{expr},D}
  \bigr),
\]
并要求
\[
  \mathcal P_D\models \mathfrak C_{\mathrm{core}},
  \qquad
  \widehat{\mathfrak R}_D^{\mathrm{IC}}
  =
  (\mathfrak R_D,\Cert_D,\Auth_D,\Conv_D)
\]
已被定义。

若再对任意不同扇区 $D_i,D_j\in\mathfrak D_*$ 给出跨扇区耦合数据
\[
  \mathfrak C_{D_i,D_j},
\]
用以记录共享变量、一致性约束、证书兼容性与动作互锁条件，
则称
\[
  \mathfrak E_{\mathrm{sys}}
  :=
  \Bigl(
    \mathfrak E_{\mathrm{core}},
    \{\mathfrak E_D\}_{D\in\mathfrak D_*},
    \{\mathfrak C_{D_i,D_j}\}_{D_i,D_j\in\mathfrak D_*},
    \{\widehat{\mathfrak R}_D^{\mathrm{IC}}\}_{D\in\mathfrak D_*}
  \Bigr)
\]
为\textbf{全局可签发方程系统}。
\end{definition}

\begin{remark}[这里的“方程系统”不是单个 $F(x)=0$]
按定义~\ref{def:tex37-ch11-signable-system}--\ref{def:tex37-ch11-core-sector}，
这里真正要写出来的对象是
\[
  \text{状态空间}
  +
  \text{障碍}
  +
  \text{修复}
  +
  \text{正规形}
  +
  \text{动作}
  +
  \text{证书}
  +
  \text{表达投影}
\]
共同组成的系统。
因此更准确的名称不是“某个单独方程”，而是
\[
  \boxed{
    \text{闭环求解系统 / 可签发的求解--裁决--表达协议}
  }.
\]
\end{remark}

\subsection{公理（降级为定理）：有限扇区族的可签发方程系统可按插件数递归编译}
\label{subsec:tex37-ch11-axiom}

\begin{theorem}[公理降级：有限扇区族的可签发方程系统可按插件数递归编译（数学归纳法证书）]
\label{thm:tex37-ch11-induction}
设
\[
  \mathfrak D_M
  :=
  \{D_1,\dots,D_M\}
  \subseteq \mathfrak D.
\]
若满足：
\begin{enumerate}
  \item 主基底方程系统
  \[
    \mathfrak E_{\mathrm{core}}
  \]
  已被定义且可签发；
  \item 对任意 $m\in\{0,1,\dots,M-1\}$，只要前 $m$ 个扇区的部分系统
  \[
    \mathfrak E_{\mathrm{sys}}^{(m)}
    :=
    \Bigl(
      \mathfrak E_{\mathrm{core}},
      \{\mathfrak E_{D_k}\}_{k=1}^{m},
      \{\mathfrak C_{D_i,D_j}\}_{1\le i,j\le m},
      \{\widehat{\mathfrak R}_{D_k}^{\mathrm{IC}}\}_{k=1}^{m}
    \Bigr)
  \]
  已被定义且可签发，就存在唯一的扩张算子
  \[
    \operatorname{Compile}_{m+1}:
    \mathfrak E_{\mathrm{sys}}^{(m)}
    \longmapsto
    \mathfrak E_{\mathrm{sys}}^{(m+1)},
  \]
  使得新加入的扇区 $D_{m+1}$ 同时带入
  \[
    \mathfrak E_{D_{m+1}},
    \qquad
    \widehat{\mathfrak R}_{D_{m+1}}^{\mathrm{IC}},
    \qquad
    \{\mathfrak C_{D_k,D_{m+1}}\}_{k=1}^{m},
  \]
  并保持整体系统可签发。
\end{enumerate}
则对任意 $m\in\{0,1,\dots,M\}$，都存在唯一的部分系统
\[
  \mathfrak E_{\mathrm{sys}}^{(m)},
\]
特别地，存在唯一的有限全局可签发方程系统
\[
  \mathfrak E_{\mathrm{sys}}^{(M)}.
\]
\end{theorem}

\begin{Certificate}[证书：基步 + 归纳步（数学归纳法证书）]
\label{cert:tex37-ch11-induction}
定义谓词
\[
  P(m)
  :
  \Longleftrightarrow
  \exists!\,\mathfrak E_{\mathrm{sys}}^{(m)}
  \ \text{使其可签发}.
\]
则数学归纳法证书记为
\[
  \pi_{\mathrm{compile}}
  =
  \bigl(
    P(0),\
    \forall m\in\{0,\dots,M-1\},\ P(m)\Rightarrow P(m+1)
  \bigr).
\]
其中归纳步的见证对象为
\[
  \operatorname{Compile}_{m+1},
  \qquad
  \mathfrak E_{D_{m+1}},
  \qquad
  \widehat{\mathfrak R}_{D_{m+1}}^{\mathrm{IC}},
  \qquad
  \{\mathfrak C_{D_k,D_{m+1}}\}_{k=1}^{m}.
\]
\end{Certificate}

\begin{proof}[证明（数学归纳法证书；基于数学符号推演逻辑的谓词因果链）]
定义谓词
\[
  P(m)
  :
  \Longleftrightarrow
  \exists!\,\mathfrak E_{\mathrm{sys}}^{(m)}
  \ \text{使其可签发}.
\]

\textbf{基步：}
当 $m=0$ 时，
\[
  \mathfrak E_{\mathrm{sys}}^{(0)}:=\mathfrak E_{\mathrm{core}}.
\]
由条件(1)，主基底方程系统已经存在且可签发，
故 $P(0)$ 成立。

\textbf{归纳步：}
设 $P(m)$ 成立，即唯一的
\[
  \mathfrak E_{\mathrm{sys}}^{(m)}
\]
已经存在且可签发。
由条件(2)，存在唯一扩张算子
\[
  \operatorname{Compile}_{m+1}
\]
把它送到
\[
  \mathfrak E_{\mathrm{sys}}^{(m+1)}.
\]
该扩张不仅加入新的扇区方程系统
\[
  \mathfrak E_{D_{m+1}},
\]
还同时加入其接口收束认证对象
\[
  \widehat{\mathfrak R}_{D_{m+1}}^{\mathrm{IC}}
\]
以及与既有扇区的全部耦合约束
\[
  \{\mathfrak C_{D_k,D_{m+1}}\}_{k=1}^{m}.
\]
由于扩张被要求保持整体系统可签发，故新系统
\[
  \mathfrak E_{\mathrm{sys}}^{(m+1)}
\]
存在。
又由于扩张算子唯一，故该系统亦唯一，于是 $P(m+1)$ 成立。

由数学归纳法，$P(m)$ 对所有 $m=0,\dots,M$ 都成立。
因此有限扇区族的可签发方程系统可按插件数递归编译。证毕。
\end{proof}

\subsection{定理：\texttt{tex36} 给出“如何求解”的结构，而本稿前文给出“应写什么系统”的结构}
\label{subsec:tex37-ch11-theorem}

\begin{theorem}[求解法--写法分离定理]
\label{thm:tex37-ch11-solve-vs-write}
在 `tex36` 的全编码求解链口径与本稿第一章、第六章到第九章的占位闭环、主权回收、扇区插件、认证富结构、接口收束口径同时成立时，可得如下严格分离：
\[
  \boxed{
    \texttt{tex36}
    \Longrightarrow
    \text{给出“如何求解”的结构}
  }
\]
\[
  \boxed{
    \text{本稿第~\ref{sec:tex37-ch1} 章与第~\ref{sec:tex37-ch6}--\ref{sec:tex37-ch9} 章}
    \Longrightarrow
    \text{给出“应写什么方程系统”的结构}
  }.
\]
更具体地说，前者固定了
\[
  x_0
  \mapsto
  \Sigma_N(x_0)
  \mapsto
  \tau^\star
  \mapsto
  a^\star
\]
的求解--裁决链，
后者固定了
\[
  \mathfrak E_{\mathrm{core}},
  \qquad
  \{\mathfrak E_D\}_{D\in\mathfrak D_*},
  \qquad
  \{\widehat{\mathfrak R}_D^{\mathrm{IC}}\}_{D\in\mathfrak D_*},
  \qquad
  \Pi_{\mathrm{expr}}
\]
应具有的结构边界。
\end{theorem}

\begin{Certificate}[证书：求解链固定 + 写法边界固定]
\label{cert:tex37-ch11-solve-vs-write}
证书登记谓词链：
\[
\begin{aligned}
  T_1 &: \texttt{tex36}\ \text{给出}\ x_0\mapsto \Sigma_N(x_0)\mapsto \operatorname{NF}_{A_n}(x_n),\\
  T_2 &: \text{第~\ref{sec:tex37-ch6} 章把上述链压成白盒求解--裁决程序链},\\
  T_3 &: \text{第~\ref{sec:tex37-ch7} 章给出}\ \mathcal P_D\ \text{的主基底服从性与扇区插件结构},\\
  T_4 &: \text{第~\ref{sec:tex37-ch8}--\ref{sec:tex37-ch9} 章给出}\ \widehat{\mathfrak R}_D^{\mathrm{IC}}\ \text{的认证与收束边界},
    \\
  T_5 &: \text{第~\ref{sec:tex37-ch1} 章给出}\ p^\star\mapsto t\ \text{的占位--解析--审计闭环},\\
  T_6 &: T_1\wedge T_2\wedge T_3\wedge T_4\wedge T_5
  \Rightarrow
  \text{“求解法”与“方程系统写法”严格分离}.
\end{aligned}
\]
\end{Certificate}

\begin{proof}[证明（基于数学符号推演逻辑的谓词因果链）]
由 $T_1$，`tex36` 已经把从初态类
\[
  x_0
\]
到全编码求解链
\[
  \Sigma_N(x_0)
\]
再到稳定正规形
\[
  \tau^\star=\operatorname{NF}_{A_n}(x_n)
\]
的求解机制固定下来。
这说明 `tex36` 的主要贡献是“如何从给定入口出发得到可接受代表与动作候选”。

由 $T_2$，第~\ref{sec:tex37-ch6} 章进一步把这一机制压成程序确定性的白盒求解--裁决链，
因此“如何解”不再只是抽象口号，而是已有的可执行对象链。

由 $T_3$，第~\ref{sec:tex37-ch7} 章已经规定了领域扩张必须写成同一主基底上的扇区插件，
而不是外挂规则包。
由 $T_4$，第~\ref{sec:tex37-ch8}--\ref{sec:tex37-ch9} 章又把这些插件收紧为可认证、可签发、可收束的对象。
由 $T_5$，第一章再把最终外显统一压到占位--解析--审计闭环中。

因此由 $T_6$，`tex36` 所完成的是求解法层，
而本稿前文所完成的是方程系统写法层。
两者叠加后，剩余工作自然不再是再发明总框架，而是把已知结构逐项实例化为可签发方程系统。证毕。
\end{proof}

\subsection{引理：剩余主任务不是再造总框架，而是把总框架压成主基底层、扇区层与接口收束层的方程族}
\label{subsec:tex37-ch11-lemma}

\begin{lemma}[三层编译引理]
\label{lem:tex37-ch11-three-layers}
在定理~\ref{thm:tex37-ch11-solve-vs-write} 的条件下，
剩余未实例化的核心对象可被压缩为如下三层：
\[
  \text{主基底层}
  \longrightarrow
  \text{扇区层}
  \longrightarrow
  \text{接口收束层}.
\]
更具体地说，待编译的对象正是
\[
  \mathfrak E_{\mathrm{core}},
  \qquad
  \{\mathfrak E_D\}_{D\in\mathfrak D_*},
  \qquad
  \{\widehat{\mathfrak R}_D^{\mathrm{IC}}\}_{D\in\mathfrak D_*},
  \qquad
  \{\mathfrak C_{D_i,D_j}\}_{D_i,D_j\in\mathfrak D_*}.
\]
因此剩余任务不是再构造新的元框架，而是把既有元框架压成具体方程族。
\end{lemma}

\begin{Certificate}[证书：核心合同已定 + 插件结构已定 + 收束协议已定]
\label{cert:tex37-ch11-three-layers}
证书登记谓词链：
\[
\begin{aligned}
  L_1 &: \mathfrak C_{\mathrm{core}}\ \text{已由第~\ref{sec:tex37-ch7} 章固定},\\
  L_2 &: \mathcal P_D\ \text{的扇区插件结构已由第~\ref{sec:tex37-ch7} 章固定},\\
  L_3 &: \widehat{\mathfrak R}_D^{\mathrm{IC}}\ \text{的认证与收束协议已由第~\ref{sec:tex37-ch8}--\ref{sec:tex37-ch9} 章固定},\\
  L_4 &: \Pi_{\mathrm{expr}}\ \text{已由第~\ref{sec:tex37-ch1} 章固定},\\
  L_5 &: L_1\wedge L_2\wedge L_3\wedge L_4
  \Rightarrow
  \text{剩余自由度只落在具体方程实例化，而不落在新的元框架选择}.
\end{aligned}
\]
\end{Certificate}

\begin{proof}[证明（基于数学符号推演逻辑的谓词因果链）]
由 $L_1$，主基底层的硬边界已经写成核心合同，
因此“状态接口、障碍接口、修复接口、动作接口、证书接口、占位接口、审计接口”这些类型边界不再开放。

由 $L_2$，扇区扩张的正确写法也已被固定为
\[
  \mathcal P_D\models \mathfrak C_{\mathrm{core}}
\]
的同基底纤维插件，
故扇区层不允许回退成随意外挂规则。

由 $L_3$，接口收束层进一步被固定为
\[
  \widehat{\mathfrak R}_D^{\mathrm{IC}}
\]
这种带签发、带细化、带闭合的富结构对象，
因此接口层也不允许以“未签发内部公理”或“无限细化而不收束”的方式悬空。

再由 $L_4$，最终表达层已经被压成占位--解析--审计闭环。
故由 $L_5$，整个元框架的骨架已经完成；
剩余工作只是在该骨架内，把主基底层、扇区层与接口收束层分别写成具体方程族。证毕。
\end{proof}

\subsection{命题：所谓“剩下的就是写方程”，其严格含义是写可解、可审计、可签发、可收束、可执行的方程系统}
\label{subsec:tex37-ch11-proposition}

\begin{proposition}[可签发方程系统完成性命题]
\label{prop:tex37-ch11-complete-equations}
对任意领域（扇区）$D\in\mathfrak D_*$，
若某个候选系统只写出形式符号
\[
  F_D(x)=0,
\]
却未同时显式给出
\[
  \mathcal X_D,
  \qquad
  \mathcal O_D,
  \qquad
  \mathcal R_D,
  \qquad
  \operatorname{NF}_D,
  \qquad
  \operatorname{Act}_D,
  \qquad
  \mathcal I_D^{\mathrm{stmt}},
  \qquad
  \Pi_{\mathrm{expr},D},
\]
则它不能被称为“扇区方程系统已写完”。
更严格地说，一个扇区插件最终至少要回答下列问题：
\[
  \text{状态是什么？}
  \quad
  \text{障碍是什么？}
  \quad
  \text{修复算子是什么？}
  \quad
  \text{正规形是什么？}
\]
\[
  \text{动作选择律是什么？}
  \quad
  \text{证书如何签发？}
  \quad
  \text{如何投影到可审计表达？}
\]
否则该系统至多是未闭合的符号草图，而不是可签发方程系统。
\end{proposition}

\begin{Certificate}[证书：可解 + 可审计 + 可签发 + 可收束 + 可执行]
\label{cert:tex37-ch11-complete-equations}
证书登记谓词链：
\[
\begin{aligned}
  P_1 &: \mathcal X_D,\mathcal O_D,\mathcal R_D\ \text{缺失}
  \Rightarrow
  \text{状态与障碍无从定义，系统不可解},\\
  P_2 &: \operatorname{NF}_D\ \text{缺失}
  \Rightarrow
  \text{系统无正规代表，无法收束到可比较结果},\\
  P_3 &: \operatorname{Act}_D\ \text{缺失}
  \Rightarrow
  \text{系统无法从解代表过渡到动作执行},\\
  P_4 &: \mathcal I_D^{\mathrm{stmt}}\ \text{的签发规则缺失}
  \Rightarrow
  \text{系统不可审计且不可签发},\\
  P_5 &: \Pi_{\mathrm{expr},D}\ \text{缺失}
  \Rightarrow
  \text{系统无法闭合到对外表达接口},\\
  P_6 &: P_1\vee P_2\vee P_3\vee P_4\vee P_5
  \Rightarrow
  \text{候选对象不是完整的可签发方程系统}.
\end{aligned}
\]
\end{Certificate}

\begin{proof}[证明（基于数学符号推演逻辑的谓词因果链）]
若只给出
\[
  F_D(x)=0
\]
而没有给出状态空间、障碍空间与修复关系，
则连“未知量是什么、阻碍是什么、通过什么算子修复”都无法回答，
于是 $P_1$ 成立，系统不可解。

即便某个候选符号可以形式求值，若没有
\[
  \operatorname{NF}_D,
\]
就无法把结果压到唯一可比较代表，故 $P_2$ 成立，系统不可收束。
若没有
\[
  \operatorname{Act}_D,
\]
则求解链无法进入执行链，故 $P_3$ 成立，系统不可执行。

再者，若接口命题没有被压入
\[
  \mathcal I_D^{\mathrm{stmt}}
\]
及其签发规则，
则无法区分“已签发定理”与“外部导入前提”，故 $P_4$ 成立，系统不可审计且不可签发。
最后，若没有
\[
  \Pi_{\mathrm{expr},D},
\]
则内部结果不能以占位--解析--审计闭环的方式投影到外显接口，故 $P_5$ 成立。

因此只要上述任何一项缺失，就只能得到未闭合的符号草图，而得不到完整的可签发方程系统。
故 $P_6$ 成立。命题得证。
\end{proof}

\subsection{推论：体系已进入系统编写/编译阶段，而下一步应先写主基底方程系统}
\label{subsec:tex37-ch11-corollary}

\begin{corollary}[系统编写阶段推论]
\label{cor:tex37-ch11-next-step}
由定理~\ref{thm:tex37-ch11-solve-vs-write}、引理~\ref{lem:tex37-ch11-three-layers} 与命题~\ref{prop:tex37-ch11-complete-equations}
  ，
可得当前阶段的最准确定位为：
\[
  \boxed{
    \text{体系已经从“求解理论整理阶段”推进到“方程系统编写/编译阶段”。}
  }
\]
并且下一步最自然的工作顺序应写成
\[
  \mathfrak E_{\mathrm{core}}
  \Longrightarrow
  \{\mathfrak E_D\}_{D\in\mathfrak D_*}
  \Longrightarrow
  \{\mathfrak C_{D_i,D_j}\}_{D_i,D_j\in\mathfrak D_*}
  \Longrightarrow
  \{\widehat{\mathfrak R}_D^{\mathrm{IC}}\}_{D\in\mathfrak D_*}.
\]
特别地，主基底方程系统
\[
  \mathfrak E_{\mathrm{core}}
\]
应被优先写出，因为它是全部扇区插件、耦合方程与签发协议的锚点。
\end{corollary}

\begin{Certificate}[证书：阶段转移 + 编译顺序 + 主基底优先]
\label{cert:tex37-ch11-next-step}
证书登记谓词链：
\[
\begin{aligned}
  C_1 &: \text{定理~\ref{thm:tex37-ch11-solve-vs-write}}\Rightarrow \text{求解法层与写法层已被严格分离},\\
  C_2 &: \text{引理~\ref{lem:tex37-ch11-three-layers}}\Rightarrow \text{剩余对象被压缩为主基底层、扇区层与接口收束层},\\
  C_3 &: \text{命题~\ref{prop:tex37-ch11-complete-equations}}\Rightarrow \text{每一层都必须写成完整可签发方程系统},\\
  C_4 &: C_1\wedge C_2\wedge C_3
  \Rightarrow
  \text{体系已进入系统编写/编译阶段},\\
  C_5 &: \mathfrak E_{\mathrm{core}}\ \text{为全部扇区插件与耦合条件的锚点}
  \Rightarrow
  \text{应先写}\ \mathfrak E_{\mathrm{core}}.
\end{aligned}
\]
\end{Certificate}

\begin{proof}[证明（基于数学符号推演逻辑的谓词因果链）]
由 $C_1$，当前已不存在“先弄清楚到底怎么求解”这一元层不确定性；
求解法层已经由 `tex36` 及本稿第六章固定。
由 $C_2$，剩余对象也不再是无限开放的，而被压缩为三层可列举对象。
由 $C_3$，这些对象还必须满足完整性、可签发性与表达闭环约束。
因此由 $C_4$，当前阶段的准确名称只能是“系统编写/编译阶段”。

再由 $C_5$，扇区插件必须服从
\[
  \mathfrak C_{\mathrm{core}}
\]
并依附于主基底；
跨扇区耦合也必须先知道主基底变量、核心障碍、核心修复链与核心动作接口如何写出，
才能定义何谓“共享变量”“兼容正规形”“一致性约束”“可签发耦合定理”。
故主基底方程系统必须优先完成。推论得证。
\end{proof}

\subsection{备注：推荐口径、任务顺序与边界}
\label{subsec:tex37-ch11-remarks}

\begin{remark}[更严格的一句话版本]
若把本章压成一句最短的严格表述，则可写成：
\[
  \boxed{
    \begin{aligned}
      &\texttt{tex36}\ \text{给了“怎么求解”，本稿前文给了“该写什么系统”，}\\
      &\text{剩下的就是把主基底与各扇区插件编译成可签发的方程系统。}
    \end{aligned}
  }
\]
\end{remark}

\begin{remark}[推荐采用的工程任务顺序]
若把剩余工作直接写成工程任务清单，则最合适的顺序是：
\[
  \boxed{
    \begin{aligned}
      &\text{先写主基底方程系统，再写扇区插件方程系统，}\\
      &\text{再写跨扇区耦合与签发协议，最后写接口收束与表达审计条件。}
    \end{aligned}
  }
\]
这不是经验排序，而是由主基底锚点性、扇区依附性与接口收束的后置闭环性共同决定的。
\end{remark}

\begin{remark}[边界：不能把“写方程”误写成“堆符号”]
本章最需要钉死的边界是：
\[
  \boxed{
    \begin{aligned}
      &\text{所谓“写方程”，不是写一串漂亮符号，}\\
      &\text{而是写出可解、可审计、可签发、可收束、可执行的系统。}
    \end{aligned}
  }
\]
若缺少状态定义、障碍定义、修复关系、正规形、动作选择律、签发规则或表达投影闭环中的任一项，
则得到的都不是可部署系统，只是尚未闭合的形式草图。
\end{remark}

\clearpage

%%%%%%%%%%%%%%%%%%%%%%%%%%%%%%%%%%%%%%%%%%%%%%%%%%%%%%%%%%%%
\section{形式逻辑：主基底方程系统与扇区插件方程系统的统一结构签名、六层闭环与规则签发后置}
\label{sec:tex37-ch12}
%%%%%%%%%%%%%%%%%%%%%%%%%%%%%%%%%%%%%%%%%%%%%%%%%%%%%%%%%%%%

\begin{remark}[核心陈述（结构签名版）]
承接第~\ref{sec:tex37-ch11} 章关于“系统编写/编译阶段”的结论，本章进一步把主基底方程系统与扇区插件方程系统的共同骨架压成如下严格口径：
\[
  \boxed{
    \text{主基底方程系统或扇区插件方程系统的本体，不是“规则表”或“规则堆积”，}
  }
\]
\[
  \boxed{
    \text{而是由算子及其自由度、纤维化状态空间、同伦扭曲/联络演化、允许路径类、}
  }
\]
\[
  \boxed{
    \text{路径积分选类、正规形收束以及定理（证书 + 证明）签发规则共同组成的层级闭环系统。}
  }
\]
在这一口径下，路径积分并不直接等于规则接口；只有在路径类经过正规形提取、证书构造与签发闭合后，
才上升为可对外调用的规则/定理接口。
\end{remark}

\begin{remark}[阅读导航：仓内引用与承接关系]
本章直接承接 `tex36` 中关于主丛联络、路径积分支撑类、稳定同伦扭曲正规形与有效作用选择的对象链 \cite{RepoTex36HomotopyTwistBijection}；
承接本稿第~\ref{sec:tex37-ch4} 章关于“问答内核对象是上下文基底上的允许路径类/修复链对象”的口径，
以及第~\ref{sec:tex37-ch7}--\ref{sec:tex37-ch9} 章关于主基底、扇区插件、认证富结构与接口收束的边界。
对于占位--解析--审计闭环，继续沿用第一章的逻辑占位投影链 \cite{RepoTex20SemanticGaugeNormalization,RepoTex21AuditedWrapper,RepoTex23SemanticBase}。

外部标准参照方面，本章继续使用纤维丛、联络、同伦/同调与范畴化接口的标准语言 \cite{MacLane1998Categories,Hatcher2002AT,Weibel1994HomologicalAlgebra,
  KobayashiNomizu1963Foundations}。
本章新增的目标不是重述“如何求解”，而是进一步回答：主基底与扇区插件的可签发方程系统究竟应按什么统一结构签名来写。
\end{remark}

\subsection{定义：主基底或扇区方程系统的统一结构签名}
\label{subsec:tex37-ch12-def1}

\begin{definition}[统一结构签名]
\label{def:tex37-ch12-signature}
对任意层级对象
\[
  \star\in \{\mathrm{base}\}\cup\mathfrak D,
\]
定义其方程系统的统一结构签名为
\[
  \mathfrak E_\star
  :=
  \Bigl(
    \mathcal B_\star,
    \pi_\star:\mathcal E_\star\to \mathcal B_\star,
    \mathcal O_\star,
    \mathcal U_\star,
    A_\star,
    F_\star,
    \Ob_\star,
    \Gamma_\star,
    \mathcal S_\star,
    \operatorname{NF}_\star,
    \operatorname{Rep}_\star,
    \Theta_\star,
    \Cert_\star,
    \Issue_\star
  \Bigr).
\]
其中：
\begin{enumerate}
  \item $\mathcal B_\star$ 是主基底或扇区基底；
  \item
  \[
    \pi_\star:\mathcal E_\star\to \mathcal B_\star
  \]
  是其纤维丛状态空间；
  \item $\mathcal O_\star,\mathcal U_\star$ 分别是算子族与自由度族；
  \item $A_\star,F_\star$ 分别是联络与曲率；
  \item
  \[
    \Ob_\star:\mathcal E_\star\to \mathcal W_\star
  \]
  是障碍函数；
  \item $\Gamma_\star$ 是允许路径类集合；
  \item
  \[
    \mathcal S_\star:\Gamma_\star\to \mathbb{R}\cup\{+\infty\}
  \]
  是作用泛函或路径积分选择准则；
  \item
  \[
    \operatorname{NF}_\star:\Gamma_\star/\!\sim\ \to \mathcal T_\star,
    \qquad
    \operatorname{Rep}_\star:\mathcal E_\star\to \mathcal T_\star
  \]
  分别是正规形提取算子与修复代表算子；
  \item
  \[
    \Theta_\star:\mathcal T_\star\to \mathcal I_\star^{\mathrm{logic}}
  \]
  是向逻辑占位或接口层的投影；
  \item
  \[
    \Cert_\star:\mathcal T_\star\to \mathcal C_\star,
    \qquad
    \Issue_\star:\mathcal C_\star\times \mathcal T_\star\to \mathcal I_\star^{\mathrm{rule}}
  \]
  分别是证书构造器与定理/规则签发映射。
\end{enumerate}
\end{definition}

\begin{remark}[结构签名的层级含义]
定义~\ref{def:tex37-ch12-signature} 说明：方程系统从来不是“先有规则，再补求解过程”，而是先给出算子、自由度与状态空间，
再给出障碍、允许路径类、路径积分选择、正规形提取与证书签发。
因此规则接口在结构顺序上位于末端，而不是位于起点。
\end{remark}

\subsection{命题：主基底或扇区插件方程系统的最精确压缩公式}
\label{subsec:tex37-ch12-prop}

\begin{proposition}[统一骨架压缩公式]
\label{prop:tex37-ch12-chain}
对任意
\[
  \star\in \{\mathrm{base}\}\cup\mathfrak D,
\]
其主基底或扇区插件方程系统的真正结构链可压缩为
\[
  \boxed{
    (\mathcal O_\star,\mathcal U_\star)
    \Longrightarrow
    (\mathcal E_\star,\pi_\star)
    \Longrightarrow
    \Gamma_\star
    \Longrightarrow
    \mathcal S_\star
    \Longrightarrow
    \operatorname{NF}_\star=\operatorname{Rep}_\star
    \Longrightarrow
    (\Cert_\star,\Issue_\star).
  }
\]
该压缩公式比“算子对应路径积分对应规则”更严格，
因为它显式分离了纤维化状态空间、允许路径类、路径积分选类、正规形收束与证书签发各层。
\end{proposition}

\begin{Certificate}[证书：状态层分离 + 路径层分离 + 签发层后置]
\label{cert:tex37-ch12-chain}
证书登记谓词链：
\[
\begin{aligned}
  P_1 &: (\mathcal O_\star,\mathcal U_\star)\ \text{决定局部变换律与可变参数},\\
  P_2 &: P_1\Rightarrow (\mathcal E_\star,\pi_\star)\ \text{被生成并获得纤维化状态解释},\\
  P_3 &: (\mathcal E_\star,\pi_\star,A_\star,F_\star,\Ob_\star)\Rightarrow \Gamma_\star\ \text{可定义},\\
  P_4 &: \Gamma_\star\Rightarrow \mathcal S_\star\ \text{只负责选取可行且近优的路径类代表},\\
  P_5 &: \mathcal S_\star\Rightarrow \operatorname{NF}_\star=\operatorname{Rep}_\star\ \text{给出唯一可比较代表},\\
  P_6 &: \operatorname{NF}_\star=\operatorname{Rep}_\star\Rightarrow (\Cert_\star,\Issue_\star)\ \text{才可合法签发规则接口}.
\end{aligned}
\]
\end{Certificate}

\begin{proof}[证明（基于数学符号推演逻辑的谓词因果链）]
由 $P_1$ 与 $P_2$，算子及其自由度并不直接给出规则接口，而是先生成局部状态对象及其纤维化解释。
由 $P_3$，只有在联络、曲率与障碍函数被并入后，
才可讨论从给定状态出发的允许路径类。
由 $P_4$，路径积分或作用泛函只承担“在允许路径类中选类”的职责，
而不直接承担规则签发职责。
由 $P_5$，路径类必须先压到正规形/修复代表，系统才获得唯一可比较的结果接口。
最后由 $P_6$，证书构造与签发发生在正规形之后。
故命题成立。证毕。
\end{proof}

\subsection{定理：算子及其自由度先决定纤维化状态空间，而不是直接决定规则}
\label{subsec:tex37-ch12-thm1}

\begin{theorem}[算子--自由度到纤维态的生成定理]
\label{thm:tex37-ch12-operators}
设
\[
  \mathcal O_\star=\{O_{\star,\alpha}\}_{\alpha\in I_\star}
\]
为算子族，
\[
  \mathcal U_\star=\{u_{\star,\beta}\}_{\beta\in J_\star}
\]
为自由度族。
若定义局部纤维态为
\[
  x_\star=
  \bigl(
    b_\star,\,
    u_{\star,\beta},\,
    O_{\star,\alpha}\cdot u_{\star,\beta},\,
    \cdots
  \bigr)
  \in \mathcal E_\star,
\]
则方程系统的第一层结构不是规则表，而是
\[
  \boxed{
    \text{算子--自由度耦合生成的纤维化状态空间。}
  }
\]
\end{theorem}

\begin{Certificate}[证书：变换律 + 参数族 + 局部态生成]
\label{cert:tex37-ch12-operators}
证书登记谓词链：
\[
\begin{aligned}
  T_1 &: \mathcal O_\star\ \text{只给出变换律，不单独给出完整状态},\\
  T_2 &: \mathcal U_\star\ \text{只给出可变参数，不单独给出作用结果},\\
  T_3 &: T_1\wedge T_2\Rightarrow
  (O_{\star,\alpha},u_{\star,\beta})\ \text{必须联立},\\
  T_4 &: T_3\Rightarrow x_\star=
  (b_\star,u_{\star,\beta},O_{\star,\alpha}\cdot u_{\star,\beta},\cdots)\in\mathcal E_\star,\\
  T_5 &: T_4\Rightarrow \text{第一层对象是纤维化状态空间而非规则接口}.
\end{aligned}
\]
\end{Certificate}

\begin{proof}[证明（基于数学符号推演逻辑的谓词因果链）]
由 $T_1$，若只给出算子族 $\mathcal O_\star$，
则系统只知道“允许怎样变换”，而不知道“具体对什么对象、在什么参数下变换”。
由 $T_2$，若只给出自由度族 $\mathcal U_\star$，
则系统只知道“有哪些可变参数”，而不知道“这些参数在什么算子作用下形成何种状态”。

因此由 $T_3$，只有把算子族与自由度族联立，才可能得到具体局部态。
由 $T_4$，这些局部态以
\[
  x_\star=
  (b_\star,u_{\star,\beta},O_{\star,\alpha}\cdot u_{\star,\beta},\cdots)
\]
的形式进入纤维对象
\[
  \mathcal E_\star.
\]
故由 $T_5$，方程系统的第一层不是“规则”，而是由算子--自由度耦合生成的纤维化状态空间。证毕。
\end{proof}

\subsection{定理：同伦扭曲是纤维丛空间上的修复/演化机制，路径积分是允许路径类的选取机制}
\label{subsec:tex37-ch12-thm2}

\begin{theorem}[同伦扭曲与路径积分的职责分离定理]
\label{thm:tex37-ch12-homotopy}
在
\[
  \mathcal E_\star
\]
上给定联络
\[
  A_\star
\]
与曲率
\[
  F_\star.
\]
定义障碍函数
\[
  \Ob_\star(x)
  :=
  \bigl(
    \lVert D_{A_\star}F_\star\rVert,\,
    \lVert J_\star\rVert,\,
    \lVert \Delta_{\mathrm{cert},\star}\rVert,\,
    \lVert \Delta_{\mathrm{ref},\star}\rVert
  \bigr),
\]
并设修复算子
\[
  \mathcal R_\star:\mathcal E_\star\to \mathcal E_\star
\]
满足
\[
  \Ob_\star(\mathcal R_\star(x))\preceq \Ob_\star(x).
\]
令
\[
  \Gamma_\star(x)
\]
为从 $x$ 出发的允许路径类集合，
作用泛函为
\[
  \mathcal S_\star:\Gamma_\star(x)\to \mathbb{R}\cup\{+\infty\}.
\]
则同伦扭曲对应的是纤维丛状态上的可修复变形机制，
而路径积分/作用泛函对应的是允许路径类中的代表选取机制。
\end{theorem}

\begin{Certificate}[证书：修复下降 + 允许路径类 + 近优选类]
\label{cert:tex37-ch12-homotopy}
证书登记谓词链：
\[
\begin{aligned}
  H_1 &: \Ob_\star(\mathcal R_\star(x))\preceq \Ob_\star(x),\\
  H_2 &: H_1\Rightarrow \mathcal R_\star\ \text{负责把状态沿可修复方向推进},\\
  H_3 &: \Gamma_\star(x)\ \text{由联络、障碍与允许变形共同决定},\\
  H_4 &: \mathcal S_\star:\Gamma_\star(x)\to \mathbb{R}\cup\{+\infty\}\ \text{只在路径类上比较代价/作用},\\
  H_5 &: H_3\wedge H_4\Rightarrow
  \gamma_\star^\ast\in
  \arg\min_{\gamma\in \Gamma_\star(x)^{\mathrm{adm}}}\mathcal S_\star(\gamma),\\
  H_6 &: H_2\wedge H_5\Rightarrow
  \text{同伦扭曲负责“怎么变”，路径积分负责“选哪条变形路径类”}.
\end{aligned}
\]
\end{Certificate}

\begin{proof}[证明（基于数学符号推演逻辑的谓词因果链）]
由 $H_1$，修复算子 $\mathcal R_\star$ 的作用是降低或不增加障碍，
故它控制的是状态在纤维丛中的修复/演化方向。
于是由 $H_2$，同伦扭曲的严格数学含义不是“最后写成哪句规则”，而是
\[
  x
  \rightsquigarrow
  \mathcal R_\star(x)
  \rightsquigarrow
  \cdots
\]
这样的可修复变形链。

由 $H_3$，允许路径类集合
\[
  \Gamma_\star(x)
\]
由联络结构、障碍边界与允许变形共同决定。
由 $H_4$，作用泛函
\[
  \mathcal S_\star
\]
是在这些路径类上比较可行性与近优性，而不是在规则文本之间直接比较。
因此由 $H_5$，路径积分/作用泛函最多只能选出近优路径类代表
\[
  \gamma_\star^\ast.
\]
综合 $H_2$ 与 $H_5$，即得 $H_6$。定理成立。证毕。
\end{proof}

\subsection{定理：规则不是直接从路径积分读出，而是从正规形与证书签发层读出}
\label{subsec:tex37-ch12-thm3}

\begin{theorem}[规则签发后置定理]
\label{thm:tex37-ch12-issue}
设正规形提取算子为
\[
  \operatorname{NF}_\star:\Gamma_\star/\!\sim\ \to \mathcal T_\star,
\]
并定义修复代表为
\[
  \operatorname{Rep}_\star(x):=\operatorname{NF}_\star([x]).
\]
再定义证书构造器与签发器
\[
  \Cert_\star:\mathcal T_\star\to \mathcal C_\star,
  \qquad
  \Issue_\star:\mathcal C_\star\times \mathcal T_\star\to \mathcal I_\star^{\mathrm{rule}}.
\]
若
\[
  \gamma_\star^\ast
  \in
  \arg\min_{\gamma\in\Gamma_\star(x)^{\mathrm{adm}}}\mathcal S_\star(\gamma),
\]
并记
\[
  \tau_\star^\ast:=\operatorname{NF}_\star([\gamma_\star^\ast]),
\]
则合法规则对象应写为
\[
  R_\star^{\mathrm{thm}}
  :=
  \Issue_\star\bigl(\Cert_\star(\tau_\star^\ast),\tau_\star^\ast\bigr),
\]
而不能写成“规则 = 任意路径积分输出”。
\end{theorem}

\begin{Certificate}[证书：选类不足 + 正规化必要 + 证书签发必要]
\label{cert:tex37-ch12-issue}
证书登记谓词链：
\[
\begin{aligned}
  Q_1 &: \gamma_\star^\ast\ \text{只是允许路径类中的近优代表},\\
  Q_2 &: Q_1\Rightarrow \gamma_\star^\ast\ \text{仍可能依赖具体表象},\\
  Q_3 &: \tau_\star^\ast:=\operatorname{NF}_\star([\gamma_\star^\ast])\Rightarrow \text{得到唯一可比较正规代表},\\
  Q_4 &: \Cert_\star(\tau_\star^\ast)\ \text{给出可审计证书包},\\
  Q_5 &: \Issue_\star(\Cert_\star(\tau_\star^\ast),\tau_\star^\ast)\ \text{才把正规代表提升为规则/定理接口},\\
  Q_6 &: Q_1\wedge Q_3\wedge Q_4\wedge Q_5
  \Rightarrow
  \text{规则接口来自正规形与证书签发层，而非直接来自路径积分层}.
\end{aligned}
\]
\end{Certificate}

\begin{proof}[证明（基于数学符号推演逻辑的谓词因果链）]
由 $Q_1$，路径积分或作用泛函只选出某个路径类代表
\[
  \gamma_\star^\ast,
\]
这仍然是路径层对象，而不是规则层对象。
由 $Q_2$，若不经过正规化，同一等价类可能仍有不同表象，
因此系统尚未得到唯一可比较结果。

由 $Q_3$，对路径类代表施加正规形提取算子后，
才得到唯一正规代表
\[
  \tau_\star^\ast.
\]
由 $Q_4$，证书构造器为该正规代表生成可审计证书包。
由 $Q_5$，只有在证书与正规代表同时送入签发器后，
系统才得到对外规则/定理接口
\[
  R_\star^{\mathrm{thm}}
  =
  \Issue_\star\bigl(\Cert_\star(\tau_\star^\ast),\tau_\star^\ast\bigr).
\]
因此由 $Q_6$，规则接口并不是从路径积分层直接读出，而是从正规形与证书签发层读出。证毕。
\end{proof}

\subsection{推论：对原判断的最严格改写}
\label{subsec:tex37-ch12-cor1}

\begin{corollary}[六层顺序推论]
\label{cor:tex37-ch12-strict}
由命题~\ref{prop:tex37-ch12-chain} 与定理~\ref{thm:tex37-ch12-operators}--\ref{thm:tex37-ch12-issue}，
可把主基底方程系统或扇区插件方程系统的结构，严格改写为：
\[
  \boxed{
    \text{算子及其自由度生成纤维丛状态空间；}
  }
\]
\[
  \boxed{
    \text{该状态空间上的同伦扭曲与联络演化生成允许路径类；}
  }
\]
\[
  \boxed{
    \text{路径积分/作用泛函从允许路径类中选出可行且近优的代表；}
  }
\]
\[
  \boxed{
    \text{再经正规形提取、证书构造与定理签发，闭合为规则接口。}
  }
\]
这才是“主基底或扇区插件方程系统的共同骨架”的最严密层次顺序。
\end{corollary}

\begin{Certificate}[证书：四步改写]
\label{cert:tex37-ch12-strict}
证书登记谓词链：
\[
\begin{aligned}
  C_1 &: \text{定理~\ref{thm:tex37-ch12-operators}}\Rightarrow \text{算子及自由度先生成纤维态},\\
  C_2 &: \text{定理~\ref{thm:tex37-ch12-homotopy}}\Rightarrow \text{同伦扭曲与路径积分职责分离},\\
  C_3 &: \text{定理~\ref{thm:tex37-ch12-issue}}\Rightarrow \text{规则签发后置于正规形层之后},\\
  C_4 &: C_1\wedge C_2\wedge C_3
  \Rightarrow
  \text{得到上述四步严格改写}.
\end{aligned}
\]
\end{Certificate}

\begin{proof}[证明（基于数学符号推演逻辑的谓词因果链）]
由 $C_1$，系统的起点不是规则，而是纤维态生成。
由 $C_2$，中间层必须区分“怎么变”与“选哪条变形路径类”。
由 $C_3$，末端规则接口又必须后置到正规形与证书之后。
因此由 $C_4$，上述四步改写即为最严格的结构压缩。推论成立。证毕。
\end{proof}

\subsection{定义：主基底方程系统与扇区插件方程系统的关系}
\label{subsec:tex37-ch12-def2}

\begin{definition}[主基底回拉与扇区细化差量]
\label{def:tex37-ch12-base-sector}
记主基底方程系统为
\[
  \mathfrak E_{\mathrm{base}}
  :=
  \mathfrak E_{\mathrm{core}},
\]
记扇区 $D$ 的插件方程系统为
\[
  \mathfrak E_D.
\]
则最自然的结构写法为
\[
  \mathfrak E_D
  =
  \rho_D^\ast \mathfrak E_{\mathrm{base}}
  \oplus
  \Delta_D,
\]
其中
\[
  \Delta_D
  :=
  \bigl(
    \Delta\mathcal O_D,\,
    \Delta\mathcal U_D,\,
    \Delta\Ob_D,\,
    \Delta\mathcal A_D,\,
    \Delta\Cert_D
  \bigr)
\]
表示扇区细化差量。
换言之，扇区插件不是另起一个宇宙，而是在同一主基底上的纤维细化方程系统。
\end{definition}

\begin{remark}[与第七章、第八章记号的对齐]
定义~\ref{def:tex37-ch12-base-sector} 中的
\[
  \Delta_D
\]
可视为第~\ref{sec:tex37-ch7} 章五类增量
\[
  \Delta_D^{\mathrm{obs}},
  \Delta_D^{\mathrm{op}},
  \Delta_D^{\mathrm{obst}},
  \Delta_D^{\mathrm{act}},
  \Delta_D^{\mathrm{cert}}
\]
的更抽象统摄记号；
而
\[
  \rho_D^\ast \mathfrak E_{\mathrm{base}}
\]
则强调：扇区方程系统来自同一主基底的回拉与细化，而不是平行本体的拼接。
\end{remark}

\subsection{公理（降级为定理）：有限层级规则候选的闭合/细化链可按层级递归构造}
\label{subsec:tex37-ch12-axiom}

\begin{theorem}[公理降级：有限层级规则候选的闭合/细化链可按层级递归构造（数学归纳法证书）]
\label{thm:tex37-ch12-induction}
设规则候选对象按层级分布为
\[
  R_\star^{(n)}\in \mathcal I_{\star}^{\mathrm{cand},(n)},
  \qquad
  n\in\NN,
\]
并定义层级秩函数
\[
  r_\star:
  \bigsqcup_{n\in\NN}\mathcal I_{\star}^{\mathrm{cand},(n)}
  \to
  \NN.
\]
若满足：
\begin{enumerate}
  \item 对任意 $R_\star^{(n)}$，若 $r_\star(R_\star^{(n)})=0$，则它可直接闭合签发，即存在
  \[
    c_\star^{(n)}=\Cert_\star(T_\star^{(n)})
  \]
  使
  \[
    \Issue_\star(c_\star^{(n)},T_\star^{(n)})
  \]
  合法定义；
  \item 对任意 $k\in\NN$，若某个规则候选满足
  \[
    r_\star(R_\star^{(n)})=k+1
  \]
  且当前层不能直接签发，
  则存在唯一细化族
  \[
    \{R_{\star,\beta}^{(n+1)}\}_{\beta\in B_n}
  \]
  满足
  \[
    \rho_{n+1\to n}(R_{\star,\beta}^{(n+1)})=R_\star^{(n)},
    \qquad
    r_\star(R_{\star,\beta}^{(n+1)})\le k.
  \]
\end{enumerate}
则对任意满足
\[
  r_\star(R)\le N
\]
的规则候选 $R$，都存在有限条“闭合签发或继续细化”的链，
并最终到达可签发规则接口。
\end{theorem}

\begin{Certificate}[证书：基步 + 归纳步（数学归纳法证书）]
\label{cert:tex37-ch12-induction}
定义谓词
\[
  P(k)
  :
  \Longleftrightarrow
  \forall R,\ r_\star(R)\le k
  \Rightarrow
  \text{$R$ 存在有限闭合/细化链并最终到达可签发接口}.
\]
则数学归纳法证书记为
\[
  \pi_{\mathrm{rule}}
  =
  \bigl(
    P(0),\
    \forall k\in\NN,\ P(k)\Rightarrow P(k+1)
  \bigr).
\]
其中归纳步的见证数据为
\[
  \Cert_\star,\ \Issue_\star,\ \rho_{n+1\to n},\ \{R_{\star,\beta}^{(n+1)}\}_{\beta\in B_n}.
\]
\end{Certificate}

\begin{proof}[证明（数学归纳法证书；基于数学符号推演逻辑的谓词因果链）]
定义谓词
\[
  P(k)
  :
  \Longleftrightarrow
  \forall R,\ r_\star(R)\le k
  \Rightarrow
  \text{$R$ 存在有限闭合/细化链并最终到达可签发接口}.
\]

\textbf{基步：}
若
\[
  r_\star(R)=0,
\]
则由条件(1)，$R$ 可直接被证书化并签发。
因此一切秩不超过 $0$ 的规则候选都存在有限闭合链，故 $P(0)$ 成立。

\textbf{归纳步：}
设 $P(k)$ 成立。取任意满足
\[
  r_\star(R)=k+1
\]
的规则候选。
若它在当前层已可直接签发，则闭合链立即存在。
若不能直接签发，则由条件(2)，存在唯一细化族
\[
  \{R_{\star,\beta}^{(n+1)}\}_{\beta\in B_n}
\]
满足
\[
  \rho_{n+1\to n}(R_{\star,\beta}^{(n+1)})=R,
  \qquad
  r_\star(R_{\star,\beta}^{(n+1)})\le k.
\]
由归纳假设 $P(k)$，每个细化子候选都存在有限闭合/细化链并最终到达可签发接口。
把
\[
  R\longmapsto \{R_{\star,\beta}^{(n+1)}\}_{\beta\in B_n}
\]
接到这些子链前端，即得原候选的有限闭合/细化链。
因此 $P(k+1)$ 成立。

由数学归纳法，对所有 $k\in\NN$，$P(k)$ 都成立。
故任意有限秩规则候选都可按层级递归收束到可签发接口。证毕。
\end{proof}

\subsection{定理：接口收束性在这里体现为“规则必须可签发，不能签发就必须继续细化”}
\label{subsec:tex37-ch12-thm4}

\begin{theorem}[规则层接口收束定理]
\label{thm:tex37-ch12-convergence}
对任意当前层级的规则候选
\[
  R_\star^{(n)},
\]
只允许两种合法去向：
\begin{enumerate}
  \item \textbf{闭合签发：}
  存在定理对象
  \[
    T_\star^{(n)}
  \]
  与证书
  \[
    c_\star^{(n)}=\Cert_\star(T_\star^{(n)})
  \]
  使
  \[
    \Issue_\star(c_\star^{(n)},T_\star^{(n)})
  \]
  合法定义；
  \item \textbf{继续细化：}
  若当前层不能签发，则必须进入
  \[
    R_\star^{(n)}
    \longmapsto
    \{R_{\star,\beta}^{(n+1)}\}_{\beta\in B_n},
    \qquad
    \rho_{n+1\to n}(R_{\star,\beta}^{(n+1)})=R_\star^{(n)}.
  \]
\end{enumerate}
因此接口收束性的严格含义是：
\[
  \boxed{
    \text{规则层必须“定理化”或“继续纤维细化”，不能停留在未闭合中间态。}
  }
\]
\end{theorem}

\begin{Certificate}[证书：可签发 or 继续细化，且无第三种合法状态]
\label{cert:tex37-ch12-convergence}
证书登记谓词链：
\[
\begin{aligned}
  K_1 &: \text{定理~\ref{thm:tex37-ch12-induction}}\Rightarrow \text{有限秩规则候选总存在闭合/细化链},\\
  K_2 &: \text{若某规则候选当前层可签发，则它属于闭合签发型},\\
  K_3 &: \text{若某规则候选当前层不可签发，则它必须进入保投影细化族},\\
  K_4 &: K_2\wedge K_3\Rightarrow \text{不存在第三种合法中间态},\\
  K_5 &: K_1\wedge K_4\Rightarrow \text{接口收束性等价于“可签发或继续细化”}.
\end{aligned}
\]
\end{Certificate}

\begin{proof}[证明（基于数学符号推演逻辑的谓词因果链）]
由 $K_1$，有限秩规则候选不允许永久悬空；
它们总必须沿某条闭合/细化链向可签发接口推进。
若某规则候选在当前层已经具备证书与定理对象，
则由 $K_2$，它属于闭合签发型。
若当前层尚不能签发，则由 $K_3$，它必须被送入下一层细化，
而不能伪装成已完成规则。

因此由 $K_4$，规则层不存在第三种合法状态：
既不能签发、又不继续细化的对象，只能是悬空接口，不属于合法白盒系统。
结合 $K_1$，即得 $K_5$。定理成立。证毕。
\end{proof}

\subsection{推论：因此“写方程”真正要写的是六层，不是一条式子}
\label{subsec:tex37-ch12-cor2}

\begin{corollary}[六层闭环推论]
\label{cor:tex37-ch12-six-layers}
由命题~\ref{prop:tex37-ch12-chain}、定理~\ref{thm:tex37-ch12-operators}、定理~\ref{thm:tex37-ch12-homotopy}、
定理~\ref{thm:tex37-ch12-issue} 与定理~\ref{thm:tex37-ch12-convergence}，
真正需要写出的并不是单个
\[
  F(x)=0,
\]
而至少是如下六层：
\[
  \boxed{
    \text{(1) 状态层：}\ (\mathcal O_\star,\mathcal U_\star)\to(\mathcal E_\star,\pi_\star)
  }
\]
\[
  \boxed{
    \text{(2) 障碍层：}\ \Ob_\star
  }
\]
\[
  \boxed{
    \text{(3) 修复层：}\ \mathcal R_\star
  }
\]
\[
  \boxed{
    \text{(4) 路径积分层：}\ \Gamma_\star,\mathcal S_\star
  }
\]
\[
  \boxed{
    \text{(5) 正规形层：}\ \operatorname{NF}_\star=\operatorname{Rep}_\star
  }
\]
\[
  \boxed{
    \text{(6) 签发层：}\ \Cert_\star,\Issue_\star
  }
\]
其中主基底写通用骨架，扇区插件写各自的细化差量。
\end{corollary}

\begin{Certificate}[证书：六层逐层不可省略]
\label{cert:tex37-ch12-six-layers}
证书登记谓词链：
\[
\begin{aligned}
  U_1 &: \text{定理~\ref{thm:tex37-ch12-operators}}\Rightarrow \text{状态层不可省略},\\
  U_2 &: \text{定理~\ref{thm:tex37-ch12-homotopy}}\Rightarrow \text{障碍层、修复层与路径积分层不可省略},\\
  U_3 &: \text{定理~\ref{thm:tex37-ch12-issue}}\Rightarrow \text{正规形层与签发层不可省略},\\
  U_4 &: \text{定义~\ref{def:tex37-ch12-base-sector}}\Rightarrow \text{扇区插件应写成主基底的细化差量},\\
  U_5 &: U_1\wedge U_2\wedge U_3\wedge U_4
  \Rightarrow
  \text{真正需要写出的是六层闭环系统，而非单个方程符号}.
\end{aligned}
\]
\end{Certificate}

\begin{proof}[证明（基于数学符号推演逻辑的谓词因果链）]
由 $U_1$，如果没有状态层，则方程系统甚至没有对象域。
由 $U_2$，如果没有障碍层、修复层与路径积分层，
则系统无法区分“当前卡在哪里”“沿什么修复机制变”“从哪些允许路径类中选代表”。
由 $U_3$，如果没有正规形层与签发层，
则系统无法把路径类结果压成唯一可比较对象，更无法把该对象提升为规则/定理接口。
由 $U_4$，扇区插件又不能另起宇宙，而必须写成主基底的细化差量。
因此由 $U_5$，真正应写出的必然是六层闭环系统，而不是单条符号等式。推论成立。证毕。
\end{proof}

\subsection{备注：统一口径、最短公式与下一步}
\label{subsec:tex37-ch12-remarks}

\begin{remark}[最值得钉死的一句]
本章最值得钉死的一句，可写成：
\[
  \boxed{
    \begin{aligned}
      &\text{主基底方程系统/扇区插件方程系统的本体，不是“规则表”，}\\
      &\text{而是“由算子--自由度--纤维丛--路径类--正规形--证书签发”}\\
      &\text{组成的层级闭环系统。}
    \end{aligned}
  }
\]
\end{remark}

\begin{remark}[最精确的共同骨架公式]
若把本章再压成一条最短公式，则可写成：
\[
  \boxed{
    (\mathcal O_\star,\mathcal U_\star)
    \to
    \mathcal E_\star
    \to
    \Gamma_\star
    \to
    \mathcal S_\star
    \to
    \operatorname{NF}_\star
    \to
    (\Cert_\star,\Issue_\star),
    \qquad
    \star\in \{\mathrm{base}\}\cup\mathfrak D.
  }
\]
这条式子基本就把主基底方程系统与扇区插件方程系统的共同骨架钉死了。
\end{remark}

\begin{remark}[下一步]
本章与第~\ref{sec:tex37-ch11} 章联立后，下一步最自然的工作已经非常明确：
\[
  \boxed{
    \text{先把 }\mathfrak E_{\mathrm{base}}=\mathfrak E_{\mathrm{core}}\text{ 写成这套统一结构签名的具体版本。}
  }
\]
因为只有主基底的六层闭环骨架先被写实，
各扇区插件的细化差量、耦合条件、证书构造器与签发规则，才有稳定锚点可依附。
\end{remark}

\clearpage

%%%%%%%%%%%%%%%%%%%%%%%%%%%%%%%%%%%%%%%%%%%%%%%%%%%%%%%%%%%%
\section{形式逻辑：离散语义基底、认证纤维塔、竖直认证路径与可收束签发}
\label{sec:tex37-ch13}
%%%%%%%%%%%%%%%%%%%%%%%%%%%%%%%%%%%%%%%%%%%%%%%%%%%%%%%%%%%%

\begin{remark}[核心陈述（四点收紧后）]
若把自然语言逻辑占位写成有限字符集上的协议内最小占位串，并把依赖限制为局部有限的良基依赖图，
则“语义纤维丛空间”可被严格拆成如下对象链：
\[
  \text{可数离散语义基底}
  \to
  \text{逐层认证纤维塔}
  \to
  \text{竖直认证路径空间}
  \to
  \text{定理/规则签发接口}.
\]
在这一口径下，必须同步收紧四点：
\begin{enumerate}
  \item “离散拓扑的流形”更安全地写成“可数离散 $0$ 维语义基底空间”；
  \item “遍历”不是通常拓扑连续路径，而是认证依赖图上的离散组合路径；
  \item “有限字符集表达无限语义”只能写成“充分编码可组合、可递归、可签发的无限语义族”；
  \item “自然双射”“同胚”“同伦等价”必须分层书写，不能在未选定一致拓扑时直接并列。
\end{enumerate}
\end{remark}

\begin{remark}[阅读导航：仓内引用与外部参照]
本章把仓内关于“判定+规范化”接口与语义规范场的整理稿 \cite{RepoTex20SemanticGaugeNormalization,RepoTex21AuditedWrapper,RepoTex23SemanticBase}，
关于可修复障碍边界与接口收束/签发塔的整理稿 \cite{RepoTex29RepairableObstruction,RepoTex35ExtendedCertificationTower}，
以及关于纤维丛空间、路径积分支撑类与同伦扭曲正规形集合等价的整理稿 \cite{RepoTex36HomotopyTwistBijection}
压缩为“离散语义基底--认证纤维塔--竖直认证路径--签发收束”这一条严格对象链。
更系统的 G 框架背景仍可参照 \cite{GaoZhengAppVol2,GaoZhengAppVol3,GaoZheng2025GFramework}。

外部标准参照方面，本章继续借用范畴语言、离散/纤维丛拓扑、同伦理论与同调代数的标准写法 \cite{MacLane1998Categories,Hatcher2002AT,Weibel1994HomologicalAlgebra,
  KobayashiNomizu1963Foundations}，
但不把外部教材结论当作本稿内部未登记的前提。
\end{remark}

\subsection{定义：有限字符集、逻辑占位、层级依赖与可数离散语义基底}
\label{subsec:tex37-ch13-defs}

\begin{definition}[有限字符集与逻辑占位]
\label{def:tex37-ch13-placeholder}
设
\[
  \Sigma=\Sigma_{\mathrm{CN}}\sqcup \Sigma_{\mathrm{Math}},
  \qquad
  |\Sigma|<\infty,
\]
其中 $\Sigma_{\mathrm{CN}}$ 表示中文字符集合，$\Sigma_{\mathrm{Math}}$ 表示数学符号集合。
记
\[
  \Sigma^\ast:=\bigcup_{m\ge 0}\Sigma^m.
\]
取子集
\[
  \mathcal P\subseteq \Sigma^\ast.
\]
称每个 $p\in\mathcal P$ 为一个\textbf{逻辑占位}，若满足：
\begin{enumerate}
  \item $p$ 在当前协议下可被解析为原子语义槽；
  \item $p$ 在当前层级下不可再被无损拆成同级占位的串联；
  \item $p$ 的合法性由占位证书接口认证，即 $\Cert(p)=\top$。
\end{enumerate}
因此，“最小”不是自然语言意义下的最短，而是
\[
  \boxed{
    \text{在当前层级与当前认证协议下不可再无损同级拆分。}
  }
\]
\end{definition}

\begin{definition}[层级、良基依赖与局部有限性]
\label{def:tex37-ch13-dependency}
定义层级函数
\[
  \ell:\mathcal P\to \NN,
\]
以及依赖关系
\[
  \prec\ \subseteq\ \mathcal P\times \mathcal P.
\]
要求：
\[
  q\prec p \Longrightarrow \ell(q)<\ell(p).
\]
并进一步要求每个占位的直接前驱集
\[
  \operatorname{Pred}(p):=\{q\in\mathcal P:q\prec p\}
\]
是有限集。于是得到局部有限的分层有向无环图
\[
  \mathcal G_{\mathrm{dep}}=(\mathcal P,E_{\mathrm{dep}}),
  \qquad
  E_{\mathrm{dep}}:=\{(q,p)\in\mathcal P\times \mathcal P:q\prec p\}.
\]
\end{definition}

\begin{definition}[可数离散语义基底空间]
\label{def:tex37-ch13-base}
在 $\mathcal P$ 上赋离散拓扑
\[
  \tau_{\mathrm{disc}}:=2^{\mathcal P}.
\]
记
\[
  \mathcal B_{\mathrm{sem}}:=(\mathcal P,\tau_{\mathrm{disc}})
\]
为\textbf{离散语义基底空间}。由于 $|\Sigma|<\infty$，故 $|\Sigma^\ast|=\aleph_0$；又 $\mathcal P\subseteq \Sigma^\ast$，
所以 $\mathcal B_{\mathrm{sem}}$ 至多可数。
\end{definition}

\begin{remark}[关于“离散流形”的安全写法]
\label{rem:tex37-ch13-discrete-manifold}
在本章口径下，最安全的表述是：
\[
  \boxed{
    \mathcal B_{\mathrm{sem}}\ \text{是一个可数离散 }0\text{ 维语义基底空间。}
  }
\]
若额外采用“每点局部同胚于单点离散空间”的 $0$ 维流形口径，
则可进一步把它称为“可数离散 $0$ 维拓扑流形”；
若不预设该口径，则不应直接把“离散拓扑”与“流形”并列写成无条件同义词。
\end{remark}

\begin{definition}[N 阶认证语义纤维塔]
\label{def:tex37-ch13-tower}
对任意 $p\in\mathcal B_{\mathrm{sem}}$，先定义
\[
  E^{\le 0}(p):=\{(p)\}.
\]
对每个 $n\ge 1$ 与每个已认证前缀状态
\[
  x_{n-1}=(p;\tau_1,\dots,\tau_{n-1})\in E^{\le n-1}(p),
\]
定义第 $n$ 阶纤维
\[
  F_n(x_{n-1})
  :=
  \left\{
    \tau_n
    \,\middle|\,
    \begin{array}{l}
      \tau_n\text{ 是建立在 }x_{n-1}\text{ 上的第 }n\text{ 阶同伦扭曲/修复状态},\\
      \tau_n\text{ 只依赖于严格低层数据与有限个低层占位前驱},\\
      \Cert_n(x_{n-1},\tau_n)=\top
    \end{array}
  \right\}.
\]
递归定义
\[
  E^{\le n}(p)
  :=
  \left\{
    (x_{n-1},\tau_n)
    \,\middle|\,
    x_{n-1}\in E^{\le n-1}(p),\ \tau_n\in F_n(x_{n-1})
  \right\}.
\]
再定义总空间
\[
  \mathcal E^{(N)}
  :=
  \bigsqcup_{p\in\mathcal B_{\mathrm{sem}}} E^{\le N}(p),
  \qquad
  \pi_0:\mathcal E^{(N)}\to \mathcal B_{\mathrm{sem}},
  \quad
  \pi_0(p;\tau_1,\dots,\tau_N)=p.
\]
称其为\textbf{$N$ 阶认证语义纤维塔}。
\end{definition}

\subsection{命题：有限字符集足以充分编码可组合、可递归、可签发的无限语义族}
\label{subsec:tex37-ch13-encoding}

\begin{proposition}[有限字符集足以充分编码可组合、可递归、可签发的无限语义族]
\label{prop:tex37-ch13-encoding}
设 $\mathfrak S_\infty$ 为一族语义对象，并满足：对每个 $s\in\mathfrak S_\infty$，
都存在可被固定序列化规约写入 $\Sigma^\ast$ 的有限描述串、有限组合规则、递归生成规则以及有限证书链，
且这些数据可完全恢复 $s$。
则存在单射
\[
  \operatorname{Enc}:\mathfrak S_\infty \hookrightarrow \Sigma^\ast.
\]
特别地，若 $\mathfrak S_\infty$ 为无限族，则它仍可被有限字符集 $\Sigma$ 充分编码。
\end{proposition}

\begin{Certificate}[证书：有限字母表上的有限串可数无限]
\label{cert:tex37-ch13-encoding}
证书登记谓词链：
\[
\begin{aligned}
  E_1 &: |\Sigma|<\infty,\\
  E_2 &: E_1\Rightarrow |\Sigma^\ast|=\aleph_0,\\
  E_3 &: \forall s\in\mathfrak S_\infty,\ \exists \mathbf d(s)\in \Sigma^\ast\ \text{使得 }\mathbf d(s)\text{ 序列化记录 }
    s\text{ 的有限描述/规则/证书链},\\
  E_4 &: \mathbf d(s)\ \text{对 }s\text{ 可完全恢复},\\
  E_5 &: E_3\wedge E_4\Rightarrow \operatorname{Enc}(s):=\mathbf d(s)\ \text{为良定义单射}.
\end{aligned}
\]
\end{Certificate}

\begin{proof}[证明（基于数学符号推演逻辑的谓词因果链）]
由 $E_1$，$\Sigma$ 是有限字母表；因此由 $E_2$，
\[
  |\Sigma^\ast|=\aleph_0.
\]
由 $E_3$，每个语义对象 $s\in\mathfrak S_\infty$ 都可以通过固定序列化规约写成某个有限串
\[
  \mathbf d(s)\in \Sigma^\ast,
\]
该串同时记录其有限描述、有限组合规则、递归生成规则与有限证书链。
再由 $E_4$，$\mathbf d(s)$ 对应的数据足以完全恢复 $s$，因而不同对象不可能共享同一完整编码串；
于是
\[
  \operatorname{Enc}(s):=\mathbf d(s)
\]
给出 $\mathfrak S_\infty$ 到 $\Sigma^\ast$ 的单射，即 $E_5$。

因此，有限字符集足以充分编码“可组合、可递归、可签发”的无限语义族。
但这里不能无条件推广为“编码任意意义上的所有无限语义”：
若对象要求不可数参数化且又不允许通过有限规则、递归规约或证书链离散化，则它超出当前框架。证毕。
\end{proof}

\subsection{引理：离散语义基底上的拓扑连续路径只能是常值路径}
\label{subsec:tex37-ch13-const-path}

\begin{lemma}[离散语义基底上的拓扑连续路径只能是常值路径]
\label{lem:tex37-ch13-const-path}
若
\[
  \gamma:[0,1]\to \mathcal B_{\mathrm{sem}}
\]
是通常拓扑意义下的连续映射，则 $\gamma$ 必为常值映射。
\end{lemma}

\begin{Certificate}[证书：连通像保持连通，而离散空间的连通子集只有单点]
\label{cert:tex37-ch13-const-path}
证书登记谓词链：
\[
\begin{aligned}
  P_1 &: [0,1]\ \text{是连通空间},\\
  P_2 &: \gamma\ \text{连续}\Rightarrow \gamma([0,1])\ \text{连通},\\
  P_3 &: \mathcal B_{\mathrm{sem}}\ \text{离散}\Rightarrow \text{其连通子集只能是单点},\\
  P_4 &: P_2\wedge P_3\Rightarrow |\gamma([0,1])|=1.
\end{aligned}
\]
\end{Certificate}

\begin{proof}[证明（基于数学符号推演逻辑的谓词因果链）]
由 $P_1$，区间 $[0,1]$ 是连通空间。
若 $\gamma$ 连续，则由连通像保持连通可得 $P_2$：
\[
  \gamma([0,1])\subseteq \mathcal B_{\mathrm{sem}}
\]
必须是连通子集。
另一方面，由定义~\ref{def:tex37-ch13-base}，$\mathcal B_{\mathrm{sem}}$ 赋离散拓扑，
故任意含有至少两个点的子集都可被拆成两个非空开闭子集，从而不连通；这正是 $P_3$。
于是由 $P_2$ 与 $P_3$ 得
\[
  |\gamma([0,1])|=1,
\]
即 $\gamma$ 的像为单点，故 $\gamma$ 为常值映射。证毕。
\end{proof}

\subsection{定义：离散认证路径、竖直认证路径空间与认证保持同伦商}
\label{subsec:tex37-ch13-paths}

\begin{definition}[基底离散认证路径]
\label{def:tex37-ch13-base-path}
定义可遍历边集
\[
  E_{\mathrm{trav}}
  :=
  \left\{
    (p,q)\in\mathcal P\times\mathcal P
    \,\middle|\,
    p\prec q\ \vee\ \Auth(p,q)=\top\ \vee\ \Tw(p,q)=\top
  \right\}.
\]
称有限序列
\[
  \gamma_{\mathrm{base}}=(p_0\to p_1\to \cdots\to p_m)
\]
为一条\textbf{基底离散认证路径}，若对每个 $0\le i<m$ 都有 $(p_i,p_{i+1})\in E_{\mathrm{trav}}$。
记所有此类路径构成的集合为 $\Gamma_{\mathrm{base}}$。
\end{definition}

\begin{definition}[竖直认证路径空间]
\label{def:tex37-ch13-vertical-path}
对任意 $p\in\mathcal B_{\mathrm{sem}}$ 与 $N\in\NN$，定义
\[
  \Pi_N^{\mathrm{vert}}(p)
  :=
  \left\{
    (x_0\to x_1\to \cdots \to x_N)
    \,\middle|\,
    \begin{array}{l}
      x_0=(p),\\
      x_k=(p;\tau_1,\dots,\tau_k)\in E^{\le k}(p)\quad (1\le k\le N),\\
      x_k=(x_{k-1},\tau_k)\ \text{且}\ \Cert_k(x_{k-1},\tau_k)=\top
    \end{array}
  \right\}.
\]
再定义
\[
  \Pi_N^{\mathrm{vert}}
  :=
  \bigsqcup_{p\in\mathcal B_{\mathrm{sem}}}\Pi_N^{\mathrm{vert}}(p).
\]
称其为\textbf{$N$ 阶竖直认证路径空间}。
\end{definition}

\begin{definition}[认证保持同伦商与路径积分口径]
\label{def:tex37-ch13-homotopy-quotient}
在 $\Pi_N^{\mathrm{vert}}$ 上定义关系 $\sim_h^{\mathrm{cert}}$：若两条路径可通过有限次局部换序、局部细化或局部压缩互达，
且在整个过程中同时保持
\[
  \pi_0,\qquad
  \text{每一阶认证条件},\qquad
  \text{终点状态}
\]
不变，则称它们\textbf{认证保持同伦等价}。
若再给定作用泛函
\[
  \mathcal S_N:\Pi_N^{\mathrm{vert}}\to \mathbb{R}\sqcup\{+\infty\},
\]
则本章的“语义路径积分”应理解为：在商空间
\[
  \Pi_N^{\mathrm{vert}}/\!\sim_h^{\mathrm{cert}}
\]
上选择可行且近优的代表。
\end{definition}

\begin{remark}[“遍历”与“双射对象”的严格区分]
\label{rem:tex37-ch13-path-distinction}
本章故意把两类路径分开：
\[
  \Gamma_{\mathrm{base}}
  \quad\text{负责基底上的离散组合遍历},\qquad
  \Pi_N^{\mathrm{vert}}
  \quad\text{负责固定基点上的逐层认证抬升}.
\]
后续双射定理只对 $\mathcal E^{(N)}$ 与 $\Pi_N^{\mathrm{vert}}$ 断言自然双射；
若把跨基点的全部离散遍历也一起并入对象端，再直接写全局双射，命题会被写得过强。
\end{remark}

\subsection{公理（降级为定理）：良基依赖使认证语义纤维塔可按阶递归构造}
\label{subsec:tex37-ch13-recursion}

\begin{theorem}[公理降级：局部有限良基依赖使有限阶认证语义纤维塔可递归构造]
\label{thm:tex37-ch13-recursion}
设定义~\ref{def:tex37-ch13-dependency} 的依赖图局部有限且良基，且对每个 $n\ge 1$ 与每个已认证前缀状态 $x_{n-1}\in E^{\le n-1}(p)$，
纤维 $F_n(x_{n-1})$ 的判定仅调用有限个严格低层占位前驱与低阶已认证数据。
则对任意 $N\in\NN$ 与任意 $p\in\mathcal B_{\mathrm{sem}}$，
族
\[
  E^{\le 0}(p),E^{\le 1}(p),\dots,E^{\le N}(p)
\]
都可按递推式唯一构造；特别地，$\mathcal E^{(N)}$ 是良定义的。
\end{theorem}

\begin{Certificate}[证书：基步 + 归纳步（数学归纳法证书）]
\label{cert:tex37-ch13-recursion}
定义谓词
\[
  Q(n):\ \forall p\in\mathcal B_{\mathrm{sem}},\ E^{\le n}(p)\ \text{已良定义，且其每个元素都可由有限个严格低层数据构造}.
\]
记数学归纳法证书为
\[
  \pi_{\mathrm{ind}}
  =
  \bigl(Q(0),\ \forall n\in\NN,\ Q(n)\Rightarrow Q(n+1)\bigr).
\]
其中归纳步的递推式登记为
\[
  E^{\le n+1}(p)
  =
  \left\{
    (x_n,\tau_{n+1})
    \,\middle|\,
    x_n\in E^{\le n}(p),\ \tau_{n+1}\in F_{n+1}(x_n)
  \right\}.
\]
\end{Certificate}

\begin{proof}[证明（数学归纳法证书；基于数学符号推演逻辑的谓词因果链）]
\textbf{基步 $Q(0)$：}
由定义~\ref{def:tex37-ch13-tower}，$E^{\le 0}(p)=\{(p)\}$，
故对每个 $p$，$E^{\le 0}(p)$ 都良定义，且其元素不含更高阶数据，$Q(0)$ 成立。

\textbf{归纳步：}
设 $Q(n)$ 成立。任取 $p\in\mathcal B_{\mathrm{sem}}$ 与 $x_n\in E^{\le n}(p)$。
根据假设，判定 $\tau_{n+1}\in F_{n+1}(x_n)$ 时，
只会调用有限个严格低层占位前驱与低阶已认证数据。
由于依赖图局部有限且满足
\[
  q\prec p \Rightarrow \ell(q)<\ell(p),
\]
不存在向上或同层回指的依赖环；因此这些被调用的数据是良基的，并且认证检查不会产生未闭合的循环引用。
于是 $F_{n+1}(x_n)$ 作为一个由确定谓词切出的集合是良定义的。

进一步，由归纳假设，所有 $x_n$ 已经良定义；
故递推式
\[
  E^{\le n+1}(p)
  =
  \left\{
    (x_n,\tau_{n+1})
    \,\middle|\,
    x_n\in E^{\le n}(p),\ \tau_{n+1}\in F_{n+1}(x_n)
  \right\}
\]
给出一个唯一的下一层对象集。
并且每个新元素 $(x_n,\tau_{n+1})$ 仍只依赖有限个严格低层数据，
故 $Q(n+1)$ 成立。

综上，$Q(n)\Rightarrow Q(n+1)$。
由数学归纳法，$\forall n\in\NN,\ Q(n)$。
因此对任意有限 $N$，$\mathcal E^{(N)}$ 都可按层递归构造且良定义。证毕。
\end{proof}

\subsection{定理：N 阶认证语义纤维态与竖直认证路径空间之间存在自然双射}
\label{subsec:tex37-ch13-bijection}

\begin{theorem}[N 阶认证语义纤维态与竖直认证路径空间自然双射]
\label{thm:tex37-ch13-bijection}
对任意 $N\in\NN$，定义
\[
  \Phi_N:\mathcal E^{(N)}\to \Pi_N^{\mathrm{vert}},
  \qquad
  \Phi_N(p;\tau_1,\dots,\tau_N)
  :=
  \bigl(
    (p)\to
    (p;\tau_1)\to\cdots\to
    (p;\tau_1,\dots,\tau_N)
  \bigr).
\]
则 $\Phi_N$ 是自然双射。其逆映射
\[
  \Psi_N:\Pi_N^{\mathrm{vert}}\to \mathcal E^{(N)}
\]
由“取竖直认证路径的终点状态”给出。
\end{theorem}

\begin{Certificate}[证书：前缀链唯一化 + 终点回收]
\label{cert:tex37-ch13-bijection}
证书登记谓词链：
\[
\begin{aligned}
  B_1 &: x=(p;\tau_1,\dots,\tau_N)\in\mathcal E^{(N)}
  \Rightarrow
  \Phi_N(x)\in \Pi_N^{\mathrm{vert}},\\
  B_2 &: \gamma=(x_0\to\cdots\to x_N)\in\Pi_N^{\mathrm{vert}}
  \Rightarrow
  x_N\in \mathcal E^{(N)},\\
  B_3 &: \Psi_N(\Phi_N(x))=x,\qquad \forall x\in\mathcal E^{(N)},\\
  B_4 &: \Phi_N(\Psi_N(\gamma))=\gamma,\qquad \forall \gamma\in \Pi_N^{\mathrm{vert}},\\
  B_5 &: B_3\wedge B_4 \Rightarrow \Phi_N\ \text{与}\ \Psi_N\ \text{互为逆映射}.
\end{aligned}
\]
\end{Certificate}

\begin{proof}[证明（基于数学符号推演逻辑的谓词因果链）]
先证 $B_1$。任取
\[
  x=(p;\tau_1,\dots,\tau_N)\in \mathcal E^{(N)}.
\]
由定义~\ref{def:tex37-ch13-tower}，
每个前缀
\[
  (p),\ (p;\tau_1),\ \dots,\ (p;\tau_1,\dots,\tau_k)
\]
都属于对应的 $E^{\le k}(p)$，且每一阶扩张都满足认证条件。
故由定义~\ref{def:tex37-ch13-vertical-path}，
\[
  \Phi_N(x)\in \Pi_N^{\mathrm{vert}}.
\]

再证 $B_2$。任取
\[
  \gamma=(x_0\to\cdots\to x_N)\in \Pi_N^{\mathrm{vert}}.
\]
根据竖直认证路径的定义，
\[
  x_k=(p;\tau_1,\dots,\tau_k)\in E^{\le k}(p)\qquad (0\le k\le N),
\]
尤其终点 $x_N\in E^{\le N}(p)\subseteq \mathcal E^{(N)}$。
故可定义
\[
  \Psi_N(\gamma):=x_N.
\]

对任意 $x=(p;\tau_1,\dots,\tau_N)\in\mathcal E^{(N)}$，
由 $\Phi_N$ 的定义，$\Phi_N(x)$ 的终点正是 $x$ 自身，
所以
\[
  \Psi_N(\Phi_N(x))=x,
\]
即 $B_3$。

反之，任取
\[
  \gamma=(x_0\to\cdots\to x_N)\in \Pi_N^{\mathrm{vert}}.
\]
由竖直认证路径定义，每个 $x_k$ 必为终点 $x_N$ 的第 $k$ 个前缀。
因此从 $\Psi_N(\gamma)=x_N$ 再施加 $\Phi_N$，
会重新得到同一条前缀链：
\[
  \Phi_N(\Psi_N(\gamma))=\gamma.
\]
故 $B_4$ 成立。

由 $B_3$ 与 $B_4$ 得 $\Phi_N$ 与 $\Psi_N$ 互为逆映射，即 $B_5$。
因此 $\Phi_N$ 是自然双射。证毕。
\end{proof}

\subsection{命题：认证证书把同伦扭曲自由度收紧为可收束自由度}
\label{subsec:tex37-ch13-convergence}

\begin{definition}[接口状态的闭合/细化二分]
\label{def:tex37-ch13-refine}
设 $\mathcal X$ 表示尚待签发的接口状态集合。给定秩函数
\[
  r:\mathcal X\to \NN.
\]
要求每个状态 $x\in\mathcal X$ 满足以下 fail-closed 二分：
\begin{enumerate}
  \item 若 $\Cert(x)=\top$，则允许签发 $\Issue(\Cert(x),x)=\top$；
  \item 若 $\Cert(x)=\bot$，则只能执行有限分支细化
  \[
    \Refine(x)=\{x_\beta'\}_{\beta\in B(x)},
  \]
  其中 $B(x)$ 为有限集，且对所有 $\beta\in B(x)$ 都有
  \[
    r(x_\beta')<r(x).
  \]
  若 $r(x)=0$ 且仍无法签发，则必须显式返回失败。
\end{enumerate}
\end{definition}

\begin{proposition}[认证证书把同伦扭曲自由度收紧为可收束自由度]
\label{prop:tex37-ch13-convergence}
在定义~\ref{def:tex37-ch13-refine} 的条件下，
任意合法认证分支
\[
  x_0\leadsto x_1\leadsto x_2\leadsto \cdots
\]
都不可能无限下降；因此每条分支都会在有限步后进入“签发”或“显式失败”二者之一。
特别地，同伦扭曲自由度不再是任意游走，而是受证书合同约束的可收束自由度。
\end{proposition}

\begin{Certificate}[证书：闭合签发或严格降秩细化]
\label{cert:tex37-ch13-convergence}
证书登记谓词链：
\[
\begin{aligned}
  C_1 &: x_k\ \text{若可签发，则}\ \Issue(\Cert(x_k),x_k)=\top,\\
  C_2 &: x_k\ \text{若不可签发，则}\ x_{k+1}\in \Refine(x_k),\\
  C_3 &: C_2 \Rightarrow r(x_{k+1})<r(x_k),\\
  C_4 &: r(\mathcal X)\subseteq \NN\ \text{不存在无限严格下降链},\\
  C_5 &: C_2\wedge C_3\wedge C_4 \Rightarrow \text{认证分支不可能无限细化}.
\end{aligned}
\]
\end{Certificate}

\begin{proof}[证明（基于数学符号推演逻辑的谓词因果链）]
若某一步满足 $C_1$，则分支在该步闭合并签发，结论成立。
故只需排除“永不签发而无限细化”的情况。

反设存在合法认证分支
\[
  x_0\leadsto x_1\leadsto x_2\leadsto \cdots
\]
永不结束。
由于每一步都未签发，根据定义~\ref{def:tex37-ch13-refine} 只能走细化分支，
故由 $C_2$ 与 $C_3$ 得
\[
  r(x_0)>r(x_1)>r(x_2)>\cdots .
\]
这给出 $\NN$ 上的一条无限严格下降链，
与 $C_4$ 矛盾。
因此无限细化不可能发生。

若某分支走到 $r(x)=0$ 仍不能签发，则根据定义必须显式返回失败；
故每条合法分支都在有限步后结束于“签发”或“显式失败”。
于是同伦扭曲自由度被证书合同收紧为可收束自由度。证毕。
\end{proof}

\subsection{推论：离散语义基底—认证纤维塔—竖直认证路径—定理签发构成良基闭环}
\label{subsec:tex37-ch13-corollary}

\begin{corollary}[纯数卷语义纤维丛空间的严格对象链]
\label{cor:tex37-ch13-chain}
在定义~\ref{def:tex37-ch13-placeholder}--\ref{def:tex37-ch13-homotopy-quotient} 与定理~\ref{thm:tex37-ch13-recursion}、\ref{thm:tex37-ch13-bijection} 及命题~\ref{prop:tex37-ch13-encoding}、\ref{prop:tex37-ch13-convergence} 的条件下，
G 框架纯数卷中的相关对象链可被严格压缩为
\[
  \boxed{
    \begin{aligned}
      &\text{有限字符集}
      \to
      \text{逻辑占位}
      \to
      \text{可数离散语义基底}
      \to
      \text{良基依赖}\\
      &\to
      \text{$N$ 阶认证语义纤维塔}
      \to
      \text{竖直认证路径空间}
      \to
      \text{正规形/定理签发}.
    \end{aligned}
  }
\]
并且其中所谓“基底自由度”应解释为 $\Gamma_{\mathrm{base}}$ 上的离散组合遍历自由度，
所谓“扭曲自由度”应解释为命题~\ref{prop:tex37-ch13-convergence} 所给的可收束自由度。
\end{corollary}

\begin{Certificate}[证书：编码可数性 + 良基递归 + 竖直双射 + 收束签发]
\label{cert:tex37-ch13-chain}
证书登记谓词链：
\[
\begin{aligned}
  K_1 &: \text{命题~\ref{prop:tex37-ch13-encoding}}\Rightarrow \mathcal P\subseteq \Sigma^\ast\ \text{可承载无限但可编码的语义族},\\
  K_2 &: \text{定义~\ref{def:tex37-ch13-base}}\Rightarrow \mathcal B_{\mathrm{sem}}\ \text{可数离散},\\
  K_3 &: \text{定理~\ref{thm:tex37-ch13-recursion}}\Rightarrow \mathcal E^{(N)}\ \text{可按层良基展开},\\
  K_4 &: \text{定理~\ref{thm:tex37-ch13-bijection}}\Rightarrow \mathcal E^{(N)}\cong \Pi_N^{\mathrm{vert}},\\
  K_5 &: \text{命题~\ref{prop:tex37-ch13-convergence}}\Rightarrow \text{认证分支可收束并可进入签发接口},\\
  K_6 &: K_1\wedge K_2\wedge K_3\wedge K_4\wedge K_5
  \Rightarrow
  \text{上述链条构成良基闭环}.
\end{aligned}
\]
\end{Certificate}

\begin{proof}[证明（基于数学符号推演逻辑的谓词因果链）]
由 $K_1$，有限字符集与逻辑占位层足以承载目标语义族的编码。
由 $K_2$，这些逻辑占位在对象层上形成可数离散语义基底，而不是连续流形上的连续路径对象。
由 $K_3$，良基依赖保证 $N$ 阶认证语义纤维塔可按层递归构造。
由 $K_4$，每个 $N$ 阶认证语义纤维态都对应一条唯一的竖直认证路径，反之亦然。
由 $K_5$，路径上的扭曲/细化不会无界游走，而会在有限步后进入签发或显式失败。

因此，
\[
  \text{编码层}
  \to
  \text{基底层}
  \to
  \text{纤维塔层}
  \to
  \text{认证路径层}
  \to
  \text{签发层}
\]
构成一个 fail-closed 的良基闭环。推论成立。证毕。
\end{proof}

\subsection{备注：同胚升级条件、推荐压缩表述与一句话结论}
\label{subsec:tex37-ch13-remarks}

\begin{remark}[何时可以写“同胚”，何时只能写“同伦等价”]
若在 $\mathcal E^{(N)}$ 上赋由有限坐标前缀生成的离散乘积拓扑，在 $\Pi_N^{\mathrm{vert}}$ 上赋对应的柱状拓扑，
则定理~\ref{thm:tex37-ch13-bijection} 给出的自然双射可加强为同胚。
但若进一步对路径空间取商
\[
  \Pi_N^{\mathrm{vert}}/\!\sim_h^{\mathrm{cert}},
\]
则更稳妥的记法通常先写成“状态空间与路径类空间存在自然对应”，
或在补齐商映射的同伦数据后再写
\[
  \mathcal E^{(N)} \simeq \Pi_N^{\mathrm{vert}}/\!\sim_h^{\mathrm{cert}}.
\]
在未补齐这些条件前，不应直接把“同胚”与“同伦等价”并列当作无条件同义命题。
\end{remark}

\begin{remark}[建议采用的压缩版表述]
以下口径可直接作为本章的压缩表述：
\[
  \boxed{
    \text{按逻辑占位的良基依赖关系分层，只允许向下依赖，不允许同层与向上依赖。}
  }
\]
\[
  \boxed{
    \text{逻辑占位由有限字符集上的有限串构成，并在当前认证协议下不可再无损同级拆分。}
  }
\]
\[
  \boxed{
    \begin{aligned}
      &\text{这些占位形成可数离散语义基底；其上“遍历”应理解为认证依赖图上的离散组合路径，}\\
      &\text{而非通常连续路径。}
    \end{aligned}
  }
\]
\[
  \boxed{
    \begin{aligned}
      &\text{建立在该基底上的 }N\text{ 阶认证语义纤维塔，}\\
      &\text{与固定基点上的竖直认证路径空间存在自然双射。}
    \end{aligned}
  }
\]
\[
  \boxed{
    \begin{aligned}
      &\text{认证证书通过“当前层闭合签发”或“进入严格降秩细化”两种机制，}\\
      &\text{使扭曲自由度具有可收束性。}
    \end{aligned}
  }
\]
\end{remark}

\begin{remark}[一句话结论]
\[
  \boxed{
    \begin{aligned}
      &\text{把“路径”改写为离散认证路径，把“自由度”改写为带证书收束的自由度后，}\\
      &\text{“离散语义基底--认证纤维塔--认证路径--定理签发”这条链就能严格立起来。}
    \end{aligned}
  }
\]
\end{remark}

\clearpage

%%%%%%%%%%%%%%%%%%%%%%%%%%%%%%%%%%%%%%%%%%%%%%%%%%%%%%%%%%%%
\section{形式逻辑：离散语义基底与认证纤维塔新增的主骨架级增益评价}
\label{sec:tex37-ch14}
%%%%%%%%%%%%%%%%%%%%%%%%%%%%%%%%%%%%%%%%%%%%%%%%%%%%%%%%%%%%

\begin{remark}[核心结论（评价严格化后）]
第~\ref{sec:tex37-ch13} 章的新增不是普通局部补丁，而是把
\[
  \text{逻辑占位}
  \to
  \text{良基层级依赖}
  \to
  \text{认证纤维塔}
  \to
  \text{离散认证路径}
  \to
  \text{证书收束}
\]
压成了可直接进入纯数卷的统一数学骨架。
因此，这部分新增的总评价应写为
\[
  \boxed{
    \text{增益极高，且属于“纯数卷主骨架级”显著增益。}
  }
\]
\end{remark}

本章不再新增第~\ref{sec:tex37-ch13} 章的对象，而是把“这一步究竟重在哪里”严格写成：
\[
  \begin{aligned}
    &\text{增益函数}
    \to
    \text{对象层提升}
    \to
    \text{良基性}
    \to
    \text{路径化}
    \to
    \text{坐标统一}\\
    &\to
    \text{证书收束}
    \to
    \text{方程系统可写性}
    \to
    \text{主骨架级判定}.
  \end{aligned}
\]

\begin{remark}[阅读导航：承接位置与引用口径]
本章直接承接第~\ref{sec:tex37-ch13} 章关于离散语义基底、认证纤维塔、竖直认证路径与可收束签发的形式化结果；
对“可签发方程系统”与“统一结构签名”的工程写法，继续承接第~\ref{sec:tex37-ch11}--\ref{sec:tex37-ch12} 章；
更早的语义规范场、审计闭环与运行时降级背景，仍参考 \cite{RepoTex20SemanticGaugeNormalization,RepoTex21AuditedWrapper,RepoTex23SemanticBase}；
对障碍边界、接口收束与路径/纤维对象语言，继续参考 \cite{RepoTex29RepairableObstruction,RepoTex35ExtendedCertificationTower,
  RepoTex36HomotopyTwistBijection,GaoZhengAppVol3,GaoZheng2025GFramework}。
\end{remark}

\subsection{定义：本次新增的综合增益函数、见证谓词与定性评级集}
\label{subsec:tex37-ch14-def}

\begin{definition}[本次新增的综合增益函数]
\label{def:tex37-ch14-gain}
定义第~\ref{sec:tex37-ch13} 章新增部分的综合增益为
\[
  \mathrm{Gain}_{\mathrm{sem\_bundle}}
  :=
  \frac{
    G_{\mathrm{base}}
    +
    G_{\mathrm{wf}}
    +
    G_{\mathrm{fiber}}
    +
    G_{\mathrm{path}}
    +
    G_{\mathrm{cert}}
    +
    G_{\mathrm{closure}}
    +
    G_{\mathrm{equation}}
  }{
    C_{\mathrm{assump}}
  },
\]
其中约定
\[
  C_{\mathrm{assump}}>0.
\]
各项含义如下：
\begin{enumerate}
  \item $G_{\mathrm{base}}$：离散语义基底建立增益；
  \item $G_{\mathrm{wf}}$：良基层级依赖增益；
  \item $G_{\mathrm{fiber}}$：认证语义纤维丛化增益；
  \item $G_{\mathrm{path}}$：认证路径空间化增益；
  \item $G_{\mathrm{cert}}$：证书认证增益；
  \item $G_{\mathrm{closure}}$：可收束自由度增益；
  \item $G_{\mathrm{equation}}$：可写成主基底方程系统骨架的增益；
  \item $C_{\mathrm{assump}}$：仍待显式补全的前提成本，包括最小语义单元协议、认证路径允许边全集、
  从自然双射升级到同胚所需拓扑、以及扇区扩展时的耦合规则等。
\end{enumerate}
\end{definition}

\begin{definition}[结构见证谓词与定性评级集]
\label{def:tex37-ch14-rate}
定义七个结构见证谓词：
\[
\begin{aligned}
  W_{\mathrm{base}}
  &: \mathcal B_{\mathrm{sem}}\ \text{已被明确写成可数离散语义基底},\\
  W_{\mathrm{wf}}
  &: \prec\ \text{满足只向下依赖且局部有限},\\
  W_{\mathrm{fiber}}
  &: \mathcal E^{(N)}\ \text{已被写成逐层认证纤维塔},\\
  W_{\mathrm{path}}
  &: \Gamma_{\mathrm{base}},\ \Pi_N^{\mathrm{vert}},\ \sim_h^{\mathrm{cert}}\ \text{已被定义},\\
  W_{\mathrm{cert}}
  &: \Cert,\ \Cert_n,\ \Issue\ \text{已被写入对象链},\\
  W_{\mathrm{closure}}
  &: r,\ \Refine\ \text{已把自由度改写为闭合/细化二分},\\
  W_{\mathrm{equation}}
  &: \text{第~\ref{sec:tex37-ch11}--\ref{sec:tex37-ch12} 章的结构签名已有语义基底前体实例}.
\end{aligned}
\]

再定义定性评级集
\[
  \mathfrak R_{\mathrm{gain}}
  :=
  \{\mathrm{局部},\mathrm{高},\mathrm{很高},\mathrm{极高},\mathrm{主骨架级}\},
\]
并约定其顺序为
\[
  \mathrm{局部}
  <
  \mathrm{高}
  <
  \mathrm{很高}
  <
  \mathrm{极高}
  <
  \mathrm{主骨架级}.
\]
定义本次新增的评级映射
\[
  \operatorname{Rate}_{\mathrm{sem\_bundle}}
  :=
  \begin{cases}
    \mathrm{主骨架级},
    &
    \bigwedge\limits_{\bullet\in\{\mathrm{base},\mathrm{wf},\mathrm{fiber},\mathrm{path},\mathrm{cert},\mathrm{closure},
      \mathrm{equation}\}}
    W_\bullet,\\[1ex]
    \mathrm{极高},
    &
    W_{\mathrm{base}}\wedge W_{\mathrm{wf}}\wedge W_{\mathrm{fiber}}\wedge W_{\mathrm{path}}\wedge W_{\mathrm{cert}}
      \wedge W_{\mathrm{closure}}
    \wedge \neg W_{\mathrm{equation}},\\[1ex]
    \mathrm{很高},
    &
    W_{\mathrm{base}}\wedge W_{\mathrm{wf}}\wedge W_{\mathrm{fiber}}\wedge W_{\mathrm{path}}
    \wedge \neg(W_{\mathrm{cert}}\wedge W_{\mathrm{closure}}\wedge W_{\mathrm{equation}}),\\[1ex]
    \mathrm{高},
    &
    W_{\mathrm{base}}\wedge W_{\mathrm{wf}}
    \wedge \neg(W_{\mathrm{fiber}}\wedge W_{\mathrm{path}}),\\[1ex]
    \mathrm{局部},
    & \text{其余情形}.
  \end{cases}
\]
\end{definition}

\begin{remark}[评价口径]
本章的判定不以“是否又多写了几个定义”为尺度，而以是否同时完成如下四件事为尺度：
\[
  \text{对象层提升}
  +
  \text{良基递归}
  +
  \text{证书收束}
  +
  \text{方程系统可写性}.
\]
若这四类条件同时成立，则应判为主骨架级增益，而不能降格为普通局部命题增益。
\end{remark}

\begin{remark}[对后续隐式解章节的直接推论]
一旦“离散语义基底--认证纤维塔--竖直认证路径--证书收束”这条主骨架成立，
则后续关于“系统究竟应该持久化什么”的答案也会随之改变。
更具体地说，第~\ref{sec:tex37-ch15} 章将证明：
\[
  \boxed{
    \begin{aligned}
      &\text{系统原则上无需穷举持久化所有显式答案，}\\
      &\text{而只需持久化可重构、可命中、可签发的解空间结构与认证机制。}
    \end{aligned}
  }
\]
因此，第~\ref{sec:tex37-ch14} 章的“主骨架级增益”不仅是对象论上的提升，
也是后续“隐式解空间持久化原则”得以成立的前提。
\end{remark}

\subsection{公理（降级为定理）：综合增益分子可按七类结构增量递归累加}
\label{subsec:tex37-ch14-axiom}

\begin{theorem}[公理降级：本次新增的综合增益分子可按七类结构增量递归累加（数学归纳法证书）]
\label{thm:tex37-ch14-decomposition}
定义结构增量
\[
  \begin{aligned}
    &\Delta_1:=G_{\mathrm{base}},\quad
    \Delta_2:=G_{\mathrm{wf}},\quad
    \Delta_3:=G_{\mathrm{fiber}},\quad
    \Delta_4:=G_{\mathrm{path}},\\
    &\Delta_5:=G_{\mathrm{cert}},\quad
    \Delta_6:=G_{\mathrm{closure}},\quad
    \Delta_7:=G_{\mathrm{equation}}.
  \end{aligned}
\]
再定义部分和
\[
  N(1):=\Delta_1,
  \qquad
  N(m+1):=N(m)+\Delta_{m+1}
  \quad (1\le m\le 6).
\]
则对所有 $m\in\{1,2,\dots,7\}$，有
\[
  N(m)=\sum_{k=1}^{m}\Delta_k.
\]
特别地，
\[
  N(7)
  =
  G_{\mathrm{base}}
  +
  G_{\mathrm{wf}}
  +
  G_{\mathrm{fiber}}
  +
  G_{\mathrm{path}}
  +
  G_{\mathrm{cert}}
  +
  G_{\mathrm{closure}}
  +
  G_{\mathrm{equation}},
\]
正是定义~\ref{def:tex37-ch14-gain} 中综合增益函数的分子。
\end{theorem}

\begin{Certificate}[证书：基步 + 归纳步（数学归纳法证书）]
\label{cert:tex37-ch14-decomposition}
定义谓词
\[
  P(m):\ N(m)=\sum_{k=1}^{m}\Delta_k.
\]
数学归纳法证书记为
\[
  \pi_{\mathrm{ind}}
  =
  \bigl(P(1),\ \forall m\in\{1,\dots,6\},\ P(m)\Rightarrow P(m+1)\bigr).
\]
其中
\[
  P(m)\Rightarrow P(m+1)
\]
只调用递推关系
\[
  N(m+1)=N(m)+\Delta_{m+1}.
\]
\end{Certificate}

\begin{proof}[证明（数学归纳法证书；基于数学符号推演逻辑的谓词因果链）]
\textbf{基步 $P(1)$：}
由定义，
\[
  N(1)=\Delta_1=\sum_{k=1}^{1}\Delta_k.
\]
故 $P(1)$ 成立。

\textbf{归纳步：}
设 $P(m)$ 成立，即
\[
  N(m)=\sum_{k=1}^{m}\Delta_k.
\]
则由递推式，
\[
  N(m+1)=N(m)+\Delta_{m+1}.
\]
代入归纳假设得
\[
  N(m+1)
  =
  \sum_{k=1}^{m}\Delta_k+\Delta_{m+1}
  =
  \sum_{k=1}^{m+1}\Delta_k.
\]
故 $P(m)\Rightarrow P(m+1)$ 成立。

由数学归纳法，$P(m)$ 对 $m=1,\dots,7$ 全部成立。
于是 $N(7)$ 就是七类结构增量的总和，也就是综合增益函数的分子。证毕。
\end{proof}

\subsection{定理：这部分新增的最大增益，是把已认证语义态压成几何对象}
\label{subsec:tex37-ch14-geometry}

\begin{theorem}[语义几何对象化定理]
\label{thm:tex37-ch14-geometry}
对任意 $N\in\NN$，定义映射
\[
  \operatorname{Geo}_N:\mathcal E^{(N)}\to \mathcal B_{\mathrm{sem}}\times \Pi_N^{\mathrm{vert}},
  \qquad
  \operatorname{Geo}_N(x):=\bigl(\pi_0(x),\Phi_N(x)\bigr).
\]
则 $\operatorname{Geo}_N$ 为单射。
因此，每个已认证 $N$ 阶语义态都可被忠实表示为
\[
  \text{基底点}
  +
  \text{竖直认证路径}
\]
的几何对象，而不再只是表述层文本。
\end{theorem}

\begin{Certificate}[证书：基底投影 + 竖直双射 + 单射恢复]
\label{cert:tex37-ch14-geometry}
证书登记谓词链：
\[
\begin{aligned}
  G_1 &: x\in\mathcal E^{(N)}\Rightarrow \Phi_N(x)\in \Pi_N^{\mathrm{vert}}
  \ \text{（定理~\ref{thm:tex37-ch13-bijection}）},\\
  G_2 &: \operatorname{Geo}_N(x)=\bigl(\pi_0(x),\Phi_N(x)\bigr)\ \text{良定义},\\
  G_3 &: \operatorname{Geo}_N(x_1)=\operatorname{Geo}_N(x_2)\Rightarrow \Phi_N(x_1)=\Phi_N(x_2),\\
  G_4 &: \Phi_N\ \text{为双射，故特别是单射},\\
  G_5 &: G_3\wedge G_4\Rightarrow x_1=x_2.
\end{aligned}
\]
\end{Certificate}

\begin{proof}[证明（基于数学符号推演逻辑的谓词因果链）]
由定理~\ref{thm:tex37-ch13-bijection}，任意
\[
  x\in \mathcal E^{(N)}
\]
都唯一对应一条竖直认证路径 $\Phi_N(x)$，故 $G_1$ 成立。
因此
\[
  \operatorname{Geo}_N(x)=\bigl(\pi_0(x),\Phi_N(x)\bigr)
\]
确为良定义映射，即 $G_2$。

若
\[
  \operatorname{Geo}_N(x_1)=\operatorname{Geo}_N(x_2),
\]
则特别有
\[
  \Phi_N(x_1)=\Phi_N(x_2),
\]
即 $G_3$。
再由 $G_4$，$\Phi_N$ 单射，于是立得
\[
  x_1=x_2,
\]
即 $G_5$。

故 $\operatorname{Geo}_N$ 为单射。
这说明已认证语义态可被忠实压成“基底点 + 路径对象”的几何数据，而不再只是文本表述。证毕。
\end{proof}

\begin{remark}[评价]
上面这一步的真正重量在于：
\[
  \boxed{
    \text{语义从“被说出来的内容”提升为“可被拓扑化、路径化、认证化处理的对象”。}
  }
\]
因此
\[
  G_{\mathrm{base}}+G_{\mathrm{fiber}}+G_{\mathrm{path}}
\]
应判为极高。
\end{remark}

\subsection{定理：只允许向下依赖这一条，带来了决定性的良基性增益}
\label{subsec:tex37-ch14-wf}

\begin{theorem}[良基层级依赖定理]
\label{thm:tex37-ch14-wf}
对任意 $N\in\NN$ 与任意 $x\in \mathcal E^{(N)}$，存在一个有限依赖证书树
\[
  \mathfrak T_x
\]
满足：
\begin{enumerate}
  \item $\mathfrak T_x$ 的根对应 $x$ 的基底占位 $\pi_0(x)$；
  \item 每条树边都对应一个真实依赖调用；
  \item 若 $u$ 是 $v$ 的子节点，则 $\ell(u)<\ell(v)$；
  \item $\mathfrak T_x$ 无有向环，且每条最大链在有限步内终止于低层占位。
\end{enumerate}
特别地，任何高层语义态都能沿依赖链回压到底层占位，而不存在循环依赖。
\end{theorem}

\begin{Certificate}[证书：局部有限依赖 + 严格降层 + 有限层数]
\label{cert:tex37-ch14-wf}
证书登记谓词链：
\[
\begin{aligned}
  W_1 &: x\in E^{\le 0}(p)\Rightarrow \mathfrak T_x\ \text{可取单点树},\\
  W_2 &: x_{n+1}=(x_n,\tau_{n+1})\Rightarrow \tau_{n+1}\ \text{只依赖有限个严格低层前驱}
  \ \text{（定理~\ref{thm:tex37-ch13-recursion}）},\\
  W_3 &: q\prec p\Rightarrow \ell(q)<\ell(p)
  \ \text{（定义~\ref{def:tex37-ch13-dependency}）},\\
  W_4 &: W_2\wedge W_3\Rightarrow \mathfrak T_{x_{n+1}}\ \text{可由有限棵严格降层子树拼接而成},\\
  W_5 &: W_4\Rightarrow \mathfrak T_x\ \text{有限且无环}.
\end{aligned}
\]
\end{Certificate}

\begin{proof}[证明（基于数学符号推演逻辑的谓词因果链）]
先看基步。若
\[
  x\in E^{\le 0}(p)=\{(p)\},
\]
则 $x$ 不再调用更高阶数据，其依赖证书树可以取为只含根节点 $p$ 的单点树，即 $W_1$。

对一般情形，设
\[
  x_{n+1}=(x_n,\tau_{n+1})\in E^{\le n+1}(p).
\]
由定理~\ref{thm:tex37-ch13-recursion}，
$\tau_{n+1}$ 的认证判定只调用有限个严格低层占位前驱与低阶已认证数据，故 $W_2$ 成立。
再由定义~\ref{def:tex37-ch13-dependency}，
每次真实依赖都满足
\[
  q\prec p\Rightarrow \ell(q)<\ell(p),
\]
即 $W_3$。

因此，在展开 $x_{n+1}$ 的构造时，
每一步只会生成有限个子调用，并且沿每条边层级严格下降。
故由 $W_2$ 与 $W_3$，可把 $x_{n+1}$ 的全部调用展开为一棵有限分支、严格降层的证书树 $\mathfrak T_{x_{n+1}}$，即 $W_4$。
严格降层又立即排除了回到旧节点的有向环；于是 $\mathfrak T_x$ 必有限且无环，即 $W_5$。

故任意高层语义态都能在有限步内回压到底层占位，循环依赖不可能发生。证毕。
\end{proof}

\begin{remark}[评价]
这一条不是修饰语，而是系统成立的基础条件。
没有它，语义纤维塔很容易退化为“术语很强，但接口无法闭合”的漂亮草图；
有了它，才真正进入可递归证明的范围。
因此
\[
  G_{\mathrm{wf}}
\]
应判为极高。
\end{remark}

\subsection{引理：把“基底遍历”改写成“离散认证路径”，增益非常大}
\label{subsec:tex37-ch14-path}

\begin{lemma}[基底路径概念的离散化纠偏]
\label{lem:tex37-ch14-path}
在第~\ref{sec:tex37-ch13} 章的口径下，
任何与语义基底有关的非平凡“遍历”都不能由通常拓扑连续路径表达，
而必须写成定义~\ref{def:tex37-ch13-base-path} 的基底离散认证路径。
因此，基底路径空间与其局部改写/同伦语言都获得了严格的组合对象解释。
\end{lemma}

\begin{Certificate}[证书：连续路径退化 + 离散路径显式定义]
\label{cert:tex37-ch14-path}
证书登记谓词链：
\[
\begin{aligned}
  P_1 &: \text{引理~\ref{lem:tex37-ch13-const-path}}\Rightarrow
  [0,1]\to \mathcal B_{\mathrm{sem}}\ \text{的连续路径只能常值},\\
  P_2 &: P_1\Rightarrow \text{基底上的非平凡运动不能用通常连续路径承载},\\
  P_3 &: \Gamma_{\mathrm{base}}\ \text{在定义~\ref{def:tex37-ch13-base-path} 中已被显式定义},\\
  P_4 &: \Pi_N^{\mathrm{vert}}/\!\sim_h^{\mathrm{cert}}\ \text{在定义~\ref{def:tex37-ch13-homotopy-quotient} 中已被显式定义},\\
  P_5 &: P_2\wedge P_3\wedge P_4\Rightarrow \text{路径与同伦语言获得严格离散对象解释}.
\end{aligned}
\]
\end{Certificate}

\begin{proof}[证明（基于数学符号推演逻辑的谓词因果链）]
由引理~\ref{lem:tex37-ch13-const-path}，从连通区间 $[0,1]$ 到离散语义基底 $\mathcal B_{\mathrm{sem}}$ 的连续路径只能是常值路径，即 $P_1$。
因此，只要讨论的“基底遍历”是非平凡的，就不可能再用通常拓扑意义上的连续路径承载，这就是 $P_2$。

另一方面，定义~\ref{def:tex37-ch13-base-path} 已把基底上的可遍历对象明确写成
\[
  \gamma_{\mathrm{base}}=(p_0\to\cdots\to p_m)\in \Gamma_{\mathrm{base}},
\]
即 $P_3$。
而定义~\ref{def:tex37-ch13-homotopy-quotient} 又把竖直认证路径的局部改写/压缩语言写成了商空间
\[
  \Pi_N^{\mathrm{vert}}/\!\sim_h^{\mathrm{cert}},
\]
即 $P_4$。

于是由 $P_2$、$P_3$ 与 $P_4$，原本可能含混的“基底运动”被严格替换为离散认证路径及其局部改写对象。
因此路径概念从隐喻变为可定义对象。证毕。
\end{proof}

\begin{remark}[评价]
这一步的意义在于，它把“连续几何直觉”正确投影到了“离散语义几何”里。
因此
\[
  G_{\mathrm{path}}
\]
应判为很高，而且是消除术语误用的关键纠偏增益。
\end{remark}

\subsection{定理：纤维态表示与认证路径表示是等价的两种坐标系}
\label{subsec:tex37-ch14-bijection}

\begin{theorem}[纤维态/路径态坐标统一定理]
\label{thm:tex37-ch14-bijection}
对任意谓词
\[
  P:\mathcal E^{(N)}\to\{\top,\bot\},
\]
定义其路径侧提升为
\[
  P^\sharp:\Pi_N^{\mathrm{vert}}\to\{\top,\bot\},
  \qquad
  P^\sharp(\gamma):=P(\Psi_N(\gamma)).
\]
则对任意 $x\in\mathcal E^{(N)}$，有
\[
  P(x)\Longleftrightarrow P^\sharp(\Phi_N(x)).
\]
对任意路径谓词
\[
  Q:\Pi_N^{\mathrm{vert}}\to\{\top,\bot\},
\]
定义其状态侧回拉为
\[
  Q^\flat:\mathcal E^{(N)}\to\{\top,\bot\},
  \qquad
  Q^\flat(x):=Q(\Phi_N(x)),
\]
则对任意 $\gamma\in \Pi_N^{\mathrm{vert}}$，有
\[
  Q(\gamma)\Longleftrightarrow Q^\flat(\Psi_N(\gamma)).
\]
因此，纤维态表示与认证路径表示是等价的两种坐标系。
\end{theorem}

\begin{Certificate}[证书：双向互逆 + 谓词搬运]
\label{cert:tex37-ch14-bijection}
证书登记谓词链：
\[
\begin{aligned}
  B_1 &: \Psi_N\circ \Phi_N=\id_{\mathcal E^{(N)}}
  \ \text{（定理~\ref{thm:tex37-ch13-bijection}）},\\
  B_2 &: \Phi_N\circ \Psi_N=\id_{\Pi_N^{\mathrm{vert}}}
  \ \text{（定理~\ref{thm:tex37-ch13-bijection}）},\\
  B_3 &: P^\sharp(\Phi_N(x))=P(\Psi_N(\Phi_N(x))),\\
  B_4 &: Q^\flat(\Psi_N(\gamma))=Q(\Phi_N(\Psi_N(\gamma))),\\
  B_5 &: B_1\wedge B_2\wedge B_3\wedge B_4
  \Rightarrow
  \text{状态谓词与路径谓词可双向搬运}.
\end{aligned}
\]
\end{Certificate}

\begin{proof}[证明（基于数学符号推演逻辑的谓词因果链）]
由定理~\ref{thm:tex37-ch13-bijection}，$\Phi_N$ 与 $\Psi_N$ 互为逆映射，
即 $B_1$ 与 $B_2$。

对任意状态谓词 $P$ 与任意 $x\in\mathcal E^{(N)}$，
由定义，
\[
  P^\sharp(\Phi_N(x))
  =
  P(\Psi_N(\Phi_N(x))).
\]
再由 $B_1$，右端化为
\[
  P(\Psi_N(\Phi_N(x)))=P(x).
\]
故
\[
  P(x)\Longleftrightarrow P^\sharp(\Phi_N(x)).
\]

同理，对任意路径谓词 $Q$ 与任意 $\gamma\in \Pi_N^{\mathrm{vert}}$，
由定义，
\[
  Q^\flat(\Psi_N(\gamma))
  =
  Q(\Phi_N(\Psi_N(\gamma))).
\]
再由 $B_2$，
\[
  Q(\Phi_N(\Psi_N(\gamma)))=Q(\gamma).
\]
故
\[
  Q(\gamma)\Longleftrightarrow Q^\flat(\Psi_N(\gamma)).
\]

因此，状态表示与路径表示之间不存在信息损失，
它们只是同一对象的两种坐标表达。证毕。
\end{proof}

\begin{remark}[评价]
这一点的重量在于，它把“静态结构”和“动态生成过程”统一到了同一语言。
因此
\[
  G_{\mathrm{fiber}}+G_{\mathrm{path}}
\]
应判为极高。
\end{remark}

\subsection{命题：认证证书把自由度从任意游走收紧为可收束自由度}
\label{subsec:tex37-ch14-cert}

\begin{proposition}[证书收束把自由度改写为受控演化机制]
\label{prop:tex37-ch14-cert}
设
\[
  x_0\leadsto x_1\leadsto \cdots \leadsto x_m
\]
是一条合法认证链，其中每一步要么满足
\[
  \Issue(\Cert(x_k),x_k)=\top,
\]
要么满足
\[
  x_{k+1}\in \Refine(x_k)
  \qquad\text{且}\qquad
  r(x_{k+1})<r(x_k).
\]
则若整条链在前 $m$ 步中始终未签发，必有
\[
  m\le r(x_0).
\]
因此，合法自由度链不可能无限游走；它至多经历 $r(x_0)$ 次真正细化，就必须进入“签发”或“显式失败”二分。
\end{proposition}

\begin{Certificate}[证书：严格降秩细化给出步数上界]
\label{cert:tex37-ch14-cert}
证书登记谓词链：
\[
\begin{aligned}
  C_1 &: x_{k+1}\in \Refine(x_k)\Rightarrow r(x_{k+1})\le r(x_k)-1,\\
  C_2 &: \text{若前 }m\text{ 步都未签发，则 } r(x_m)\le r(x_0)-m,\\
  C_3 &: r(x_m)\in\NN\Rightarrow r(x_m)\ge 0,\\
  C_4 &: C_2\wedge C_3\Rightarrow m\le r(x_0),\\
  C_5 &: C_4\Rightarrow \text{合法自由度链不可能无限细化}.
\end{aligned}
\]
\end{Certificate}

\begin{proof}[证明（基于数学符号推演逻辑的谓词因果链）]
若某一步已经满足
\[
  \Issue(\Cert(x_k),x_k)=\top,
\]
则链条在该步闭合并签发，无需再证。
因此只需考虑“前 $m$ 步都未签发”的情形。

在这种情形下，每一步都只能进入细化支路。
由定义~\ref{def:tex37-ch13-refine}，一旦
\[
  x_{k+1}\in \Refine(x_k),
\]
就有
\[
  r(x_{k+1})<r(x_k).
\]
由于 $r$ 取自然数值，故进一步有
\[
  r(x_{k+1})\le r(x_k)-1,
\]
即 $C_1$。
迭代应用 $C_1$，可得
\[
  r(x_m)\le r(x_0)-m,
\]
即 $C_2$。
另一方面，$r(x_m)\in\NN$，故
\[
  r(x_m)\ge 0,
\]
即 $C_3$。
于是由 $C_2$ 与 $C_3$ 得
\[
  0\le r(x_m)\le r(x_0)-m,
\]
从而
\[
  m\le r(x_0),
\]
即 $C_4$。

故真正的细化步数有显式上界，合法自由度链不可能无限游走。
因此自由度已经被改写为受证书合同约束的可收束演化机制。证毕。
\end{proof}

\begin{remark}[评价]
这一步把“自由度大”改写成了
\[
  \boxed{
    \text{自由但不失控，开放但可收束。}
  }
\]
因此
\[
  G_{\mathrm{cert}}+G_{\mathrm{closure}}
\]
应判为极高。
\end{remark}

\subsection{定理：这部分新增直接把主基底语义方程系统推进到了可下笔状态}
\label{subsec:tex37-ch14-equation}

\begin{theorem}[主基底语义前体签名定理]
\label{thm:tex37-ch14-equation}
定义
\[
  \Gamma_{\mathrm{sem}}^{(N)}
  :=
  \Pi_N^{\mathrm{vert}}/\!\sim_h^{\mathrm{cert}},
\]
\[
  \operatorname{Rep}_{\mathrm{sem}}^{(N)}:\mathcal E^{(N)}\to \Gamma_{\mathrm{sem}}^{(N)},
  \qquad
  \operatorname{Rep}_{\mathrm{sem}}^{(N)}(x):=[\Phi_N(x)],
\]
\[
  \Theta_{\mathrm{sem}}^{(N)}:\Gamma_{\mathrm{sem}}^{(N)}\to \mathcal P,
  \qquad
  \Theta_{\mathrm{sem}}^{(N)}([\gamma]):=\pi_0(\Psi_N(\gamma)).
\]
再定义语义主基底前体签名
\[
  \operatorname{PreSig}_{\mathrm{sem}}^{(N)}
  :=
  \Bigl(
    \mathcal B_{\mathrm{sem}},
    \pi_0:\mathcal E^{(N)}\to \mathcal B_{\mathrm{sem}},
    \Gamma_{\mathrm{base}},
    \Gamma_{\mathrm{sem}}^{(N)},
    \operatorname{Rep}_{\mathrm{sem}}^{(N)},
    \Theta_{\mathrm{sem}}^{(N)},
    r,
    \Cert,
    \Issue
  \Bigr).
\]
则 $\operatorname{PreSig}_{\mathrm{sem}}^{(N)}$ 良定义，
并为定义~\ref{def:tex37-ch12-signature} 中统一结构签名的基底层、状态层、路径层、代表层、投影层与签发层
提供了第一套可直接下笔的语义基底实例。
\end{theorem}

\begin{Certificate}[证书：基底/状态/路径/代表/投影/签发六层就位]
\label{cert:tex37-ch14-equation}
证书登记谓词链：
\[
\begin{aligned}
  E_1 &: \mathcal B_{\mathrm{sem}},\ \mathcal E^{(N)},\ \pi_0\ \text{在定义~\ref{def:tex37-ch13-base} 与定义~\ref{def:tex37-ch13-tower} 中已被写出},\\
  E_2 &: \Gamma_{\mathrm{base}},\ \Gamma_{\mathrm{sem}}^{(N)}\ \text{在定义~\ref{def:tex37-ch13-base-path} 与定义~\ref{def:tex37-ch13-homotopy-quotient} 中已被写出},\\
  E_3 &: \operatorname{Rep}_{\mathrm{sem}}^{(N)}(x)=[\Phi_N(x)]\ \text{良定义}
  \ \text{（定理~\ref{thm:tex37-ch13-bijection}）},\\
  E_4 &: \Theta_{\mathrm{sem}}^{(N)}\ \text{良定义，因为 }\sim_h^{\mathrm{cert}}\ \text{保持}\ \pi_0\ \text{与终点状态不变},\\
  E_5 &: r,\Cert,\Issue\ \text{在定义~\ref{def:tex37-ch13-refine} 与命题~\ref{prop:tex37-ch13-convergence} 中已进入收束/签发链},\\
  E_6 &: E_1\wedge E_2\wedge E_3\wedge E_4\wedge E_5
  \Rightarrow
  \operatorname{PreSig}_{\mathrm{sem}}^{(N)}\ \text{为统一结构签名的语义前体实例}.
\end{aligned}
\]
\end{Certificate}

\begin{proof}[证明（基于数学符号推演逻辑的谓词因果链）]
由定义~\ref{def:tex37-ch13-base} 与定义~\ref{def:tex37-ch13-tower}，
基底
\[
  \mathcal B_{\mathrm{sem}}
\]
与状态空间
\[
  \pi_0:\mathcal E^{(N)}\to\mathcal B_{\mathrm{sem}}
\]
都已经被显式写出，即 $E_1$。
由定义~\ref{def:tex37-ch13-base-path} 与定义~\ref{def:tex37-ch13-homotopy-quotient}，
基底路径层与竖直路径商层也都已给出，即 $E_2$。

再由定理~\ref{thm:tex37-ch13-bijection}，任意 $x\in \mathcal E^{(N)}$ 都唯一对应一条竖直认证路径 $\Phi_N(x)$，
故
\[
  \operatorname{Rep}_{\mathrm{sem}}^{(N)}(x):=[\Phi_N(x)]
\]
良定义，即 $E_3$。

对投影层，若 $[\gamma_1]=[\gamma_2]$，则根据定义~\ref{def:tex37-ch13-homotopy-quotient}，
认证保持同伦等价同时保持 $\pi_0$ 与终点状态不变。
因此
\[
  \pi_0(\Psi_N(\gamma_1))=\pi_0(\Psi_N(\gamma_2)),
\]
从而
\[
  \Theta_{\mathrm{sem}}^{(N)}([\gamma])=\pi_0(\Psi_N(\gamma))
\]
良定义，即 $E_4$。

最后，由定义~\ref{def:tex37-ch13-refine} 与命题~\ref{prop:tex37-ch13-convergence}，
$r,\Cert,\Issue$ 已经构成可收束、可签发的接口层，即 $E_5$。
于是由 $E_1$--$E_5$，
\[
  \operatorname{PreSig}_{\mathrm{sem}}^{(N)}
\]
已经把统一结构签名中的基底层、状态层、路径层、代表层、投影层与签发层具体化。
这说明主基底语义方程系统不再停留在抽象口号，而已有第一套可直接下笔的前体签名。证毕。
\end{proof}

\begin{remark}[评价]
这一步的关键不在于“方程已经全部写完”，而在于：
\[
  \boxed{
    \text{主基底方程系统至少已经拿到了第一套可下笔的语义基底签名。}
  }
\]
因此
\[
  G_{\mathrm{equation}}
\]
应判为极高。
\end{remark}

\subsection{推论：这部分新增应判为“纯数卷主骨架级”显著增益}
\label{subsec:tex37-ch14-cor}

\begin{corollary}[主骨架级增益判定]
\label{cor:tex37-ch14-main}
在第~\ref{sec:tex37-ch13} 章与本章前述结果条件下，
\[
  W_{\mathrm{base}}\wedge W_{\mathrm{wf}}\wedge W_{\mathrm{fiber}}\wedge W_{\mathrm{path}}\wedge W_{\mathrm{cert}}\wedge
    W_{\mathrm{closure}}\wedge W_{\mathrm{equation}}
\]
全部成立，从而
\[
  \operatorname{Rate}_{\mathrm{sem\_bundle}}=\mathrm{主骨架级}.
\]
换言之，这部分新增的评价应严格写成：
\[
  \boxed{
    \text{这是一次“纯数卷主骨架级”的显著增益，重量明显高于一般局部命题增益。}
  }
\]
\end{corollary}

\begin{Certificate}[证书：七类见证谓词全部成立]
\label{cert:tex37-ch14-main}
证书登记谓词链：
\[
\begin{aligned}
  M_1 &: \text{定理~\ref{thm:tex37-ch14-geometry}}\Rightarrow W_{\mathrm{base}}\wedge W_{\mathrm{fiber}}\wedge
    W_{\mathrm{path}},\\
  M_2 &: \text{定理~\ref{thm:tex37-ch14-wf}}\Rightarrow W_{\mathrm{wf}},\\
  M_3 &: \text{引理~\ref{lem:tex37-ch14-path}} \Rightarrow W_{\mathrm{path}},\\
  M_4 &: \text{定理~\ref{thm:tex37-ch14-bijection}}\Rightarrow W_{\mathrm{fiber}}\wedge W_{\mathrm{path}},\\
  M_5 &: \text{命题~\ref{prop:tex37-ch14-cert}}\Rightarrow W_{\mathrm{cert}}\wedge W_{\mathrm{closure}},\\
  M_6 &: \text{定理~\ref{thm:tex37-ch14-equation}}\Rightarrow W_{\mathrm{equation}},\\
  M_7 &: M_1\wedge M_2\wedge M_3\wedge M_4\wedge M_5\wedge M_6
  \Rightarrow
  \bigwedge_{\bullet} W_\bullet,\\
  M_8 &: \bigwedge_{\bullet} W_\bullet
  \Rightarrow
  \operatorname{Rate}_{\mathrm{sem\_bundle}}=\mathrm{主骨架级}
  \ \text{（定义~\ref{def:tex37-ch14-rate}）}.
\end{aligned}
\]
\end{Certificate}

\begin{proof}[证明（基于数学符号推演逻辑的谓词因果链）]
由定理~\ref{thm:tex37-ch14-geometry}，
已认证语义态已经被忠实压成基底点与路径对象的几何数据，
故 $W_{\mathrm{base}}$、$W_{\mathrm{fiber}}$ 与 $W_{\mathrm{path}}$ 成立，即 $M_1$。
由定理~\ref{thm:tex37-ch14-wf}，良基层级依赖与有限依赖证书树已经被写实，故 $W_{\mathrm{wf}}$ 成立，即 $M_2$。
由引理~\ref{lem:tex37-ch14-path}，路径对象不再是隐喻而是离散认证路径对象，故 $M_3$ 成立。
由定理~\ref{thm:tex37-ch14-bijection}，纤维态坐标与路径态坐标可双向搬运，故 $M_4$ 成立。
由命题~\ref{prop:tex37-ch14-cert}，自由度已经被收紧为可收束的受控演化机制，故 $W_{\mathrm{cert}}$ 与 $W_{\mathrm{closure}}$ 成立，即 $M_5$。
由定理~\ref{thm:tex37-ch14-equation}，主基底语义前体签名已经就位，故 $W_{\mathrm{equation}}$ 成立，即 $M_6$。

综上得到 $M_7$：
\[
  W_{\mathrm{base}}\wedge W_{\mathrm{wf}}\wedge W_{\mathrm{fiber}}\wedge W_{\mathrm{path}}\wedge W_{\mathrm{cert}}\wedge
    W_{\mathrm{closure}}\wedge W_{\mathrm{equation}}.
\]
再由定义~\ref{def:tex37-ch14-rate}，这正是评级映射取到最高档
\[
  \operatorname{Rate}_{\mathrm{sem\_bundle}}=\mathrm{主骨架级}
\]
的充足条件，即 $M_8$。
故推论成立。证毕。
\end{proof}

\subsection{备注：分层总评、最重价值、边界与一句话结论}
\label{subsec:tex37-ch14-remarks}

\begin{remark}[分层总评]
若把本次新增按层次压成总评，则最稳妥的写法是：
\[
  \begin{aligned}
    \text{纯数学骨架增益} &= \text{极高},\\
    \text{严格性增益} &= \text{极高},\\
    \text{可写性增益} &= \text{极高},\\
    \text{跨卷桥接增益} &= \text{很高}.
  \end{aligned}
\]
原因在于：语义对象已被正式基底化、纤维塔化、路径空间化与证书收束化，
并且开始与第~\ref{sec:tex37-ch11}--\ref{sec:tex37-ch12} 章的可签发方程系统共享同一底层骨架。
\end{remark}

\begin{remark}[这部分新增最重的价值，不是“又多了一个定义”，而是“纯数卷终于找到语义基底”]
过去“语义”很容易停留在解释层；
现在它被压成了
\[
  \text{离散基底}
  +
  \text{认证纤维塔}
  +
  \text{路径类}
  +
  \text{证书闭合}.
\]
这意味着“语义如何进入 G 框架纯数卷”已经不再只是强直觉，而有了可下笔的正式骨架。
\end{remark}

\begin{remark}[边界：仍有若干前提需要继续显式化]
本次评价虽已可判为主骨架级，但仍有若干前提需要继续展开：
\begin{enumerate}
  \item “最小语义单元”的判定准则需进一步写成显式协议；
  \item “认证路径”的允许边类型仍应补成完整清单；
  \item 从自然双射升级到同胚所需的拓扑仍需固定；
  \item 扇区扩展时，基底共享与扇区特化的耦合规则仍需继续写清。
\end{enumerate}
这些属于继续展开的工作，不影响当前“主骨架级增益”的判断。
\end{remark}

\begin{remark}[一句话结论]
\[
  \boxed{
    \begin{aligned}
      &\text{上述新增，等于把“语义如何进入 G 框架纯数卷”这件事，}\\
      &\text{从强直觉推进到了可写定理、可写签名、可写方程系统前体的正式骨架。}
    \end{aligned}
  }
\]
\end{remark}

\clearpage

%%%%%%%%%%%%%%%%%%%%%%%%%%%%%%%%%%%%%%%%%%%%%%%%%%%%%%%%%%%%
\section{形式逻辑：隐式解空间、上下文条件化命中与解空间持久化原则}
\label{sec:tex37-ch15}
%%%%%%%%%%%%%%%%%%%%%%%%%%%%%%%%%%%%%%%%%%%%%%%%%%%%%%%%%%%%

\begin{remark}[核心陈述（严格收紧后）]
本章把“启发命中解”收紧为如下严格口径：
\[
  \boxed{
    \begin{aligned}
      &\text{很多解不是“事先枚举并存库”的显式解，}\\
      &\text{而是原本就存在于可达支撑类中的隐式解。}
    \end{aligned}
  }
\]
\[
  \boxed{
    \begin{aligned}
      &\text{上下文的作用不是凭空创造解，}\\
      &\text{而是对隐式解空间做条件化、加权、筛选与命中。}
    \end{aligned}
  }
\]
更严格地说：
\[
  \boxed{
    \begin{aligned}
      &\text{解原本就作为可认证的同伦扭曲正规形、路径积分支撑类}\\
      &\text{与稳定修复序列存在于求解空间中；上下文只负责把其中一部分命中为显式正规形。}
    \end{aligned}
  }
\]
\end{remark}

\begin{remark}[阅读导航：承接位置与引用口径]
本章直接承接 `tex36` 关于纤维终点空间、路径积分支撑类、同伦扭曲正规形与稳定修复序列的对象链 \cite{RepoTex36HomotopyTwistBijection}，
并承接本稿第~\ref{sec:tex37-ch4} 章关于“上下文基底上的允许路径类才是内核求解对象”的结论，
以及第~\ref{sec:tex37-ch11}--\ref{sec:tex37-ch12} 章关于“真正应持久化的是可签发方程系统与统一结构签名”的结论，
再结合第~\ref{sec:tex37-ch13} 章与第~\ref{sec:tex37-ch14} 章关于离散语义基底、认证纤维塔、证书收束与主骨架级增益的结果，
把“隐式解存在、上下文命中、机制持久化”写成一套严格命题链。
\end{remark}

\subsection{定义：隐式解空间、显式库项、上下文权重与命中解}
\label{subsec:tex37-ch15-def}

\begin{definition}[隐式解空间与显式库项]
\label{def:tex37-ch15-implicit}
固定一个问题态 $x$ 与一组作用/联络数据 $A$。沿用 \cite{RepoTex36HomotopyTwistBijection} 的记号，
记
\[
  \mathcal F_A(x)
\]
为纤维终点编码集合，
\[
  \mathcal{PI}_A(x)
\]
为路径积分支撑类编码集合，
\[
  \mathcal T_A(x)
\]
为同伦扭曲正规形编码集合，
\[
  \mathcal R_A^{\mathrm{stab}}(x)
\]
为稳定修复序列集合。
假定存在编码等价
\[
  \mathcal F_A(x)
  \cong
  \mathcal{PI}_A(x)
  \cong
  \mathcal T_A(x)
  \cong
  \mathcal R_A^{\mathrm{stab}}(x).
\]
定义问题态 $x$ 的\textbf{隐式解空间}为
\[
  \mathfrak S_{\mathrm{imp}}(x)
  :=
  \mathcal{PI}_A(x)
  \cong
  \mathcal T_A(x)
  \cong
  \mathcal R_A^{\mathrm{stab}}(x).
\]

再定义\textbf{显式库项集合}
\[
  \mathfrak L_{\mathrm{exp}}(x)\subseteq \mathfrak S_{\mathrm{imp}}(x),
\]
其元素表示已被显式物化、缓存或持久化的正规形代表。
\end{definition}

\begin{definition}[结构支撑权、上下文条件化权与命中解]
\label{def:tex37-ch15-weight}
设当前上下文为
\[
  c=(\kappa_1,\dots,\kappa_m),
  \qquad
  m\in\NN.
\]
定义结构支撑权
\[
  \Omega_x:\mathfrak S_{\mathrm{imp}}(x)\to \mathbb R_{\ge 0},
\]
其值综合记录路径积分支撑、同伦扭曲可达性与微观隧穿正支撑；
定义上下文条件化权
\[
  W_c:\mathfrak S_{\mathrm{imp}}(x)\to \mathbb R_{\ge 0}.
\]
定义总命中权
\[
  H_{x,c}(s):=\Omega_x(s)\,W_c(s).
\]

再定义在上下文 $c$ 下的命中候选集
\[
  \mathfrak S_{\mathrm{hit}}(x,c)
  :=
  \left\{
    s\in \mathfrak S_{\mathrm{imp}}(x)
    \,\middle|\,
    \begin{array}{l}
      H_{x,c}(s)>0,\\
      \Norm(s)\neq \bot,\\
      \Issue(\Cert(\Norm(s)),\Norm(s))=\top
    \end{array}
  \right\}.
\]
称任意
\[
  s^\ast(x,c)\in \mathfrak S_{\mathrm{hit}}(x,c)
\]
为一个\textbf{显式命中解}。若进一步给定选择算子
\[
  \Select_c:\mathcal P(\mathfrak S_{\mathrm{hit}}(x,c))\to \mathcal P(\mathfrak S_{\mathrm{hit}}(x,c)),
\]
则可要求
\[
  s^\ast(x,c)\in \Select_c\bigl(\mathfrak S_{\mathrm{hit}}(x,c)\bigr).
\]
更严格地，定义
\[
  \Select_c\bigl(\mathfrak S_{\mathrm{hit}}(x,c)\bigr)
  :=
  \operatorname*{ArgSel}_{s\in \mathfrak S_{\mathrm{hit}}(x,c)}
  H_{x,c}(s),
\]
其中
\[
  \operatorname*{ArgSel}
\]
表示在证书完备、代价可接受与收束约束满足的前提下，对总命中权 $H_{x,c}$ 进行的选择算子。
若命中代表唯一，则其显式正规形可写为
\[
  n^\ast(x,c):=\Norm\bigl(s^\ast(x,c)\bigr),
  \qquad
  s^\ast(x,c)\in \operatorname*{ArgSel}_{s\in \mathfrak S_{\mathrm{hit}}(x,c)} H_{x,c}(s).
\]
\end{definition}

\begin{remark}[“隐式”的严格含义]
本章所谓“隐式解”不是说解不存在，而是说：
\[
  \boxed{
    \begin{aligned}
      &\text{解已经作为支撑类、正规形或稳定修复序列存在于求解空间中，}\\
      &\text{但尚未被某个具体上下文命中为当前会话中的显式输出。}
    \end{aligned}
  }
\]
\end{remark}

\begin{remark}[“启发”的严格替代写法]
在本章语境下，“启发”若不加限制，容易被误读为经验性生成。
因此更严格的术语应统一写为
\[
  \boxed{
    \text{上下文条件化}
    +
    \text{上下文诱导加权}
    +
    \text{上下文选择算子}.
  }
\]
特别地，$\Select_c$ 不是生成新解的算子，而是在
\[
  \mathfrak S_{\mathrm{hit}}(x,c)\subseteq \mathfrak S_{\mathrm{imp}}(x)
\]
上选择显式代表的算子。
\end{remark}

\subsection{公理（降级为定理）：有限上下文条件化权与命中候选集可递归构造}
\label{subsec:tex37-ch15-axiom}

\begin{theorem}[公理降级：有限上下文条件化权与命中候选集可按上下文长度递归构造（数学归纳法证书）]
\label{thm:tex37-ch15-conditioning}
对每个上下文子句 $\kappa_i$，设存在局部权因子
\[
  w_{\kappa_i}:\mathfrak S_{\mathrm{imp}}(x)\to \mathbb R_{\ge 0},
\]
以及局部可行谓词
\[
  \operatorname{Adm}_{\kappa_i}:\mathfrak S_{\mathrm{imp}}(x)\to\{\top,\bot\}.
\]
对上下文前缀
\[
  c^{(m)}:=(\kappa_1,\dots,\kappa_m)
\]
递归定义
\[
  W_{c^{(1)}}:=w_{\kappa_1},
  \qquad
  W_{c^{(m+1)}}(s):=W_{c^{(m)}}(s)\,w_{\kappa_{m+1}}(s),
\]
以及命中候选集
\[
  \mathfrak S_{\mathrm{hit}}(x,c^{(1)})
  :=
  \left\{
    s\in\mathfrak S_{\mathrm{imp}}(x)
    \,\middle|\,
    \begin{array}{l}
      \Omega_x(s)W_{c^{(1)}}(s)>0,\\
      \operatorname{Adm}_{\kappa_1}(s)=\top,\\
      \Norm(s)\neq\bot,\\
      \Issue(\Cert(\Norm(s)),\Norm(s))=\top
    \end{array}
  \right\},
\]
\[
  \mathfrak S_{\mathrm{hit}}(x,c^{(m+1)})
  :=
  \left\{
    s\in \mathfrak S_{\mathrm{hit}}(x,c^{(m)})
    \,\middle|\,
    \begin{array}{l}
      w_{\kappa_{m+1}}(s)>0,\\
      \operatorname{Adm}_{\kappa_{m+1}}(s)=\top
    \end{array}
  \right\}.
\]
则对任意有限 $m\in\NN_{>0}$，
\[
  W_{c^{(m)}}\quad\text{与}\quad \mathfrak S_{\mathrm{hit}}(x,c^{(m)})
\]
都可唯一递归构造。
\end{theorem}

\begin{Certificate}[证书：基步 + 归纳步（数学归纳法证书）]
\label{cert:tex37-ch15-conditioning}
定义谓词
\[
  P(m):\ \text{“}W_{c^{(m)}}\text{ 与 }\mathfrak S_{\mathrm{hit}}(x,c^{(m)})\text{已唯一构造完成”}.
\]
数学归纳法证书记为
\[
  \pi_{\mathrm{ind}}
  =
  \bigl(P(1),\ \forall m\in\NN_{>0},\ P(m)\Rightarrow P(m+1)\bigr).
\]
其中归纳步的显式递推式为
\[
  W_{c^{(m+1)}}(s)=W_{c^{(m)}}(s)\,w_{\kappa_{m+1}}(s),
\]
\[
  \mathfrak S_{\mathrm{hit}}(x,c^{(m+1)})
  =
  \left\{
    s\in \mathfrak S_{\mathrm{hit}}(x,c^{(m)})
    \,\middle|\,
    w_{\kappa_{m+1}}(s)>0,\ \operatorname{Adm}_{\kappa_{m+1}}(s)=\top
  \right\}.
\]
\end{Certificate}

\begin{proof}[证明（数学归纳法证书；基于数学符号推演逻辑的谓词因果链）]
\textbf{基步 $P(1)$：}
由局部权因子 $w_{\kappa_1}$ 与局部可行谓词 $\operatorname{Adm}_{\kappa_1}$ 的给定，
\[
  W_{c^{(1)}}:=w_{\kappa_1}
\]
被唯一确定；
于是
\[
  \mathfrak S_{\mathrm{hit}}(x,c^{(1)})
\]
只是从既有集合 $\mathfrak S_{\mathrm{imp}}(x)$ 中按显式谓词切出的子集，也被唯一确定。
故 $P(1)$ 成立。

\textbf{归纳步：}
设 $P(m)$ 成立，即
\[
  W_{c^{(m)}}
  \quad\text{与}\quad
  \mathfrak S_{\mathrm{hit}}(x,c^{(m)})
\]
已经唯一构造完成。
加入新子句 $\kappa_{m+1}$ 后，由递推式
\[
  W_{c^{(m+1)}}(s)=W_{c^{(m)}}(s)\,w_{\kappa_{m+1}}(s)
\]
可知新的条件化权函数被唯一确定。
再由
\[
  \mathfrak S_{\mathrm{hit}}(x,c^{(m+1)})
  =
  \left\{
    s\in \mathfrak S_{\mathrm{hit}}(x,c^{(m)})
    \,\middle|\,
    w_{\kappa_{m+1}}(s)>0,\ \operatorname{Adm}_{\kappa_{m+1}}(s)=\top
  \right\},
\]
新的命中候选集也是从旧集合中按新增上下文条件切出的子集，故同样唯一确定。

因此 $P(m)\Rightarrow P(m+1)$。
由数学归纳法，任意有限上下文前缀的条件化权与命中候选集都可唯一递归构造。证毕。
\end{proof}

\subsection{定理：很多解可以是隐式存在的，而不是必须预先持久化为显式库项}
\label{subsec:tex37-ch15-implicit}

\begin{theorem}[隐式可达解存在定理]
\label{thm:tex37-ch15-implicit}
设 $x$ 为一个问题态。若存在
\[
  s\in \mathfrak S_{\mathrm{imp}}(x)
\]
满足
\[
  \Omega_x(s)>0
  \qquad\text{且}\qquad
  \exists c,\ W_c(s)>0,
  \qquad\text{且}\qquad
  \Issue(\Cert(\Norm(s)),\Norm(s))=\top,
\]
则该解 $s$ 可以是一个\textbf{隐式可达解}：
即使
\[
  s\notin \mathfrak L_{\mathrm{exp}}(x),
\]
它仍属于可被命中、可被正规化、可被签发的解类。
因此，问题的解不必先以“显式答案表”形式预存；它可以原本就作为支撑类/正规形/稳定修复序列存在于求解空间中。
\end{theorem}

\begin{Certificate}[证书：三重编码等价 + 稳定修复序列等价 + 微观隧穿正支撑]
\label{cert:tex37-ch15-implicit}
证书登记谓词链：
\[
\begin{aligned}
  I_1 &: \mathcal F_A(x)\cong \mathcal{PI}_A(x)\cong \mathcal T_A(x)
  \ \text{（承接 \cite{RepoTex36HomotopyTwistBijection}）},\\
  I_2 &: \mathcal T_A(x)\cong \mathcal R_A^{\mathrm{stab}}(x)
  \ \text{（承接 \cite{RepoTex36HomotopyTwistBijection}）},\\
  I_3 &: \Omega_x(s)>0\Rightarrow s\ \text{属于具有正支撑的可达解类}
  \ \text{（承接 \cite{RepoTex36HomotopyTwistBijection} 的微观隧穿命中口径）},\\
  I_4 &: \exists c,\ W_c(s)>0
  \Rightarrow
  s\ \text{在某个具体上下文下满足条件化命中前提},\\
  I_5 &: I_1\wedge I_2\Rightarrow s\ \text{可被视为支撑类/正规形/稳定修复序列的同一编码对象},\\
  I_6 &: I_3\wedge I_4\wedge \Issue(\Cert(\Norm(s)),\Norm(s))=\top
  \Rightarrow
  s\ \text{可被命中且可被签发},\\
  I_7 &: I_5\wedge I_6\Rightarrow s\notin \mathfrak L_{\mathrm{exp}}(x)\ \text{并不妨碍其作为可达隐式解存在}.
\end{aligned}
\]
\end{Certificate}

\begin{proof}[证明（基于数学符号推演逻辑的谓词因果链）]
由 $I_1$ 与 $I_2$，问题态 $x$ 的求解对象并不只是一张显式答案表；
它同样可以编码为路径积分支撑类、同伦扭曲正规形或稳定修复序列。
因此任意
\[
  s\in \mathfrak S_{\mathrm{imp}}(x)
\]
都可以被视为这三类结构对象的同一编码代表，即 $I_5$。

若再满足
\[
  \Omega_x(s)>0,
\]
则由 $I_3$，$s$ 落在具有正支撑的可达解类中。
若进一步存在某个上下文 $c$ 使
\[
  W_c(s)>0,
\]
则由 $I_4$，该结构代表在该上下文下满足条件化命中前提。
若同时
\[
  \Issue(\Cert(\Norm(s)),\Norm(s))=\top,
\]
则说明该解不仅可达，而且可被提取为正规形并进入签发接口，即 $I_6$。

因此，$s$ 是否已经进入显式库项集合 $\mathfrak L_{\mathrm{exp}}(x)$，
并不是它是否属于可达解类的必要条件。
即使
\[
  s\notin \mathfrak L_{\mathrm{exp}}(x),
\]
它仍可作为支撑类/正规形/稳定修复序列的结构对象存在，并在合适上下文下被命中。
故隐式可达解存在定理成立。证毕。
\end{proof}

\subsection{引理：上下文不是创造解，而是命中解}
\label{subsec:tex37-ch15-hit}

\begin{lemma}[上下文命中引理]
\label{lem:tex37-ch15-hit}
对任意问题态 $x$ 与任意上下文 $c$，都有
\[
  \mathfrak S_{\mathrm{hit}}(x,c)\subseteq \mathfrak S_{\mathrm{imp}}(x).
\]
特别地，任意显式命中解
\[
  s^\ast(x,c)\in \mathfrak S_{\mathrm{hit}}(x,c)
\]
本来就属于隐式解空间，而不是由上下文凭空创造出来的新对象。
\end{lemma}

\begin{Certificate}[证书：命中候选集按定义从隐式解空间中切出]
\label{cert:tex37-ch15-hit}
证书登记谓词链：
\[
\begin{aligned}
  H_1 &: \mathfrak S_{\mathrm{hit}}(x,c)=\\
  &\{s\in \mathfrak S_{\mathrm{imp}}(x)\mid H_{x,c}(s)>0,\ \Norm(s)\neq\bot,\\
  &\qquad \Issue(\Cert(\Norm(s)),\Norm(s))=\top\},\\
  H_2 &: H_1\Rightarrow \forall s\in \mathfrak S_{\mathrm{hit}}(x,c),\ s\in \mathfrak S_{\mathrm{imp}}(x),\\
  H_3 &: s^\ast(x,c)\in \mathfrak S_{\mathrm{hit}}(x,c)\Rightarrow s^\ast(x,c)\in \mathfrak S_{\mathrm{imp}}(x).
\end{aligned}
\]
\end{Certificate}

\begin{proof}[证明（基于数学符号推演逻辑的谓词因果链）]
由定义~\ref{def:tex37-ch15-weight}，
\[
  \mathfrak S_{\mathrm{hit}}(x,c)
\]
本身就是从
\[
  \mathfrak S_{\mathrm{imp}}(x)
\]
中按
\[
  H_{x,c}(s)>0,\qquad \Norm(s)\neq\bot,\qquad \Issue(\Cert(\Norm(s)),\Norm(s))=\top
\]
三类条件切出的子集，即 $H_1$。
于是每个
\[
  s\in \mathfrak S_{\mathrm{hit}}(x,c)
\]
自动都属于 $\mathfrak S_{\mathrm{imp}}(x)$，即 $H_2$。
因此任意显式命中解
\[
  s^\ast(x,c)\in \mathfrak S_{\mathrm{hit}}(x,c)
\]
也自动属于 $\mathfrak S_{\mathrm{imp}}(x)$，即 $H_3$。
故上下文只是在既有隐式解空间中做条件化命中，而不是创造新解。证毕。
\end{proof}

\subsection{定理：因此没有必要把所有显式解穷举持久化到架构中}
\label{subsec:tex37-ch15-persist}

\begin{theorem}[非穷举持久化定理]
\label{thm:tex37-ch15-persist}
定义系统架构
\[
  \mathfrak A
  :=
  \Bigl(
    \mathfrak E_{\mathrm{core}},
    \{\mathfrak E_D\}_{D\in\mathfrak D_*},
    \Norm,
    \Theta,
    \Cert,
    \Issue,
    \Omega,
    \{W_c\}_c,
    \{\Select_c\}_c
  \Bigr).
\]
若对任意问题态 $x$ 与任意上下文 $c$，系统都能：
\begin{enumerate}
  \item 从 $\mathfrak E_{\mathrm{core}}$ 与 $\{\mathfrak E_D\}_{D\in\mathfrak D_*}$ 恢复允许解域；
  \item 在 $\mathfrak S_{\mathrm{imp}}(x)$ 上重构结构支撑权 $\Omega_x$ 与上下文条件化权 $W_c$；
  \item 由 $\Select_c$ 在 $\mathfrak S_{\mathrm{hit}}(x,c)$ 中重现显式命中解；
  \item 通过 $\Norm,\Cert,\Issue,\Theta$ 完成正规化、签发与占位/文本投影；
\end{enumerate}
则架构正确性并不要求持久化
\[
  \bigcup_{x,c}\mathfrak S_{\mathrm{hit}}(x,c)
\]
的全部显式元素。
换言之，原则上无需穷举持久化所有显式解。
\end{theorem}

\begin{Certificate}[证书：可重构 + 可复现 + 可签发]
\label{cert:tex37-ch15-persist}
证书登记谓词链：
\[
\begin{aligned}
  P_1 &: \mathfrak E_{\mathrm{core}},\ \{\mathfrak E_D\}_{D\in\mathfrak D_*}
  \ \text{给出可签发解空间结构}
  \ \text{（第~\ref{sec:tex37-ch11}--\ref{sec:tex37-ch12} 章）},\\
  P_2 &: \mathfrak S_{\mathrm{imp}}(x)\ \text{中的显式命中解可由}\ \Omega_x,\ W_c,\ \Select_c\ \text{重构},\\
  P_3 &: \Norm,\ \Cert,\ \Issue,\ \Theta\ \text{足以把命中解重现为可签发正规形及其投影},\\
  P_4 &: P_2\wedge P_3\Rightarrow \text{显式解的重现不依赖于其事先全量持久化},\\
  P_5 &: P_1\wedge P_4\Rightarrow \bigcup_{x,c}\mathfrak S_{\mathrm{hit}}(x,c)\ \text{无需被穷举持久化}.
\end{aligned}
\]
\end{Certificate}

\begin{proof}[证明（基于数学符号推演逻辑的谓词因果链）]
由 $P_1$，系统的核心资产首先是可签发解空间结构本身，而不是某个固定答案表；
这部分结构已经在第~\ref{sec:tex37-ch11}--\ref{sec:tex37-ch12} 章被写成主基底/扇区插件方程系统与统一结构签名。

若对给定问题态 $x$ 与上下文 $c$，系统能够在隐式解空间 $\mathfrak S_{\mathrm{imp}}(x)$ 上重构
\[
  \Omega_x,\qquad W_c,\qquad \Select_c,
\]
则显式命中解可由结构机制再现，即 $P_2$。
若再由
\[
  \Norm,\ \Cert,\ \Issue,\ \Theta
\]
把该命中解正规化、签发并投影到显式占位/文本层，则系统还能够重现该解的可审计显式形态，即 $P_3$。

因此，由 $P_2$ 与 $P_3$，
显式解是否已经事先作为库项驻留，并不是系统正确给出该解的必要条件；只要可重构、可复现、可签发，这些显式解就可以按需再现，即 $P_4$。
再结合 $P_1$，得到 $P_5$：
系统不需要穷举持久化全部显式解，而只需持久化足以重构和签发它们的结构机制。证毕。
\end{proof}

\subsection{命题：所谓“保存在同伦扭曲/路径积分空间中”，应收紧为“以可达支撑类形式存在”}
\label{subsec:tex37-ch15-support}

\begin{proposition}[隐式解的结构存在性命题]
\label{prop:tex37-ch15-support}
设
\[
  s\in \mathfrak S_{\mathrm{imp}}(x).
\]
若
\[
  \Omega_x(s)>0,
\]
则说 $s$ “存在于同伦扭曲/路径积分/稳定修复序列空间中”的严格含义应是：
\[
  \boxed{
    \begin{aligned}
      &s\text{ 以可达支撑类、可命中正规形、非零权重有效组合}\\
      &\text{或稳定修复序列的形式存在于该结构表示空间中。}
    \end{aligned}
  }
\]
而不能理解为：
\[
  \boxed{
    \text{对应的自然语言答案已经作为逐条数据库记录驻留其中。}
  }
\]
\end{proposition}

\begin{Certificate}[证书：结构表示空间与文本缓存空间类型不同]
\label{cert:tex37-ch15-support}
证书登记谓词链：
\[
\begin{aligned}
  S_1 &: \mathfrak S_{\mathrm{imp}}(x)\cong \mathcal{PI}_A(x)\cong \mathcal T_A(x)\cong \mathcal R_A^{\mathrm{stab}}(x),
    \\
  S_2 &: \mathcal{PI}_A(x),\ \mathcal T_A(x),\ \mathcal R_A^{\mathrm{stab}}(x)\ \text{都是结构表示空间},\\
  S_3 &: \text{显式文本输出位于占位/文本投影层，而非结构表示层}
  \ \text{（第~\ref{sec:tex37-ch4} 章）},\\
  S_4 &: S_1\wedge S_2\wedge S_3\Rightarrow \text{“存在于该空间中”只能理解为结构表示意义上的存在}.
\end{aligned}
\]
\end{Certificate}

\begin{proof}[证明（基于数学符号推演逻辑的谓词因果链）]
由 $S_1$，隐式解空间与路径积分支撑类、同伦扭曲正规形、稳定修复序列只是同一求解对象的不同编码。
由 $S_2$，这些空间的元素类型是“支撑类/正规形/序列”，而不是“自然语言句子记录”。
另一方面，由第~\ref{sec:tex37-ch4} 章的对象链，
显式文本只位于
\[
  \text{路径类}
  \to
  \text{规范语义代表}
  \to
  \text{逻辑占位}
  \to
  \text{表面文本}
\]
这条投影链的最末端，即 $S_3$。

因此，当说某个解“存在于同伦扭曲/路径积分/稳定修复序列空间中”时，
其严格含义只能是：该解在结构表示空间中具有一个可达代表，而不是说相应自然语言答案已经作为文本库项驻留。
故命题成立。证毕。
\end{proof}

\subsection{推论：架构应持久化“解空间与签发机制”，而不是“答案全集”}
\label{subsec:tex37-ch15-cor}

\begin{corollary}[隐式解空间持久化原则]
\label{cor:tex37-ch15-cor}
在定理~\ref{thm:tex37-ch15-persist} 与命题~\ref{prop:tex37-ch15-support} 的条件下，
系统更严格的持久化原则应写为
\[
  \boxed{
    \begin{aligned}
      &\operatorname{Persist}(\mathfrak A)
      =
      \operatorname{Persist}(\mathfrak E_{\mathrm{core}})
      \oplus
      \operatorname{Persist}\bigl(\{\mathfrak E_D\}_{D\in\mathfrak D_*}\bigr)\\
      &\oplus
      \operatorname{Persist}(\Norm,\Theta,\Cert,\Issue,\Omega,\{W_c\}_c,\{\Select_c\}_c)
      \oplus
      \operatorname{Persist}(\mathcal C_{\mathrm{NF}},\mathcal C_{\mathrm{Cert}},\mathcal C_{\mathrm{NoGo}}),
    \end{aligned}
  }
\]
而不必写成
\[
  \operatorname{Persist}\!\left(
    \bigcup_{x,c}\mathfrak S_{\mathrm{hit}}(x,c)
  \right).
\]
因此，原则上应持久化“解空间结构、命中机制、正规化/签发机制与必要缓存”，
而不是穷举持久化所有显式答案。
\end{corollary}

\begin{Certificate}[证书：机制持久化足够 + 选择性缓存允许]
\label{cert:tex37-ch15-cor}
证书登记谓词链：
\[
\begin{aligned}
  C_1 &: \text{定理~\ref{thm:tex37-ch15-persist}}\Rightarrow \text{全量显式解持久化不是正确性的必要条件},\\
  C_2 &: \text{命题~\ref{prop:tex37-ch15-support}}\Rightarrow \text{隐式解的主要存在方式是结构表示，而非文本记录},\\
  C_3 &: \mathcal C_{\mathrm{NF}},\ \mathcal C_{\mathrm{Cert}},\ \mathcal C_{\mathrm{NoGo}}
  \ \text{只构成选择性缓存，不改变主原则},\\
  C_4 &: C_1\wedge C_2\wedge C_3
  \Rightarrow
  \text{应持久化机制而非答案全集}.
\end{aligned}
\]
\end{Certificate}

\begin{proof}[证明（基于数学符号推演逻辑的谓词因果链）]
由定理~\ref{thm:tex37-ch15-persist}，只要架构保存了可重构、可复现、可签发的机制，
系统就不需要依赖全部显式命中解的事先穷举持久化，即 $C_1$。
由命题~\ref{prop:tex37-ch15-support}，隐式解的主要存在方式是结构表示空间中的可达代表，
而不是文本库项，即 $C_2$。

另一方面，出于代价与频率考虑，正规形缓存 $\mathcal C_{\mathrm{NF}}$、
证书缓存 $\mathcal C_{\mathrm{Cert}}$ 与反例/失败缓存 $\mathcal C_{\mathrm{NoGo}}$
都可以被保留，但它们只属于选择性缓存层，而不改变“主存的是机制不是答案全集”这一原则，即 $C_3$。
因此由 $C_1$、$C_2$ 与 $C_3$ 得 $C_4$，推论成立。证毕。
\end{proof}

\subsection{备注：架构意义、缓存边界、推荐压缩表述与一句话结论}
\label{subsec:tex37-ch15-remarks}

\begin{remark}[这一步对架构特别重要]
这一章把系统从
\[
  \text{检索巨型答案库}
\]
改写成
\[
  \text{在白盒解空间中按上下文命中并签发正规形}.
\]
因此系统的核心资产不再是
\[
  \text{“记住了多少答案”},
\]
而是
\[
  \boxed{
    \text{“是否掌握了可重构、可命中、可签发的解空间与机制”。}
  }
\]
\end{remark}

\begin{remark}[边界：不能绝对化成“完全不需要持久化任何解”]
本章主张的不是“绝不缓存”，而是“不必穷举持久化答案全集”。
通常仍有三类对象值得选择性缓存：
\begin{enumerate}
  \item 代价高且高频命中的正规形缓存
  \[
    \mathcal C_{\mathrm{NF}};
  \]
  \item 已签发且复核代价高的证书链缓存
  \[
    \mathcal C_{\mathrm{Cert}};
  \]
  \item 已知不可行分支、障碍模式与反例路径缓存
  \[
    \mathcal C_{\mathrm{NoGo}}.
  \]
\end{enumerate}
因此更精确的原则是：
\[
  \boxed{
    \begin{aligned}
      &\text{原则上不需要穷举持久化所有显式解；}\\
      &\text{只需按代价与频度选择性缓存正规形、证书与反例。}
    \end{aligned}
  }
\]
\end{remark}

\begin{remark}[建议采用的压缩版表述]
以下口径可直接作为本章的压缩表述：
\[
  \boxed{
    \begin{aligned}
      &\text{在该框架下，大量解以隐式方式存在于同伦扭曲正规形空间、}\\
      &\text{路径积分支撑类空间与稳定修复序列空间中。}
    \end{aligned}
  }
\]
\[
  \boxed{
    \begin{aligned}
      &\text{这些解并非必须事先作为显式占位或显式文本持久化到架构内；}\\
      &\text{它们可以仅以可达支撑类、可命中正规形与非零权重有效组合的形式存在。}
    \end{aligned}
  }
\]
\[
  \boxed{
    \begin{aligned}
      &\text{上下文的作用不是创造解，而是对隐式解空间施加条件化权重，}\\
      &\text{从而命中满足当前依赖、证书、代价与收束约束的显式正规形代表。}
    \end{aligned}
  }
\]
\[
  \boxed{
    \begin{aligned}
      &\text{因此，系统原则上只需持久化主基底方程系统、扇区插件方程系统、}\\
      &\text{正规形提取算子、证书签发接口与必要缓存，而不必穷举持久化所有显式解。}
    \end{aligned}
  }
\]
\end{remark}

\begin{remark}[一句话结论]
\[
  \boxed{
    \begin{aligned}
      &\text{很多解可以被看作原本就隐式存在于}\\
      &\text{同伦扭曲/路径积分/稳定修复序列空间中的可达正规形；}\\
      &\text{上下文只负责把它们命中并显式化，}\\
      &\text{所以架构保存“解空间与签发机制”即可，而不必保存“答案全集”。}
    \end{aligned}
  }
\]
\end{remark}

%%%%%%%%%%%%%%%%%%%%%%%%%%%%%%%%%%%%%%%%%%%%%%%%%%%%%%%%%%%%
\section{形式逻辑：逻辑性度量、占位提升与持久化判据}
\label{sec:tex37-ch16}
%%%%%%%%%%%%%%%%%%%%%%%%%%%%%%%%%%%%%%%%%%%%%%%%%%%%%%%%%%%%

\begin{remark}[核心陈述（严格化后）]
本章把“哪些被命中的占位值得进入架构本体”收紧为如下严格口径：
\[
  \boxed{
    \begin{aligned}
      &\text{不是所有被命中的占位都值得进入架构本体；}\\
      &\text{只有高逻辑性、可签发、可复用、可闭合的占位，才应被提升为持久化对象。}
    \end{aligned}
  }
\]
因此，第~\ref{sec:tex37-ch15} 章解决的是“哪些解不必事先显式存库”，
而本章要解决的是：
\[
  \boxed{
    \text{哪些已命中的占位应从隐式候选层升格为主基底/扇区插件中的持久化对象。}
  }
\]
\end{remark}

\begin{remark}[阅读导航：承接位置与引用口径]
本章直接承接第~\ref{sec:tex37-ch15} 章关于“隐式解空间、上下文条件化命中与解空间持久化原则”的结论，
并把仓内 PureMath 主稿中“逻辑占位与逻辑性度量”的总背景 \cite{GaoZheng2025GFramework}
收紧到可签发占位口径；
对认证纤维塔、结构表示与隐式可达支撑类，继续承接 \cite{RepoTex35ExtendedCertificationTower,RepoTex36HomotopyTwistBijection}；
对优先权、归因链、时间戳与证据包的仓内 formalization，参考 \cite{RepoTex23PriorityGovernance}。
\end{remark}

\subsection{定义：逻辑性度量、三类占位、提升算子与证据包}
\label{subsec:tex37-ch16-def}

\begin{definition}[逻辑占位全集与逻辑性度量]
\label{def:tex37-ch16-logicity}
为避免与幂集记号 $\mathcal P(\cdot)$ 混淆，定义逻辑占位全集为
\[
  \mathfrak P
  :=
  \mathsf{PL}_{\mathrm{NL}}/\!\equiv_{\mathrm{core}},
\]
其中
\[
  p_1\equiv_{\mathrm{core}}p_2
\]
表示二者具有相同的量词骨架、谓词依赖、证书槽、引用接口与失败口径，只允许表面措辞不同。

再定义六个结构分量
\[
  \operatorname{Nec},\operatorname{Gen},\operatorname{Trans},
  \operatorname{Comp},\operatorname{CertV},\operatorname{Vol}
  :\mathfrak P\to\mathbb R_{\ge 0},
\]
并定义\textbf{逻辑性度量}
\[
  \mathfrak L:\mathfrak P\to \mathbb R_{\ge 0}
\]
为
\[
  \mathfrak L(p)
  :=
  \frac{
    \alpha_1\operatorname{Nec}(p)
    +
    \alpha_2\operatorname{Gen}(p)
    +
    \alpha_3\operatorname{Trans}(p)
    +
    \alpha_4\operatorname{Comp}(p)
    +
    \alpha_5\operatorname{CertV}(p)
  }{
    1+\alpha_6\operatorname{Vol}(p)
  },
  \qquad
  \alpha_i>0.
\]
其中：
\begin{enumerate}
  \item $\operatorname{Nec}(p)$ 衡量删除 $p$ 后对主基底、接口闭环与定理链的破坏强度；
  \item $\operatorname{Gen}(p)$ 衡量 $p$ 对下游占位、定理、插件与规则族的生成/支撑能力；
  \item $\operatorname{Trans}(p)$ 衡量 $p$ 的跨扇区迁移能力；
  \item $\operatorname{Comp}(p)$ 衡量 $p$ 带来的证明压缩、接口压缩与搜索压缩收益；
  \item $\operatorname{CertV}(p)$ 衡量 $p$ 的可认证、可签发、可复审程度；
  \item $\operatorname{Vol}(p)$ 衡量 $p$ 对局部上下文的依赖波动，即离开当前问题后快速失效的程度。
\end{enumerate}
因此，$\mathfrak L(p)$ 不是文本长度，也不是表面复杂度，而是 $p$ 的结构重量、持久化必要性与签发价值。
\end{definition}

\begin{definition}[三类占位与持久化提升算子]
\label{def:tex37-ch16-promote}
固定阈值
\[
  0<\theta_{\mathrm{sec}}<\theta_{\mathrm{arch}}.
\]
定义三类占位：
\[
  \mathfrak P_{\mathrm{arch}}
  :=
  \left\{
    p\in\mathfrak P
    \,\middle|\,
    \mathfrak L(p)\ge \theta_{\mathrm{arch}}
    \ \wedge\
    \Issue(\Cert(p),p)=\top
  \right\},
\]
\[
  \mathfrak P_{\mathrm{sec}}
  :=
  \left\{
    p\in\mathfrak P
    \,\middle|\,
    \theta_{\mathrm{sec}}\le \mathfrak L(p)<\theta_{\mathrm{arch}}
    \ \wedge\
    \Issue(\Cert(p),p)=\top
  \right\},
\]
\[
  \mathfrak P_{\mathrm{imp}}
  :=
  \left\{
    p\in\mathfrak P
    \,\middle|\,
    \mathfrak L(p)<\theta_{\mathrm{sec}}
    \ \vee\
    \Issue(\Cert(p),p)\neq\top
  \right\}.
\]
它们分别表示：
\begin{enumerate}
  \item \textbf{架构级占位：}必须进入主基底或共享核心接口；
  \item \textbf{扇区级占位：}应进入对应扇区插件，但不一定上升到核心主基底；
  \item \textbf{隐式/瞬时占位：}可停留在隐式解空间、命中空间与支撑类空间中，不必进入架构本体。
\end{enumerate}

再定义提升算子
\[
  \Promote:\mathfrak P\to\{\mathrm{arch},\mathrm{sector},\mathrm{implicit}\}
\]
为
\[
  \Promote(p)
  :=
  \begin{cases}
    \mathrm{arch},
    &
    \mathfrak L(p)\ge \theta_{\mathrm{arch}}
    \ \wedge\
    \Issue(\Cert(p),p)=\top,\\[4pt]
    \mathrm{sector},
    &
    \theta_{\mathrm{sec}}\le \mathfrak L(p)<\theta_{\mathrm{arch}}
    \ \wedge\
    \Issue(\Cert(p),p)=\top,\\[4pt]
    \mathrm{implicit},
    &
    \text{其余情形}.
  \end{cases}
\]
\end{definition}

\begin{definition}[高逻辑性占位的技术证据包]
\label{def:tex37-ch16-evidence}
对任意满足
\[
  \Promote(p)\in\{\mathrm{arch},\mathrm{sector}\}
\]
的占位 $p$，定义其技术证据包为
\[
  \Pi_{\mathrm{evid}}(p)
  :=
  \Bigl(
    \operatorname{Author}(p),
    \operatorname{Time}(p),
    \operatorname{Ctx}(p),
    \operatorname{Dep}(p),
    \mathfrak L(p),
    \Cert(p),
    \Issue(\Cert(p),p),
    \operatorname{Version}(p)
  \Bigr).
\]
其中各分量依次记录：提出者、时间戳、提出时上下文、依赖链、逻辑性度量、证书链、签发状态与版本演化信息。
\end{definition}

\begin{remark}[逻辑性强的占位不等于文本更长]
本章所谓“逻辑性强”并不意味着文本更长，而意味着：
\[
  \boxed{
    \text{删掉会明显破坏闭环、留下能继续生长、签发后可稳定复用。}
  }
\]
\end{remark}

\subsection{公理（降级为定理）：有限逻辑性度量可按分量数递归构造}
\label{subsec:tex37-ch16-axiom}

\begin{theorem}[公理降级：有限逻辑性度量可按分量数递归构造（数学归纳法证书）]
\label{thm:tex37-ch16-aggregation}
记
\[
  \Phi_1:=\operatorname{Nec},\quad
  \Phi_2:=\operatorname{Gen},\quad
  \Phi_3:=\operatorname{Trans},\quad
  \Phi_4:=\operatorname{Comp},\quad
  \Phi_5:=\operatorname{CertV}.
\]
对任意 $p\in\mathfrak P$ 递归定义部分和
\[
  \Lambda_1(p):=\alpha_1\Phi_1(p),
  \qquad
  \Lambda_{n+1}(p):=\Lambda_n(p)+\alpha_{n+1}\Phi_{n+1}(p)
  \quad (1\le n\le 4).
\]
再定义
\[
  \mathfrak L_n(p)
  :=
  \frac{\Lambda_n(p)}{1+\alpha_6\operatorname{Vol}(p)}
  \qquad (1\le n\le 5).
\]
则对任意 $n\in\{1,2,3,4,5\}$，
\[
  \Lambda_n(p)=\sum_{k=1}^{n}\alpha_k\Phi_k(p),
  \qquad
  \mathfrak L_n(p)=
  \frac{\sum_{k=1}^{n}\alpha_k\Phi_k(p)}{1+\alpha_6\operatorname{Vol}(p)}.
\]
特别地，
\[
  \mathfrak L_5(p)=\mathfrak L(p).
\]
\end{theorem}

\begin{Certificate}[证书：基步 + 归纳步（数学归纳法证书）]
\label{cert:tex37-ch16-aggregation}
定义谓词
\[
  P(n):\
  \Lambda_n(p)=\sum_{k=1}^{n}\alpha_k\Phi_k(p)
  \ \wedge\
  \mathfrak L_n(p)=
  \frac{\sum_{k=1}^{n}\alpha_k\Phi_k(p)}{1+\alpha_6\operatorname{Vol}(p)}.
\]
数学归纳法证书记为
\[
  \pi_{\mathrm{ind}}
  =
  \bigl(
    P(1),\
    \forall n\in\{1,2,3,4\},\ P(n)\Rightarrow P(n+1)
  \bigr).
\]
\end{Certificate}

\begin{proof}[证明（数学归纳法证书；基于数学符号推演逻辑的谓词因果链）]
\textbf{基步 $P(1)$：}
由定义，
\[
  \Lambda_1(p)=\alpha_1\Phi_1(p)
  =
  \sum_{k=1}^{1}\alpha_k\Phi_k(p).
\]
于是
\[
  \mathfrak L_1(p)
  =
  \frac{\Lambda_1(p)}{1+\alpha_6\operatorname{Vol}(p)}
  =
  \frac{\sum_{k=1}^{1}\alpha_k\Phi_k(p)}{1+\alpha_6\operatorname{Vol}(p)}.
\]
故 $P(1)$ 成立。

\textbf{归纳步：}
设 $P(n)$ 成立，即
\[
  \Lambda_n(p)=\sum_{k=1}^{n}\alpha_k\Phi_k(p).
\]
则由递推式
\[
  \Lambda_{n+1}(p)=\Lambda_n(p)+\alpha_{n+1}\Phi_{n+1}(p)
\]
可得
\[
  \Lambda_{n+1}(p)
  =
  \sum_{k=1}^{n}\alpha_k\Phi_k(p)+\alpha_{n+1}\Phi_{n+1}(p)
  =
  \sum_{k=1}^{n+1}\alpha_k\Phi_k(p).
\]
再把上式除以
\[
  1+\alpha_6\operatorname{Vol}(p)>0
\]
即得
\[
  \mathfrak L_{n+1}(p)
  =
  \frac{\sum_{k=1}^{n+1}\alpha_k\Phi_k(p)}{1+\alpha_6\operatorname{Vol}(p)}.
\]
故 $P(n)\Rightarrow P(n+1)$。
由数学归纳法，$P(n)$ 对 $n=1,\dots,5$ 全部成立，特别有
\[
  \mathfrak L_5(p)=\mathfrak L(p).
\]
证毕。
\end{proof}

\subsection{定理：逻辑性度量给出占位持久化的正式判据}
\label{subsec:tex37-ch16-theorem-criterion}

\begin{theorem}[逻辑性度量持久化判据定理]
\label{thm:tex37-ch16-criterion}
设架构持久化策略满足下列准则：
\begin{enumerate}
  \item[\textnormal{(H1)}] 若删除某占位 $p$ 会使主基底、共享接口或跨扇区闭环破坏超过容忍阈值，则 $p$ 必须被持久化到某一层级；
  \item[\textnormal{(H2)}] 生成能力、迁移能力、压缩收益与可签发性越高的占位，应被优先提升到更高持久化层级；
  \item[\textnormal{(H3)}] 上下文波动性高、不可签发或高度局部化的占位，不得无条件进入核心主基底。
\end{enumerate}
则存在阈值
\[
  0<\theta_{\mathrm{sec}}<\theta_{\mathrm{arch}}
\]
使得对任意 $p\in\mathfrak P$，
\[
  \Promote(p)=\mathrm{arch}
  \iff
  p\in\mathfrak P_{\mathrm{arch}},
\]
\[
  \Promote(p)=\mathrm{sector}
  \iff
  p\in\mathfrak P_{\mathrm{sec}},
\]
\[
  \Promote(p)=\mathrm{implicit}
  \iff
  p\in\mathfrak P_{\mathrm{imp}}.
\]
因此，在上述策略下，逻辑性度量与签发谓词恰好刻画了“哪些占位应进入架构本体、哪些占位只应保留在扇区层、哪些占位可停留于隐式空间”。
\end{theorem}

\begin{Certificate}[证书：高结构重量应持久化 + 高波动对象不得直接入核]
\label{cert:tex37-ch16-criterion}
证书登记谓词链：
\[
\begin{aligned}
  T_1 &: \text{第~\ref{sec:tex37-ch15} 章已证：不是所有命中解都需要持久化},\\
  T_2 &: \textnormal{(H1)}\Rightarrow \operatorname{Nec}(p)\ \text{高的对象必须至少进入某个持久化层级},\\
  T_3 &: \textnormal{(H2)}\Rightarrow \operatorname{Gen}(p),\operatorname{Trans}(p),\operatorname{Comp}(p),
    \operatorname{CertV}(p)\ \text{越高，提升层级越高},\\
  T_4 &: \textnormal{(H3)}\Rightarrow \operatorname{Vol}(p)\ \text{高或不可签发对象不得直接进入核心主基底},\\
  T_5 &: T_2\wedge T_3\wedge T_4\Rightarrow \mathfrak L(p)\ \text{可被阈值化为三层提升规则},\\
  T_6 &: T_5\wedge \text{定义~\ref{def:tex37-ch16-promote}}
  \Rightarrow
  \Promote\ \text{恰好给出持久化判据}.
\end{aligned}
\]
\end{Certificate}

\begin{proof}[证明（基于数学符号推演逻辑的谓词因果链）]
由第~\ref{sec:tex37-ch15} 章，系统不应穷举持久化所有命中对象，故 $T_1$ 成立。
因此，真正需要补出的不是“再多存一点”，而是“什么对象值得被持久化”。

由 \textnormal{(H1)}，若某占位的删除会显著破坏主基底、共享接口或跨扇区闭环，
则其 $\operatorname{Nec}(p)$ 必然应被赋予高权重，并且该对象必须被持久化到至少一个层级，即 $T_2$。
由 \textnormal{(H2)}，若某占位具有高生成、高迁移、高压缩收益与高可签发性，
则它不是局部技巧，而是长期可复用的结构支点，故应被提升到更高层级，即 $T_3$。
由 \textnormal{(H3)}，若某占位的上下文波动性高，或者尚未完成签发，
则即使它曾被某次命中，也不应无条件进入核心主基底，即 $T_4$。

把 $T_2$、$T_3$ 与 $T_4$ 合并，可知一个自然的三层判据必须同时奖励
\[
  \operatorname{Nec}(p),\operatorname{Gen}(p),\operatorname{Trans}(p),\operatorname{Comp}(p),\operatorname{CertV}(p)
\]
并惩罚
\[
  \operatorname{Vol}(p).
\]
定义~\ref{def:tex37-ch16-logicity} 中的 $\mathfrak L(p)$ 正是这样一个加权比值。
因此存在阈值
\[
  \theta_{\mathrm{sec}}<\theta_{\mathrm{arch}}
\]
把对象划分为“核心主基底值得持久化”“扇区插件值得持久化”“只适合停留于隐式空间”三类，即 $T_5$。
再由定义~\ref{def:tex37-ch16-promote}，$\Promote$ 正是这一阈值化规则的显式写法，故 $T_6$ 成立。
因此逻辑性度量与签发谓词恰好给出了占位持久化的正式判据。证毕。
\end{proof}

\subsection{引理：提升算子给出互斥且穷尽的三分层}
\label{subsec:tex37-ch16-lemma}

\begin{lemma}[占位提升三分层引理]
\label{lem:tex37-ch16-trichotomy}
对任意 $p\in\mathfrak P$，恰有且仅有一种情形成立：
\[
  p\in \mathfrak P_{\mathrm{arch}},
  \qquad
  p\in \mathfrak P_{\mathrm{sec}},
  \qquad
  p\in \mathfrak P_{\mathrm{imp}}.
\]
等价地，
\[
  \mathfrak P
  =
  \mathfrak P_{\mathrm{arch}}
  \sqcup
  \mathfrak P_{\mathrm{sec}}
  \sqcup
  \mathfrak P_{\mathrm{imp}},
\]
并且
\[
  \Promote(p)
\]
总是良定义。
\end{lemma}

\begin{Certificate}[证书：阈值互斥 + 剩余情形并入隐式层]
\label{cert:tex37-ch16-trichotomy}
证书登记谓词链：
\[
\begin{aligned}
  L_1 &: 0<\theta_{\mathrm{sec}}<\theta_{\mathrm{arch}},\\
  L_2 &: \mathfrak L(p)\ge \theta_{\mathrm{arch}}
  \Rightarrow
  \neg\bigl(\theta_{\mathrm{sec}}\le \mathfrak L(p)<\theta_{\mathrm{arch}}\bigr),\\
  L_3 &: \theta_{\mathrm{sec}}\le \mathfrak L(p)<\theta_{\mathrm{arch}}
  \Rightarrow
  \neg\bigl(\mathfrak L(p)\ge \theta_{\mathrm{arch}}\bigr),\\
  L_4 &: \neg\bigl(p\in\mathfrak P_{\mathrm{arch}}\cup\mathfrak P_{\mathrm{sec}}\bigr)
  \Rightarrow
  p\in\mathfrak P_{\mathrm{imp}},\\
  L_5 &: L_2\wedge L_3\wedge L_4
  \Rightarrow
  \mathfrak P_{\mathrm{arch}},\mathfrak P_{\mathrm{sec}},\mathfrak P_{\mathrm{imp}}\ \text{互斥且穷尽}.
\end{aligned}
\]
\end{Certificate}

\begin{proof}[证明（基于数学符号推演逻辑的谓词因果链）]
由 $L_1$，两个阈值满足严格不等式
\[
  \theta_{\mathrm{sec}}<\theta_{\mathrm{arch}}.
\]
因此若
\[
  \mathfrak L(p)\ge \theta_{\mathrm{arch}},
\]
则不可能同时满足
\[
  \theta_{\mathrm{sec}}\le \mathfrak L(p)<\theta_{\mathrm{arch}},
\]
即 $L_2$。
反过来，若
\[
  \theta_{\mathrm{sec}}\le \mathfrak L(p)<\theta_{\mathrm{arch}},
\]
则显然不可能再满足
\[
  \mathfrak L(p)\ge \theta_{\mathrm{arch}},
\]
即 $L_3$。

另一方面，若某占位既不属于 $\mathfrak P_{\mathrm{arch}}$，也不属于 $\mathfrak P_{\mathrm{sec}}$，
则按定义~\ref{def:tex37-ch16-promote}，它或者度量低于 $\theta_{\mathrm{sec}}$，
或者尚未签发，因而自动并入 $\mathfrak P_{\mathrm{imp}}$，即 $L_4$。
故由 $L_2$、$L_3$ 与 $L_4$ 可知，三类集合互斥且穷尽，$\Promote$ 总是良定义。证毕。
\end{proof}

\subsection{命题：科研过程可被写成“从低度量候选云中提炼并固化高度量占位”}
\label{subsec:tex37-ch16-proposition}

\begin{proposition}[科研过程的持久化演化命题]
\label{prop:tex37-ch16-research}
设第 $t$ 阶段的候选占位集合为
\[
  Q_t\subseteq \mathfrak P,
\]
已持久化集合为
\[
  A_t\subseteq \mathfrak P_{\mathrm{arch}}\cup\mathfrak P_{\mathrm{sec}}.
\]
若研究演化由
\[
  Q_t
  \xrightarrow{\text{探索/命中/修复}}
  Q_{t+1}
\]
推进，则持久化更新可严格写为
\[
  A_{t+1}
  :=
  A_t
  \cup
  \left\{
    p\in Q_{t+1}
    \,\middle|\,
    \Promote(p)\in\{\mathrm{arch},\mathrm{sector}\}
  \right\}.
\]
因此，科研过程不是“把所有中间想法都存起来”，而是：
\[
  \boxed{
    \text{从低逻辑性候选云中，发现、证明、签发并固化少量高逻辑性占位。}
  }
\]
\end{proposition}

\begin{Certificate}[证书：候选云扩张 + 提升判据筛选 + 持久化并集更新]
\label{cert:tex37-ch16-research}
证书登记谓词链：
\[
\begin{aligned}
  P_1 &: Q_t\xrightarrow{\text{探索/命中/修复}}Q_{t+1},\\
  P_2 &: p\in Q_{t+1}\ \wedge\ \Promote(p)\in\{\mathrm{arch},\mathrm{sector}\}
  \Rightarrow
  p\ \text{应进入持久化层},\\
  P_3 &: p\in Q_{t+1}\ \wedge\ \Promote(p)=\mathrm{implicit}
  \Rightarrow
  p\ \text{可继续停留在隐式候选层},\\
  P_4 &: P_2\wedge P_3
  \Rightarrow
  A_{t+1}
  =
  A_t\cup
  \{p\in Q_{t+1}\mid \Promote(p)\in\{\mathrm{arch},\mathrm{sector}\}\}.
\end{aligned}
\]
\end{Certificate}

\begin{proof}[证明（基于数学符号推演逻辑的谓词因果链）]
由 $P_1$，研究过程的每一步都会产生新的候选占位集合 $Q_{t+1}$。
这些对象并不天然都值得进入架构本体，而必须经过第~\ref{sec:tex37-ch16} 章的提升判据筛选。

若某个 $p\in Q_{t+1}$ 满足
\[
  \Promote(p)\in\{\mathrm{arch},\mathrm{sector}\},
\]
则由定义~\ref{def:tex37-ch16-promote}，它已经具备足够高的逻辑性且通过了签发接口，
因此应并入持久化集合，即 $P_2$。
反之，若
\[
  \Promote(p)=\mathrm{implicit},
\]
则它要么逻辑性不足，要么尚未签发，因而仍可停留在隐式层继续等待后续发现、细化或淘汰，即 $P_3$。

把 $P_2$ 与 $P_3$ 合并，就得到唯一合理的更新规则
\[
  A_{t+1}
  =
  A_t
  \cup
  \left\{
    p\in Q_{t+1}
    \,\middle|\,
    \Promote(p)\in\{\mathrm{arch},\mathrm{sector}\}
  \right\}.
\]
故科研过程的严格写法不是“保存所有中间命中”，而是“从候选云中筛出并固化高逻辑性占位”。证毕。
\end{proof}

\subsection{定理：该持久化判据与第 15 章的隐式解空间完全兼容}
\label{subsec:tex37-ch16-theorem-compatible}

\begin{theorem}[逻辑性度量与隐式解空间兼容定理]
\label{thm:tex37-ch16-compatible}
设问题态 $x$ 在上下文 $c$ 下命中一个显式正规形
\[
  s^\ast(x,c)\in \mathfrak S_{\mathrm{hit}}(x,c),
\]
并令
\[
  p^\ast(x,c):=\Theta\bigl(\Norm(s^\ast(x,c))\bigr)\in\mathfrak P
\]
表示其逻辑占位代表。
则：
\begin{enumerate}
  \item 若 $\Promote\bigl(p^\ast(x,c)\bigr)=\mathrm{implicit}$，则 $p^\ast(x,c)$ 可继续仅以隐式占位代表的方式停留于支撑类/正规形/稳定修复序列空间，
    不要求持久化进架构本体；
  \item 若 $\Promote\bigl(p^\ast(x,c)\bigr)=\mathrm{sector}$，则 $p^\ast(x,c)$ 应被签发并提升到对应扇区插件；
  \item 若 $\Promote\bigl(p^\ast(x,c)\bigr)=\mathrm{arch}$，则 $p^\ast(x,c)$ 应被签发并提升到主基底或共享核心接口。
\end{enumerate}
因此：
\[
  \boxed{
    \begin{aligned}
      &\text{隐式解空间负责保留可达可能性，}\\
      &\text{逻辑性度量负责决定哪些可能性应升格为持久化架构对象。}
    \end{aligned}
  }
\]
\end{theorem}

\begin{Certificate}[证书：第 15 章给出隐式保留原则 + 本章给出提升判据]
\label{cert:tex37-ch16-compatible}
证书登记谓词链：
\[
\begin{aligned}
  K_1 &: \text{第~\ref{sec:tex37-ch15} 章已证：低价值或未签发对象可停留在隐式解空间},\\
  K_2 &: p^\ast(x,c)=\Theta(\Norm(s^\ast(x,c)))\ \text{把命中正规形投影为逻辑占位代表},\\
  K_3 &: \Promote(p^\ast)=\mathrm{implicit}
  \Rightarrow
  p^\ast\in\mathfrak P_{\mathrm{imp}},\\
  K_4 &: \Promote(p^\ast)=\mathrm{sector}
  \Rightarrow
  p^\ast\in\mathfrak P_{\mathrm{sec}},\\
  K_5 &: \Promote(p^\ast)=\mathrm{arch}
  \Rightarrow
  p^\ast\in\mathfrak P_{\mathrm{arch}},\\
  K_6 &: K_1\wedge K_3\wedge K_4\wedge K_5
  \Rightarrow
  \text{隐式保留与持久化提升构成上下游兼容关系}.
\end{aligned}
\]
\end{Certificate}

\begin{proof}[证明（基于数学符号推演逻辑的谓词因果链）]
由第~\ref{sec:tex37-ch15} 章，系统不需要穷举持久化所有显式命中对象；
大量对象完全可以只以支撑类、正规形或稳定修复序列的形式保留在隐式空间中，即 $K_1$。
由 $K_2$，每个命中正规形都可经由占位投影算子 $\Theta$ 被压成一个逻辑占位代表 $p^\ast(x,c)$。

若
\[
  \Promote\bigl(p^\ast(x,c)\bigr)=\mathrm{implicit},
\]
则由定义~\ref{def:tex37-ch16-promote}，它不是高逻辑性且已签发的持久化对象，故仍可停留在隐式层，即 $K_3$。
若
\[
  \Promote\bigl(p^\ast(x,c)\bigr)=\mathrm{sector},
\]
则它已满足扇区层面的逻辑性与签发条件，应进入对应扇区插件，即 $K_4$。
若
\[
  \Promote\bigl(p^\ast(x,c)\bigr)=\mathrm{arch},
\]
则它已满足核心层面的逻辑性与签发条件，应进入主基底或共享接口，即 $K_5$。

因此，第~\ref{sec:tex37-ch15} 章并不是在否定持久化，而是在否定“无差别地持久化所有命中对象”。
本章则进一步给出：哪些对象应继续停留在隐式层，哪些对象应被提升到扇区或架构层。
故二者不是冲突，而是上下游兼容关系，即 $K_6$。证毕。
\end{proof}

\subsection{定理：高逻辑性、已签发占位可形成稳定的技术归因与优先权证据链}
\label{subsec:tex37-ch16-theorem-attribution}

\begin{theorem}[技术归因与优先权证据链定理]
\label{thm:tex37-ch16-attribution}
对任意满足
\[
  \Promote(p)\in\{\mathrm{arch},\mathrm{sector}\}
  \quad\text{且}\quad
  \Issue(\Cert(p),p)=\top
\]
的占位 $p$，其技术证据包
\[
  \Pi_{\mathrm{evid}}(p)
\]
给出一条可回放、可审计、可复核的结构归因链。
更具体地说，它能够同时回答：
\begin{enumerate}
  \item 谁提出了该结构对象；
  \item 它在何时、何种上下文下首次被明确提出；
  \item 它依赖哪些上游对象，并支撑哪些下游对象；
  \item 它具有多高的结构重量与持久化价值；
  \item 它是否已经完成签发，以及版本如何演化。
\end{enumerate}
因此，$\Pi_{\mathrm{evid}}(p)$ 是 AI 时代技术归因、优先权记录与结构性创新确证的强技术证据基础。
但应严格说明：
\[
  \boxed{
    \text{它首先是技术证据链与结构归因链，而不是对最终法律权利的自动裁决。}
  }
\]
\end{theorem}

\begin{Certificate}[证书：作者/时间/依赖/度量/签发/版本六链闭合]
\label{cert:tex37-ch16-attribution}
证书登记谓词链：
\[
\begin{aligned}
  A_1 &: \operatorname{Author}(p),\operatorname{Time}(p),\operatorname{Version}(p)
  \ \text{给出对象的提出者、时间戳与版本链},\\
  A_2 &: \operatorname{Ctx}(p),\operatorname{Dep}(p)
  \ \text{给出提出场景与上下游依赖链},\\
  A_3 &: \mathfrak L(p)
  \ \text{给出对象的结构重量、持久化必要性与复用价值},\\
  A_4 &: \Cert(p),\Issue(\Cert(p),p)
  \ \text{给出对象是否已签发及其证书链},\\
  A_5 &: A_1\wedge A_2\wedge A_3\wedge A_4
  \Rightarrow
  \Pi_{\mathrm{evid}}(p)\ \text{形成可回放的技术归因与优先权证据链},\\
  A_6 &: A_5\Rightarrow \text{可区分“结构贡献”与“表面措辞”},\\
  A_7 &: A_6\Rightarrow \text{该证据链可服务于归因、优先权记录与创新估值，但不直接替代法域裁决}.
\end{aligned}
\]
\end{Certificate}

\begin{proof}[证明（基于数学符号推演逻辑的谓词因果链）]
由 $A_1$，证据包首先记录了提出者、时间戳与版本演化信息，
因此它能回答“谁在何时以哪个版本提出了该结构对象”。
由 $A_2$，证据包进一步记录提出时上下文与依赖链，
因此它能回答“这个对象依赖什么、又支撑什么”。
由 $A_3$，证据包还记录了逻辑性度量 $\mathfrak L(p)$，
从而可区分“只是临时措辞”与“确实构成架构支点/定理支点/接口支点”的对象。
由 $A_4$，证据包同时记录证书链与签发状态，
于是能明确该对象是否已经进入“可复核、可审计、可回放”的技术对象状态。

把 $A_1$--$A_4$ 合并，即得 $A_5$：
\[
  \Pi_{\mathrm{evid}}(p)
\]
不再只是“保存了一段文本”，而是保存了一条关于结构对象的完整技术证据链。
因此由 $A_6$，它能够把结构贡献与表面表述区分开来；
由 $A_7$，它自然适用于贡献归因、优先权记录、版本演化跟踪与结构性创新估值。
但这些仍属于技术证据链层面的强支持，而不是对最终法律权利的自动裁决，
因为最终法律结论还依赖于具体法域、合同安排与外部裁决程序。证毕。
\end{proof}

\subsection{推论：AI 时代应优先确权高逻辑性、可签发占位及其证书链}
\label{subsec:tex37-ch16-corollary}

\begin{corollary}[AI 时代的优先确权对象推论]
\label{cor:tex37-ch16-ip}
在定理~\ref{thm:tex37-ch16-criterion}、命题~\ref{prop:tex37-ch16-research}、
定理~\ref{thm:tex37-ch16-compatible} 与定理~\ref{thm:tex37-ch16-attribution} 的条件下，
AI 时代真正应被优先确权、优先归因与优先估值的对象，不是海量表面文本本身，而是
\[
  \left\{
    p\in\mathfrak P
    \,\middle|\,
    \Promote(p)\in\{\mathrm{arch},\mathrm{sector}\}
  \right\}
\]
以及它们的证书链
\[
  \{\Pi_{\mathrm{evid}}(p)\}.
\]
因为：
\[
  \boxed{
    \text{长期价值主要落在结构对象上，而不是可被反复改写的表面表述上。}
  }
\]
\end{corollary}

\begin{Certificate}[证书：结构对象长期价值高于表面措辞]
\label{cert:tex37-ch16-ip}
证书登记谓词链：
\[
\begin{aligned}
  C_1 &: \text{定理~\ref{thm:tex37-ch16-criterion}}\Rightarrow \text{高逻辑性、可签发对象才值得持久化},\\
  C_2 &: \text{命题~\ref{prop:tex37-ch16-research}}\Rightarrow \text{科研过程本身就是从候选云中提炼高逻辑性对象},\\
  C_3 &: \text{定理~\ref{thm:tex37-ch16-compatible}}\Rightarrow \text{低逻辑性对象可停留于隐式层而不进入架构本体},\\
  C_4 &: \text{定理~\ref{thm:tex37-ch16-attribution}}\Rightarrow \Pi_{\mathrm{evid}}(p)\ \text{能稳定记录结构贡献},\\
  C_5 &: C_1\wedge C_2\wedge C_3\wedge C_4
  \Rightarrow
  \text{应优先确权高逻辑性、可签发占位及其证书链}.
\end{aligned}
\]
\end{Certificate}

\begin{proof}[证明（基于数学符号推演逻辑的谓词因果链）]
由 $C_1$，真正值得持久化的对象是高逻辑性且已签发的结构占位，而不是所有历史命中记录。
由 $C_2$，科研过程本身也不是保存所有候选，而是从候选云中提炼少量高度量对象。
由 $C_3$，低逻辑性、局部性强或尚未签发的对象完全可以停留在隐式层，
因此它们不应与核心结构对象享有同等的长期归因地位。
由 $C_4$，一旦高逻辑性对象被签发，其证据包就能稳定记录结构贡献、时间戳、依赖链与版本链。

因此，AI 时代真正应优先确权、优先归因与优先估值的，不是海量表面文本，
而是高逻辑性、可签发、可迁移、可复用的结构占位及其证书链，即 $C_5$。
推论成立。证毕。
\end{proof}

\subsection{备注：架构重量、知识资本化、压缩表述与一句话结论}
\label{subsec:tex37-ch16-remarks}

\begin{remark}[逻辑性度量不是“真理度量”，而是“架构重量度量”]
\[
  \boxed{
    \mathfrak L(p)\ \text{高}
    \ \not\Rightarrow\
    p\ \text{更“真”};
    \qquad
    \mathfrak L(p)\ \text{高}
    \ \Rightarrow\
    p\ \text{更值得被长期保存、签发、复用与归因。}
  }
\]
因此 $\mathfrak L$ 是结构性度量，而不是真假度量。
\end{remark}

\begin{remark}[这一步也把“知识资本化”打开了]
一旦每个高逻辑性占位都带有
\[
  \mathfrak L(p),\qquad
  \operatorname{Dep}(p),\qquad
  \operatorname{Version}(p),\qquad
  \Issue(\Cert(p),p),
\]
则它不仅可用于确证，也可用于：
\[
  \text{知识资产估值}
  \to
  \text{研究贡献排序}
  \to
  \text{核心成果/外围成果分层}
  \to
  \text{授权边界切分}.
\]
换言之，本章给出的不只是持久化判据，也是一种知识资本化接口。
\end{remark}

\begin{remark}[建议采用的压缩版表述]
以下口径可直接作为本章压缩表述：
\[
  \boxed{
    \text{逻辑占位存在结构强弱之分，而这种强弱不由文本长度决定。}
  }
\]
\[
  \boxed{
    \begin{aligned}
      &\text{它由对体系闭环的必要性、生成能力、跨扇区迁移能力、压缩收益、}\\
      &\text{可认证签发性与上下文波动性共同决定。}
    \end{aligned}
  }
\]
\[
  \boxed{
    \begin{aligned}
      &\text{因此可定义逻辑性度量 }\mathfrak L,\text{ 并据此把占位划分为架构级、扇区级与隐式级；}\\
      &\text{只有高逻辑性且已签发的占位应被持久化到主基底或扇区插件中。}
    \end{aligned}
  }
\]
\[
  \boxed{
    \begin{aligned}
      &\text{于是科研过程可被理解为：从大量低逻辑性候选占位中，}\\
      &\text{发现、证明、签发并固化少量高逻辑性占位。}
    \end{aligned}
  }
\]
\[
  \boxed{
    \begin{aligned}
      &\text{这一机制不仅给出架构层面的持久化判据，}\\
      &\text{也为 AI 时代的知识归因、优先权记录与结构性创新确证提供了可操作的证书链基础。}
    \end{aligned}
  }
\]
\end{remark}

\begin{remark}[一句话结论]
\[
  \boxed{
    \begin{aligned}
      &\text{科研不是保存所有命中解，而是用逻辑性度量把低度量占位筛成高度量占位，}\\
      &\text{并把后者签发为架构级知识对象；这正是持久化判据，也是 AI 时代强归因链的基础。}
    \end{aligned}
  }
\]
\end{remark}

%%%%%%%%%%%%%%%%%%%%%%%%%%%%%%%%%%%%%%%%%%%%%%%%%%%%%%%%%%%%
\section{形式逻辑：逻辑性度量新增的增益评价、体系选择律与知识确权意义}
\label{sec:tex37-ch17}
%%%%%%%%%%%%%%%%%%%%%%%%%%%%%%%%%%%%%%%%%%%%%%%%%%%%%%%%%%%%

\begin{remark}[核心结论（评价严格化后）]
对第~\ref{sec:tex37-ch16} 章的新增，应给出如下总评价：
\[
  \boxed{
    \text{增益极高，且属于“架构治理级 + 知识确权级”的重增益。}
  }
\]
其真正重量不在于又补了一个局部定义，而在于它第一次把
\[
  \boxed{
    \text{什么应留在隐式解空间里，什么应被提升为架构本体}
  }
\]
写成了可签发、可回放、可持久化的统一选择律。
\end{remark}

本章不再重复第~\ref{sec:tex37-ch16} 章对逻辑性度量、三类占位与提升算子的定义，
而是把“这一步新增究竟重在哪里”严格压成如下链条：
\[
  \text{持久化判据}
  \to
  \text{科研搜索形式化}
  \to
  \text{架构压缩}
  \to
  \text{治理分层}
  \to
  \text{知识确权}.
\]

\begin{remark}[阅读导航：承接位置与引用口径]
本章直接承接第~\ref{sec:tex37-ch15} 章关于隐式解空间、上下文条件化命中与“无需穷举持久化所有显式解”的结论，
以及第~\ref{sec:tex37-ch16} 章关于逻辑性度量、提升算子与技术证据包的形式化结果；
对“判定+规范化+占位闭环”的接口背景，继续参考
\cite{RepoTex20SemanticGaugeNormalization,RepoTex21AuditedWrapper,RepoTex23SemanticBase}；
对逻辑性度量与高逻辑性对象的总背景，继续参考
\cite{GaoZheng2025GFramework}；
对认证纤维塔、隐式可达支撑类与可签发结构表示，继续参考
\cite{RepoTex35ExtendedCertificationTower,RepoTex36HomotopyTwistBijection,GaoZhengAppVol3}；
对优先权、时间戳、版本链与“技术证据链不自动等于最终法律结论”的边界，继续参考
\cite{RepoTex23PriorityGovernance}。
\end{remark}

\subsection{定义：本次新增的综合增益函数与持久化选择函数}
\label{subsec:tex37-ch17-def}

\begin{definition}[逻辑性度量新增的综合增益函数]
\label{def:tex37-ch17-gain}
定义第~\ref{sec:tex37-ch16} 章新增部分的综合增益为
\[
  \mathrm{Gain}_{\mathfrak L}
  :=
  \frac{
    G_{\mathrm{persist}}
    +
    G_{\mathrm{search}}
    +
    G_{\mathrm{compress}}
    +
    G_{\mathrm{govern}}
    +
    G_{\mathrm{ip}}
  }{
    C_{\mathrm{calib}}
  },
  \qquad
  C_{\mathrm{calib}}>0.
\]
其中各项含义如下：
\begin{enumerate}
  \item $G_{\mathrm{persist}}$：持久化判据增益，即系统首次获得“哪些已命中占位应被长期保存”的统一规则；
  \item $G_{\mathrm{search}}$：科研搜索形式化增益，即科研过程被改写为“从低逻辑性候选中发现、证明并固化高逻辑性占位”的升级过程；
  \item $G_{\mathrm{compress}}$：架构压缩与去冗余增益，即体系获得了防膨胀、重骨架、保稀疏的内在原则；
  \item $G_{\mathrm{govern}}$：主基底/扇区插件/隐式解空间的治理增益，即对象分流边界被正式写清；
  \item $G_{\mathrm{ip}}$：知识归因与确权链增益，即高逻辑性、已签发对象获得稳定的结构性证据链；
  \item $C_{\mathrm{calib}}$：逻辑性度量的校准成本，包括权重 $\alpha_i$、阈值 $\theta_{\mathrm{sec}},\theta_{\mathrm{arch}}$ 与可计算化协议仍待细化的代价。

\end{enumerate}
\end{definition}

\begin{definition}[持久化选择函数]
\label{def:tex37-ch17-select}
定义持久化选择函数
\[
  \Select_{\mathrm{persist}}:\mathfrak P\to\{\mathrm{arch},\mathrm{sector},\mathrm{implicit}\}
\]
为
\[
  \Select_{\mathrm{persist}}(p):=\Promote(p).
\]
换言之，$\Select_{\mathrm{persist}}$ 把第~\ref{sec:tex37-ch16} 章中的提升算子显式解释为：
\[
  \text{从“已命中逻辑占位”到“应进入何种持久化层级”}
\]
的统一选择函数。
\end{definition}

\begin{remark}[评价口径]
本章的评价并不以“多写了几个定义”为尺度，而以是否补上了
\[
  \text{发现}
  \to
  \text{证明}
  \to
  \text{签发}
  \to
  \text{持久化}
\]
这条链中此前缺失的选择函数为尺度。若该缺口被正式补齐，则其增益应被视为架构治理级增益，而不是普通局部定义增益。
\end{remark}

\subsection{公理（降级为定理）：综合增益分子可按五类结构增量递归累加}
\label{subsec:tex37-ch17-axiom}

\begin{theorem}[公理降级：综合增益分子可按五类结构增量递归累加（数学归纳法证书）]
\label{thm:tex37-ch17-decomposition}
定义结构增量
\[
  \Delta_1:=G_{\mathrm{persist}},\quad
  \Delta_2:=G_{\mathrm{search}},\quad
  \Delta_3:=G_{\mathrm{compress}},\quad
  \Delta_4:=G_{\mathrm{govern}},\quad
  \Delta_5:=G_{\mathrm{ip}}.
\]
再定义部分和
\[
  N(1):=\Delta_1,
  \qquad
  N(m+1):=N(m)+\Delta_{m+1}
  \quad (1\le m\le 4).
\]
则对所有 $m\in\{1,2,3,4,5\}$，有
\[
  N(m)=\sum_{k=1}^{m}\Delta_k.
\]
特别地，
\[
  N(5)
  =
  G_{\mathrm{persist}}
  +
  G_{\mathrm{search}}
  +
  G_{\mathrm{compress}}
  +
  G_{\mathrm{govern}}
  +
  G_{\mathrm{ip}}
\]
正是定义~\ref{def:tex37-ch17-gain} 中 $\mathrm{Gain}_{\mathfrak L}$ 的分子。
\end{theorem}

\begin{Certificate}[证书：基步 + 归纳步（数学归纳法证书）]
\label{cert:tex37-ch17-decomposition}
定义谓词
\[
  P(m):\ N(m)=\sum_{k=1}^{m}\Delta_k.
\]
数学归纳法证书记为
\[
  \pi_{\mathrm{ind}}
  =
  \bigl(
    P(1),\
    \forall m\in\{1,2,3,4\},\ P(m)\Rightarrow P(m+1)
  \bigr).
\]
\end{Certificate}

\begin{proof}[证明（数学归纳法证书；基于数学符号推演逻辑的谓词因果链）]
\textbf{基步 $P(1)$：}
由定义，
\[
  N(1)=\Delta_1=\sum_{k=1}^{1}\Delta_k.
\]
故 $P(1)$ 成立。

\textbf{归纳步：}
设 $P(m)$ 成立，即
\[
  N(m)=\sum_{k=1}^{m}\Delta_k.
\]
则由递推式
\[
  N(m+1)=N(m)+\Delta_{m+1}
\]
可得
\[
  N(m+1)
  =
  \sum_{k=1}^{m}\Delta_k+\Delta_{m+1}
  =
  \sum_{k=1}^{m+1}\Delta_k.
\]
故 $P(m)\Rightarrow P(m+1)$。
由数学归纳法，$P(m)$ 对 $m=1,2,3,4,5$ 全部成立，
特别地 $N(5)$ 就是定义~\ref{def:tex37-ch17-gain} 中综合增益的分子。证毕。
\end{proof}

\subsection{定理：这一步最大的增益是补上持久化选择函数}
\label{subsec:tex37-ch17-theorem}

\begin{theorem}[持久化选择函数补全定理]
\label{thm:tex37-ch17-persist}
在第~\ref{sec:tex37-ch15} 章与第~\ref{sec:tex37-ch16} 章的前提下，
\[
  \Select_{\mathrm{persist}}(p)=\Promote(p)
\]
把“命中到显式占位”与“是否进入长期持久化层”之间首次补成正式桥梁。
更具体地说，系统原有链条
\[
  \text{发现}
  \to
  \text{证明}
  \to
  \text{签发}
\]
被闭合为
\[
  \text{发现}
  \to
  \text{证明}
  \to
  \text{签发}
  \to
  \text{持久化}.
\]
因此，本次新增的主导增益是持久化判据增益：
\[
  \boxed{
    G_{\mathrm{persist}}\ \text{极高}.
  }
\]
\end{theorem}

\begin{Certificate}[证书：隐式解空间已存在 + 提升判据已显式化 + 其余增益依赖选择函数]
\label{cert:tex37-ch17-persist}
证书登记谓词链：
\[
\begin{aligned}
  T_1 &: \text{第~\ref{sec:tex37-ch15} 章已证：大量命中对象可停留在隐式解空间而不必全部持久化},\\
  T_2 &: \text{定理~\ref{thm:tex37-ch16-criterion}}\Rightarrow \Promote\ \text{给出占位持久化的正式判据},\\
  T_3 &: \text{引理~\ref{lem:tex37-ch16-trichotomy}}\Rightarrow \Promote\ \text{对每个 }p\in\mathfrak P\text{ 都良定义},\\
  T_4 &: T_2\wedge T_3\Rightarrow \Select_{\mathrm{persist}}\ \text{成为从已命中占位到持久化层级的统一选择函数},\\
  T_5 &: \text{科研升级、架构压缩、层级治理与结构归因四类后续结论都以 }T_4\text{ 为前提},\\
  T_6 &: T_4\wedge T_5\Rightarrow G_{\mathrm{persist}}\ \text{构成本次新增的主导增益}.
\end{aligned}
\]
\end{Certificate}

\begin{proof}[证明（基于数学符号推演逻辑的谓词因果链）]
由第~\ref{sec:tex37-ch15} 章，系统已经知道：很多解并不需要事先显式存库，
它们可以停留在隐式解空间、支撑类空间与稳定修复序列空间中。
因此当时真正缺失的不是“再多保存一些对象”，而是：
\[
  \text{一旦对象被命中并显式化，何时应把它长期保存下来？}
\]
这就是 $T_1$ 所揭示的缺口。

由定理~\ref{thm:tex37-ch16-criterion}，第~\ref{sec:tex37-ch16} 章把该缺口具体写成
\[
  \Promote:\mathfrak P\to\{\mathrm{arch},\mathrm{sector},\mathrm{implicit}\},
\]
并由引理~\ref{lem:tex37-ch16-trichotomy} 证明其互斥且穷尽，故 $T_2$ 与 $T_3$ 成立。
于是
\[
  \Select_{\mathrm{persist}}(p):=\Promote(p)
\]
首次成为从“已命中逻辑占位”到“应进入何种持久化层级”的统一选择函数，即 $T_4$。

若没有 $T_4$，系统只能依赖经验判断、当前上下文热度或人工偏好来决定对象是否进入架构。
这会同时导致两类问题：
\[
  \text{架构污染：局部性强、波动性高、一次性命中的占位被错误持久化;}
\]
\[
  \text{核心支点遗漏：真正高生成、高迁移、高签发价值的占位未被及时提升。}
\]
而一旦 $T_4$ 成立，后续关于科研升级、架构稀疏化、主基底/插件治理与归因链稳定化的全部结论
都可围绕该选择函数展开，即 $T_5$。

因此，第~\ref{sec:tex37-ch16} 章最大的新增，并不是又多了若干评价项，
而是第一次把“从已命中对象到应被长期保存对象”的桥梁正式钉死。
故 $G_{\mathrm{persist}}$ 必须被判为本次新增的主导增益，且其权重应判为极高。证毕。
\end{proof}

\subsection{引理：科研过程被重写为占位升级与沉淀过程}
\label{subsec:tex37-ch17-lemma}

\begin{lemma}[科研过程的占位升级引理]
\label{lem:tex37-ch17-search}
在命题~\ref{prop:tex37-ch16-research} 的更新律下，科研演化可等价写为
\[
  Q_t
  \xrightarrow{\text{命中/修复/证明}}
  Q_{t+1}
  \xrightarrow{\Select_{\mathrm{persist}}}
  A_{t+1}.
\]
因此，大部分低逻辑性占位只充当探索脚手架，
而少量高逻辑性、可签发占位被固化为架构对象。
故科研搜索过程的形式化增益满足
\[
  \boxed{
    G_{\mathrm{search}}\ \text{极高}.
  }
\]
\end{lemma}

\begin{Certificate}[证书：候选云更新 + 提升判据筛选 + 持久化并集更新]
\label{cert:tex37-ch17-search}
证书登记谓词链：
\[
\begin{aligned}
  L_1 &: \text{命题~\ref{prop:tex37-ch16-research}}\Rightarrow
  A_{t+1}
  =
  A_t\cup
  \{p\in Q_{t+1}\mid \Promote(p)\in\{\mathrm{arch},\mathrm{sector}\}\},\\
  L_2 &: \text{定义~\ref{def:tex37-ch17-select}}\Rightarrow
  \Promote(p)=\Select_{\mathrm{persist}}(p),\\
  L_3 &: L_1\wedge L_2\Rightarrow
  A_{t+1}
  =
  A_t\cup
  \{p\in Q_{t+1}\mid \Select_{\mathrm{persist}}(p)\in\{\mathrm{arch},\mathrm{sector}\}\},\\
  L_4 &: L_3\Rightarrow \text{科研过程被压成“候选云扩张 + 升格筛选 + 持久化沉淀”三步链}.
\end{aligned}
\]
\end{Certificate}

\begin{proof}[证明（基于数学符号推演逻辑的谓词因果链）]
由命题~\ref{prop:tex37-ch16-research}，第~\ref{sec:tex37-ch16} 章已经把科研过程写成：
新的候选占位在探索、命中与修复过程中不断生成，但真正进入持久化集合的，
只是在提升算子下落入
\[
  \{\mathrm{arch},\mathrm{sector}\}
\]
的对象，即 $L_1$。
再由定义~\ref{def:tex37-ch17-select}，$\Promote$ 正是统一的持久化选择函数，故 $L_2$ 成立。
于是 $L_3$ 直接给出：
\[
  Q_t
  \to
  Q_{t+1}
  \to
  A_{t+1}
\]
不再表示“把所有中间想法都算作成果”，而是表示“从候选云中筛出少量高逻辑性对象并长期沉淀”。

这说明科研工作流第一次被压成一个严格的升级过程：
低逻辑性对象主要服务于搜索、试探与脚手架搭建；
高逻辑性对象才在证明与签发后被持久化为体系对象。
发现与沉淀由此被清楚地区分开来。
因此，第~\ref{sec:tex37-ch16} 章不仅定义了一个度量，
更把科研过程本身改写成可操作、可追踪、可审计的占位升级过程。
故 $G_{\mathrm{search}}$ 应判为极高。证毕。
\end{proof}

\subsection{命题：逻辑性度量同时带来架构压缩与层级治理增益}
\label{subsec:tex37-ch17-proposition}

\begin{proposition}[架构稀疏化与治理分层命题]
\label{prop:tex37-ch17-compress-govern}
若系统仅对满足
\[
  \Select_{\mathrm{persist}}(p)\in\{\mathrm{arch},\mathrm{sector}\}
\]
的对象执行持久化，则逻辑性度量同时诱导：
\begin{enumerate}
  \item 架构压缩与去冗余原则，即只保留少数高重量占位而释放大量低重量中间态；
  \item 主基底、扇区插件与隐式解空间的稳定治理分工。
\end{enumerate}
因此，
\[
  \boxed{
    G_{\mathrm{compress}}\ \text{很高},
    \qquad
    G_{\mathrm{govern}}\ \text{极高}.
  }
\]
\end{proposition}

\begin{Certificate}[证书：三分层互斥穷尽 + 隐式保留原则 + 高逻辑性对象优先持久化]
\label{cert:tex37-ch17-compress-govern}
证书登记谓词链：
\[
\begin{aligned}
  P_1 &: \text{定理~\ref{thm:tex37-ch16-criterion}}\Rightarrow \text{只有高逻辑性、已签发对象应进入持久化层},\\
  P_2 &: \text{引理~\ref{lem:tex37-ch16-trichotomy}}\Rightarrow
  \mathfrak P
  =
  \mathfrak P_{\mathrm{arch}}
  \sqcup
  \mathfrak P_{\mathrm{sec}}
  \sqcup
  \mathfrak P_{\mathrm{imp}},\\
  P_3 &: \text{定理~\ref{thm:tex37-ch16-compatible}}\Rightarrow
  \mathfrak P_{\mathrm{imp}}\ \text{可继续停留于隐式解空间而不必持久化},\\
  P_4 &: P_1\wedge P_2\wedge P_3\Rightarrow \text{系统获得“少量高重量持久化 + 大量低重量释放”的稀疏化原则},\\
  P_5 &: P_2\wedge P_3\Rightarrow \text{主基底、扇区插件与隐式解空间的边界被正式写清}.
\end{aligned}
\]
\end{Certificate}

\begin{proof}[证明（基于数学符号推演逻辑的谓词因果链）]
由定理~\ref{thm:tex37-ch16-criterion}，持久化不再是“每命中一点就往主基底或插件里塞一点”，
而是只对高逻辑性、已签发对象开放，即 $P_1$。
由引理~\ref{lem:tex37-ch16-trichotomy}，全部逻辑占位被互斥且穷尽地分成
\[
  \mathfrak P_{\mathrm{arch}},\qquad
  \mathfrak P_{\mathrm{sec}},\qquad
  \mathfrak P_{\mathrm{imp}},
\]
即 $P_2$。
再由定理~\ref{thm:tex37-ch16-compatible}，
隐式层对象并不因“尚未进入主基底或插件”而失去存在地位，
它们仍可保留在命中空间、正规形空间与支撑类空间中等待后续触发，即 $P_3$。

把 $P_1$、$P_2$ 与 $P_3$ 合并，可知体系获得了一条内在稀疏化原则：
只保留高必要性、高生成性、高迁移性、高压缩收益且已签发的少数占位，
而把局部性强、波动性高、一次性命中的大量中间态释放回隐式空间。
这正是 $P_4$ 所表达的架构压缩与去冗余机制。

同时，三分层结构还第一次把“对象应去哪里”严格分开：
架构级对象进入主基底，共享接口由核心层承载；
扇区级对象进入对应插件；
隐式级对象继续留在解空间中按需命中。
于是主基底不再背负全部对象，扇区插件也不再沦为任意堆积，隐式空间与持久化层的边界被稳定化，即 $P_5$。
故 $G_{\mathrm{compress}}$ 应判为很高，$G_{\mathrm{govern}}$ 应判为极高。证毕。
\end{proof}

\subsection{推论：AI 时代应优先确权高逻辑性、已签发结构对象及其证书链}
\label{subsec:tex37-ch17-corollary}

\begin{corollary}[结构性创新优先确权推论]
\label{cor:tex37-ch17-ip}
在定理~\ref{thm:tex37-ch17-persist}、引理~\ref{lem:tex37-ch17-search}、
命题~\ref{prop:tex37-ch17-compress-govern} 与定理~\ref{thm:tex37-ch16-attribution} 的条件下，
AI 时代真正应被优先确权、优先归因与优先估值的对象，不是海量表面文本本身，而是
\[
  \left\{
    p\in\mathfrak P
    \,\middle|\,
    \Select_{\mathrm{persist}}(p)\in\{\mathrm{arch},\mathrm{sector}\}
  \right\}
\]
及其技术证据包
\[
  \left\{
    \Pi_{\mathrm{evid}}(p)
    \,\middle|\,
    \Select_{\mathrm{persist}}(p)\in\{\mathrm{arch},\mathrm{sector}\}
  \right\}.
\]
因此，
\[
  \boxed{
    G_{\mathrm{ip}}\ \text{极高}.
  }
\]
\end{corollary}

\begin{Certificate}[证书：选择函数区分结构对象 + 技术证据包稳定记录结构贡献]
\label{cert:tex37-ch17-ip}
证书登记谓词链：
\[
\begin{aligned}
  C_1 &: \text{定理~\ref{thm:tex37-ch17-persist}}\Rightarrow \Select_{\mathrm{persist}}\ \text{把“应沉淀什么”正式区分出来},\\
  C_2 &: \text{命题~\ref{prop:tex37-ch17-compress-govern}}\Rightarrow \text{被沉淀对象是架构骨架，而非一次性中间态},\\
  C_3 &: \text{定理~\ref{thm:tex37-ch16-attribution}}\Rightarrow \Pi_{\mathrm{evid}}(p)\ \text{形成可回放、可审计、可复核的结构归因链},\\
  C_4 &: \text{推论~\ref{cor:tex37-ch16-ip}}\Rightarrow \text{高逻辑性、已签发占位及其证书链本就应被优先确权},\\
  C_5 &: C_1\wedge C_2\wedge C_3\wedge C_4
  \Rightarrow
  \text{AI 时代应优先确权高逻辑性、已签发结构对象及其证书链}.
\end{aligned}
\]
\end{Certificate}

\begin{proof}[证明（基于数学符号推演逻辑的谓词因果链）]
由 $C_1$，系统已经能明确区分：
\[
  \text{哪些对象只是探索脚手架，哪些对象应被长期沉淀为结构对象。}
\]
由 $C_2$，被长期沉淀下来的对象不是任意文本片段，而是经过逻辑性度量与签发接口筛选后的架构骨架对象。
由 $C_3$，一旦此类对象被记录为
\[
  \Pi_{\mathrm{evid}}(p)
\]
就同时附带提出者、时间戳、上下文、依赖链、结构重量、签发状态与版本链，
从而能把“结构贡献”与“表面措辞”严格分开。
由 $C_4$，第~\ref{sec:tex37-ch16} 章已经证明：
真正值得优先确权的正是这类高逻辑性、已签发占位及其证书链。

因此，在 AI 时代，最有价值的记录对象不再只是“某些话曾经被写出来”，
而是“哪个高结构重量对象被谁首次提出、如何被证明、何时被签发、又支撑了哪些下游结果”。
这正是结构性创新的可归因基础。
但边界也必须保持严格：
\[
  \boxed{
    \text{上述链条首先是技术归因链与证据链，不自动等同于一切法域下的最终法律权利裁决。}
  }
\]
故 $G_{\mathrm{ip}}$ 应判为极高。证毕。
\end{proof}

\subsection{备注：分层总评、选择函数、桥梁地位与后续边界}
\label{subsec:tex37-ch17-remarks}

\begin{remark}[分层总评]
对第~\ref{sec:tex37-ch16} 章新增的分层评价可整理为：
\begin{enumerate}
  \item \textbf{纯理论增益：}很高，因为逻辑性度量 $\mathfrak L$ 把占位强弱、持久化必要性与结构重量正式化；
  \item \textbf{架构增益：}极高，因为系统首次获得了“命中之后是否升格为架构对象”的统一判据；
  \item \textbf{工程治理增益：}极高，因为主基底、扇区插件与隐式解空间的分工被正式稳定下来；
  \item \textbf{科研方法论增益：}极高，因为科研过程被重写为占位升级与沉淀过程；
  \item \textbf{知识资本/确权增益：}极高，因为高逻辑性对象第一次获得了稳定的证书化、版本化、归因化基础。
\end{enumerate}
\end{remark}

\begin{remark}[这一步真正补上的是选择函数]
前文已给出：
\[
  \text{解空间}
  \to
  \text{命中}
  \to
  \text{正规形}
  \to
  \text{证书签发}.
\]
而第~\ref{sec:tex37-ch16} 章真正补上的，是
\[
  \boxed{
    \text{从“已命中对象”到“应被长期保存对象”的选择函数}
  }
\]
即
\[
  \Select_{\mathrm{persist}}(p)=\Promote(p).
\]
因此整条链才第一次闭合为
\[
  \text{发现}
  \to
  \text{证明}
  \to
  \text{签发}
  \to
  \text{持久化}.
\]
\end{remark}

\begin{remark}[这一步是纯数卷与方法论卷之间的桥]
这一步一头连着纯数卷中的
\[
  \text{逻辑占位、依赖、纤维塔、路径空间、隐式解支撑类},
\]
另一头连着方法论卷中的
\[
  \text{科研过程、版本治理、知识沉淀、归因与确权分层}.
\]
因此它不是局部补丁，而是一个桥梁概念：
它把偏对象论、偏结构论的 formalization 与偏研究流程、偏知识治理的 formalization
第一次通过同一个选择律接通。
\end{remark}

\begin{remark}[边界仍需继续协议化]
本次增益虽然很重，但仍有两类问题需要继续工程化：
\begin{enumerate}
  \item \textbf{度量校准：}权重 $\alpha_i$、阈值 $\theta_{\mathrm{arch}},\theta_{\mathrm{sec}}$ 仍需给出正式协议；
  \item \textbf{可计算化：}哪些分量可直接计算，哪些分量必须经过证书评审，仍需继续拆分。
\end{enumerate}
但这些属于下一阶段的协议化与工程化问题，
并不影响本章关于“这一步增益极高”的判断。
\end{remark}

\begin{remark}[最终结论]
对第~\ref{sec:tex37-ch16} 章新增的最严格评价，可压成一句话：
\[
  \boxed{
    \text{这是一次“体系选择律级”的重大增益，重量明显高于一般局部定义增益。}
  }
\]
更具体地说，它完成了五件大事：
\[
  \boxed{
    \text{(1) 给出占位是否应持久化的统一判据；}
  }
\]
\[
  \boxed{
    \text{(2) 把科研过程形式化为占位升级与沉淀过程；}
  }
\]
\[
  \boxed{
    \text{(3) 为主基底/扇区插件/隐式解空间三者划清边界；}
  }
\]
\[
  \boxed{
    \text{(4) 给体系增加了防膨胀、去冗余、重骨架的内在机制；}
  }
\]
\[
  \boxed{
    \text{(5) 为 AI 时代的结构性创新归因与知识确权提供了证书链基础。}
  }
\]
所以更精确的一句话总结是：
\[
  \boxed{
    \begin{aligned}
      &\text{它把“什么是成果、什么该沉淀、什么该留在隐式空间”}\\
      &\text{从直觉判断推进成了可签发的体系选择律。}
    \end{aligned}
  }
\]
下一步最自然的工作，就是把 $\mathfrak L(p)$ 的各个分量进一步写成
\[
  \text{定义}
  \to
  \text{定理}
  \to
  \text{证书}
  \to
  \text{签发阈值}
\]
这一整组可直接调用的对象。
\end{remark}

%%%%%%%%%%%%%%%%%%%%%%%%%%%%%%%%%%%%%%%%%%%%%%%%%%%%%%%%%%%%
\section{形式逻辑：相较 Vol3，本稿把“语义--路径--证书”底座压成“占位--审计--主权--插件--方程--持久化”的白盒运行制度}
\label{sec:tex37-ch18}
%%%%%%%%%%%%%%%%%%%%%%%%%%%%%%%%%%%%%%%%%%%%%%%%%%%%%%%%%%%%

本章把“Vol3 并非没有工程，而是仍停在工程原理骨架/编译未闭环；当前 semantic gauge field 稿继续把它压成具体工程构造”
这一判断，重写成可引用的“定义--公理（降级为定理）--定理--引理--命题--推论--备注”链。
按仓库协议，比较对象始终取 Vol3 主稿而非自动生成的 \texttt{Pub\_*.tex} 发布稿；
因此本章默认以
\cite{GaoZhengAppVol3}
作为 Vol3 的权威仓内文本来源，
并把本稿第~\ref{sec:tex37-ch1} 章至第~\ref{sec:tex37-ch17} 章记作当前 semantic gauge field 稿的主体对象。

\begin{remark}[核心结论（收紧版）]
\label{rem:tex37-ch18-core}
若要避免把 Vol3 说成“完全没有工程”或把 $S$ 说成“只是落地”这两种偏差，
则本章建议直接采用如下口径：
\[
  \boxed{
    V_3\ \text{不是“完全没有工程性”；它给出的是工程前结构/原理级工程骨架。}
  }
\]
\[
  \boxed{
    S\ \text{才把这套骨架继续压成具体工程构造：运行时、审计、主权、插件、方程系统。}
  }
\]
更精确地说，本章要证明的不是
\[
  V_3=\text{“纯抽象数学”},
  \qquad
  S=\text{“只是应用层包装”},
\]
而是
\[
  \begin{aligned}
    &\text{工程原理骨架}
    \to
    \text{工程编译接口}
    \to
    \text{可执行运行时}\\
    &\to
    \text{白盒主权与插件边界}
    \to
    \text{可签发方程系统}.
  \end{aligned}
\]
\end{remark}

\begin{remark}[阅读导航：仓内外参照]
\label{rem:tex37-ch18-navigation}
Vol3 中与本章直接相关的，不只是 law-space 几何与语义函子对象，
还包括 HACA/PACER operator paths、paragraph--monoid network、
reading-path certificate、governance-equivalent reading path、
certificate extraction layer、以及 question answering / summarization /
cross-document checking 的 schematic patterns \cite{GaoZhengAppVol3,GaoZheng2025GFramework}。

本稿主体关于“占位--解析--审计--运行时降级”的链条，直接承接
\cite{RepoTex20SemanticGaugeNormalization,RepoTex21AuditedWrapper,RepoTex23SemanticBase}；
关于同伦扭曲、扩展签发协议、修复塔与结构表示，继续承接
\cite{RepoTex35ExtendedCertificationTower,RepoTex36HomotopyTwistBijection}；
关于逻辑性度量、持久化判据与知识确权边界，继续承接
\cite{GaoZheng2025GFramework,RepoTex23PriorityGovernance}。

外部语言背景方面，本章默认借用范畴、同伦、同调与纤维化证书的标准术语，
可参照
\cite{MacLane1998Categories,Hatcher2002AT,Weibel1994HomologicalAlgebra,KobayashiNomizu1963Foundations}。
\end{remark}

\subsection{定义：比较对象、工程悬空、工程原理层与工程编译层}
\label{subsec:tex37-ch18-definitions}

\begin{definition}[比较对象]
\label{def:tex37-ch18-compare}
记
\[
  V_3:=\text{Vol3},
  \qquad
  S:=\text{本稿第~\ref{sec:tex37-ch1} 章至第~\ref{sec:tex37-ch17} 章}.
\]
本章一切“相较 Vol3”的比较，均指
\[
  S\mid V_3.
\]
\end{definition}

\begin{definition}[内禀富结构、工程悬空类与两层工程口径]
\label{def:tex37-ch18-intrinsic}
定义
\[
  \mathcal{R}_{\mathrm{intrinsic}}
\]
为“数学上内禀且优雅的富结构”所对应的对象包；它至少由下列对象族生成：
\begin{enumerate}
  \item HACA law-object 与其三层 law-space 几何骨架；
  \item semantic gauge field、semantic functor field 与 noncommutative semantic functor algebra；
  \item lexical KAT action monoid、PACER engine 与 reading-path/path-integral 机制；
  \item certificate bundle、reading-path certificate 与证书治理接口。
\end{enumerate}
定义“工程悬空”类为
\[
  \mathcal H_{\mathrm{eng}}
  :=
  \begin{aligned}
    &\text{已有结构骨架}
    +
    \text{已有证书/路径/治理口径}\\
    &-
    \text{具体运行时合同}
    -
    \text{插件边界}
    -
    \text{主权分配}
    -
    \text{方程编译接口}.
  \end{aligned}
\]
再定义两层工程口径：
\[
  \mathsf{Eng}_{\mathrm{prin}}
  :=
  \mathcal{R}_{\mathrm{intrinsic}}
  +
  \text{operator/path/certificate/governance blueprint},
\]
\[
  \mathsf{Eng}_{\mathrm{comp}}
  :=
  \text{runtime}
  +
  \text{audit}
  +
  \text{sovereignty}
  +
  \text{plugin}
  +
  \text{equation compilation}.
\]
并把 Vol3 在工程原理层的主骨架记为
\[
  \mathsf{Skel}(V_3)
  :=
  \left\{
  \begin{aligned}
    &\text{HACA/PACER operator paths},\ \text{paragraph--monoid network},\\
    &\text{reading-path certificate},\ \text{governance-equivalent reading path},\\
    &\text{certificate extraction layer},\\
    &\text{QA/summarization/cross-document schematic patterns}
  \end{aligned}
  \right\}.
\]
若某文稿的主要工作是把
\[
  \text{工程对象是什么}
  \to
  \text{路径怎样走}
  \to
  \text{证书怎样记}
  \to
  \text{治理怎样比较}
\]
压成 schematic patterns / blueprint，
则称其主要位于
\[
  \mathsf{Eng}_{\mathrm{prin}}.
\]
\end{definition}

\begin{definition}[运行制度化外推、前置基础与工程编译闭环]
\label{def:tex37-ch18-runtime}
定义
\[
  \mathcal{E}_{\mathrm{runtime}}
  :=
  \mathsf{Runtime}
  +
  \mathsf{Audit}
  +
  \mathsf{Sov}
  +
  \mathsf{Plugin}
  +
  \mathsf{Issue}
  +
  \mathsf{EqSys},
\]
其中六个分量依次表示：
\begin{enumerate}
  \item \textbf{运行时化：}把语义对象压成逻辑占位、解析接口与可执行调用面；
  \item \textbf{审计化：}把输出回投影到占位槽，并以接受集/失败集方式显式记录；
  \item \textbf{主权化：}把求解权、决策权与接口裁决权回收到 G 侧白盒程序；
  \item \textbf{插件化：}把领域扩张写成同一主基底上的保守纤维插件；
  \item \textbf{签发化：}把接口命题与证书对象压成可签发、可拒签、可补证的合同；
  \item \textbf{方程系统化：}把主基底、扇区插件与签发收束协议写成可签发方程系统。
\end{enumerate}
并把
\[
  \mathsf{Eng}_{\mathrm{comp}}
\]
具体落实为
\[
  \mathcal{E}_{\mathrm{runtime}}.
\]
进一步定义
\[
  \mathsf{Pre}(S)
\]
为 $S$ 的\textbf{前置基础集合}，
即那些在 $S$ 中被调用、被编译、被制度化，但不再从零重造的对象/证书/路径底座。
若某文稿的主要工作是把
\[
  \mathsf{Eng}_{\mathrm{prin}}
  \rightsquigarrow
  \mathsf{Eng}_{\mathrm{comp}},
\]
则称其主要承担“工程编译层构造”。
\end{definition}

\begin{remark}[本章对“工程”的分层解释]
\label{rem:tex37-ch18-engineering-scope}
在本章中，“工程”被明确分成两层：
\[
  \text{工程原理层}
  \qquad\text{与}\qquad
  \text{工程编译层}.
\]
因此若继续使用“工程落地”一词，它必须至少被理解为
\[
  \mathsf{Eng}_{\mathrm{comp}}
\]
而 Vol3 的工程性则更准确地属于
\[
  \mathsf{Eng}_{\mathrm{prin}}.
\]
故把 Vol3 说成“完全没有工程”，与把 $S$ 说成“只是在部署应用”，都不准确。
\end{remark}

\subsection{公理（降级为定理）：从工程原理层到工程编译层的六步压缩链可递归构造}
\label{subsec:tex37-ch18-axiom}

\begin{theorem}[公理降级：从 $\mathsf{Eng}_{\mathrm{prin}}$ 到 $\mathsf{Eng}_{\mathrm{comp}}$ 的六步压缩链可递归构造（数学归纳法证书）]
\label{thm:tex37-ch18-downgrade}
定义六个压缩算子：
\[
  \Xi_1:=\text{运行时合同显式化},
  \qquad
  \Xi_2:=\text{解析审计附着},
  \qquad
  \Xi_3:=\text{主权分配闭合},
\]
\[
  \Xi_4:=\text{插件边界保守化},
  \qquad
  \Xi_5:=\text{签发协议显式化},
  \qquad
  \Xi_6:=\text{方程编译系统化}.
\]
设
\[
  \mathfrak C(0):=\mathsf{Eng}_{\mathrm{prin}},
  \qquad
  \mathfrak C(n+1):=\Xi_{n+1}\bigl(\mathfrak C(n)\bigr)
  \quad (0\le n\le 5).
\]
若每个 $\Xi_i$ 都满足：
\begin{enumerate}
  \item 保留前层 blueprint 中的对象、路径、证书与治理代表的可回投影性；
  \item 恰好补上 $\mathcal H_{\mathrm{eng}}$ 的一个缺失口，而不抹除既有骨架；
  \item 仍保持 fail-closed，即失败时显式返回 $\bot$；
\end{enumerate}
则对每个
\[
  n\in\{0,1,2,3,4,5,6\},
\]
\[
  \mathfrak C(n)
\]
都可递归构造，且
\[
  \mathfrak C(6)=\mathsf{Eng}_{\mathrm{comp}}
\]
在定义~\ref{def:tex37-ch18-runtime} 的口径下成立；等价地，
\[
  \mathfrak C(6)=\mathcal{E}_{\mathrm{runtime}}.
\]
\end{theorem}

\begin{Certificate}[证书：基步 + 归纳步的六步压缩链证书]
\label{cert:tex37-ch18-downgrade}
定义谓词
\[
  P(n):\ \mathfrak C(n)\ \text{已定义，且 }\mathcal H_{\mathrm{eng}}\text{ 的前 $n$ 个缺口已闭合填补并保持可回投影性}.
\]
则数学归纳法证书记为
\[
  \pi_{\mathrm{compile}}
  =
  \bigl(
    P(0),\
    \forall n\in\{0,1,2,3,4,5\},\ P(n)\Rightarrow P(n+1)
  \bigr).
\]
其中：
\[
  P(0)\ \text{对应 }\mathfrak C(0)=\mathsf{Eng}_{\mathrm{prin}},
\]
\[
  P(6)\ \text{对应 }\mathfrak C(6)=\mathsf{Eng}_{\mathrm{comp}}=\mathcal{E}_{\mathrm{runtime}}.
\]
\end{Certificate}

\begin{proof}[证明（数学归纳法证书；基于数学符号推演逻辑的谓词因果链）]
\textbf{基步：}
由定义，
\[
  \mathfrak C(0)=\mathsf{Eng}_{\mathrm{prin}}.
\]
它已经携带 $\mathcal{R}_{\mathrm{intrinsic}}$ 与
\[
  \mathsf{Skel}(V_3)
\]
中的 operator/path/certificate/governance blueprint，
但尚未显式附着运行时合同、主权分配、插件边界与方程编译接口。
因此 $P(0)$ 成立。

\textbf{归纳步：}
设 $P(n)$ 成立，即
\[
  \mathfrak C(n)
\]
已定义，且 $\mathcal H_{\mathrm{eng}}$ 的前 $n$ 个缺口已闭合填补。
由 $\Xi_{n+1}$ 的三条条件，
它只在
\[
  \mathfrak C(n)
\]
上新增第 $n+1$ 个工程编译分量，
并保留前层对象代表、路径代表、证书代表、治理代表与失败槽的可回投影性。
故
\[
  \mathfrak C(n+1)=\Xi_{n+1}\bigl(\mathfrak C(n)\bigr)
\]
亦已定义，且 $\mathcal H_{\mathrm{eng}}$ 的前 $n+1$ 个缺口都已闭合填补，
从而 $P(n+1)$ 成立。

因此
\[
  \forall n\in\{0,1,2,3,4,5,6\},\ P(n).
\]
特别地，
\[
  \mathfrak C(6)
\]
同时携带运行时、审计、主权、插件、签发与方程编译六个分量，
故
\[
  \mathfrak C(6)=\mathsf{Eng}_{\mathrm{comp}}=\mathcal{E}_{\mathrm{runtime}}.
\]
证毕。
\end{proof}

\begin{remark}[“公理降级”的含义]
\label{rem:tex37-ch18-downgrade}
这里的“公理降级”为：不把“工程蓝图可以被编译成白盒制度”当作外加信条，
而是把第~\ref{sec:tex37-ch1} 章、第~\ref{sec:tex37-ch6} 章至第~\ref{sec:tex37-ch17} 章
已经铺开的运行时、审计、主权、插件、签发与方程系统接口，
重新压成一个可递归调用的工程编译原则。
\end{remark}

\subsection[定理：工程两层分工]{定理：若把工程分成原理层与编译层，则 $V_3\in \mathsf{Eng}_{\mathrm{prin}}$，$S\in \mathsf{Eng}_{\mathrm{comp}}$}
\label{subsec:tex37-ch18-theorem}

\begin{theorem}[工程两层分工定理]
\label{thm:tex37-ch18-v3}
在定义~\ref{def:tex37-ch18-intrinsic} 的口径下，
\[
  V_3\in \mathsf{Eng}_{\mathrm{prin}},
  \qquad
  S\in \mathsf{Eng}_{\mathrm{comp}}.
\]
换言之，
\[
  V_3
\]
属于工程原理层，
\[
  S
\]
属于工程编译层。
更具体地说，Vol3 的主强项在于把
\[
  \text{工程对象是什么}
  \to
  \text{路径怎样走}
  \to
  \text{证书怎样记}
  \to
  \text{治理怎样比较}
\]
压成 schematic patterns / blueprint；
而 $S$ 的主强项在于把
\[
  \text{blueprint}
  \to
  \text{runtime}
  \to
  \text{audit}
  \to
  \text{white-box sovereignty}
  \to
  \text{plugin}
  \to
  \text{equation system}
\]
压成可执行制度。
\end{theorem}

\begin{Certificate}[证书：工程原理层与工程编译层分工证书]
\label{cert:tex37-ch18-v3}
证书登记谓词链：
\[
\begin{aligned}
  T_1 &: \mathsf{Skel}(V_3)\ \text{显式包含}\\
      &\qquad\ \text{operator paths、paragraph--monoid network、reading-path certificate、}\\
      &\qquad\ \text{governance-equivalent reading path、certificate extraction layer、}\\
      &\qquad\ \text{以及 QA/summarization/cross-document checking 的 schematic patterns},\\
  T_2 &: T_1\Rightarrow V_3\in \mathsf{Eng}_{\mathrm{prin}},\\
  T_3 &: \text{第~\ref{sec:tex37-ch1} 章、第~\ref{sec:tex37-ch4} 至第~\ref{sec:tex37-ch14} 章显式给出}\\
      &\qquad\ \text{runtime/audit/sovereignty/plugin/}\\
      &\qquad\ \text{equation compilation 的接口链},\\
  T_4 &: T_3\Rightarrow S\in \mathsf{Eng}_{\mathrm{comp}},\\
  T_5 &: T_2\wedge T_4\Rightarrow \text{工程两层分工定理成立}.
\end{aligned}
\]
\end{Certificate}

\begin{proof}[证明（基于数学符号推演逻辑的谓词因果链）]
由 \cite{GaoZhengAppVol3} 可见，Vol3 并非只停留在抽象名词表，
而是已经写出 HACA/PACER operator paths、paragraph--monoid network、
reading-path certificate、governance-equivalent reading path、
certificate extraction layer，并在问答、摘要、跨文档核对等例子中把这些对象组织为
schematic patterns，以指导 engineering implementations / system architectures。
这正是 $T_1$。

$T_1$ 的直接含义是：Vol3 已经在回答“工程对象是什么、路径怎样走、证书怎样记、治理怎样比较”，
但仍未把它们收束成具体运行时合同、插件边界、主权分配与方程编译接口。
因此它对应的是工程原理层，而非工程编译层，故 $T_2$ 成立。

另一方面，由第~\ref{sec:tex37-ch1} 章，
$S$ 把 LLM 收紧为可审计表达运行时；
由第~\ref{sec:tex37-ch4} 至第~\ref{sec:tex37-ch9} 章，
$S$ 把真正求解对象、白盒主权与纤维插件边界写成确定性合同接口；
由第~\ref{sec:tex37-ch11} 至第~\ref{sec:tex37-ch14} 章，
$S$ 又把主基底、扇区插件与签发协议推进到可签发方程系统。
这正是 $T_3$。

而 $T_3$ 所描述的恰好就是定义~\ref{def:tex37-ch18-runtime} 中
\[
  \mathsf{Eng}_{\mathrm{comp}}
\]
的内容，因此 $T_4$ 成立。
综上，$T_2\wedge T_4$ 推出 $T_5$，即
\[
  V_3\in \mathsf{Eng}_{\mathrm{prin}},
  \qquad
  S\in \mathsf{Eng}_{\mathrm{comp}}.
\]
证毕。
\end{proof}

\subsection[引理：工程悬空引理]{引理：$V_3$ 属于工程悬空类，但该“悬空”仅指编译未闭环}
\label{subsec:tex37-ch18-lemma}

\begin{lemma}[工程悬空引理]
\label{lem:tex37-ch18-pre}
在定义~\ref{def:tex37-ch18-intrinsic} 的口径下，
\[
  V_3\in \mathcal H_{\mathrm{eng}}.
\]
也就是说，Vol3 已经给出结构骨架、路径/证书/治理口径，
但仍缺少具体运行时合同、插件边界、主权分配与方程编译接口；
因此其“工程悬空”不是“完全无工程”，而是“工程编译未闭环”。
\end{lemma}

\begin{Certificate}[证书：工程悬空证书]
\label{cert:tex37-ch18-pre}
证书登记谓词链：
\[
\begin{aligned}
  L_1 &: \mathsf{Skel}(V_3)\ \text{已包含结构骨架、路径构造、证书登记与治理比较口径},\\
  L_2 &: L_1\Rightarrow V_3\ \text{不是“完全没有工程内容”},\\
  L_3 &: V_3\ \text{尚未把上述骨架闭合成 runtime contract、plugin boundary、}\\
      &\qquad\ \text{sovereignty allocation、equation compilation interface},\\
  L_4 &: L_2\wedge L_3\Rightarrow V_3\in \mathcal H_{\mathrm{eng}},\\
  L_5 &: L_4\Rightarrow \text{“工程悬空”应解释为“编译未闭环”，而非“无工程”}.
\end{aligned}
\]
\end{Certificate}

\begin{proof}[证明（基于数学符号推演逻辑的谓词因果链）]
由 $L_1$，Vol3 已经明确给出 engineering blueprint 所需的四类要件：
对象骨架、operator/path 机制、reading-path certificate，以及 governance-equivalent 比较口径。
因此 $L_2$ 成立，即它不能被描述为“完全没有工程内容”。

另一方面，由 $L_3$，Vol3 并未把这些骨架直接收束为
可执行运行时合同、白盒主权裁决接口、统一插件边界与可签发方程编译系统。
这意味着它的缺口不在于“无结构”，而在于“编译接口尚未闭合”。

故由 $L_2\wedge L_3$ 推出 $L_4$，
即
\[
  V_3\in \mathcal H_{\mathrm{eng}}.
\]
进一步，$L_4$ 的语义解释正是 $L_5$。
故引理成立。证毕。
\end{proof}

\subsection[命题：结构蓝图与编译规约]{命题：就语义规范场作为 AI/逻辑引擎架构主线而言，$V_3$ 给出结构蓝图，$S$ 给出编译规约}
\label{subsec:tex37-ch18-proposition}

\begin{proposition}[结构蓝图与编译规约命题]
\label{prop:tex37-ch18-runtime}
在定义~\ref{def:tex37-ch18-intrinsic} 与定义~\ref{def:tex37-ch18-runtime} 的口径下，
\[
  V_3
\]
给出语义工程的结构蓝图，
\[
  S
\]
给出语义工程的编译规约。
更准确地说，$V_3$ 的对象链应读作
\[
  \text{law-space object}
  \to
  \text{path}
  \to
  \text{certificate}
  \to
  \text{governance},
\]
而 $S$ 的对象链应读作
\[
  \text{逻辑占位}
  \to
  \text{解析}
  \to
  \text{审计}
  \to
  \text{运行时}
  \to
  \text{白盒主权}
  \to
  \text{领域插件}
  \to
  \text{可签发方程系统}.
\]
\end{proposition}

\begin{Certificate}[证书：结构蓝图与编译规约证书]
\label{cert:tex37-ch18-runtime}
证书登记谓词链：
\[
\begin{aligned}
  P_1 &: \cite{GaoZhengAppVol3}\Rightarrow
  V_3\ \text{把 law-space object、path、certificate、governance 串成统一骨架},\\
  P_2 &: P_1\Rightarrow V_3\ \text{给出的是结构蓝图而非运行时规约},\\
  P_3 &: \text{第~\ref{sec:tex37-ch1} 章至第~\ref{sec:tex37-ch14} 章}\Rightarrow
  S\ \text{显式组织 placeholder、parse、audit、runtime、}\\
      &\qquad\ \text{white-box sovereignty、plugin、equation system 的合同链},\\
  P_4 &: P_3\Rightarrow S\ \text{给出的是工程编译规约而非仅仅对象目录},\\
  P_5 &: P_2\wedge P_4\Rightarrow \text{结构蓝图与编译规约命题成立}.
\end{aligned}
\]
\end{Certificate}

\begin{proof}[证明（基于数学符号推演逻辑的谓词因果链）]
由 $P_1$，Vol3 的主链条先把语义对象压进 law-space，
再把 law-space 上的 operator/path 结构压成 reading-path certificate，
最后接入 governance-equivalent 的比较与抽取接口。
因此 $P_2$ 成立：这套工作回答的是“蓝图怎样搭”，而不是“接口怎样部署”。

由 $P_3$，$S$ 则把真正求解对象压到逻辑占位、解析与审计链，
再通过白盒主权与插件边界把接口收束到可签发方程系统。
因此 $P_4$ 成立：这套工作回答的是“蓝图怎样被真正写成工程制度”。

于是 $P_2\wedge P_4$ 推出 $P_5$。
故命题成立。证毕。
\end{proof}

\subsection[命题：工程悬空弱表述纠偏]{命题：若说“仅 $V_3$ 是工程悬空的”，则方向对但措辞偏重}
\label{subsec:tex37-ch18-weak}

\begin{proposition}[“仅 $V_3$ 是工程悬空的”弱表述命题]
\label{prop:tex37-ch18-weak}
若说
\[
  \text{“仅 }V_3\text{ 是工程悬空的”},
\]
则该表述\textbf{方向对，但措辞偏重}。
它只有在“工程悬空”被理解为“工程原理已在、工程编译未闭环”时才成立；
若把“工程悬空”理解成“完全没有工程内容”，则对 $V_3$ 不成立。
\end{proposition}

\begin{Certificate}[证书：弱表述纠偏证书]
\label{cert:tex37-ch18-weak}
证书登记谓词链：
\[
\begin{aligned}
  W_1 &: \mathsf{Skel}(V_3)\ \text{包含任务示例、路径构造、证书记录、治理等价与实现导向 schematic patterns},\\
  W_2 &: W_1\Rightarrow V_3\ \text{并非“无工程”},\\
  W_3 &: V_3\ \text{仍缺少 runtime contract、plugin boundary、sovereignty allocation、}\\
      &\qquad\ \text{equation compilation interface},\\
  W_4 &: W_3\Rightarrow V_3\ \text{在工程编译层意义下仍属 }\mathcal H_{\mathrm{eng}},\\
  W_5 &: W_2\wedge W_4\Rightarrow \text{“仅 }V_3\text{ 是工程悬空的”只有在“编译未闭环”意义下才成立}.
\end{aligned}
\]
\end{Certificate}

\begin{proof}[证明（基于数学符号推演逻辑的谓词因果链）]
由 $W_1$，Vol3 已经含有工程导向的任务示例、路径机制与证书/治理记录，
故 $W_2$ 成立，即它并不能被判定为“纯粹无工程”。

另一方面，由 $W_3$，Vol3 又确实没有把这些内容闭合成
运行时合同、插件边界、主权分配与方程编译接口；
故 $W_4$ 成立，即它在工程编译层意义下仍属
\[
  \mathcal H_{\mathrm{eng}}.
\]

因此，$W_2\wedge W_4$ 推出 $W_5$：
“仅 $V_3$ 是工程悬空的”这句话只能在“原理骨架已在、编译闭环未成”的意义下成立。
若把它理解成“Vol3 完全没有工程内容”，则该表述为假。
故命题成立。证毕。
\end{proof}

\subsection{推论：严格改写、压缩一句话与更强口径}
\label{subsec:tex37-ch18-corollaries}

\begin{corollary}[严格改写推论]
\label{cor:tex37-ch18-rewrite}
对本章结论的严格改写可写成：
\[
  \boxed{
    V_3\ \text{并非没有工程，而是工程上仍偏“原理悬空/编译未闭环”；}
  }
\]
\[
  \boxed{
    S\ \text{填补了 }V_3\text{ 在具体工程构造上的关键空白。}
  }
\]
\end{corollary}

\begin{Certificate}[证书：工程两层分工定理 + 结构蓝图与编译规约命题]
\label{cert:tex37-ch18-rewrite}
证书登记谓词链：
\[
\begin{aligned}
  C_1 &: \text{定理~\ref{thm:tex37-ch18-v3}}\Rightarrow
  V_3\in \mathsf{Eng}_{\mathrm{prin}},\ S\in \mathsf{Eng}_{\mathrm{comp}},\\
  C_2 &: \text{命题~\ref{prop:tex37-ch18-runtime}}\Rightarrow
  V_3\ \text{给出结构蓝图，而 }S\ \text{给出编译规约},\\
  C_3 &: C_1\wedge C_2\Rightarrow \text{严格改写推论成立}.
\end{aligned}
\]
\end{Certificate}

\begin{proof}[证明（基于数学符号推演逻辑的谓词因果链）]
由 $C_1$，两份文稿已被分配到工程原理层与工程编译层；
由 $C_2$，这种分工又可进一步读成“蓝图”与“编译规约”的差异。
因此最严格的改写，不再是“有无工程”的二值判断，
而是“Vol3 的工程性停在原理骨架，$S$ 继续把它闭合成具体工程构造”。
故由 $C_3$，严格改写推论成立。证毕。
\end{proof}

\begin{corollary}[一句话压缩推论]
\label{cor:tex37-ch18-one-line}
若继续压缩成一句话，则可写成：
\[
  \boxed{
    V_3\ \text{是工程蓝图；}
    \qquad
    S\ \text{是工程编译层。}
  }
\]
\end{corollary}

\begin{Certificate}[证书：严格改写推论 + 弱表述纠偏命题]
\label{cert:tex37-ch18-one-line}
证书登记谓词链：
\[
\begin{aligned}
  C'_1 &: \text{推论~\ref{cor:tex37-ch18-rewrite}}\Rightarrow
  V_3\ \text{的工程性停在 blueprint 层，而 }S\ \text{补上 compile 层},\\
  C'_2 &: \text{命题~\ref{prop:tex37-ch18-weak}}\Rightarrow
  \text{“工程悬空”必须解释为“编译未闭环”而非“无工程”},\\
  C'_3 &: C'_1\wedge C'_2\Rightarrow \text{一句话压缩推论成立}.
\end{aligned}
\]
\end{Certificate}

\begin{proof}[证明（基于数学符号推演逻辑的谓词因果链）]
由 $C'_1$，Vol3 与 $S$ 的差异已被压成
\[
  \text{blueprint}
  \to
  \text{compile}
\]
的分工；
由 $C'_2$，其中的 blueprint 又不是“空白”，而是“未闭环”。
因此进一步把它们压成“工程蓝图”与“工程编译层”不会丢失实质信息。
故 $C'_3$ 成立。证毕。
\end{proof}

\begin{corollary}[更强口径推论]
\label{cor:tex37-ch18-position}
若再加强半步，则可写成：
\[
  \boxed{
    V_3\ \text{回答“这套语义富结构可以怎样用于工程”；}
  }
\]
\[
  \boxed{
    S\ \text{回答“这套语义富结构应当怎样被真正写成工程系统”。}
  }
\]
\end{corollary}

\begin{Certificate}[证书：工程两层分工定理 + 结构蓝图与编译规约命题]
\label{cert:tex37-ch18-position}
证书登记谓词链：
\[
\begin{aligned}
  C''_1 &: \text{定理~\ref{thm:tex37-ch18-v3}}\Rightarrow
  V_3\in \mathsf{Eng}_{\mathrm{prin}},\ S\in \mathsf{Eng}_{\mathrm{comp}},\\
  C''_2 &: \text{命题~\ref{prop:tex37-ch18-runtime}}\Rightarrow
  V_3\ \text{给出用法蓝图，而 }S\ \text{给出写法规约},\\
  C''_3 &: C''_1\wedge C''_2\Rightarrow \text{更强口径推论成立}.
\end{aligned}
\]
\end{Certificate}

\begin{proof}[证明（基于数学符号推演逻辑的谓词因果链）]
由 $C''_1$，Vol3 与 $S$ 已经被区分为工程原理层与工程编译层；
由 $C''_2$，这种区分进一步对应于“可以怎样用于工程”与“应当怎样写成工程”。
故 $C''_3$ 成立。证毕。
\end{proof}

\subsection{备注：最终建议口径}
\label{subsec:tex37-ch18-remarks}

\begin{remark}[边界补充：完整闭环还依赖 tex36 的求解链]
还应补一个边界，避免把
\[
  S
\]
说得过满。最完整的闭环应写成：
\[
  V_3:\ \text{语义富结构与工程蓝图},
\]
\[
  \texttt{tex36}:\ \text{求解链},
\]
\[
  S:\ \text{运行时合同 + 主权分配 + 插件结构 + 方程系统写法}.
\]
其中 \texttt{tex36} 对应
\cite{RepoTex36HomotopyTwistBijection}，
它提供“如何求解”的结构桥梁；因此只有 $V_3$、\texttt{tex36} 与 $S$ 三者联立时，才构成完整闭环。
\end{remark}

\begin{remark}[最终建议采用的总口径]
若把本章结论压成最后一版对外口径，我建议写成：
\[
  \boxed{
    \begin{aligned}
      &\text{您的判断基本成立；但应改成：}\\
      &\text{“}V_3\text{ 工程上并非空白，而是尚未编译闭环；}\\
      &S\text{ 补上了具体工程构造空白”。}
    \end{aligned}
  }
\]
或者再展开半句：
\[
  \boxed{
    \begin{aligned}
      &V_3\ \text{给出工程所需的抽象结构与原理骨架，但仍偏“悬空”；}\\
      &S\ \text{填补了它在具体工程构造层面的关键空白。}
    \end{aligned}
  }
\]
这比把 Vol3 说成“无工程”，或者把 $S$ 只压成“落地”，都更准确。
\end{remark}

\clearpage

%%%%%%%%%%%%%%%%%%%%%%%%%%%%%%%%%%%%%%%%%%%%%%%%%%%%%%%%%%%%
\section{形式逻辑：$S$ 使 $V_3$ 的语义规范场富结构在开放系统博弈中获得动态成立}
\label{sec:tex37-ch19}
%%%%%%%%%%%%%%%%%%%%%%%%%%%%%%%%%%%%%%%%%%%%%%%%%%%%%%%%%%%%

沿用定义~\ref{def:tex37-ch18-compare} 的记号，仍记
\[
  V_3:=\text{Vol3},
  \qquad
  S:=\text{本稿第~\ref{sec:tex37-ch1} 章至第~\ref{sec:tex37-ch17} 章}.
\]
本章把“$S$ 是否使 $V_3$ 在开放系统博弈下成立”这一判断，继续收紧为一个更精确的问题：
\[
  \text{它究竟是在“本体首创”意义下成立，还是在“动态求解闭环”意义下成立？}
\]

\begin{remark}[核心结论（最稳版本）]
本章将证明：若把“成立”理解为“把 Vol3 的语义规范场富结构嵌入开放系统博弈的上下文延拓--障碍修复--可审计投影闭环”，
则该判断成立；若把“成立”理解为“$S$ 从零独立发明了该富结构所需的一切本体”，则不宜说满。
\end{remark}

\begin{remark}[阅读导航：仓内外参照]
Vol3 提供 law-space、HACA/PACER、semantic gauge field、property-layer Jacobiator-gap、
certificate bundle 与治理背景 \cite{GaoZhengAppVol3,GaoZheng2025GFramework}；
\texttt{tex23} 中关于开放系统的非平凡 holonomy / 未知--未知口径，可作为“开放性”判据的抽象背景
\cite{RepoTex23OpenSystemGame}；
\texttt{tex36} 提供修复塔、路径积分支撑类、同伦扭曲正规形与一般开放系统博弈的求解桥
\cite{RepoTex36HomotopyTwistBijection,RepoTex35ExtendedCertificationTower}；
\texttt{tex38} 进一步把“比安基闭合是修复后的目标态，而不是原始表示的先验事实”压成离散证书
\cite{RepoTex38JacobiBianchiRepair}。

外部语言背景方面，本章继续借用联络、曲率、holonomy、obstruction 与 homotopy repair 的标准术语，
可参照 \cite{MacLane1998Categories,Hatcher2002AT,Weibel1994HomologicalAlgebra,KobayashiNomizu1963Foundations}。
\end{remark}

\subsection{定义：两种“成立”、Vol3 富结构底座与开放系统博弈}
\label{subsec:tex37-ch19-definitions}

\begin{definition}[Vol3 的语义规范场富结构底座]
\label{def:tex37-ch19-rv3}
定义
\[
  \mathcal R_{V_3}
\]
为 Vol3 在本章中的\textbf{语义规范场富结构底座}，它至少包含下列对象链：
\[
  \text{law-space}
  \to
  \text{HACA/PACER}
  \to
  \text{semantic gauge field}
  \to
  \text{Jacobi gap}
  \to
  \text{certificate bundle}.
\]
若需要把治理与应用语境也并入，则再附加
\[
  \text{reading-path certificate}
  \to
  \text{value geometry}
  \to
  \text{multi-agent/open-system governance}.
\]
该底座的核心职责是回答“富结构是什么、层次如何组织、障碍如何出现、证书如何束化”。
\end{definition}

\begin{definition}[开放系统博弈的本章口径]
\label{def:tex37-ch19-osg}
定义
\[
  \mathsf{OSG}
\]
为满足下列条件的问题语境类：
\begin{enumerate}
  \item \textbf{上下文持续延拓：}新问题 $q$ 进入时，不是直接在旧上下文 $b$ 上静态读值，而是先形成
  \[
    b^+=\operatorname{Ext}_q(b);
  \]
  \item \textbf{约束持续变化：}上下文延拓后，联络、曲率、证书依赖与引用接口可被重写；
  \item \textbf{障碍持续改写：}比安基缺陷、雅可比障碍与接口残差需要在新基点上重新计算，即
  \[
    \operatorname{Ob}_{b^+}\neq \operatorname{Ob}_b
  \]
  一般允许成立；
  \item \textbf{解答不能静态读值：}对外回答必须经过
  \[
    \text{延拓}
    \to
    \text{障碍更新}
    \to
    \text{修复塔选路}
    \to
    \text{可审计投影}
  \]
  的闭环，而不是把旧状态直接当作答案。
\end{enumerate}
因此，本章并不把“开放系统博弈”粗糙地理解为“参与者很多”，而把它收紧为
\[
  \text{上下文可变}
  +
  \text{障碍会重写}
  +
  \text{解必须经修复生成}.
\]
\end{definition}

\begin{definition}[两种“成立”]
\label{def:tex37-ch19-meanings}
定义两种不同的成立谓词：
\[
  \mathsf{Hold}^{\mathrm{ont}}(S,V_3)
  :\Longleftrightarrow
  S\ \text{首次独立发明了开放系统博弈所需的修复塔、障碍与雅可比间隙本体};
\]
\[
  \mathsf{Hold}^{\mathrm{dyn}}(S,V_3)
  :\Longleftrightarrow
  S\ \text{把 }\mathcal R_{V_3}\text{ 嵌入}
\]
\[
  \operatorname{Ext}_q
  \to
  \operatorname{Ob}
  \to
  \Gamma^{\mathrm{adm}}
  \to
  \Sigma_n
  \to
  \Theta_n
  \to
  \mathsf{PL}_{\mathrm{NL}}^{(0)}
  \to
  \Render^{\mathrm{LLM}}_\xi
\]
这一动态求解闭环，使其在开放系统博弈中成为可运行、可修复、可审计的求解机制。
\end{definition}

\begin{remark}[本章要证明的究竟是哪一种“成立”]
本章只主张
\[
  \mathsf{Hold}^{\mathrm{dyn}}(S,V_3)
\]
成立，而不主张
\[
  \mathsf{Hold}^{\mathrm{ont}}(S,V_3)
\]
成立。换言之，本章证明的是“动态成立”，不是“本体首创”。
\end{remark}

\subsection{公理（降级为定理）：开放系统博弈中的动态成立链可按回合递归构造}
\label{subsec:tex37-ch19-axiom}

\begin{theorem}[公理降级：动态成立链的递归构造定理（数学归纳法证书）]
\label{thm:tex37-ch19-dynamic-chain}
固定初始上下文基点
\[
  b_0\in \mathcal{B}_{\mathrm{ctx}},
\]
问题序列
\[
  q_0,q_1,\dots,q_N,
\]
以及输出模式序列
\[
  \xi_0,\xi_1,\dots,\xi_N\in \Xi.
\]
设对每个
\[
  k\in\{0,1,\dots,N\}
\]
都满足：
\begin{enumerate}
  \item 上下文延拓
  \[
    b_k^+=\operatorname{Ext}_{q_k}(b_k)
  \]
  良定义；
  \item 存在某个层级
  \[
    N_k\in\NN
  \]
  使可接受路径类集合
  \[
    \Gamma_{b_k^+}^{\mathrm{adm},(N_k)}
  \]
  非空；
  \item 每个
  \[
    \Gamma_{b_k^+}^{\mathrm{adm},(N_k)}
  \]
  都配有固定全序
  \[
    \prec_{b_k^+},
  \]
  从而唯一代表
  \[
    [\gamma_k^\star]
    :=
    \min\nolimits_{\prec_{b_k^+}}
    \Gamma_{b_k^+}^{\mathrm{adm},(N_k)}
  \]
  存在；
  \item 代表链按
  \[
    r_k:=\Sigma_{N_k}([\gamma_k^\star]),
    \qquad
    p_k:=\operatorname{pr}_{N_k\to 0}\bigl(\Theta_{N_k}(r_k)\bigr),
    \qquad
    t_k:=\Render^{\mathrm{LLM}}_{\xi_k}(p_k)
  \]
  定义；
  \item 当 $k<N$ 时，下一回合基点取为
  \[
    b_{k+1}:=\gamma_k^\star(1).
  \]
\end{enumerate}
则对每个
\[
  m\in\{0,1,\dots,N\},
\]
第 $m$ 回合的动态成立链
\[
  D_m:
  b_m
  \xrightarrow{\operatorname{Ext}_{q_m}}
  b_m^+
  \rightsquigarrow
  [\gamma_m^\star]
  \mapsto
  r_m
  \mapsto
  p_m
  \mapsto
  t_m
\]
都存在且唯一确定。
因此，开放系统博弈中的“上下文延拓--障碍修复--逻辑占位--审计表达”闭环，可以按回合数递归构造。
\end{theorem}

\begin{Certificate}[证书：基步 + 归纳步的动态成立链证书]
\label{cert:tex37-ch19-dynamic-chain}
定义谓词
\[
  P(m):
  \text{“前 $m+1$ 个回合的动态成立链 }
  D_0,D_1,\dots,D_m
  \text{ 已唯一构造完成”}.
\]
则数学归纳法证书记为
\[
  \pi_{\mathrm{osg}}
  =
  \bigl(
    P(0),\
    \forall m\in\{0,1,\dots,N-1\},\ P(m)\Rightarrow P(m+1)
  \bigr).
\]
其中归纳步显式调用的递推式为
\[
  b_{m+1}:=\gamma_m^\star(1),
  \qquad
  b_{m+1}^+=\operatorname{Ext}_{q_{m+1}}(b_{m+1}),
\]
\[
  [\gamma_{m+1}^\star]
  :=
  \min\nolimits_{\prec_{b_{m+1}^+}}
  \Gamma_{b_{m+1}^+}^{\mathrm{adm},(N_{m+1})},
\]
\[
  r_{m+1}:=\Sigma_{N_{m+1}}([\gamma_{m+1}^\star]),
  \qquad
  p_{m+1}:=\operatorname{pr}_{N_{m+1}\to 0}\bigl(\Theta_{N_{m+1}}(r_{m+1})\bigr),
  \qquad
  t_{m+1}:=\Render^{\mathrm{LLM}}_{\xi_{m+1}}(p_{m+1}).
\]
\end{Certificate}

\begin{proof}[证明（数学归纳法证书；基于数学符号推演逻辑的谓词因果链）]
\textbf{基步 $P(0)$：}
由条件(1)--(4)，第 $0$ 回合的各个对象
\[
  b_0^+,\quad
  [\gamma_0^\star],\quad
  r_0,\quad
  p_0,\quad
  t_0
\]
都被唯一确定。
因此
\[
  D_0:
  b_0
  \xrightarrow{\operatorname{Ext}_{q_0}}
  b_0^+
  \rightsquigarrow
  [\gamma_0^\star]
  \mapsto
  r_0
  \mapsto
  p_0
  \mapsto
  t_0
\]
存在且唯一，故 $P(0)$ 成立。

\textbf{归纳步：}
设 $P(m)$ 成立，其中
\[
  0\le m<N.
\]
则第 $m$ 回合动态成立链已唯一构造完成。
由条件(5)，下一回合起点
\[
  b_{m+1}:=\gamma_m^\star(1)
\]
唯一确定。
再由条件(1)，
\[
  b_{m+1}^+=\operatorname{Ext}_{q_{m+1}}(b_{m+1})
\]
良定义。
由条件(2) 与条件(3)，非空集合
\[
  \Gamma_{b_{m+1}^+}^{\mathrm{adm},(N_{m+1})}
\]
在固定全序
\[
  \prec_{b_{m+1}^+}
\]
下具有唯一最小元
\[
  [\gamma_{m+1}^\star].
\]
由条件(4)，
\[
  r_{m+1}:=\Sigma_{N_{m+1}}([\gamma_{m+1}^\star]),
  \qquad
  p_{m+1}:=\operatorname{pr}_{N_{m+1}\to 0}\bigl(\Theta_{N_{m+1}}(r_{m+1})\bigr),
  \qquad
  t_{m+1}:=\Render^{\mathrm{LLM}}_{\xi_{m+1}}(p_{m+1})
\]
也唯一确定。
故第 $m+1$ 回合链
\[
  D_{m+1}:
  b_{m+1}
  \xrightarrow{\operatorname{Ext}_{q_{m+1}}}
  b_{m+1}^+
  \rightsquigarrow
  [\gamma_{m+1}^\star]
  \mapsto
  r_{m+1}
  \mapsto
  p_{m+1}
  \mapsto
  t_{m+1}
\]
存在且唯一，从而
\[
  P(m)\Rightarrow P(m+1).
\]

由数学归纳法，
\[
  \forall m\in\{0,1,\dots,N\},\ P(m).
\]
故所有回合的动态成立链都可递归构造。证毕。
\end{proof}

\subsection{定理：若按动态成立义理解，则您的判断成立}
\label{subsec:tex37-ch19-theorem}

\begin{theorem}[动态成立定理]
\label{thm:tex37-ch19-dynamic-hold}
在定义~\ref{def:tex37-ch19-meanings} 的口径下，
\[
  \mathsf{Hold}^{\mathrm{dyn}}(S,V_3)
\]
成立。
等价地，
\[
  \boxed{
    \begin{gathered}
      S\ \text{确实使 }V_3\text{ 的语义规范场、}\\
      \text{语义纤维丛与证书束富结构，}\\
      \text{在开放系统博弈口径下得到动态成立。}
    \end{gathered}
  }
\]
更准确地说，
\[
  \boxed{
    \begin{gathered}
      V_3\ \text{给出语义规范场的富结构底座；}\\
      S\ \text{给出该富结构在开放系统博弈中的}\\
      \text{上下文化修复机制。}
    \end{gathered}
  }
\]
\end{theorem}

\begin{Certificate}[证书：结构底座 + 动态链嵌入证书]
\label{cert:tex37-ch19-dynamic-hold}
证书登记为下列对象链与谓词链：
\[
  V_3:
  \quad
  \text{law-space}
  \to
  \text{HACA/PACER}
  \to
  \text{semantic gauge field}
  \to
  \text{Jacobi gap}
  \to
  \text{certificate bundle};
\]
\[
  S:
  \quad
  (b,q)
  \to
  b^+=\operatorname{Ext}_q(b)
  \to
  \operatorname{Ob}_{b^+},J_{b^+}
  \to
  \Gamma_{b^+}^{\mathrm{adm}}
  \to
  \Sigma_n
  \to
  \Theta_n
  \to
  \mathsf{PL}_{\mathrm{NL}}^{(0)}
  \to
  \Render^{\mathrm{LLM}}_\xi.
\]
并登记五段谓词：
\[
\begin{aligned}
  T_1 &: V_3\ \text{显式给出 }\mathcal R_{V_3}\ \text{这一富结构底座},\\
  T_2 &: S\ \text{显式给出 }\operatorname{Ext}_q,\ \operatorname{Ob},\ \Gamma^{\mathrm{adm}},\ \Sigma_n,\ \Theta_n,\
    \Render^{\mathrm{LLM}}_\xi,\\
  T_3 &: \mathsf{OSG}\ \text{要求新问题进入时上下文与障碍被重写，而非静态读值},\\
  T_4 &: T_2\wedge T_3\Rightarrow \mathcal R_{V_3}\ \text{被嵌入“延拓 }\to\text{ 修复 }\to\text{ 投影”的动态闭环},\\
  T_5 &: T_1\wedge T_4\Rightarrow \mathsf{Hold}^{\mathrm{dyn}}(S,V_3).
\end{aligned}
\]
\end{Certificate}

\begin{proof}[证明（基于数学符号推演逻辑的谓词因果链）]
\textbf{第一步：$V_3$ 给的是富结构底座，但主重心仍偏结构论。}

由 \cite{GaoZhengAppVol3,GaoZheng2025GFramework}，Vol3 已经明确给出 law-space、
HACA/PACER、semantic gauge field、property-layer Jacobiator-gap、
certificate bundle，以及开放治理与多主体背景。
这说明 $T_1$ 成立：
\[
  V_3\ \text{并不缺富结构底座。}
\]
但这些对象的主要语义功能，仍是回答
\[
  \text{对象是什么}
  \to
  \text{层次如何组织}
  \to
  \text{障碍如何出现}
  \to
  \text{证书如何登记},
\]
即它首先确立
\[
  \mathcal R_{V_3}
\]
的内禀成立。

\textbf{第二步：$S$ 把“开放系统博弈”具体压成“上下文延拓--障碍重写--修复选路”。}

由第四章的定义~\ref{def:tex37-ch4-base}、
定义~\ref{def:tex37-ch4-obstruction}、
定义~\ref{def:tex37-ch4-action} 与
定义~\ref{def:tex37-ch4-projection}，
$S$ 明确写出了
\[
  b\xrightarrow{\operatorname{Ext}_q} b^+,
\]
\[
  \operatorname{Ob}_{b^+},
  \qquad
  \Gamma_{b^+}^{\mathrm{adm}},
  \qquad
  \Sigma_n,
  \qquad
  \Theta_n,
  \qquad
  \Render^{\mathrm{LLM}}_\xi.
\]
这正是 $T_2$。
再由定义~\ref{def:tex37-ch19-osg}，开放系统博弈的最小特征不是参与者个数，而是
\[
  \text{上下文可被延拓}
  +
  \text{障碍会被改写}
  +
  \text{解必须经修复生成},
\]
即 $T_3$。
于是 $T_2\wedge T_3$ 给出：$S$ 已把开放性压成严格动力学，而不是抽象口号。

\textbf{第三步：在 $S$ 中，雅可比间隙从“结构量”推进成“求解驱动量”。}

由第四章引理~\ref{lem:tex37-ch4-bianchi-not-prior}，
上下文延拓并不先验保持比安基闭合；
由定义~\ref{def:tex37-ch4-action}，
可接受路径类由
\[
  \operatorname{Ob}_{\gamma(1)}
  \preceq_{\mathrm{lex}}
  \varepsilon_b
\]
决定，而可行性作用量
\[
  \mathcal{S}_{\mathrm{feas},b}
\]
直接对
\[
  |\mathcal{B}\!i|,
  \ |J|,
  \ |\Delta_{\mathrm{cert}}|,
  \ |\Delta_{\mathrm{ref}}|
\]
积分。
因此，雅可比间隙与总障碍不再只是“结构缺口的说明书”，而成为“为何需要修复、为何必须选路、为何不能直接读值”的驱动力。

\textbf{第四步：修复塔把 Vol3 富结构第一次压成开放系统中的可解对象。}

由第四章定理~\ref{thm:tex37-ch4-object}、
定理~\ref{thm:tex37-ch4-answer-as-repair} 与
推论~\ref{cor:tex37-ch4-llm-projection}，
真正被求解的对象不是表面文本，而是修复塔上的允许路径类
\[
  [\gamma]\in \Gamma_b^{\mathrm{adm},(n)}.
\]
其后才经由
\[
  [\gamma]
  \to
  \Sigma_n([\gamma])
  \to
  \Theta_n
  \to
  \mathsf{PL}_{\mathrm{NL}}^{(0)}
  \to
  \Render^{\mathrm{LLM}}_\xi
\]
投影成可审计文本。
这就得到 $T_4$：
\[
  \mathcal R_{V_3}
\]
已被嵌入
\[
  \text{延拓}
  \to
  \text{障碍更新}
  \to
  \text{修复塔选路}
  \to
  \text{逻辑占位}
  \to
  \text{审计表达}
\]
这一动态求解闭环。

\textbf{第五步：因此成立的是“动态成立”，而不是“本体首创”。}

由前四步，$T_1$ 与 $T_4$ 同时成立，故 $T_5$ 成立，即
\[
  \mathsf{Hold}^{\mathrm{dyn}}(S,V_3).
\]
也就是说，$S$ 的贡献不在于替代 Vol3 去发明全部本体，
而在于把 Vol3 的富结构推进为开放系统博弈中的上下文驱动障碍修复系统。
故定理成立。证毕。
\end{proof}

\subsection{引理：$V_3$ 单独给出的是富结构底座，而非完整的上下文驱动求解链}
\label{subsec:tex37-ch19-lemma}

\begin{lemma}[Vol3 单独并不推出开放系统博弈下的动态求解闭环]
\label{lem:tex37-ch19-static-vs-dynamic}
仅由 $V_3$ 中的
\[
  \mathcal R_{V_3}
\]
以及开放治理背景，不能直接推出对任意新问题 $q$ 都存在
\[
  b
  \xrightarrow{\operatorname{Ext}_q}
  b^+
  \to
  \operatorname{Ob}_{b^+}
  \to
  \Gamma_{b^+}^{\mathrm{adm}}
  \to
  \Sigma_n
  \to
  \Theta_n
  \to
  \mathsf{PL}_{\mathrm{NL}}^{(0)}
  \to
  \Render^{\mathrm{LLM}}_\xi
\]
这一完整求解链。
更严格地说，$V_3$ 主要确立的是富结构底座与工程蓝图；把该蓝图压成上下文驱动的动态求解闭环，则是
\texttt{tex36} 与 $S$ 的继续工作。
\end{lemma}

\begin{Certificate}[证书：结构底座存在，但动态闭环仍需补链]
\label{cert:tex37-ch19-static-vs-dynamic}
证书登记谓词链：
\[
\begin{aligned}
  L_1 &: V_3\ \text{已给出 }\mathcal R_{V_3}\ \text{及其治理/证书背景},\\
  L_2 &: \text{开放系统博弈下的动态求解闭环需要 }\operatorname{Ext}_q,\ \Gamma^{\mathrm{adm}},\ \Sigma_n,\ \Theta_n,\
    \Render^{\mathrm{LLM}}_\xi,\\
  L_3 &: L_2\ \text{中的上下文延拓、修复塔选路与投影链，在 }S\ \text{第四章才被显式写成接口},\\
  L_4 &: \texttt{tex36}\ \text{补出“路径积分 }\leftrightarrow\text{ 同伦扭曲正规形 }\leftrightarrow\text{ 修复塔”求解桥},\\
  L_5 &: L_1\wedge L_3\wedge L_4\Rightarrow V_3\ \text{单独不足以推出完整动态求解闭环}.
\end{aligned}
\]
\end{Certificate}

\begin{proof}[证明（基于数学符号推演逻辑的谓词因果链）]
由 \cite{GaoZhengAppVol3} 与第~\ref{sec:tex37-ch18} 章定理~\ref{thm:tex37-ch18-v3}，
$V_3$ 已经位于
\[
  \mathsf{Eng}_{\mathrm{prin}}
\]
层，故 $L_1$ 成立。
这意味着 Vol3 已经给出结构对象、路径口径、证书登记与治理比较的蓝图。

然而，若要把新问题
\[
  q
\]
送入开放系统博弈的动态求解核，就必须显式给出
\[
  \operatorname{Ext}_q,
  \qquad
  \Gamma_{b^+}^{\mathrm{adm}},
  \qquad
  \Sigma_n,
  \qquad
  \Theta_n,
  \qquad
  \Render^{\mathrm{LLM}}_\xi,
\]
即 $L_2$。
这些对象在本稿第四章才被以“上下文延拓--障碍重写--修复塔选路--可审计投影”的链条写成严格接口，
故 $L_3$ 成立。

另一方面，由 \cite{RepoTex36HomotopyTwistBijection,RepoTex35ExtendedCertificationTower}，
\texttt{tex36} 已把路径积分支撑类、同伦扭曲正规形与修复塔的求解桥压成显式可调用结构，故 $L_4$ 成立。
于是
\[
  L_1\wedge L_3\wedge L_4
\]
推出 $L_5$：
$V_3$ 单独给出的是富结构底座，而把该底座闭合为开放系统博弈下的动态求解链，还需要 \texttt{tex36} 与 $S$ 的继续工作。
故引理成立。证毕。
\end{proof}

\subsection{命题：在 $S$ 中，雅可比间隙与总障碍从结构量升级为求解驱动量}
\label{subsec:tex37-ch19-proposition}

\begin{proposition}[Jacobi gap 的求解驱动化命题]
\label{prop:tex37-ch19-jacobi-driver}
在定义~\ref{def:tex37-ch4-obstruction} 与定义~\ref{def:tex37-ch4-action} 的口径下，
\[
  J_b,
  \qquad
  \operatorname{Ob}_b
\]
不再只是描述结构不闭合的量，而是决定修复算子
\[
  \mathcal R,
\]
可接受路径类
\[
  \Gamma_b^{\mathrm{adm}},
\]
以及作用泛函
\[
  \mathcal S_b
\]
的核心驱动量。
因此，雅可比间隙在 $S$ 中从“结构证据”升级为“开放系统求解核心”。
\end{proposition}

\begin{Certificate}[证书：障碍进入修复算子、允许类与作用泛函]
\label{cert:tex37-ch19-jacobi-driver}
证书登记谓词链：
\[
\begin{aligned}
  P_1 &: V_3\ \text{中的 Jacobi gap 首先是结构不闭合/不相容的证据量},\\
  P_2 &: \Gamma_b^{\mathrm{adm}}
  =
  \{[\gamma]\in\Gamma_b\mid \operatorname{Ob}_{\gamma(1)}\preceq_{\mathrm{lex}}\varepsilon_b\},\\
  P_3 &: \operatorname{Ob}(\mathcal R(x))\preceq_{\mathrm{lex}}\operatorname{Ob}(x)\ \text{表明修复算子直接受障碍驱动},\\
  P_4 &: \mathcal{S}_{\mathrm{feas},b}([\gamma])=\int_\gamma(\alpha_1|\mathcal{B}
    \!i|+\alpha_2|J|+\alpha_3|\Delta_{\mathrm{cert}}|+\alpha_4|\Delta_{\mathrm{ref}}|)\,\mathrm{d}s,\\
  P_5 &: P_2\wedge P_3\wedge P_4\Rightarrow J_b,\operatorname{Ob}_b\ \text{从结构量升级为求解驱动量}.
\end{aligned}
\]
\end{Certificate}

\begin{proof}[证明（基于数学符号推演逻辑的谓词因果链）]
由 \cite{GaoZhengAppVol3}，Vol3 中的 Jacobi gap 主要承担“记录何处失配、何处需高阶填充”的结构证据角色，
故 $P_1$ 成立。

再由定义~\ref{def:tex37-ch4-action}，
一条路径类是否属于
\[
  \Gamma_b^{\mathrm{adm}}
\]
并不是自由指定的，而是由终点障碍是否压到可接受边界内决定，即 $P_2$。
同一一定义还规定修复算子满足
\[
  \operatorname{Ob}(\mathcal R(x))
  \preceq_{\mathrm{lex}}
  \operatorname{Ob}(x),
\]
故 $P_3$ 成立：修复为何必要、修复朝哪个方向推进，都由障碍向量直接驱动。

进一步，由 $P_4$，
可行性作用量直接对
\[
  |\mathcal{B}\!i|,
  \ |J|,
  \ |\Delta_{\mathrm{cert}}|,
  \ |\Delta_{\mathrm{ref}}|
\]
积分。
这意味着路径积分的排序准则本身就由障碍量控制，而不是在障碍之外另设一套无关指标。

因此，
\[
  P_2\wedge P_3\wedge P_4
\]
推出 $P_5$：
在 $S$ 中，雅可比间隙与总障碍已经进入“允许类生成 + 修复算子收缩 + 作用量排序”三条核心机制，
故它们不再只是静态描述量，而是开放系统求解核心。证毕。
\end{proof}

\subsection{推论：严格改写、压缩一句话与开放系统链翻译}
\label{subsec:tex37-ch19-corollaries}

\begin{corollary}[严格改写推论]
\label{cor:tex37-ch19-rewrite}
您的原判断可被严格改写为：
\[
  \boxed{
    \begin{aligned}
      &S\ \text{使 }V_3\text{ 的语义规范场/语义纤维丛富结构，}\\
      &\text{在开放系统博弈中不再只是静态优雅对象，}\\
      &\text{而成为由“上下文延拓 }\to\text{ 雅可比障碍 }\to\text{ 修复塔选路”驱动的动态可解系统。}
    \end{aligned}
  }.
\]
\end{corollary}

\begin{Certificate}[证书：动态成立定理 + Jacobi 驱动化命题]
\label{cert:tex37-ch19-rewrite}
证书登记谓词链：
\[
\begin{aligned}
  C_1 &: \text{定理~\ref{thm:tex37-ch19-dynamic-hold}}\Rightarrow
  S\ \text{已把 }\mathcal R_{V_3}\text{ 嵌入开放系统博弈的动态闭环},\\
  C_2 &: \text{命题~\ref{prop:tex37-ch19-jacobi-driver}}\Rightarrow
  \text{该动态闭环由雅可比障碍与总障碍直接驱动},\\
  C_3 &: C_1\wedge C_2\Rightarrow \text{严格改写推论成立}.
\end{aligned}
\]
\end{Certificate}

\begin{proof}[证明（基于数学符号推演逻辑的谓词因果链）]
由 $C_1$，$S$ 已经把 $V_3$ 的富结构从“对象底座”推进到“开放系统中的动态闭环”；
由 $C_2$，这个闭环的驱动核不是抽象修辞，而是上下文延拓后重写的雅可比障碍与总障碍。
于是 $C_1\wedge C_2$ 推出 $C_3$，即推论成立。证毕。
\end{proof}

\begin{corollary}[一句话压缩推论]
\label{cor:tex37-ch19-one-line}
若继续压缩成一句话，则可写成：
\[
  \boxed{
    \begin{aligned}
      &V_3\ \text{说明“富结构是什么”；}\\
      &S\ \text{说明“在开放系统博弈里，这个富结构怎样失闭合、怎样被修复、怎样产生解”。}
    \end{aligned}
  }
\]
\end{corollary}

\begin{Certificate}[证书：富结构底座定义 + 动态成立定理 + 上下文驱动引理]
\label{cert:tex37-ch19-one-line}
证书登记谓词链：
\[
\begin{aligned}
  C'_1 &: \text{定义~\ref{def:tex37-ch19-rv3}}\Rightarrow V_3\ \text{主要回答“富结构是什么”},\\
  C'_2 &: \text{定理~\ref{thm:tex37-ch19-dynamic-hold}}\Rightarrow S\ \text{把该富结构送入动态求解闭环},\\
  C'_3 &: \text{引理~\ref{lem:tex37-ch19-static-vs-dynamic}}\Rightarrow
  \text{该闭环并非 }V_3\text{ 单独就已显式写成},\\
  C'_4 &: C'_1\wedge C'_2\wedge C'_3\Rightarrow \text{一句话压缩推论成立}.
\end{aligned}
\]
\end{Certificate}

\begin{proof}[证明（基于数学符号推演逻辑的谓词因果链）]
由 $C'_1$，Vol3 的直接长项在于富结构底座本身；
由 $C'_2$，$S$ 的直接长项在于把该底座放进开放系统中的动态求解闭环；
由 $C'_3$，这种闭环化恰恰不是 Vol3 已经显式完成的部分。
故 $C'_1\wedge C'_2\wedge C'_3$ 推出 $C'_4$，即一句话压缩推论成立。证毕。
\end{proof}

\begin{corollary}[开放系统链翻译推论]
\label{cor:tex37-ch19-osg-chain}
若把“开放系统博弈”进一步翻译成本文体系语言，则可严格写成：
\[
  \boxed{
    \begin{aligned}
      &\text{新问题进入}
      \Longrightarrow
      \text{上下文基底延拓}
      \Longrightarrow
      \text{局部联络与曲率更新}\\
      &\Longrightarrow
      \text{比安基缺陷/雅可比障碍显现}
    \end{aligned}
  }
\]
\[
  \boxed{
    \begin{aligned}
      &\text{修复塔上允许路径类生成}
      \Longrightarrow
      \text{路径积分在允许类中选代表}\\
      &\Longrightarrow
      \text{逻辑占位投影}
      \Longrightarrow
      \text{审计表达输出。}
    \end{aligned}
  }
\]
而这条链，正是 $S$ 对 $V_3$ 的最大补完之一。
\end{corollary}

\begin{Certificate}[证书：动态成立链递归构造 + 第四章求解链]
\label{cert:tex37-ch19-osg-chain}
证书登记谓词链：
\[
\begin{aligned}
  C''_1 &: \text{定理~\ref{thm:tex37-ch19-dynamic-chain}}\Rightarrow
  \text{各回合存在 }b_k\to b_k^+\to[\gamma_k^\star]\to r_k\to p_k\to t_k\ \text{链},\\
  C''_2 &: \text{第四章定理~\ref{thm:tex37-ch4-answer-as-repair} 与命题~\ref{prop:tex37-ch4-path-integral-selector}}
    \Rightarrow\\
         &\qquad \text{“修复塔允许类 + 路径积分选代表”是求解核，而非表面表达},\\
  C''_3 &: C''_1\wedge C''_2\Rightarrow \text{开放系统链翻译推论成立}.
\end{aligned}
\]
\end{Certificate}

\begin{proof}[证明（基于数学符号推演逻辑的谓词因果链）]
由 $C''_1$，开放系统中的每个问题回合都可以被压成
\[
  b_k
  \xrightarrow{\operatorname{Ext}_{q_k}}
  b_k^+
  \rightsquigarrow
  [\gamma_k^\star]
  \mapsto
  r_k
  \mapsto
  p_k
  \mapsto
  t_k.
\]
把其中的
\[
  b_k^+
\]
按定义~\ref{def:tex37-ch4-obstruction} 展开，就得到“联络/曲率更新与障碍显现”；
把其中的
\[
  [\gamma_k^\star]
\]
按第四章命题~\ref{prop:tex37-ch4-path-integral-selector} 展开，就得到“修复塔允许类生成与路径积分选代表”；
把其中的
\[
  p_k\to t_k
\]
按第一章与第四章的运行时结论展开，就得到“逻辑占位投影与审计表达输出”。
因此 $C''_1\wedge C''_2$ 推出 $C''_3$，即推论成立。证毕。
\end{proof}

\subsection{备注：边界、完整闭环与最终建议口径}
\label{subsec:tex37-ch19-remarks}

\begin{remark}[若把“成立”理解成“本体首创”，则不宜说满]
若把您的原话理解为
\[
  \mathsf{Hold}^{\mathrm{ont}}(S,V_3),
\]
即
\[
  S\ \text{第一次独立创造了“修复塔 + 雅可比障碍 + 开放系统博弈”这整套本体},
\]
那就不宜说满。
因为由 \cite{GaoZhengAppVol3,RepoTex23OpenSystemGame,RepoTex36HomotopyTwistBijection,RepoTex38JacobiBianchiRepair} 可见，
law-space、Jacobi gap、开放系统判据、修复塔、路径积分、比安基/雅可比障碍与正规形对象链，都已经在前置链条中存在。
\end{remark}

\begin{remark}[若把“成立”理解成“动态成立”，则判断非常准确]
若把您的原话理解为
\[
  \mathsf{Hold}^{\mathrm{dyn}}(S,V_3),
\]
即
\[
  S\ \text{把 }V_3\text{ 的富结构真正放进开放系统博弈的求解内核里},
\]
那我认为这是非常准确的。
因为 $S$ 的关键增量，不是再定义一套新的顶层名词表，
而是把
\[
  \text{上下文延拓}
  \to
  \text{障碍重写}
  \to
  \text{修复塔选路}
  \to
  \text{可审计投影}
\]
这一动态闭环写成了严格接口。
\end{remark}

\begin{remark}[完整闭环版本与最终建议句]
若强调最完整闭环，则更稳妥的总口径应写成：
\[
  \boxed{
    \begin{aligned}
      &V_3\ \text{给出语义规范场的内禀富结构；}\\
      &\texttt{tex36}\ \text{给出“如何求解”的结构桥。}
    \end{aligned}
  }
\]
\[
  \boxed{
    \begin{aligned}
      &S\ \text{给出该富结构在开放系统博弈中的}\\
      &\text{上下文延拓、雅可比障碍、修复塔选路与审计投影机制。}
    \end{aligned}
  }
\]
若只保留一句最凝练且最稳的表述，则建议写成：
\[
  \boxed{
    \begin{aligned}
      &V_3\ \text{给出语义规范场的内禀富结构；}\\
      &S\ \text{给出该富结构在开放系统博弈中的上下文化修复机制。}
    \end{aligned}
  }
\]
这一句比泛泛地说“工程落地”更准确，因为它直接点明：
\[
  \text{开放性来自上下文延拓，}
  \qquad
  \text{求解性来自障碍修复。}
\]
\end{remark}

\section{形式逻辑：\texorpdfstring{$S$}{S} 把 \texorpdfstring{$V_3$}{V3} 的协变语义宇宙编码为上下文索引的同伦修复塔}
\label{sec:tex37-ch20}
%%%%%%%%%%%%%%%%%%%%%%%%%%%%%%%%%%%%%%%%%%%%%%%%%%%%%%%%%%%%

沿用定义~\ref{def:tex37-ch18-compare}、定义~\ref{def:tex37-ch19-meanings} 的记号，仍记
\[
  V_3:=\text{Vol3},
  \qquad
  S:=\text{本稿第~\ref{sec:tex37-ch1} 章至第~\ref{sec:tex37-ch19} 章}.
\]
本章不再使用“工程落地”“边缘算子命中”“无条件全局最优”这类宽口径，
而把您的原判断继续收紧为如下链条：
\[
  \text{协变语义宇宙}
  \to
  \text{上下文索引修复塔}
  \to
  \text{总障碍向量}
  \to
  \text{同伦扭曲下降}
  \to
  \text{允许路径类选类}.
\]

\begin{remark}[核心结论（先给出最稳口径）]
本章将证明：更准确的说法不是“$S$ 让 $V_3$ 命中某个工程边缘算子并直接给出无条件全局最优解”，
而是：
\[
  \boxed{
    \begin{aligned}
      &V_3\ \text{负责建立基于主纤维丛版广义非交换李代数投影的}\\
      &\text{内禀协变富结构语义宇宙；}\\
      &S\ \text{负责把这一语义宇宙编码为以上下文为基底的同伦修复塔，}\\
      &\text{并把上下文延拓诱导的比安基缺陷、雅可比障碍、证书断裂与引用断裂}\\
      &\text{统一写成总障碍向量；}\\
      &\text{随后通过修复算子诱导的同伦扭曲下降、允许路径类约束与作用泛函选类，}\\
      &\text{把“上下文修复”实现为对可行且近优路径类代表的收敛选取。}
    \end{aligned}
  }.
\]
若再附加强极小代表存在条件，则该近优代表才可升级为可行域上的全局最优代表。
\end{remark}

\begin{remark}[阅读导航：承接的仓内外文献]
本章直接承接 Vol3 的对象底座
\[
  \text{law-space}
  \to
  \text{HACA/PACER}
  \to
  \text{semantic gauge field}
  \to
  \text{Jacobi gap}
  \to
  \text{certificate bundle}
\]
及其治理背景 \cite{GaoZhengAppVol3,GaoZheng2025GFramework}；
承接第~\ref{sec:tex37-ch4} 章关于上下文基底、总障碍向量、修复算子、允许路径类与作用泛函的严格接口；
承接第~\ref{sec:tex37-ch19} 章关于“动态成立”的总判断；
并调用 \texttt{tex36} 的修复塔/路径积分/同伦扭曲正规形桥
\cite{RepoTex36HomotopyTwistBijection,RepoTex35ExtendedCertificationTower}，
以及 \texttt{tex38} 关于“比安基闭合是修复后目标态而非原始表示先验事实”的离散证书
\cite{RepoTex38JacobiBianchiRepair}。

外部术语背景方面，本章仍借用范畴、同伦、同调与联络几何的标准语言，
可参照 \cite{MacLane1998Categories,Hatcher2002AT,Weibel1994HomologicalAlgebra,KobayashiNomizu1963Foundations}。
\end{remark}

\subsection{定义：协变语义宇宙、上下文索引修复塔与术语收紧}
\label{subsec:tex37-ch20-definitions}

\begin{definition}[两份文稿的职责对象]
\label{def:tex37-ch20-objects}
定义
\[
  \mathcal{U}_{\mathrm{sem}}^{\mathrm{cov}}
  :=
  \mathcal R_{V_3}
\]
为 Vol3 的\textbf{协变语义宇宙}。按定义~\ref{def:tex37-ch19-rv3}，它至少由下列对象链生成：
\[
  \text{law-space}
  \to
  \text{HACA/PACER}
  \to
  \text{semantic gauge field}
  \to
  \text{Jacobi gap}
  \to
  \text{certificate bundle}.
\]
若把治理语境一并并入，则再附加
\[
  \text{reading-path certificate}
  \to
  \text{value geometry}
  \to
  \text{open-system governance}.
\]

再定义
\[
  \mathcal{T}_{\mathrm{ctx}}^{(N)}
  :=
  \bigl(
    \mathcal{B}_{\mathrm{ctx}},
    \pi_N:\mathcal{E}_{\mathrm{ctx}}^{(N)}\to \mathcal{B}_{\mathrm{ctx}},
    \operatorname{Ext}_q,
    \operatorname{Ob},
    \mathcal R,
    \Gamma^{\mathrm{adm}},
    \mathcal S,
    \Sigma_n,
    \Theta_n,
    \operatorname{pr}_{n\to 0},
    \Render^{\mathrm{LLM}}_\xi
  \bigr)
\]
为 $S$ 中的\textbf{上下文索引 $N$ 阶同伦修复塔系统}，
其中各个构件分别由定义~\ref{def:tex37-ch4-base}、定义~\ref{def:tex37-ch4-obstruction}、
定义~\ref{def:tex37-ch4-action}、定义~\ref{def:tex37-ch4-projection} 与第一章的高阶同伦扭曲双射给出。
\end{definition}

\begin{remark}[术语收紧：本章的“上下文修复”]
本章中的“上下文修复”一律收紧为
\[
  \text{上下文延拓后的障碍修复}.
\]
也就是说，新问题 $q$ 进入时先发生
\[
  b\mapsto b^+=\operatorname{Ext}_q(b),
\]
随后才重新计算
\[
  \mathcal{B}\!i_b,\quad
  J_b,\quad
  \Delta_{\mathrm{cert},b},\quad
  \Delta_{\mathrm{ref},b},
\]
并在修复塔上进行下降、选类与投影。
因此，本章不把“下文修复”理解为对旧上下文状态的静态读值，而理解为对延拓后障碍态的动态修复。
\end{remark}

\subsection{公理（降级为定理）：从协变语义宇宙到上下文索引修复塔的编码可按阶递归构造}
\label{subsec:tex37-ch20-axiom}

\begin{theorem}[公理降级：$\mathcal{U}_{\mathrm{sem}}^{\mathrm{cov}}$ 向 $\mathcal{T}_{\mathrm{ctx}}^{(N)}$ 的编码可按阶递归构造（数学归纳法证书）
  ]
\label{thm:tex37-ch20-encoding}
固定 $N\in\NN$。若满足下列条件：
\begin{enumerate}
  \item 第一章定理~\ref{thm:tex37-axiom-downgrade} 已给出逐阶兼容的高阶同伦扭曲双射
  \[
    \Theta_n:\mathsf{NF}_G^{(n)}\to \mathsf{PL}_{\mathrm{NL}}^{(n)};
  \]
  \item 第四章定义~\ref{def:tex37-ch4-base} 已给出上下文基底上的 $N$ 阶同伦修复塔纤维丛
  \[
    \pi_N:\mathcal E_{\mathrm{ctx}}^{(N)}\to \mathcal B_{\mathrm{ctx}};
  \]
  \item 第四章定理~\ref{thm:tex37-ch4-tower-induction} 已保证各阶纤维与允许路径类可递归构造；
  \item 第四章定义~\ref{def:tex37-ch4-projection} 已给出
  \[
    \Sigma_n,\qquad
    \Theta_n,\qquad
    \operatorname{pr}_{n\to 0}
  \]
  组成的投影链。
\end{enumerate}
则对每个
\[
  n\in\{0,1,\dots,N\},
\]
都存在唯一的阶段编码
\[
  \mathfrak T_n(b)
  :=
  \bigl(
    \mathcal E_b^{(n)},
    \operatorname{Ob}_b,
    \Gamma_b^{\mathrm{adm},(n)},
    \Sigma_n,
    \Theta_n,
    \operatorname{pr}_{n\to 0}
  \bigr),
\]
并由这些阶段编码拼成上下文索引修复塔系统
\[
  \mathcal U_{\mathrm{sem}}^{\mathrm{cov}}
  \rightsquigarrow
  \mathcal T_{\mathrm{ctx}}^{(N)}.
\]
因此，把 Vol3 的协变富结构编码进上下文索引修复塔，不是额外口号，而是可递归构造的定理性对象。
\end{theorem}

\begin{Certificate}[证书：基步 + 归纳步的编码链证书]
\label{cert:tex37-ch20-encoding}
定义谓词
\[
  P(n):
  \left\{
  \begin{aligned}
    &\text{对每个 }b\in\mathcal B_{\mathrm{ctx}},\ \text{阶段编码 }\mathfrak T_n(b)\text{ 已定义；}\\
    &\text{且它保留 Vol3 富结构的规范语义/证书/障碍投影链。}
  \end{aligned}
  \right.
\]
则数学归纳法证书记为
\[
  \pi_{\mathrm{enc}}
  =
  \bigl(
    P(0),\
    \forall n\in\{0,1,\dots,N-1\},\ P(n)\Rightarrow P(n+1)
  \bigr).
\]
其中：
\[
  P(0)
  \ \text{调用第一章基阶双射 }\Theta_0\text{ 与第四章基阶纤维 }\mathcal E_b^{(0)};
\]
\[
  \begin{aligned}
    P(n)\Rightarrow P(n+1)
    \ \text{调用第四章修复塔递归构造与}\\
    \text{第一章逐阶兼容式 }
    \Theta_{n+1}(x,\eta)=\bigl(\Theta_n(x),\Psi_x^{(n)}(\eta)\bigr).
  \end{aligned}
\]
\end{Certificate}

\begin{proof}[证明（数学归纳法证书；基于数学符号推演逻辑的谓词因果链）]
\textbf{基步 $P(0)$：}
由定义~\ref{def:tex37-ch4-base}，
\[
  \mathcal E_b^{(0)}
\]
已经编码了上下文基点 $b$ 上允许的规范形、证书依赖与展开层级；
由第一章定理~\ref{thm:tex37-axiom-downgrade}，
\[
  \Theta_0:\mathsf{NF}_G^{(0)}\to \mathsf{PL}_{\mathrm{NL}}^{(0)}
\]
为双射。
因此，基阶对象
\[
  \mathfrak T_0(b)
  =
  \bigl(
    \mathcal E_b^{(0)},
    \operatorname{Ob}_b,
    \Gamma_b^{\mathrm{adm},(0)},
    \Sigma_0,
    \Theta_0,
    \operatorname{pr}_{0\to 0}
  \bigr)
\]
良定义，故 $P(0)$ 成立。

\textbf{归纳步：}
设 $P(n)$ 成立。则第 $n$ 阶编码
\[
  \mathfrak T_n(b)
\]
已存在。
由第四章定理~\ref{thm:tex37-ch4-tower-induction}，
\[
  \mathcal E_b^{(n+1)}
  \qquad\text{与}\qquad
  \Gamma_b^{\mathrm{adm},(n+1)}
\]
可由第 $n$ 阶纤维与路径类递归构造。
再由第一章定理~\ref{thm:tex37-axiom-downgrade}，
\[
  \Theta_{n+1}
\]
由
\[
  \Theta_{n+1}(x,\eta)
  =
  \bigl(
    \Theta_n(x),
    \Psi_x^{(n)}(\eta)
  \bigr)
\]
唯一确定。
于是
\[
  \mathfrak T_{n+1}(b)
  :=
  \bigl(
    \mathcal E_b^{(n+1)},
    \operatorname{Ob}_b,
    \Gamma_b^{\mathrm{adm},(n+1)},
    \Sigma_{n+1},
    \Theta_{n+1},
    \operatorname{pr}_{n+1\to 0}
  \bigr)
\]
良定义且唯一。
故
\[
  P(n)\Rightarrow P(n+1).
\]

由数学归纳法，
\[
  \forall n\in\{0,1,\dots,N\},\ P(n).
\]
因此各阶编码都存在且唯一，并拼成
\[
  \mathcal T_{\mathrm{ctx}}^{(N)}.
\]
故定理成立。证毕。
\end{proof}

\subsection{命题：您原判断的可成立版本}
\label{subsec:tex37-ch20-proposition}

\begin{proposition}[原判断的可成立版本]
\label{prop:tex37-ch20-rewriteable}
您的原判断可严格改写为：
\[
  \boxed{
    V_3\ \text{负责建立语义宇宙的内禀协变富结构。}
  }
\]
\[
  \boxed{
    S\ \text{负责把该富结构编码为上下文索引的同伦修复塔，}
  }
\]
\[
  \boxed{
    \begin{aligned}
      &\text{并把上下文延拓所诱导的比安基缺陷、雅可比障碍、}\\
      &\text{证书断裂与引用断裂统一写成可修复总障碍，}
    \end{aligned}
  }
\]
\[
  \boxed{
    \begin{aligned}
      &\text{再以同伦扭曲下降 + 允许路径类约束 + 作用泛函选类的方式，}\\
      &\text{给出“上下文延拓后的障碍修复”机制。}
    \end{aligned}
  }
\]
\end{proposition}

\begin{Certificate}[证书：$V_3$ 的职责链]
\label{cert:tex37-ch20-v3-role}
Vol3 的主责可压成对象链
\[
  \text{law-space}
  \to
  \text{HACA/PACER}
  \to
  \text{semantic gauge field}
  \to
  \text{Jacobi gap}
  \to
  \text{certificate bundle}.
\]
换言之，
\[
  V_3:
  \qquad
  \mathcal U_{\mathrm{sem}}^{\mathrm{cov}}
  \qquad
  \text{协变语义对象宇宙的建立}.
\]
\end{Certificate}

\begin{Certificate}[证书：$S$ 的职责链]
\label{cert:tex37-ch20-s-role}
$S$ 的关键结构链可压成
\[
  b\in\mathcal B_{\mathrm{ctx}}
  \quad\mapsto\quad
  \pi_N:\mathcal E_{\mathrm{ctx}}^{(N)}\to\mathcal B_{\mathrm{ctx}},
\]
其中
\[
  \mathcal E_b^{(N)}=\pi_N^{-1}(b)
\]
表示在上下文基点 $b$ 上全部允许的规范形、扭曲数据、证书依赖、修复分支与展开层级。
随后，对新问题 $q$ 定义
\[
  \operatorname{Ext}_q:\mathcal B_{\mathrm{ctx}}\to \mathcal B_{\mathrm{ctx}},
  \qquad
  b\mapsto b^+,
\]
并在 $b$ 上定义
\[
  A_b,\qquad
  F_b,\qquad
  J_b,\qquad
  \mathcal{B}\!i_b:=D_{A_b}F_b.
\]
再把总障碍统一写成
\[
  \operatorname{Ob}_b
  =
  \bigl(
    |\mathcal{B}\!i_b|,
    |J_b|,
    |\Delta_{\mathrm{cert},b}|,
    |\Delta_{\mathrm{ref},b}|
  \bigr),
\]
把修复算子写成
\[
  \mathcal R:\mathcal E_{\mathrm{ctx}}^{(N)}\to \mathcal E_{\mathrm{ctx}}^{(N)},
  \qquad
  \operatorname{Ob}(\mathcal R(x))
  \preceq_{\mathrm{lex}}
  \operatorname{Ob}(x),
\]
把允许路径类写成
\[
  \Gamma_b^{\mathrm{adm}}
  =
  \bigl\{
    [\gamma]\in\Gamma_b
    \mid
    \operatorname{Ob}_{\gamma(1)}
    \preceq_{\mathrm{lex}}
    \varepsilon_b
  \bigr\},
\]
并把作用泛函写成
\[
  \mathcal S_b([\gamma])
  =
  \mathcal S_{\mathrm{feas},b}([\gamma])
  +
  \lambda\,\mathcal S_{\mathrm{opt},b}([\gamma]).
\]
\end{Certificate}

\begin{proof}[证明（基于数学符号推演逻辑的谓词因果链）]
由证书~\ref{cert:tex37-ch20-v3-role}，
$V_3$ 的主轴可压成
\[
  \text{law-space}
  \to
  \text{HACA/PACER}
  \to
  \text{semantic gauge field}
  \to
  \text{Jacobi gap}
  \to
  \text{certificate bundle},
\]
这说明它首先在建立
\[
  \mathcal U_{\mathrm{sem}}^{\mathrm{cov}}.
\]

由证书~\ref{cert:tex37-ch20-s-role}，
$S$ 并不是停留在“结构对象已经存在”的层面，
而是进一步引入
\[
  \mathcal B_{\mathrm{ctx}},
  \quad
  \operatorname{Ext}_q,
  \quad
  \mathcal E_{\mathrm{ctx}}^{(N)},
  \quad
  \operatorname{Ob}_b,
  \quad
  \mathcal R,
  \quad
  \Gamma_b^{\mathrm{adm}},
  \quad
  \mathcal S_b.
\]
于是一个新问题 $q$ 不再是“对既有语义宇宙直接读值”，
而是先诱导
\[
  b\mapsto b^+,
\]
再在 $b^+$ 上检查并修复
\[
  \mathcal{B}\!i_{b^+},\quad
  J_{b^+},\quad
  \Delta_{\mathrm{cert},b^+},\quad
  \Delta_{\mathrm{ref},b^+}.
\]

再由 \texttt{tex38} 的结论 \cite{RepoTex38JacobiBianchiRepair}，
比安基闭合更稳妥地应被理解为修复后的目标态，而不是上下文延拓的先验事实。
因此，本命题中的“上下文修复”必须收紧为“上下文延拓后的障碍修复”。
命题成立。证毕。
\end{proof}

\subsection{定理：\texorpdfstring{$S$}{S} 确实把 \texorpdfstring{$V_3$}{V3} 的协变富结构编码为上下文驱动的同伦修复塔求解系统}
\label{subsec:tex37-ch20-theorem}

\begin{theorem}[协变富结构的修复塔编码定理]
\label{thm:tex37-ch20-encoded-solver}
\[
  \boxed{
    S\ \text{确实把 }V_3\text{ 的协变富结构，编码成了一个“上下文驱动的同伦修复塔求解系统”。}
  }
\]
更具体地说，一个新问题的求解链被压成
\[
  b
  \mapsto
  b^+
  \rightsquigarrow
  [\gamma]
  \to
  \Sigma_n([\gamma])
  \to
  \Theta_n(\Sigma_n([\gamma]))
  \to
  \operatorname{pr}_{n\to 0}
  \to
  \Render^{\mathrm{LLM}}_\xi,
\]
其中真正的内核对象是
\[
  [\gamma]\in \Gamma_b^{\mathrm{adm},(n)},
\]
而表面文本只是不变量经投影后的显示壳。
\end{theorem}

\begin{Certificate}[证书：协变富结构 $\Rightarrow$ 上下文化求解链]
\label{cert:tex37-ch20-encoded-solver}
证书登记为下列五段谓词链：
\[
\begin{aligned}
  T_1 &: V_3\ \text{已给出 }\mathcal U_{\mathrm{sem}}^{\mathrm{cov}}\ \text{这一协变富结构底座},\\
  T_2 &: S\ \text{已给出 }\mathcal B_{\mathrm{ctx}},\ \operatorname{Ext}_q,\ \mathcal E_{\mathrm{ctx}}^{(N)},\
    \operatorname{Ob}_b,\ \Gamma_b^{\mathrm{adm}},\ \mathcal S_b,\\
  T_3 &: T_2\Rightarrow \text{新问题 }q\text{ 必先诱导 }b\mapsto b^+\text{，而非在旧基点静态读值},\\
  T_4 &: T_3\Rightarrow \text{求解对象被压成修复塔上的允许路径类 }[\gamma],\\
  T_5 &: T_1\wedge T_4\Rightarrow S\ \text{把 }V_3\text{ 的协变富结构编码为上下文驱动求解系统}.
\end{aligned}
\]
\end{Certificate}

\begin{proof}[证明（基于数学符号推演逻辑的谓词因果链）]
由定义~\ref{def:tex37-ch20-objects} 与证书~\ref{cert:tex37-ch20-v3-role}，
$V_3$ 已经给出
\[
  \mathcal U_{\mathrm{sem}}^{\mathrm{cov}},
\]
即 $T_1$。

由定义~\ref{def:tex37-ch20-objects}、定义~\ref{def:tex37-ch4-base}、
定义~\ref{def:tex37-ch4-obstruction}、定义~\ref{def:tex37-ch4-action} 与定义~\ref{def:tex37-ch4-projection}，
$S$ 明确引入了
\[
  \mathcal B_{\mathrm{ctx}},
  \quad
  \operatorname{Ext}_q,
  \quad
  \mathcal E_{\mathrm{ctx}}^{(N)},
  \quad
  \operatorname{Ob}_b,
  \quad
  \Gamma_b^{\mathrm{adm}},
  \quad
  \mathcal S_b,
\]
故 $T_2$ 成立。

由 $T_2$，新问题进入时并不直接输出文本，而是先发生
\[
  b\mapsto b^+,
\]
即 $T_3$。
再由第四章定理~\ref{thm:tex37-ch4-object} 与定理~\ref{thm:tex37-ch4-answer-as-repair}，
真正被求解的对象不是表面文本，而是
\[
  [\gamma]\in \Gamma_b^{\mathrm{adm},(n)},
\]
随后才经由
\[
  [\gamma]
  \to
  \Sigma_n([\gamma])
  \to
  \Theta_n(\Sigma_n([\gamma]))
  \to
  \operatorname{pr}_{n\to 0}
  \to
  \Render^{\mathrm{LLM}}_\xi
\]
投影为文本，故 $T_4$ 成立。

于是
\[
  T_1\wedge T_4
\]
推出 $T_5$：$S$ 确实把 $V_3$ 的协变富结构编码成上下文驱动的同伦修复塔求解系统。
定理成立。证毕。
\end{proof}

\subsection{引理：上下文缺陷被统一写成“几何 + 代数 + 证书 + 引用”的总障碍向量}
\label{subsec:tex37-ch20-lemma-obstruction}

\begin{lemma}[上下文缺陷的四分量统一写法]
\label{lem:tex37-ch20-obstruction}
\[
  \boxed{
    \begin{aligned}
      &S\ \text{把上下文缺陷正式写成了}\\
      &\text{“几何缺陷 + 代数缺陷 + 证书缺陷 + 引用缺陷”的统一障碍向量。}
    \end{aligned}
  }
\]
更准确地说，
\[
  \operatorname{Ob}_b
  =
  \bigl(
    |\mathcal{B}\!i_b|,
    |J_b|,
    |\Delta_{\mathrm{cert},b}|,
    |\Delta_{\mathrm{ref},b}|
  \bigr)
\]
就是 $S$ 中上下文缺陷的正式数学编码。
\end{lemma}

\begin{Certificate}[证书：四类缺陷统一入障碍向量]
\label{cert:tex37-ch20-obstruction}
证书登记谓词链：
\[
\begin{aligned}
  L_1 &: |\mathcal{B}\!i_b|\ \text{衡量联络--曲率--上下文延拓是否闭合},\\
  L_2 &: |J_b|\ \text{衡量语义括号--推理拼接--结构依赖是否出现雅可比残差},\\
  L_3 &: |\Delta_{\mathrm{cert},b}|\ \text{衡量证书链是否断裂},\\
  L_4 &: |\Delta_{\mathrm{ref},b}|\ \text{衡量引用接口是否断裂},\\
  L_5 &: L_1\wedge L_2\wedge L_3\wedge L_4\Rightarrow \operatorname{Ob}_b\ \text{统一编码四类上下文缺陷}.
\end{aligned}
\]
\end{Certificate}

\begin{proof}[证明（基于数学符号推演逻辑的谓词因果链）]
由定义~\ref{def:tex37-ch4-obstruction}，
\[
  \operatorname{Ob}_b
  =
  \bigl(
    |\mathcal{B}\!i_b|,
    |J_b|,
    |\Delta_{\mathrm{cert},b}|,
    |\Delta_{\mathrm{ref},b}|
  \bigr).
\]
其中第一分量正是联络、曲率与上下文延拓的闭合性残差，即 $L_1$；
第二分量记录括号结构与依赖拼接的雅可比残差，即 $L_2$；
后两项分别记录证书链与引用接口的断裂程度，即 $L_3$、$L_4$。

因此，
\[
  L_1\wedge L_2\wedge L_3\wedge L_4
\]
推出 $L_5$：
在 $S$ 中，“上下文缺陷”不是口头描述，而是被正式编码为四分量总障碍向量。
引理成立。证毕。
\end{proof}

\subsection{引理：同伦扭曲负责变形机制，路径积分负责选类机制}
\label{subsec:tex37-ch20-lemma-twist}

\begin{lemma}[同伦扭曲与路径积分的职责分离]
\label{lem:tex37-ch20-twist-vs-select}
\[
  \boxed{
    \begin{aligned}
      &\text{在 }S\text{ 中，同伦扭曲负责“怎样沿修复方向变形”，}\\
      &\text{路径积分/作用泛函负责“从允许路径类中选哪条近优代表”。}
    \end{aligned}
  }
\]
因此，严格地说，
\[
  \text{同伦扭曲}
  \neq
  \text{直接给最终规则}，
\]
而是
\[
  \text{同伦扭曲负责变形机制，路径积分负责选类机制。}
\]
\end{lemma}

\begin{Certificate}[证书：修复算子与作用泛函的职责分离]
\label{cert:tex37-ch20-twist-vs-select}
证书登记谓词链：
\[
\begin{aligned}
  L_1 &: \operatorname{Ob}(\mathcal R(x))\preceq_{\mathrm{lex}}\operatorname{Ob}(x)\ \text{表明 }\mathcal R\text{ 控制修复下降方向}
    ,\\
  L_2 &: x\rightsquigarrow \mathcal R(x)\rightsquigarrow \mathcal R^2(x)\rightsquigarrow\cdots\ \text{构成可修复变形链},\\
  L_3 &: \mathcal S_b:\Gamma_b^{\mathrm{adm}}\to\mathbb R\cup\{+\infty\}\ \text{只在允许路径类上比较代价},\\
  L_4 &: L_1\wedge L_2\wedge L_3\Rightarrow \text{同伦扭曲负责变形，路径积分负责选类}.
\end{aligned}
\]
\end{Certificate}

\begin{proof}[证明（基于数学符号推演逻辑的谓词因果链）]
由定义~\ref{def:tex37-ch4-action}，
\[
  \operatorname{Ob}(\mathcal R(x))
  \preceq_{\mathrm{lex}}
  \operatorname{Ob}(x),
\]
故 $\mathcal R$ 的职责是把状态沿着“障碍下降”的方向推进，这给出 $L_1$。
把 $\mathcal R$ 反复作用，即得到
\[
  x\rightsquigarrow \mathcal R(x)\rightsquigarrow \mathcal R^2(x)\rightsquigarrow\cdots,
\]
这正是修复变形链，即 $L_2$。

同一一定义还给出
\[
  \mathcal S_b:\Gamma_b^{\mathrm{adm}}\to\mathbb R\cup\{+\infty\},
\]
它只在允许路径类上比较可行性代价与优化代价，而不负责产生修复方向本身，这就是 $L_3$。
因此，
\[
  L_1\wedge L_2\wedge L_3
\]
推出 $L_4$：
同伦扭曲负责变形机制，路径积分负责选类机制。
引理成立。证毕。
\end{proof}

\subsection{命题：“边缘算子命中”的更稳收紧写法}
\label{subsec:tex37-ch20-prop-operator}

\begin{proposition}[“同伦扭曲命中边缘算子”的收紧写法]
\label{prop:tex37-ch20-operator-phrase}
若要与 $S$ 的主记号完全对齐，则
\[
  \text{“同伦扭曲命中边缘算子”}
\]
这一说法应收紧为：
\[
  \boxed{
    \text{同伦扭曲通过修复算子 }\mathcal R\text{ 降低总障碍，}
    \text{并在允许路径类 }\Gamma_b^{\mathrm{adm}}\text{ 中诱导可行代表。}
  }
\]
也就是说，$S$ 的最稳主词是
\[
  \mathcal R,\qquad
  \Gamma_b^{\mathrm{adm}},\qquad
  \mathcal S_b,\qquad
  \Sigma_n,\qquad
  \Theta_n,
\]
而不是单独的“边缘算子命中”。
\end{proposition}

\begin{Certificate}[证书：$S$ 的主记号在修复算子/允许类/选类链上]
\label{cert:tex37-ch20-operator-phrase}
证书登记谓词链：
\[
\begin{aligned}
  P_1 &: \mathcal R\ \text{直接控制障碍下降},\\
  P_2 &: \Gamma_b^{\mathrm{adm}}\ \text{给出可接受路径类边界},\\
  P_3 &: \mathcal S_b\ \text{只在允许路径类上比较代表},\\
  P_4 &: P_1\wedge P_2\wedge P_3\Rightarrow \text{更稳的写法应以 }\mathcal R/\Gamma_b^{\mathrm{adm}}/\mathcal S_b\text{ 为主词}.
\end{aligned}
\]
\end{Certificate}

\begin{proof}[证明（基于数学符号推演逻辑的谓词因果链）]
由引理~\ref{lem:tex37-ch20-twist-vs-select}，
$S$ 已经把“变形机制”与“选类机制”严格分离：
\[
  \mathcal R
  \quad\text{负责下降，}\qquad
  \mathcal S_b
  \quad\text{负责在 }\Gamma_b^{\mathrm{adm}}\text{ 内比较代表}.
\]
因此 $P_1$、$P_2$、$P_3$ 同时成立。

既然 $S$ 的主求解链实际写成
\[
  \mathcal R
  \to
  \Gamma_b^{\mathrm{adm}}
  \to
  \mathcal S_b
  \to
  \Sigma_n
  \to
  \Theta_n,
\]
那么把整个链压缩成“边缘算子命中”就会遗漏修复、允许类与选类这三层核心接口。
故更稳妥的表述应以 $\mathcal R/\Gamma_b^{\mathrm{adm}}/\mathcal S_b$ 为主词。
命题成立。证毕。
\end{proof}

\subsection{命题：“全局最优解”的更稳收紧写法}
\label{subsec:tex37-ch20-prop-global}

\begin{proposition}[“逼近变分收敛（全局最优解）”的收紧写法]
\label{prop:tex37-ch20-global-opt}
“逼近变分收敛到全局最优解”这一说法方向是对的，但必须加条件。
在 $S$ 的最稳口径下，直接由第四章定理~\ref{thm:tex37-ch4-dual-convergence} 得到的是
\[
  \boxed{
    \text{可行性收敛 + 最优性收敛}
  }
\]
以及
\[
  \mathcal S_b([\gamma_k])
  \to
  \inf_{[\gamma]\in\Gamma_b^{\mathrm{adm}}}\mathcal S_b([\gamma]),
\]
或至少存在达到该下确界的极小化子序列。
因此，更稳的写法是：
\[
  \boxed{
    \text{逼近变分收敛到可行域上的下确界/近优路径类代表。}
  }
\]
只有在附加强极小代表存在条件 \textnormal{(H3)} 时，才可把它升级为：
\[
  \boxed{
    \text{可行域上的全局最优路径类代表。}
  }
\]
\end{proposition}

\begin{Certificate}[证书：双性收敛 + 极小代表存在条件]
\label{cert:tex37-ch20-global-opt}
证书登记谓词链：
\[
\begin{aligned}
  G_1 &: \text{第四章定理~\ref{thm:tex37-ch4-dual-convergence}}\Rightarrow
  \mathcal S_b([\gamma_k])\to \inf_{[\gamma]\in\Gamma_b^{\mathrm{adm}}}\mathcal S_b([\gamma])\\
      &\qquad\text{或至少存在达到该下确界的极小化子序列},\\
  G_2 &: \textnormal{(H3)}\Rightarrow \Gamma_b^{\mathrm{adm}}\ \text{上存在极小代表，或至少存在极小化序列},\\
  G_3 &: \neg\textnormal{(H3)}\ \text{的强形式}\Rightarrow \text{最稳结论只能写成下确界逼近/近优代表},\\
  G_4 &: G_1\wedge G_2\Rightarrow \text{在附加条件下才可升级为全局最优代表}.
\end{aligned}
\]
\end{Certificate}

\begin{proof}[证明（基于数学符号推演逻辑的谓词因果链）]
由第四章定理~\ref{thm:tex37-ch4-dual-convergence}，
在 \textnormal{(H1)}--\textnormal{(H4)} 条件下，修复序列给出的最稳结论是
\[
  \mathcal S_b([\gamma_k])
  \to
  \inf_{[\gamma]\in\Gamma_b^{\mathrm{adm}}}\mathcal S_b([\gamma]),
\]
或至少存在达到该下确界的极小化子序列，这就是 $G_1$。

同一定理中的 \textnormal{(H3)} 只保证“存在极小代表，或至少存在极小化序列”。
因此，若没有 \textnormal{(H3)} 的强形式，就只能说“收敛到可行域下确界”或“存在近优/极小化子序列”，这正是 $G_3$。
只有当 \textnormal{(H3)} 以强形式给出实际极小代表时，
下确界才可被某个路径类真正达到，于是才能升级为全局最优代表，即 $G_4$。

故“全局最优解”不能被写成 $S$ 的无条件普遍结论。
命题成立。证毕。
\end{proof}

\subsection{定理：您的整句判断的最严格改写}
\label{subsec:tex37-ch20-final}

\begin{theorem}[最严格改写定理]
\label{thm:tex37-ch20-final}
综合定理~\ref{thm:tex37-ch20-encoded-solver}、引理~\ref{lem:tex37-ch20-obstruction}、
引理~\ref{lem:tex37-ch20-twist-vs-select}、命题~\ref{prop:tex37-ch20-operator-phrase} 与
命题~\ref{prop:tex37-ch20-global-opt}，
您的原句最严格、最稳、且与 $V_3+\texttt{tex36}+S+\texttt{tex38}$ 同时对齐的写法是：
\[
  \boxed{
    \begin{aligned}
      &V_3\ \text{负责建立基于主纤维丛版广义非交换李代数投影的}\\
      &\text{内禀协变富结构语义宇宙；}\\
      &S\ \text{负责把这一语义宇宙编码为以语义上下文为基底的 }N\text{ 阶同伦修复塔，}\\
      &\text{并把新问题引发的上下文延拓缺陷统一写成}\\
      &\text{比安基缺陷、雅可比障碍、证书断裂与引用断裂的总障碍向量；}\\
      &\text{随后通过修复算子诱导的同伦扭曲下降、允许路径类约束与作用泛函选类，}\\
      &\text{把“上下文修复”实现为对可行且近优路径类代表的收敛选取；}\\
      &\text{在附加强极小代表存在条件下，该近优代表可升级为可行域上的全局最优代表。}
    \end{aligned}
  }.
\]
\end{theorem}

\begin{Certificate}[证书：编码定理 + 四分量障碍 + 变形/选类分离 + 最优性条件]
\label{cert:tex37-ch20-final}
证书登记谓词链：
\[
\begin{aligned}
  F_1 &: \text{定理~\ref{thm:tex37-ch20-encoded-solver}}\Rightarrow
  S\ \text{已把 }V_3\text{ 的协变富结构编码成上下文驱动修复塔求解系统},\\
  F_2 &: \text{引理~\ref{lem:tex37-ch20-obstruction}}\Rightarrow
  \text{上下文缺陷被正式写成四分量总障碍向量},\\
  F_3 &: \text{引理~\ref{lem:tex37-ch20-twist-vs-select}}\Rightarrow
  \text{同伦扭曲负责变形，路径积分负责选类},\\
  F_4 &: \text{命题~\ref{prop:tex37-ch20-operator-phrase}}\Rightarrow
  \text{稳妥主词应是 }\mathcal R/\Gamma_b^{\mathrm{adm}}/\mathcal S_b,\\
  F_5 &: \text{命题~\ref{prop:tex37-ch20-global-opt}}\Rightarrow
  \text{“全局最优”只能在附加极小代表存在条件下成立},\\
  F_6 &: F_1\wedge F_2\wedge F_3\wedge F_4\wedge F_5\Rightarrow \text{最严格改写定理成立}.
\end{aligned}
\]
\end{Certificate}

\begin{proof}[证明（基于数学符号推演逻辑的谓词因果链）]
由 $F_1$，$S$ 已经把 $V_3$ 的协变富结构嵌入
\[
  \text{上下文延拓}
  \to
  \text{障碍重写}
  \to
  \text{修复塔选路}
  \to
  \text{占位投影}
  \to
  \text{审计表达}
\]
这条求解闭环。
由 $F_2$，所谓“上下文缺陷”被精确写成
\[
  \bigl(
    |\mathcal{B}\!i_b|,
    |J_b|,
    |\Delta_{\mathrm{cert},b}|,
    |\Delta_{\mathrm{ref},b}|
  \bigr),
\]
而不再是含混的口头概念。

由 $F_3$ 与 $F_4$，求解链内部的职责分工被严格拆开：
\[
  \mathcal R
  \quad\text{负责下降，}\qquad
  \Gamma_b^{\mathrm{adm}}
  \quad\text{负责可行边界，}\qquad
  \mathcal S_b
  \quad\text{负责选类}.
\]
因此，“边缘算子命中”并不是本章最稳的中心术语。

最后，由 $F_5$，
最稳结论只能写成“对可行且近优路径类代表的收敛选取”；
只有在极小代表存在的附加条件下，才可升级为“全局最优路径类代表”。
于是
\[
  F_1\wedge F_2\wedge F_3\wedge F_4\wedge F_5
\]
推出 $F_6$，即定理成立。证毕。
\end{proof}

\subsection{推论：压缩版表述}
\label{subsec:tex37-ch20-corollaries}

\begin{corollary}[一句话压缩推论]
\label{cor:tex37-ch20-one-line}
若压缩成一句话，则可写为：
\[
  \boxed{
    V_3\ \text{给出协变语义富结构宇宙；}
    \qquad
    S\ \text{给出该宇宙在开放上下文中的障碍修复动力学。}
  }
\]
\end{corollary}

\begin{Certificate}[证书：协变宇宙 + 障碍修复动力学]
\label{cert:tex37-ch20-one-line}
证书登记谓词链：
\[
\begin{aligned}
  C_1 &: \text{定义~\ref{def:tex37-ch20-objects}}\Rightarrow
  V_3\ \text{的长项是 }\mathcal U_{\mathrm{sem}}^{\mathrm{cov}},\\
  C_2 &: \text{定理~\ref{thm:tex37-ch20-encoded-solver}}\Rightarrow
  S\ \text{的长项是 }\mathcal T_{\mathrm{ctx}}^{(N)}\text{ 所给出的上下文化求解链},\\
  C_3 &: C_1\wedge C_2\Rightarrow \text{一句话压缩推论成立}.
\end{aligned}
\]
\end{Certificate}

\begin{proof}[证明（基于数学符号推演逻辑的谓词因果链）]
由 $C_1$，$V_3$ 的直接贡献是协变语义宇宙；
由 $C_2$，$S$ 的直接贡献是把这一宇宙推进成上下文延拓后的障碍修复动力学。
因此
\[
  C_1\wedge C_2
\]
推出 $C_3$，即推论成立。证毕。
\end{proof}

\begin{corollary}[开放系统博弈语义下的压缩推论]
\label{cor:tex37-ch20-osg}
若要突出开放系统博弈语义，则可写为：
\[
  \boxed{
    \begin{aligned}
      &V_3\ \text{负责“宇宙建模”；}\\
      &S\ \text{负责“在上下文博弈中，宇宙如何因缺陷而失闭合、又如何被修复并选路”。}
    \end{aligned}
  }
\]
\end{corollary}

\begin{Certificate}[证书：宇宙建模 + 上下文化修复选路]
\label{cert:tex37-ch20-osg}
证书登记谓词链：
\[
\begin{aligned}
  C'_1 &: \text{定理~\ref{thm:tex37-ch20-encoded-solver}}\Rightarrow
  \text{新问题的进入方式是 }b\mapsto b^+\rightsquigarrow [\gamma],\\
  C'_2 &: \text{引理~\ref{lem:tex37-ch20-obstruction}}\Rightarrow
  \text{失闭合被正式记录为总障碍向量},\\
  C'_3 &: \text{引理~\ref{lem:tex37-ch20-twist-vs-select}}\Rightarrow
  \text{修复与选路是两段不同但连续的机制},\\
  C'_4 &: C'_1\wedge C'_2\wedge C'_3\Rightarrow \text{开放系统博弈语义下的压缩推论成立}.
\end{aligned}
\]
\end{Certificate}

\begin{proof}[证明（基于数学符号推演逻辑的谓词因果链）]
由 $C'_1$，新问题的进入并不会静态读值，而会触发上下文延拓与路径类求解；
由 $C'_2$，失闭合被正式写成总障碍向量；
由 $C'_3$，修复与选路构成连续闭环。
故
\[
  C'_1\wedge C'_2\wedge C'_3
\]
推出 $C'_4$，即推论成立。证毕。
\end{proof}

\begin{corollary}[保留“同伦扭曲”措辞时的最稳压缩版]
\label{cor:tex37-ch20-twist}
若要保留“同伦扭曲”这一措辞，但又尽量贴近 $S$ 的主记号，则建议写成：
\[
  \boxed{
    S\ \text{将 }V_3\text{ 的协变语义宇宙，重写为“上下文缺陷驱动的同伦扭曲--修复选路系统”。}
  }
\]
\end{corollary}

\begin{Certificate}[证书：同伦扭曲保留 + 术语主词收紧]
\label{cert:tex37-ch20-twist}
证书登记谓词链：
\[
\begin{aligned}
  C''_1 &: \text{引理~\ref{lem:tex37-ch20-twist-vs-select}}\Rightarrow
  \text{“同伦扭曲”可以保留，但必须与“选类”分开表述},\\
  C''_2 &: \text{命题~\ref{prop:tex37-ch20-operator-phrase}}\Rightarrow
  \text{系统的稳妥主词是修复算子、允许类与作用泛函},\\
  C''_3 &: C''_1\wedge C''_2\Rightarrow \text{该压缩版成立}.
\end{aligned}
\]
\end{Certificate}

\begin{proof}[证明（基于数学符号推演逻辑的谓词因果链）]
由 $C''_1$，可以保留“同伦扭曲”这一词，但必须同时承认它只负责变形机制；
由 $C''_2$，系统的稳妥主词仍应落在修复算子、允许路径类与作用泛函上。
因此
\[
  C''_1\wedge C''_2
\]
推出 $C''_3$，即该压缩版成立。证毕。
\end{proof}

\subsection{备注：边界、术语与最终建议口径}
\label{subsec:tex37-ch20-remarks}

\begin{remark}[“全局最优”不能无条件直写]
$S$ 最稳的结论是：
\[
  \text{可行且近优}
  \qquad\text{或}\qquad
  \inf\text{-逼近}.
\]
只有在极小代表存在这类附加条件下，才升级成
\[
  \arg\min\text{ 级全局最优路径类代表}.
\]
因此，若写论文摘要或正式章节标题，最好避免把“全局最优”当作无条件普遍结论。
\end{remark}

\begin{remark}[“边缘算子”更像前序解释口径，不是 $S$ 的最稳主词]
若只看 $S$ 本身，最贴文的主词应是
\[
  \mathcal R,\qquad
  \Gamma^{\mathrm{adm}},\qquad
  \mathcal S,\qquad
  \Sigma_n,\qquad
  \Theta_n,\qquad
  \operatorname{Issue}/\operatorname{Cert}.
\]
因此，若要写成正式章节标题或摘要，建议把“边缘算子命中”收紧为“修复算子/允许路径类/正规形投影链”。
\end{remark}

\begin{remark}[本章的最终建议口径]
因此，您这句话最稳的最终版，不是
\[
  \text{$S$ 负责把宇宙编码为修复塔并命中边缘算子得到全局最优解，}
\]
而是
\[
  \boxed{
    \begin{aligned}
      &S\ \text{负责把 }V_3\text{ 的协变语义宇宙编码为上下文索引的同伦修复塔，}\\
      &\text{并把比安基/雅可比/证书/引用缺陷统一为障碍函数，}\\
      &\text{再通过同伦扭曲下降与变分选类，实现对上下文修复结果的可行且近优收敛。}
    \end{aligned}
  }.
\]
这一版最贴近本稿现有记号，也最便于继续挂接第~\ref{sec:tex37-ch4} 章与第~\ref{sec:tex37-ch19} 章的既有证明链。
\end{remark}

\clearpage

\section*{参考文献}
\begin{thebibliography}{99}

\bibitem{RepoTex20SemanticGaugeNormalization}
【仓内 TeX】语义函子场与语义规范场：词包判定--规范化接口的形式系统整理稿：
\code{NoteBook/notebook1/tex20_pub/gframework_semantic_functor_field_gauge_field_normalization_interface_formal_system.t
  ex}。

\bibitem{RepoTex21AuditedWrapper}
【仓内 TeX】语义函子正则化、验证与鲁棒性整理稿：
\code{NoteBook/notebook1/tex21_pub/gframework_semantic_functor_regex_normalization_validation_robustness_formal_system.t
  ex}。

\bibitem{RepoTex23SemanticBase}
【仓内 TeX】PTO/ETO 分解、语义规范场基座与审计包装整理稿：
\code{NoteBook/notebook1/tex23_pub/gframework_ptopb_etopb_operator_decomposition_formal_system.tex}。

\bibitem{RepoTex23PriorityGovernance}
【仓内 TeX】优先权固化、Priority Pack、Zenodo DOI 证据链与信息危害分级整理稿：
\code{NoteBook/notebook1/tex23_pub/gframework_ptopb_etopb_operator_decomposition_formal_system.tex}。

\bibitem{RepoTex23OpenSystemGame}
【仓内 TeX】开放系统、回合同伦类、UU/KU 障碍与工程设计开放博弈判据整理稿：
\code{NoteBook/notebook1/tex23_pub/gframework_ptopb_etopb_operator_decomposition_formal_system.tex}。

\bibitem{RepoTex36HomotopyTwistBijection}
【仓内 TeX】主丛联络的纤维丛空间、路径积分支撑类与同伦扭曲正规形的集合等价整理稿：
\code{NoteBook/notebook2/tex36_pub/gframework_principal_connection_fiber_space_path_integral_homotopy_twist_equivalence_
  formal_system.tex}。

\bibitem{RepoTex35ExtendedCertificationTower}
【仓内 TeX】基于扩展签发协议的主丛联络--路径积分--同伦修复塔与结构表示定理：
\code{NoteBook/notebook2/tex35_pub/gframework_extended_certification_tower.tex}。

\bibitem{RepoTex38JacobiBianchiRepair}
【仓内 TeX】离散拓扑基底上的 Jacobi--Bianchi 障碍修复、可行路径空间、离散路径积分与终点正规形整理稿：
\code{NoteBook/notebook2/tex38_pub/gframework_discrete_jacobi_bianchi_repair_tower_path_integral_terminal_normal_form_fo
  rmal_system.tex}。

\bibitem{RepoTex29RepairableObstruction}
【仓内 TeX】可修复障碍同伦类与广义分形：全集范畴与严格边界：
\code{NoteBook/notebook2/tex29_pub/gframework_repairable_obstruction_homotopy_class_vs_generalized_fractal_formal_system
  .tex}。

\bibitem{GaoZhengAppVol2}
Gao, Z. (2025).
\newblock GaoZheng G-Framework and GaoZheng G-Algebra: Applied Mathematics Volume II — LBOPB, Life-Science Monoids, and
  Generative Precision Medicine.
\newblock In \emph{GaoZheng G-Framework and GaoZheng G-Algebra: Applied Mathematics Volume II — LBOPB, Life-Science
  Monoids, and Generative Precision Medicine (v1.1)}.
\newblock Zenodo.
\newblock \url{https://doi.org/10.5281/zenodo.17744672}.


\bibitem{GaoZhengAppVol3}
Gao, Z. (2025).
\newblock GaoZheng G-Framework and GaoZheng G-Algebra: Applied Mathematics Volume III – HACA, PACER, and
  Certificate-Based AI.
\newblock In \emph{GaoZheng G-Framework and GaoZheng G-Algebra: Applied Mathematics Volume III – HACA, PACER, and
  Certificate-Based AI (v1.0)}.
\newblock Zenodo.
\newblock \url{https://doi.org/10.5281/zenodo.17762295}.


\bibitem{GaoZheng2025GFramework}
Gao, Z. (2025).
\newblock Meta-Mathematical Theory based on Pan-Logic Analysis and
Pan-Iterative Analysis (GaoZheng G-Framework) and the
Principal-Bundle-Based Generalized Noncommutative Lie Algebra
(GaoZheng G-Algebra).
\newblock In \emph{Meta-Mathematical Theory based on Pan-Logic Analysis and
Pan-Iterative Analysis (GaoZheng G-Framework) and the
Principal-Bundle-Based Generalized Noncommutative Lie Algebra
(GaoZheng G-Algebra): An Integrated Construction (v1.0)}.
\newblock Zenodo.
\newblock \url{https://doi.org/10.5281/zenodo.17651584}.


\bibitem{MacLane1998Categories}
Mac Lane, S. (1998).
\newblock \emph{Categories for the Working Mathematician} (2nd ed.).
\newblock Springer.

\bibitem{Hatcher2002AT}
Hatcher, A. (2002).
\newblock \emph{Algebraic Topology}.
\newblock Cambridge University Press.

\bibitem{Weibel1994HomologicalAlgebra}
Weibel, C. A. (1994).
\newblock \emph{An Introduction to Homological Algebra}.
\newblock Cambridge University Press.

\bibitem{KobayashiNomizu1963Foundations}
Kobayashi, S., \& Nomizu, K. (1963).
\newblock \emph{Foundations of Differential Geometry, Volume I}.
\newblock Interscience Publishers.

\end{thebibliography}

\end{CJK*}
\end{document}
